% Options for packages loaded elsewhere
\PassOptionsToPackage{unicode}{hyperref}
\PassOptionsToPackage{hyphens}{url}
\PassOptionsToPackage{dvipsnames,svgnames,x11names}{xcolor}
%
\documentclass[
  letterpaper,
  DIV=11,
  numbers=noendperiod]{scrreprt}

\usepackage{amsmath,amssymb}
\usepackage{iftex}
\ifPDFTeX
  \usepackage[T1]{fontenc}
  \usepackage[utf8]{inputenc}
  \usepackage{textcomp} % provide euro and other symbols
\else % if luatex or xetex
  \usepackage{unicode-math}
  \defaultfontfeatures{Scale=MatchLowercase}
  \defaultfontfeatures[\rmfamily]{Ligatures=TeX,Scale=1}
\fi
\usepackage{lmodern}
\ifPDFTeX\else  
    % xetex/luatex font selection
\fi
% Use upquote if available, for straight quotes in verbatim environments
\IfFileExists{upquote.sty}{\usepackage{upquote}}{}
\IfFileExists{microtype.sty}{% use microtype if available
  \usepackage[]{microtype}
  \UseMicrotypeSet[protrusion]{basicmath} % disable protrusion for tt fonts
}{}
\makeatletter
\@ifundefined{KOMAClassName}{% if non-KOMA class
  \IfFileExists{parskip.sty}{%
    \usepackage{parskip}
  }{% else
    \setlength{\parindent}{0pt}
    \setlength{\parskip}{6pt plus 2pt minus 1pt}}
}{% if KOMA class
  \KOMAoptions{parskip=half}}
\makeatother
\usepackage{xcolor}
\setlength{\emergencystretch}{3em} % prevent overfull lines
\setcounter{secnumdepth}{5}
% Make \paragraph and \subparagraph free-standing
\ifx\paragraph\undefined\else
  \let\oldparagraph\paragraph
  \renewcommand{\paragraph}[1]{\oldparagraph{#1}\mbox{}}
\fi
\ifx\subparagraph\undefined\else
  \let\oldsubparagraph\subparagraph
  \renewcommand{\subparagraph}[1]{\oldsubparagraph{#1}\mbox{}}
\fi

\usepackage{color}
\usepackage{fancyvrb}
\newcommand{\VerbBar}{|}
\newcommand{\VERB}{\Verb[commandchars=\\\{\}]}
\DefineVerbatimEnvironment{Highlighting}{Verbatim}{commandchars=\\\{\}}
% Add ',fontsize=\small' for more characters per line
\usepackage{framed}
\definecolor{shadecolor}{RGB}{241,243,245}
\newenvironment{Shaded}{\begin{snugshade}}{\end{snugshade}}
\newcommand{\AlertTok}[1]{\textcolor[rgb]{0.68,0.00,0.00}{#1}}
\newcommand{\AnnotationTok}[1]{\textcolor[rgb]{0.37,0.37,0.37}{#1}}
\newcommand{\AttributeTok}[1]{\textcolor[rgb]{0.40,0.45,0.13}{#1}}
\newcommand{\BaseNTok}[1]{\textcolor[rgb]{0.68,0.00,0.00}{#1}}
\newcommand{\BuiltInTok}[1]{\textcolor[rgb]{0.00,0.23,0.31}{#1}}
\newcommand{\CharTok}[1]{\textcolor[rgb]{0.13,0.47,0.30}{#1}}
\newcommand{\CommentTok}[1]{\textcolor[rgb]{0.37,0.37,0.37}{#1}}
\newcommand{\CommentVarTok}[1]{\textcolor[rgb]{0.37,0.37,0.37}{\textit{#1}}}
\newcommand{\ConstantTok}[1]{\textcolor[rgb]{0.56,0.35,0.01}{#1}}
\newcommand{\ControlFlowTok}[1]{\textcolor[rgb]{0.00,0.23,0.31}{#1}}
\newcommand{\DataTypeTok}[1]{\textcolor[rgb]{0.68,0.00,0.00}{#1}}
\newcommand{\DecValTok}[1]{\textcolor[rgb]{0.68,0.00,0.00}{#1}}
\newcommand{\DocumentationTok}[1]{\textcolor[rgb]{0.37,0.37,0.37}{\textit{#1}}}
\newcommand{\ErrorTok}[1]{\textcolor[rgb]{0.68,0.00,0.00}{#1}}
\newcommand{\ExtensionTok}[1]{\textcolor[rgb]{0.00,0.23,0.31}{#1}}
\newcommand{\FloatTok}[1]{\textcolor[rgb]{0.68,0.00,0.00}{#1}}
\newcommand{\FunctionTok}[1]{\textcolor[rgb]{0.28,0.35,0.67}{#1}}
\newcommand{\ImportTok}[1]{\textcolor[rgb]{0.00,0.46,0.62}{#1}}
\newcommand{\InformationTok}[1]{\textcolor[rgb]{0.37,0.37,0.37}{#1}}
\newcommand{\KeywordTok}[1]{\textcolor[rgb]{0.00,0.23,0.31}{#1}}
\newcommand{\NormalTok}[1]{\textcolor[rgb]{0.00,0.23,0.31}{#1}}
\newcommand{\OperatorTok}[1]{\textcolor[rgb]{0.37,0.37,0.37}{#1}}
\newcommand{\OtherTok}[1]{\textcolor[rgb]{0.00,0.23,0.31}{#1}}
\newcommand{\PreprocessorTok}[1]{\textcolor[rgb]{0.68,0.00,0.00}{#1}}
\newcommand{\RegionMarkerTok}[1]{\textcolor[rgb]{0.00,0.23,0.31}{#1}}
\newcommand{\SpecialCharTok}[1]{\textcolor[rgb]{0.37,0.37,0.37}{#1}}
\newcommand{\SpecialStringTok}[1]{\textcolor[rgb]{0.13,0.47,0.30}{#1}}
\newcommand{\StringTok}[1]{\textcolor[rgb]{0.13,0.47,0.30}{#1}}
\newcommand{\VariableTok}[1]{\textcolor[rgb]{0.07,0.07,0.07}{#1}}
\newcommand{\VerbatimStringTok}[1]{\textcolor[rgb]{0.13,0.47,0.30}{#1}}
\newcommand{\WarningTok}[1]{\textcolor[rgb]{0.37,0.37,0.37}{\textit{#1}}}

\providecommand{\tightlist}{%
  \setlength{\itemsep}{0pt}\setlength{\parskip}{0pt}}\usepackage{longtable,booktabs,array}
\usepackage{calc} % for calculating minipage widths
% Correct order of tables after \paragraph or \subparagraph
\usepackage{etoolbox}
\makeatletter
\patchcmd\longtable{\par}{\if@noskipsec\mbox{}\fi\par}{}{}
\makeatother
% Allow footnotes in longtable head/foot
\IfFileExists{footnotehyper.sty}{\usepackage{footnotehyper}}{\usepackage{footnote}}
\makesavenoteenv{longtable}
\usepackage{graphicx}
\makeatletter
\def\maxwidth{\ifdim\Gin@nat@width>\linewidth\linewidth\else\Gin@nat@width\fi}
\def\maxheight{\ifdim\Gin@nat@height>\textheight\textheight\else\Gin@nat@height\fi}
\makeatother
% Scale images if necessary, so that they will not overflow the page
% margins by default, and it is still possible to overwrite the defaults
% using explicit options in \includegraphics[width, height, ...]{}
\setkeys{Gin}{width=\maxwidth,height=\maxheight,keepaspectratio}
% Set default figure placement to htbp
\makeatletter
\def\fps@figure{htbp}
\makeatother

\KOMAoption{captions}{tableheading}
\makeatletter
\@ifpackageloaded{bookmark}{}{\usepackage{bookmark}}
\makeatother
\makeatletter
\@ifpackageloaded{caption}{}{\usepackage{caption}}
\AtBeginDocument{%
\ifdefined\contentsname
  \renewcommand*\contentsname{Table of contents}
\else
  \newcommand\contentsname{Table of contents}
\fi
\ifdefined\listfigurename
  \renewcommand*\listfigurename{List of Figures}
\else
  \newcommand\listfigurename{List of Figures}
\fi
\ifdefined\listtablename
  \renewcommand*\listtablename{List of Tables}
\else
  \newcommand\listtablename{List of Tables}
\fi
\ifdefined\figurename
  \renewcommand*\figurename{Figure}
\else
  \newcommand\figurename{Figure}
\fi
\ifdefined\tablename
  \renewcommand*\tablename{Table}
\else
  \newcommand\tablename{Table}
\fi
}
\@ifpackageloaded{float}{}{\usepackage{float}}
\floatstyle{ruled}
\@ifundefined{c@chapter}{\newfloat{codelisting}{h}{lop}}{\newfloat{codelisting}{h}{lop}[chapter]}
\floatname{codelisting}{Listing}
\newcommand*\listoflistings{\listof{codelisting}{List of Listings}}
\makeatother
\makeatletter
\makeatother
\makeatletter
\@ifpackageloaded{caption}{}{\usepackage{caption}}
\@ifpackageloaded{subcaption}{}{\usepackage{subcaption}}
\makeatother
\ifLuaTeX
  \usepackage{selnolig}  % disable illegal ligatures
\fi
\IfFileExists{bookmark.sty}{\usepackage{bookmark}}{\usepackage{hyperref}}
\IfFileExists{xurl.sty}{\usepackage{xurl}}{} % add URL line breaks if available
\urlstyle{same} % disable monospaced font for URLs
\hypersetup{
  pdftitle={Data Analytics with R},
  pdfauthor={Filip Edström, Joakim Wallmark},
  colorlinks=true,
  linkcolor={blue},
  filecolor={Maroon},
  citecolor={Blue},
  urlcolor={Blue},
  pdfcreator={LaTeX via pandoc}}

\title{Data Analytics with R}
\usepackage{etoolbox}
\makeatletter
\providecommand{\subtitle}[1]{% add subtitle to \maketitle
  \apptocmd{\@title}{\par {\large #1 \par}}{}{}
}
\makeatother
\subtitle{Lecture notes}
\author{Filip Edström, Joakim Wallmark}
\date{2025-11-10}

\begin{document}
\maketitle

\renewcommand*\contentsname{Table of contents}
{
\hypersetup{linkcolor=}
\setcounter{tocdepth}{2}
\tableofcontents
}
\bookmarksetup{startatroot}

\chapter*{Preface}\label{preface}
\addcontentsline{toc}{chapter}{Preface}

\markboth{Preface}{Preface}

These lecture notes are supposed to cover all the relevant content for
the course, along with the course literature, and a little bit more. The
notes are divided into three parts: Preliminaries, where R is introduced
along some basic statistics; Exploratory Data Analysis (EDA), which
gives reading instructions for the course literature
\href{https://r4ds.hadley.nz/}{R for Data Science} (R4DS in short); and
Statistical Inference and Prediction, which introduces how to gain
insight from data to both understand and predict. In addition to these
chapters there are solutions to the recommended exercises at the end of
the book. The main focus of the course lies with EDA and inference and
prediction, but the preliminaries are good (necessary) to be familiar
with to be able to appreciate the rest of the course.

The notes take a pragmatic approach focusing on giving an intuition for
concepts in data analysis rather than a rigorous theoretical account.
This aligns with the scope of this course, to introduce you to Data
Analytics in a practical way to give you the tools to learn more.

{Sometimes we will allude to more advanced concepts by writing a
paragraph in gray. The idea is to leave a small mental note for you that
there is something more to a concept, but that it is beyond the scope of
the course.}

We hope you enjoy the course!

\section*{Version 0.05}\label{version-0.05}
\addcontentsline{toc}{section}{Version 0.05}

\markright{Version 0.05}

Version notes:

\begin{itemize}
\tightlist
\item
  Added chapter on missing values and imputation.
\end{itemize}

These lecture notes are a work in progress. If you have any questions,
spot any errors, or have any suggestions fro improvement in the lecture
notes, please let us know by sending an email to one of the teachers on
the course with the tag {[}DAwR-LN{]} in the subject.

\section*{Disclaimer}\label{disclaimer}
\addcontentsline{toc}{section}{Disclaimer}

\markright{Disclaimer}

These lecture notes are used in the Data Analytics with R course at Umeå
University and is not intended to be spread outside of the context of
the course.

\part{Preliminaries}

In these two chapters we will introduce some preliminary knowledge that
is foundational to the course. You don't have to know everything in
these two chapters in detail, but it is good if you are familiar with
them. Both chapters starts with a ``TL;DR'', or ``Too Long; Didn't
Read'', which summarizes the chapters. This gives you a glance at the
chapter, so you get an idea if you are familiar with the concepts or if
you need to spend more time on it.

\chapter{Introduction to R}\label{introduction-to-r}

In this course, we'll be using R, a powerful programming language
designed specifically for statistical computations and data
visualization. While R is not suited for building websites or games, it
excels at handling data and performing statistical operations. In this
course we will mostly focus on applying R for the end goal of data
analysis and this chapter doesn't cover all the basics of R and
programming, but it provides a foundation for the rest of the course.

You may want to read the chapter Workflow: basics in
\href{https://r4ds.hadley.nz/}{R4DS} as it also covers some basic R.

Before getting started, you'll need to set up R and RStudio, a popular
integrated development environment (IDE) for R. You can download R from
\url{https://www.r-project.org/about.html} and RStudio from
\url{https://posit.co/products/open-source/rstudio/}. Once you have
these installed, you're ready to go!

\section{TL;DR}\label{tldr}

In this chapter we cover how to perform basic calculations; load
libraries (packages) into R; what objects (variables), functions,
vectors, and data frames are; and how logical operators and conditionals
work.

\section{Basic calculations}\label{basic-calculations}

Let's start by treating R as a giant calculator. You can perform basic
arithmetic operations like addition, subtraction, multiplication, and
division:

\begin{Shaded}
\begin{Highlighting}[]
\CommentTok{\# Addition}
\DecValTok{4} \SpecialCharTok{+} \DecValTok{2}
\end{Highlighting}
\end{Shaded}

\begin{verbatim}
[1] 6
\end{verbatim}

\begin{Shaded}
\begin{Highlighting}[]
\CommentTok{\# Subtraction}
\DecValTok{4} \SpecialCharTok{{-}} \DecValTok{2}
\end{Highlighting}
\end{Shaded}

\begin{verbatim}
[1] 2
\end{verbatim}

\begin{Shaded}
\begin{Highlighting}[]
\CommentTok{\# Multiplication}
\DecValTok{4} \SpecialCharTok{*} \DecValTok{2}
\end{Highlighting}
\end{Shaded}

\begin{verbatim}
[1] 8
\end{verbatim}

\begin{Shaded}
\begin{Highlighting}[]
\CommentTok{\# Division}
\DecValTok{4} \SpecialCharTok{/} \DecValTok{2}
\end{Highlighting}
\end{Shaded}

\begin{verbatim}
[1] 2
\end{verbatim}

Here we used \texttt{\#} to comment our code. Any text following a
\texttt{\#} on a line will not be executed by the program. Instead
comments are used to explain the code for others (and for future you!),
which comes in handy as our code gets more complicated.

\section{Libraries}\label{libraries}

Instead of writing our own code for everything, we use packages (or
libraries). Packages contain pre-written code and functions. One
essential package we'll use is \href{https://www.tidyverse.org/}{the
tidyverse package}, a collection of packages designed for data science.
It simplifies data loading, transformation, and visualization, making
our data analysis tasks much more manageable.

To install and load tidyverse use

\begin{Shaded}
\begin{Highlighting}[]
\FunctionTok{install.packages}\NormalTok{(}\StringTok{"tidyverse"}\NormalTok{)}
\FunctionTok{library}\NormalTok{(tidyverse)}
\end{Highlighting}
\end{Shaded}

If you need to install and/or load another package, just replace
\texttt{tidyverse} above with the package name.

You will be learning much more about tidyverse and how to use it when
you start working in R4DS, written by the creator of tidyverse, Hadley
Wickham, among others. There you will also learn all you need about
visualizing your data.

\section{Exercises}\label{exercises}

\begin{enumerate}
\def\labelenumi{\arabic{enumi}.}
\tightlist
\item
  Install the \texttt{tidyverse} package using the code above. You will
  need it for later exercises in this chapter.
\end{enumerate}

\section{Objects}\label{objects}

While R can be used for simple calculations, we are often more
interested in performing more complicated tasks. To achieve this we use
\emph{objects}, also called variables, to store our \emph{values}, the
results of our computations. Think of objects as labels for your data.
You can store values in objects using the assignment operator,
\texttt{\textless{}-}, like this

\begin{Shaded}
\begin{Highlighting}[]
\NormalTok{x }\OtherTok{\textless{}{-}} \DecValTok{4} \SpecialCharTok{*} \DecValTok{2}
\CommentTok{\# Printing the value of x}
\NormalTok{x}
\end{Highlighting}
\end{Shaded}

\begin{verbatim}
[1] 8
\end{verbatim}

\subsection{Exercises}\label{exercises-1}

\begin{enumerate}
\def\labelenumi{\arabic{enumi}.}
\tightlist
\item
  Copy the code from above and try using \texttt{=} instead of
  \texttt{\textless{}-}, what happens?
\end{enumerate}

\section{Functions}\label{functions}

Functions are like machines that take inputs and produce outputs. You
can think of a function like a vending machine: you input coins and a
number for a product and your desired product comes out. The function
\emph{transforms} your input into output. Inputs are also called
arguments or parameters, i.e.~the coins and number are the arguments of
the vending machine. Below we create a simple function called
\texttt{multiply4} that takes \(x\) as argument, multiplies it with 4,
and then returns the result.

\begin{Shaded}
\begin{Highlighting}[]
\CommentTok{\# Define function}
\NormalTok{multiply4 }\OtherTok{\textless{}{-}} \ControlFlowTok{function}\NormalTok{(x) \{}
\NormalTok{  x }\SpecialCharTok{*} \DecValTok{4}
\NormalTok{\}}
\CommentTok{\# Use the function two times with different inputs}
\FunctionTok{multiply4}\NormalTok{(}\DecValTok{1}\NormalTok{)}
\end{Highlighting}
\end{Shaded}

\begin{verbatim}
[1] 4
\end{verbatim}

\begin{Shaded}
\begin{Highlighting}[]
\FunctionTok{multiply4}\NormalTok{(}\DecValTok{10}\NormalTok{)}
\end{Highlighting}
\end{Shaded}

\begin{verbatim}
[1] 40
\end{verbatim}

You don't need to create your own functions in this course, but you'll
use plenty of them, and of course the above function is pretty useless.
We can just run \texttt{1*4} or \texttt{10*4} directly in R instead of
over-complicating things with a function, but it serves a good example
of how functions work.

We can also save the result of our function in our own variable. Below
we take the result of our function, save it in the
\texttt{two\_times\_four} variable, and print the result plus 2

\begin{Shaded}
\begin{Highlighting}[]
\NormalTok{two\_times\_four }\OtherTok{\textless{}{-}} \FunctionTok{multiply4}\NormalTok{(}\DecValTok{2}\NormalTok{)}
\FunctionTok{print}\NormalTok{(two\_times\_four }\SpecialCharTok{+} \DecValTok{2}\NormalTok{)}
\end{Highlighting}
\end{Shaded}

\begin{verbatim}
[1] 10
\end{verbatim}

\subsection{Exercises}\label{exercises-2}

\begin{enumerate}
\def\labelenumi{\arabic{enumi}.}
\tightlist
\item
  The code below is an extension of the \texttt{multiply4} function
  above, taking two arguments and multiplies them together. Copy the
  code and use the function to compute \(36 \cdot 52\)
\end{enumerate}

\begin{Shaded}
\begin{Highlighting}[]
\NormalTok{multiply }\OtherTok{\textless{}{-}} \ControlFlowTok{function}\NormalTok{(x, y) \{}
\NormalTok{  x }\SpecialCharTok{*}\NormalTok{ y}
\NormalTok{\}}
\end{Highlighting}
\end{Shaded}

\begin{enumerate}
\def\labelenumi{\arabic{enumi}.}
\setcounter{enumi}{1}
\tightlist
\item
  Write your own function called \texttt{subtract} that subtracts two
  numbers instead of multiplying them, and then use the function to
  compute \(36 - 52\).
\item
  In R it is possible to assign the parameters values by explicitly
  stating the parameter in the function call. Try using the
  \texttt{subtract} function by calling
  \texttt{subtract(y\ =\ 52,\ x\ =\ 36)}. Is the result different?
\item
  Redo the previous exercise, now calling
  \texttt{subtract(y\ \textless{}-\ 52,\ x\ \textless{}-\ 36)}. What
  happens?
\end{enumerate}

\section{Vectors}\label{vectors}

Sometimes we want to work with multiple items of the same type together.
To achieve this in R, we can use vectors. A vector is declared using the
\texttt{c} function (for combine). The arguments to \texttt{c} are the
items we want in the vector as separate arguments. For example, if we
want to store a sequence of different numbers, we can use:

\begin{Shaded}
\begin{Highlighting}[]
\NormalTok{vec }\OtherTok{\textless{}{-}} \FunctionTok{c}\NormalTok{(}\DecValTok{4}\NormalTok{, }\DecValTok{2}\NormalTok{, }\DecValTok{453}\NormalTok{)}
\FunctionTok{print}\NormalTok{(vec)}
\end{Highlighting}
\end{Shaded}

\begin{verbatim}
[1]   4   2 453
\end{verbatim}

We can also store objects in vectors. Here we store three text variables
(usually called \emph{strings} in a programming context) in the same
vector:

\begin{Shaded}
\begin{Highlighting}[]
\NormalTok{x }\OtherTok{\textless{}{-}} \StringTok{"string 1"}
\NormalTok{y }\OtherTok{\textless{}{-}} \StringTok{"string 2"}
\NormalTok{z }\OtherTok{\textless{}{-}} \StringTok{"string 3"}
\NormalTok{vec }\OtherTok{\textless{}{-}} \FunctionTok{c}\NormalTok{(x, y, z)}
\FunctionTok{print}\NormalTok{(vec)}
\end{Highlighting}
\end{Shaded}

\begin{verbatim}
[1] "string 1" "string 2" "string 3"
\end{verbatim}

If we want to access a specific element in a vector, we can use square
brackets \texttt{{[}{]}}. The first element in a vector is accessed by
\texttt{vec{[}1{]}}, the second by \texttt{vec{[}2{]}}, and so on. Below
we access the first element in the vector \texttt{vec}:

\begin{Shaded}
\begin{Highlighting}[]
\FunctionTok{print}\NormalTok{(vec[}\DecValTok{1}\NormalTok{])}
\end{Highlighting}
\end{Shaded}

\begin{verbatim}
[1] "string 1"
\end{verbatim}

\begin{Shaded}
\begin{Highlighting}[]
\FunctionTok{print}\NormalTok{(vec[}\DecValTok{2}\NormalTok{])}
\end{Highlighting}
\end{Shaded}

\begin{verbatim}
[1] "string 2"
\end{verbatim}

\begin{Shaded}
\begin{Highlighting}[]
\FunctionTok{print}\NormalTok{(vec[}\DecValTok{3}\NormalTok{])}
\end{Highlighting}
\end{Shaded}

\begin{verbatim}
[1] "string 3"
\end{verbatim}

{In other programming languages, e.g.~Python, vector indices commonly
start from 0. Watch out for this when moving between languages as it is
a common mistake! Starting indices from 0 is more true to how a computer
works, while starting from 1 is common in mathematics.}

In this course, we'll often use vectors implicitly, especially when
dealing with data frames.

\subsection{Exercises}\label{exercises-3}

\begin{enumerate}
\def\labelenumi{\arabic{enumi}.}
\tightlist
\item
  Try the basic calculations on the two vectors
  \texttt{vec1\ \textless{}-\ c(1,\ 2)} and
  \texttt{vec2\ \textless{}-\ c(3,\ 4)}. What happens?
\item
  Run the following code and try to understand what happens
\end{enumerate}

\begin{Shaded}
\begin{Highlighting}[]
\NormalTok{vec1 }\OtherTok{\textless{}{-}} \FunctionTok{c}\NormalTok{(}\DecValTok{1}\NormalTok{, }\DecValTok{2}\NormalTok{)}
\NormalTok{vec2 }\OtherTok{\textless{}{-}} \FunctionTok{c}\NormalTok{(}\DecValTok{3}\NormalTok{, }\DecValTok{4}\NormalTok{)}
\NormalTok{vec3 }\OtherTok{\textless{}{-}} \FunctionTok{c}\NormalTok{(}\DecValTok{3}\NormalTok{, }\StringTok{"four"}\NormalTok{)}\StringTok{\textasciigrave{}}

\AttributeTok{vec1 + vec2}
\AttributeTok{vec2 + vec3}
\end{Highlighting}
\end{Shaded}

\section{Matrices}\label{matrices}

A vector stores objects in one dimension. Sometimes we encounter
problems when we need to store data in two dimensions, for example when
we are dealing with correlations between different variables. In this
case, we can use matrices. A matrix is a two-dimensional array that
stores data in rows and columns. We can create a matrix using the
\texttt{matrix} function. The first argument to the function is the data
we want to store in the matrix, the second argument is the number of
rows, and the third argument is the number of columns. We can also
specify if we want to fill the matrix by row or by column. Below we
create a 3x3 matrix filled by row:

\begin{Shaded}
\begin{Highlighting}[]
\CommentTok{\# initialize matrix with ones on the diagonal}
\NormalTok{m }\OtherTok{\textless{}{-}} \FunctionTok{matrix}\NormalTok{(}
  \FunctionTok{c}\NormalTok{(}\DecValTok{1}\NormalTok{, }\DecValTok{2}\NormalTok{, }\DecValTok{3}\NormalTok{,}
    \DecValTok{2}\NormalTok{, }\DecValTok{1}\NormalTok{, }\DecValTok{2}\NormalTok{,}
    \DecValTok{4}\NormalTok{, }\DecValTok{2}\NormalTok{, }\DecValTok{1}\NormalTok{),}
  \AttributeTok{nrow =} \DecValTok{3}\NormalTok{,}
  \AttributeTok{byrow =} \ConstantTok{TRUE}\NormalTok{)}
\end{Highlighting}
\end{Shaded}

We can access a specific element in a matrix using square brackets
\texttt{{[}{]}}. The first argument is the row index, and the second
argument is the column index. We can also access a whole row or column
by leaving that index blank. Below are some examples.

\begin{Shaded}
\begin{Highlighting}[]
\CommentTok{\# Select first element of matrix}
\NormalTok{m[}\DecValTok{1}\NormalTok{, }\DecValTok{1}\NormalTok{]}
\end{Highlighting}
\end{Shaded}

\begin{verbatim}
[1] 1
\end{verbatim}

\begin{Shaded}
\begin{Highlighting}[]
\CommentTok{\# Select first row of matrix}
\NormalTok{m[}\DecValTok{1}\NormalTok{, ]}
\end{Highlighting}
\end{Shaded}

\begin{verbatim}
[1] 1 2 3
\end{verbatim}

\begin{Shaded}
\begin{Highlighting}[]
\CommentTok{\# Select first column of matrix}
\NormalTok{m[, }\DecValTok{1}\NormalTok{]}
\end{Highlighting}
\end{Shaded}

\begin{verbatim}
[1] 1 2 4
\end{verbatim}

Later in this chapter we will see that we may also use conditional
statements to select elements in a matrix, see Section
Section~\ref{sec-conditional-indexing}

\section{Data frames}\label{data-frames}

Tables are essential in data analysis. A table stores data in a
row/column format:

\begin{longtable}[]{@{}lll@{}}
\toprule\noalign{}
Name & Age & Favorite Color \\
\midrule\noalign{}
\endhead
\bottomrule\noalign{}
\endlastfoot
Alice & 25 & Blue \\
Bob & 30 & Red \\
Charlie & 22 & Green \\
\end{longtable}

where Name, Age, and Favorite Color are the columns and Alice, 25, Blue
is the first row. We can create this table in R using the
\texttt{data.frame} function:

\begin{Shaded}
\begin{Highlighting}[]
\NormalTok{df }\OtherTok{\textless{}{-}} \FunctionTok{data.frame}\NormalTok{(}
  \AttributeTok{Name =} \FunctionTok{c}\NormalTok{(}\StringTok{"Alice"}\NormalTok{, }\StringTok{"Bob"}\NormalTok{, }\StringTok{"Charlie"}\NormalTok{),}
  \AttributeTok{Age =} \FunctionTok{c}\NormalTok{(}\DecValTok{25}\NormalTok{, }\DecValTok{30}\NormalTok{, }\DecValTok{22}\NormalTok{),}
  \AttributeTok{Favorite\_Color =} \FunctionTok{c}\NormalTok{(}\StringTok{"Blue"}\NormalTok{, }\StringTok{"Red"}\NormalTok{, }\StringTok{"Green"}\NormalTok{)}
\NormalTok{)}
\FunctionTok{print}\NormalTok{(df)}
\end{Highlighting}
\end{Shaded}

\begin{verbatim}
     Name Age Favorite_Color
1   Alice  25           Blue
2     Bob  30            Red
3 Charlie  22          Green
\end{verbatim}

Or alternatively, by using the \texttt{tibble} function. Tibbles are
similar to data frames, but with some extended functionality for
printing and other things which we will not cover in detail. To create a
tibble, we need first need to load the \texttt{tidyverse} package:

\begin{Shaded}
\begin{Highlighting}[]
\FunctionTok{library}\NormalTok{(tidyverse)}
\NormalTok{tib }\OtherTok{\textless{}{-}} \FunctionTok{tibble}\NormalTok{(}
  \AttributeTok{Name =} \FunctionTok{c}\NormalTok{(}\StringTok{"Alice"}\NormalTok{, }\StringTok{"Bob"}\NormalTok{, }\StringTok{"Charlie"}\NormalTok{),}
  \AttributeTok{Age =} \FunctionTok{c}\NormalTok{(}\DecValTok{25}\NormalTok{, }\DecValTok{30}\NormalTok{, }\DecValTok{22}\NormalTok{),}
  \AttributeTok{Favorite\_Color =} \FunctionTok{c}\NormalTok{(}\StringTok{"Blue"}\NormalTok{, }\StringTok{"Red"}\NormalTok{, }\StringTok{"Green"}\NormalTok{)}
\NormalTok{)}
\end{Highlighting}
\end{Shaded}

If you want to access one of the columns in a data frame or tibble you
can use the dollar operator, \texttt{\$}, after the data frame or
tibbles name. Below we access the \texttt{Name} column and print out the
values in it.

\begin{Shaded}
\begin{Highlighting}[]
\NormalTok{tib}\SpecialCharTok{$}\NormalTok{Name}
\end{Highlighting}
\end{Shaded}

\begin{verbatim}
[1] "Alice"   "Bob"     "Charlie"
\end{verbatim}

There are a number of useful functions that you can use with data frames
and tibbles, in the table below you can see some of them and a
description of what they do.

\begin{longtable}[]{@{}
  >{\raggedright\arraybackslash}p{(\columnwidth - 2\tabcolsep) * \real{0.2737}}
  >{\raggedright\arraybackslash}p{(\columnwidth - 2\tabcolsep) * \real{0.7263}}@{}}
\toprule\noalign{}
\begin{minipage}[b]{\linewidth}\raggedright
Feature
\end{minipage} & \begin{minipage}[b]{\linewidth}\raggedright
Description
\end{minipage} \\
\midrule\noalign{}
\endhead
\bottomrule\noalign{}
\endlastfoot
glimpse() & Display a concise overview of the data frame's structure. \\
summary() & Provide summary statistics for numeric objects in the data
frame. \\
head() & Display the top rows of the data frame. \\
tail() & Display the bottom rows of the data frame. \\
nrow() & Show the length of the data frame. \\
ncol() & Show the number of columns in the data frame. \\
colnames() or names() & Display column names (as character objects). \\
\end{longtable}

Here is an example of how to use them all and their output for the
tibble \texttt{tib} that we defined before.

\begin{Shaded}
\begin{Highlighting}[]
\FunctionTok{glimpse}\NormalTok{(tib)}
\end{Highlighting}
\end{Shaded}

\begin{verbatim}
Rows: 3
Columns: 3
$ Name           <chr> "Alice", "Bob", "Charlie"
$ Age            <dbl> 25, 30, 22
$ Favorite_Color <chr> "Blue", "Red", "Green"
\end{verbatim}

\begin{Shaded}
\begin{Highlighting}[]
\FunctionTok{summary}\NormalTok{(tib)}
\end{Highlighting}
\end{Shaded}

\begin{verbatim}
     Name                Age        Favorite_Color    
 Length:3           Min.   :22.00   Length:3          
 Class :character   1st Qu.:23.50   Class :character  
 Mode  :character   Median :25.00   Mode  :character  
                    Mean   :25.67                     
                    3rd Qu.:27.50                     
                    Max.   :30.00                     
\end{verbatim}

\begin{Shaded}
\begin{Highlighting}[]
\FunctionTok{head}\NormalTok{(tib)}
\end{Highlighting}
\end{Shaded}

\begin{verbatim}
# A tibble: 3 x 3
  Name      Age Favorite_Color
  <chr>   <dbl> <chr>         
1 Alice      25 Blue          
2 Bob        30 Red           
3 Charlie    22 Green         
\end{verbatim}

\begin{Shaded}
\begin{Highlighting}[]
\FunctionTok{tail}\NormalTok{(tib)}
\end{Highlighting}
\end{Shaded}

\begin{verbatim}
# A tibble: 3 x 3
  Name      Age Favorite_Color
  <chr>   <dbl> <chr>         
1 Alice      25 Blue          
2 Bob        30 Red           
3 Charlie    22 Green         
\end{verbatim}

\begin{Shaded}
\begin{Highlighting}[]
\FunctionTok{nrow}\NormalTok{(tib)}
\end{Highlighting}
\end{Shaded}

\begin{verbatim}
[1] 3
\end{verbatim}

\begin{Shaded}
\begin{Highlighting}[]
\FunctionTok{ncol}\NormalTok{(tib)}
\end{Highlighting}
\end{Shaded}

\begin{verbatim}
[1] 3
\end{verbatim}

\begin{Shaded}
\begin{Highlighting}[]
\FunctionTok{names}\NormalTok{(tib)}
\end{Highlighting}
\end{Shaded}

\begin{verbatim}
[1] "Name"           "Age"            "Favorite_Color"
\end{verbatim}

As data frames and tibbles are very similar, we will be using the words
data frame and tibble interchangeably.

\subsection{Exercises}\label{exercises-4}

\begin{enumerate}
\def\labelenumi{\arabic{enumi}.}
\tightlist
\item
  Install and load the \texttt{palmerpenguins} package using the code
  below and look at the output from calling the data \texttt{penguins}.
  Try running \texttt{?penguins} in the console to learn more about the
  dataset. We will be using the \texttt{penguins} data throughout the
  course, so have a look at it, but there will be more time to study it
  in detail later.
\end{enumerate}

\begin{Shaded}
\begin{Highlighting}[]
\FunctionTok{install.packages}\NormalTok{(}\StringTok{"palmerpenguins"}\NormalTok{)}
\FunctionTok{library}\NormalTok{(palmerpenguins)}
\NormalTok{penguins}
\end{Highlighting}
\end{Shaded}

\begin{enumerate}
\def\labelenumi{\arabic{enumi}.}
\setcounter{enumi}{1}
\tightlist
\item
  Access the \texttt{flipper\_length\_mm} column of the penguins dataset
  using the dollar operator, \texttt{\$}.
\item
  Type the tibbles name (\texttt{penguins}) and \texttt{\$} without any
  space and then press tab. What do you see?
\item
  Load the \texttt{tidyverse} package and apply all the above functions
  on \texttt{penguins}, you can copy the code and replace \texttt{tib}
  with \texttt{penguins}. Again, don't spend too much time, but try to
  understand the functions and what they do.
\end{enumerate}

\section{Logical operators and
conditionals}\label{logical-operators-and-conditionals}

Logicals and conditionals are used to make decisions in our code. As
their names suggest, logicals lets us define logical expressions that
can either be \texttt{TRUE} or \texttt{FALSE}, while conditionals lets
us define what should happen in the code if a logical is either
\texttt{TRUE} or \texttt{FALSE}. For example, if it is raining when you
leave your house you may want to take an umbrella, but if it's not
raining you'd rather leave the umbrella at home. We can express this as
a logical and a condition: IF(\(raining\)) THEN take umbrella ELSE leave
umbrella. Here IF represents a conditional, and \(raining\) represents a
logical.

\subsection{Logical operators in R}\label{logical-operators-in-r}

In R, you can use logical operators to compare values. When a logical
condition is true, R returns the value\texttt{TRUE}, and when it is
false \texttt{FALSE}. There are several different operators to compare
values (or objects) in R:

\begin{itemize}
\tightlist
\item
  \textbf{Equality (\texttt{==}) and Inequality (\texttt{!=}):} Check if
  two objects are equal or not equal using \texttt{==} and \texttt{!=},
  respectively. For example:
\end{itemize}

\begin{Shaded}
\begin{Highlighting}[]
\DecValTok{4} \SpecialCharTok{==} \DecValTok{4}
\end{Highlighting}
\end{Shaded}

\begin{verbatim}
[1] TRUE
\end{verbatim}

\begin{Shaded}
\begin{Highlighting}[]
\DecValTok{4} \SpecialCharTok{==} \DecValTok{2}
\end{Highlighting}
\end{Shaded}

\begin{verbatim}
[1] FALSE
\end{verbatim}

\begin{Shaded}
\begin{Highlighting}[]
\DecValTok{4} \SpecialCharTok{!=} \DecValTok{2}
\end{Highlighting}
\end{Shaded}

\begin{verbatim}
[1] TRUE
\end{verbatim}

\begin{itemize}
\tightlist
\item
  \textbf{Greater Than (\texttt{\textgreater{}}), Less Than
  (\texttt{\textless{}}):} Compare values to check if one is greater or
  less than another:
\end{itemize}

\begin{Shaded}
\begin{Highlighting}[]
\DecValTok{4} \SpecialCharTok{\textgreater{}} \DecValTok{2}
\end{Highlighting}
\end{Shaded}

\begin{verbatim}
[1] TRUE
\end{verbatim}

\begin{Shaded}
\begin{Highlighting}[]
\DecValTok{4} \SpecialCharTok{\textless{}} \DecValTok{2}
\end{Highlighting}
\end{Shaded}

\begin{verbatim}
[1] FALSE
\end{verbatim}

\begin{itemize}
\tightlist
\item
  \textbf{Greater Or Equal (\texttt{\textgreater{}=}), Less Or Equal
  (\texttt{\textless{}=}):} Compare values to check if one is greater or
  equal or less or equal than another:
\end{itemize}

\begin{Shaded}
\begin{Highlighting}[]
\DecValTok{4} \SpecialCharTok{\textgreater{}} \DecValTok{4}
\end{Highlighting}
\end{Shaded}

\begin{verbatim}
[1] FALSE
\end{verbatim}

\begin{Shaded}
\begin{Highlighting}[]
\DecValTok{4} \SpecialCharTok{\textgreater{}=} \DecValTok{4}
\end{Highlighting}
\end{Shaded}

\begin{verbatim}
[1] TRUE
\end{verbatim}

\begin{itemize}
\tightlist
\item
  \textbf{Logical AND (\texttt{\&}) and Logical OR
  (\texttt{\textbar{}}):} Combine multiple conditions using \texttt{\&}
  for AND and \texttt{\textbar{}} for OR:
\end{itemize}

\begin{Shaded}
\begin{Highlighting}[]
\NormalTok{is\_raining }\OtherTok{\textless{}{-}} \ConstantTok{TRUE}
\NormalTok{has\_umbrella }\OtherTok{\textless{}{-}} \ConstantTok{TRUE}
\CommentTok{\# Testing if is raining and has umbrella}
\NormalTok{is\_raining }\SpecialCharTok{\&}\NormalTok{ has\_umbrella }
\end{Highlighting}
\end{Shaded}

\begin{verbatim}
[1] TRUE
\end{verbatim}

\begin{Shaded}
\begin{Highlighting}[]
\NormalTok{has\_raincoat }\OtherTok{\textless{}{-}} \ConstantTok{FALSE}
\CommentTok{\# Testing if has umbrella or raincoat}
\NormalTok{has\_umbrella }\SpecialCharTok{|}\NormalTok{ has\_raincoat}
\end{Highlighting}
\end{Shaded}

\begin{verbatim}
[1] TRUE
\end{verbatim}

\subsection{Conditional indexing}\label{sec-conditional-indexing}

We can use logical operators to select elements in a vector or matrix.
When we test a logical condition on a vector, R returns a logical vector
of the same length as the original vector. This logical vector can be
used to index the original vector. For example, if we have a vector of
numbers and we want to select only the numbers that are greater than 5,
we can use the following code:

\begin{Shaded}
\begin{Highlighting}[]
\NormalTok{vec }\OtherTok{\textless{}{-}} \FunctionTok{c}\NormalTok{(}\DecValTok{1}\NormalTok{, }\DecValTok{6}\NormalTok{, }\DecValTok{3}\NormalTok{, }\DecValTok{8}\NormalTok{, }\DecValTok{2}\NormalTok{, }\DecValTok{9}\NormalTok{)}
\NormalTok{vec }\SpecialCharTok{\textgreater{}} \DecValTok{5}
\end{Highlighting}
\end{Shaded}

\begin{verbatim}
[1] FALSE  TRUE FALSE  TRUE FALSE  TRUE
\end{verbatim}

\begin{Shaded}
\begin{Highlighting}[]
\NormalTok{vec[vec }\SpecialCharTok{\textgreater{}} \DecValTok{5}\NormalTok{]}
\end{Highlighting}
\end{Shaded}

\begin{verbatim}
[1] 6 8 9
\end{verbatim}

This can also be applied to matrices,

\begin{Shaded}
\begin{Highlighting}[]
\CommentTok{\# initialize matrix with ones on the diagonal}
\NormalTok{m }\OtherTok{\textless{}{-}} \FunctionTok{matrix}\NormalTok{(}
  \FunctionTok{c}\NormalTok{(}\DecValTok{1}\NormalTok{, }\DecValTok{2}\NormalTok{, }\DecValTok{3}\NormalTok{,}
    \DecValTok{2}\NormalTok{, }\DecValTok{1}\NormalTok{, }\DecValTok{2}\NormalTok{,}
    \DecValTok{4}\NormalTok{, }\DecValTok{2}\NormalTok{, }\DecValTok{1}\NormalTok{),}
  \AttributeTok{nrow =} \DecValTok{3}\NormalTok{,}
  \AttributeTok{byrow =} \ConstantTok{TRUE}\NormalTok{)}
\CommentTok{\# Test the elementwise condition m is equal to one}
\NormalTok{m }\SpecialCharTok{==} \DecValTok{1}
\end{Highlighting}
\end{Shaded}

\begin{verbatim}
      [,1]  [,2]  [,3]
[1,]  TRUE FALSE FALSE
[2,] FALSE  TRUE FALSE
[3,] FALSE FALSE  TRUE
\end{verbatim}

\begin{Shaded}
\begin{Highlighting}[]
\CommentTok{\# Selecting elements in m based on a condition}
\NormalTok{m[m }\SpecialCharTok{\textgreater{}} \DecValTok{1}\NormalTok{]}
\end{Highlighting}
\end{Shaded}

\begin{verbatim}
[1] 2 4 2 2 3 2
\end{verbatim}

\begin{Shaded}
\begin{Highlighting}[]
\CommentTok{\# Selecting elements in m based on multiple conditions}
\NormalTok{m[m }\SpecialCharTok{\textgreater{}} \DecValTok{1} \SpecialCharTok{\&}\NormalTok{ m }\SpecialCharTok{\textless{}} \DecValTok{3}\NormalTok{]}
\end{Highlighting}
\end{Shaded}

\begin{verbatim}
[1] 2 2 2 2
\end{verbatim}

Testing a condition returns another matrix of the same size, showing
where the condition is true. Selecting elements returns the elements
matching the statement.

\subsection{Conditional statements in
R}\label{conditional-statements-in-r}

Logical operators are useful to compare objects but the real power of
logical operators shows when combined with conditional statements.
Conditional statements allow you to execute different parts of the code
based on logical conditions. In R, you can use \texttt{ifelse} for
conditional assignment:

\begin{Shaded}
\begin{Highlighting}[]
\NormalTok{is\_raining }\OtherTok{\textless{}{-}} \ConstantTok{TRUE}
\CommentTok{\# IF(is\_raining) THEN bring umbrella ELSE leave umbrella}
\FunctionTok{ifelse}\NormalTok{(is\_raining, }\StringTok{"bring umbrella"}\NormalTok{, }\StringTok{"leave umbrella"}\NormalTok{)}
\end{Highlighting}
\end{Shaded}

\begin{verbatim}
[1] "bring umbrella"
\end{verbatim}

\begin{Shaded}
\begin{Highlighting}[]
\NormalTok{is\_raining }\OtherTok{\textless{}{-}} \ConstantTok{FALSE}
\FunctionTok{ifelse}\NormalTok{(is\_raining, }\StringTok{"bring umbrella"}\NormalTok{, }\StringTok{"leave umbrella"}\NormalTok{)}
\end{Highlighting}
\end{Shaded}

\begin{verbatim}
[1] "leave umbrella"
\end{verbatim}

\texttt{ifelse} is useful when transforming objects, as you will learn
more about later in the course.

In this course we will often talk about functions, inputs, parameters
and arguments, and outputs, so having an idea of these concepts is
useful, even if you don't need to write them yourself (you are of course
free to try!).

\chapter{Basic statistics}\label{sec-basicstat}

This chapter aims to cover basic statistical concepts which are crucial
in data analysis. It will cover variables and their distributions and
how descriptive statistics can tell us more about them. We will be using
R and visualization techniques throughout the chapter to illustrate. We
don't expect you to understand all the code in this chapter, especially
for the visualizations, but it can be useful later in the course as a
resource. Therefore some of the code is hidden by default, but you can
press the line that says {Code} to see the code!

The examples will be using the
\href{https://www.rdocumentation.org/packages/palmerpenguins/versions/0.1.1}{\texttt{palmerpenguins}
dataset} (short \texttt{penguins}) which you can install using

\begin{Shaded}
\begin{Highlighting}[]
\FunctionTok{install.packages}\NormalTok{(}\StringTok{"palmerpenguins"}\NormalTok{)}
\FunctionTok{library}\NormalTok{(palmerpenguins)}
\end{Highlighting}
\end{Shaded}

\section{TL;DR}\label{tldr-1}

This chapter covers numerical and categorical variables and their
(empirical) distributions, descriptive statistics (mean, median,
standard deviation, percentiles), and correlation between two numerical
variables.

\section{Variables}\label{variables}

In data analysis we encounter various types of variables. These
variables store information about different objects or quantities, like
in the \texttt{penguins} dataset where variables represent different
species of penguins and their flipper lengths. For each observed penguin
we have a \emph{value} of each variable. Variables can assume different
types which dictate the kind of analysis that can be performed based on
them. Understanding the different variable types and being able to
classify variables as a certain type is very important to perform
correct and sound analysis. We will look at two types of variables in
this introduction: numerical and categorical variables.

\subsection{Numerical}\label{numerical}

Variables that represents numbers. They can come in two forms:

\begin{itemize}
\tightlist
\item
  \textbf{Integer}: Represents discrete numerical values, e.g., number
  of items in a cart or a year. A rule of thumb is that if you can count
  it, then it is an integer.
\item
  \textbf{Continuous}: Represents continuous numerical values, e.g., the
  weight of an item or a distance.
\end{itemize}

To be able to handle variables correctly, each type is represented in R,
but sometimes under a different name. An integer is called
\texttt{integer} in R, and a continuous variable is called a
\texttt{numeric} or \texttt{double}.

\subsection{Categorical}\label{sec-basicstat-categorical}

Represents a finite set of categories, e.g., types of fruits. What sets
categorical variables apart from numerical variables are that there is
no \emph{natural ordering} between the different values: an apple is not
worth more than a pear! Sometimes we can have \emph{ordered} categorical
variables, these have some internal order, but instead there is no
natural \emph{distance} between the values. You can think of customer
satisfaction: a customer may be dissatisfied, neutral, or satisfied.
Would you say that the distance between dissatisfied and satisfied is
twice the distance between dissatisfied and neutral? Even if you could
come up with some way of thinking about a distance, there will not be a
natural distance that everyone can agree upon. Instead we can only order
the variable into ordered categories.

The most simple case of a categorical variable is a \emph{binary}
variable. It consists only of two categories, like the variable
\texttt{sex} in the \texttt{penguins} dataset that you will soon get
familiar with.

In R, a categorical variable is called \texttt{factor}, or if it is
ordered \texttt{ordered\ factor}.

\subsection{\texorpdfstring{Variables in
\texttt{penguins}}{Variables in penguins}}\label{variables-in-penguins}

Now we will take a look at the \texttt{penguins} dataset

\begin{longtable}[]{@{}
  >{\raggedright\arraybackslash}p{(\columnwidth - 6\tabcolsep) * \real{0.1652}}
  >{\raggedright\arraybackslash}p{(\columnwidth - 6\tabcolsep) * \real{0.6000}}
  >{\raggedright\arraybackslash}p{(\columnwidth - 6\tabcolsep) * \real{0.1304}}
  >{\raggedright\arraybackslash}p{(\columnwidth - 6\tabcolsep) * \real{0.1043}}@{}}
\toprule\noalign{}
\begin{minipage}[b]{\linewidth}\raggedright
Variable
\end{minipage} & \begin{minipage}[b]{\linewidth}\raggedright
Description
\end{minipage} & \begin{minipage}[b]{\linewidth}\raggedright
Variable Type
\end{minipage} & \begin{minipage}[b]{\linewidth}\raggedright
In R
\end{minipage} \\
\midrule\noalign{}
\endhead
\bottomrule\noalign{}
\endlastfoot
species & Penguin species (Adélie, Chinstrap, Gentoo) & Categorical &
\texttt{factor} \\
island & Island in Palmer Archipelago, Antarctica (Biscoe, Dream,
Torgersen) & Categorical & \texttt{factor} \\
bill\_length\_mm & Bill length in millimeters & Continuous &
\texttt{numeric} \\
bill\_depth\_mm & Bill depth in millimeters & Continuous &
\texttt{numeric} \\
flipper\_length\_mm & Flipper length in millimeters & Integer &
\texttt{integer} \\
body\_mass\_g & Body mass in grams & Integer & \texttt{integer} \\
sex & Penguin sex (female, male) & Categorical & \texttt{factor} \\
year & Study year (2007, 2008, 2009) & Integer & \texttt{integer} \\
\end{longtable}

We see that the dataset contains integer, continuous, and categorical
variables, but some variables doesn't seem to add up with what we would
expect. The variables \texttt{flipper\_length\_mm} and
\texttt{body\_mass\_g} are stored as integers but should be continuous,
as flipper lengths and body mass can assume any continuous value.
Sometimes this happens when variables are recorded using \emph{rounded}
numbers. Most of the time the difference between \texttt{integer} and
\texttt{numeric} doesn't affect our analysis since R is smart to figure
out the \emph{encoding} by itself, but for other variable types it is
important to \emph{code} them correctly. We may \emph{recode} the
variables using functions such as \texttt{as.integer} to code something
as an integer, \texttt{as.numeric} to continuous, or \texttt{as.factor}
to categorical. Recoding is especially important if a categorical
variable is stored as a \texttt{integer} or \texttt{numeric} instead of
a \texttt{factor}, or vice-versa, as it will change what analysis you
can do with the variable.

Now that we know the variables in our dataset we can look at them in R
by using the \texttt{glimpse} function.

\begin{Shaded}
\begin{Highlighting}[]
\FunctionTok{library}\NormalTok{(tidyverse)}
\CommentTok{\#install.packages("palmerpenguins")}
\FunctionTok{library}\NormalTok{(palmerpenguins)}
\FunctionTok{glimpse}\NormalTok{(penguins)}
\end{Highlighting}
\end{Shaded}

\begin{verbatim}
Rows: 344
Columns: 8
$ species           <fct> Adelie, Adelie, Adelie, Adelie, Adelie, Adelie, Adel~
$ island            <fct> Torgersen, Torgersen, Torgersen, Torgersen, Torgerse~
$ bill_length_mm    <dbl> 39.1, 39.5, 40.3, NA, 36.7, 39.3, 38.9, 39.2, 34.1, ~
$ bill_depth_mm     <dbl> 18.7, 17.4, 18.0, NA, 19.3, 20.6, 17.8, 19.6, 18.1, ~
$ flipper_length_mm <int> 181, 186, 195, NA, 193, 190, 181, 195, 193, 190, 186~
$ body_mass_g       <int> 3750, 3800, 3250, NA, 3450, 3650, 3625, 4675, 3475, ~
$ sex               <fct> male, female, female, NA, female, male, female, male~
$ year              <int> 2007, 2007, 2007, 2007, 2007, 2007, 2007, 2007, 2007~
\end{verbatim}

Now we can see the variables: in the first column we see the variable
name, in the second an abbreviation of the variable type, and then
observed values. The first penguin in our dataset has the values Adelie
species, observed on Torgersen, bill length of 39.1 mm, and so on. We
can also see that some values are marked \texttt{NA}. These are
\emph{missing values}, where a measurement is missing. This is common in
real data, and you will encounter it and have to deal with it many
times.

Now that we have seen the different variable types and an example, we
can start asking questions about how these variables behave. Which is
the most common species in our data set? What is the average flipper
length? These are all questions about the \emph{distributions} of the
variables.

\subsection{Exercises}\label{exercises-5}

\begin{enumerate}
\def\labelenumi{\arabic{enumi}.}
\tightlist
\item
  Use the \texttt{summary} function to learn more about the variables in
  the \texttt{penguins} dataset. How does the \texttt{summary} output
  differ for categorical and numerical variables? How many
  \texttt{NA\textquotesingle{}s} are there for each variable (Hint: Look
  at the last line below each variable). Don't worry if you don't
  understand every detail of the output yet, just note the differences.
\item
  Run the code below to recode the variables
  \texttt{flipper\_length\_mm} and \texttt{body\_mass\_g}. Run the
  functions \texttt{glimpse} and \texttt{summary} on the
  \texttt{penguins} dataset to check that the variables have indeed
  changed.
\end{enumerate}

\begin{Shaded}
\begin{Highlighting}[]
\NormalTok{penguins}\SpecialCharTok{$}\NormalTok{flipper\_length\_mm }\OtherTok{\textless{}{-}} \FunctionTok{as.numeric}\NormalTok{(penguins}\SpecialCharTok{$}\NormalTok{flipper\_length\_mm)}
\NormalTok{penguins}\SpecialCharTok{$}\NormalTok{body\_mass\_g }\OtherTok{\textless{}{-}} \FunctionTok{as.numeric}\NormalTok{(penguins}\SpecialCharTok{$}\NormalTok{body\_mass\_g)}
\end{Highlighting}
\end{Shaded}

\section{Distributions}\label{distributions}

Distributions describe how likely different values are for a variable.
How are the values \emph{distributed}? We have two types of
distributions: discrete and continuous.

\subsection{Discrete distributions}\label{sec-basicstat-discdist}

Discrete distributions deal with categorical variables. In the
\texttt{penguins} dataset we have a several variables that have discrete
distributions, one of them is \texttt{species}. We can visualize a
discrete distribution using a \emph{bar chart}:

\begin{Shaded}
\begin{Highlighting}[]
\NormalTok{penguins }\SpecialCharTok{|\textgreater{}}
  \FunctionTok{ggplot}\NormalTok{(}\FunctionTok{aes}\NormalTok{(}\AttributeTok{x =}\NormalTok{ species, }\AttributeTok{fill =}\NormalTok{ species)) }\SpecialCharTok{+}
  \FunctionTok{geom\_bar}\NormalTok{() }\SpecialCharTok{+}
  \FunctionTok{labs}\NormalTok{(}
    \AttributeTok{x =} \StringTok{"Species"}\NormalTok{,}
    \AttributeTok{y =} \StringTok{"Count"}\NormalTok{,}
    \AttributeTok{fill =} \StringTok{"Species"}\NormalTok{)}
\end{Highlighting}
\end{Shaded}

\includegraphics{basic_statistics_files/figure-pdf/unnamed-chunk-3-1.pdf}

In our bar chart we can see the different species that are observed in
the dataset, i.e.~what value the variable \texttt{species} can take, and
how many of each species was observed, i.e.~which values are more
common. We see that the Adelie species is the most common in our
dataset.

\subsection{Continuous distributions}\label{continuous-distributions}

All variables representing measurements of the penguins bodies are
continuous variables, like \texttt{flipper\_length\_mm}. To get an idea
of how the flipper lengths are distributed in our data, we can use a
\emph{histogram}, as shown below. A histogram is a series of bars, where
each bar represents a specific range of values, and the height of the
bar shows the frequency or count of that value. That is, the more common
it is for a flipper length to fall within a certain interval, the higher
the bar.

\begin{Shaded}
\begin{Highlighting}[]
\CommentTok{\#install.packages("patchwork") \# Uncomment to install patchwork}
\FunctionTok{library}\NormalTok{(patchwork)}

\NormalTok{penguins }\SpecialCharTok{|\textgreater{}}
  \FunctionTok{ggplot}\NormalTok{(}\FunctionTok{aes}\NormalTok{(}\AttributeTok{x =}\NormalTok{ flipper\_length\_mm)) }\SpecialCharTok{+}
  \FunctionTok{geom\_histogram}\NormalTok{(}\AttributeTok{col =} \StringTok{"white"}\NormalTok{) }\SpecialCharTok{+}
  \FunctionTok{xlim}\NormalTok{(}\DecValTok{170}\NormalTok{, }\DecValTok{232}\NormalTok{) }\SpecialCharTok{+}
  \FunctionTok{labs}\NormalTok{(}
    \AttributeTok{x =} \StringTok{"Flipper length (mm)"}\NormalTok{,}
    \AttributeTok{y =} \StringTok{"Count"}\NormalTok{)}
\end{Highlighting}
\end{Shaded}

\includegraphics{basic_statistics_files/figure-pdf/unnamed-chunk-4-1.pdf}

In the histogram we can see that flipper lengths vary from around
170-230 mm and that flipper lengths around 190 mm and 215 mm seem
common. But what is the average flipper length? And can we say something
about the variation in flipper lengths? For this we can use descriptive
statistics.

\section{Descriptive statistics}\label{descriptive-statistics}

Descriptive statistics tells us about a variable: the \emph{range} of
the variable, i.e.~its minimum and maximum value, the average or the
\emph{mean} value, and the \emph{variation}. The bar chart and histogram
offers a qualitative way of asserting these values, but descriptive
statistics gives us numbers and doesn't rely on our own (subjective)
observations.

Some of the most common descriptive statistics are

\begin{itemize}
\tightlist
\item
  \textbf{Counts}: The number of observations of each variable in a
  category. For example the number of Adelie, Chinstrap, and Gentoo
  penguins in the \texttt{penguins} dataset. We already saw this in our
  bar chart in section Section~\ref{sec-basicstat-discdist}.
\item
  \textbf{Mean}: The average value of a variable, counted by summing the
  value of each observation and dividing by the total number of
  observations. For example summing all the flipper lengths in
  \texttt{penguins} (the values) and dividing it by the total number of
  penguins observed in \texttt{penguins} (number of observations).
\item
  \textbf{Median}: The middle value after sorting a variable in
  increasing order, i.e.~the value that is bigger than 50\% of all the
  values of the variable. For example if you line up all the 344
  penguins in \texttt{penguins} in order of flipper length, the median
  value is the value in between the 172nd and 173rd penguin, where 172
  is 50\% of the penguins.
\item
  \textbf{Standard Deviation/Variance}: Measures the deviation of data
  from its mean and tells us about how much a variable varies. For
  example the variation of flipper lengths within the Adelie penguins is
  smaller than the variation of the flipper lengths of all the penguins,
  as you will see later. Standard deviation and variance are important
  statistical concepts, but we will not spend too much time on them in
  this course.
\item
  \textbf{Range, Quartiles, Quantiles, Percentiles}: These help in
  understanding the spread and distribution of the data in different
  segments. The range is the minimum and maximum value, while quartiles,
  quantiles, and percentiles divide the data into proportions.
\end{itemize}

We will now take a look at the flipper lengths using these descriptive
statistics.

\subsection{Mean and median of flipper
length}\label{mean-and-median-of-flipper-length}

Both the mean and the median are measures of the center of the
distribution. To calculate these values in R we use the \texttt{mean}
and \texttt{median} functions. Note that here we use \texttt{na.rm} to
handle the \texttt{NA}s, the missing values.

\begin{Shaded}
\begin{Highlighting}[]
\CommentTok{\# Mean}
\FunctionTok{mean}\NormalTok{(penguins}\SpecialCharTok{$}\NormalTok{flipper\_length\_mm, }\AttributeTok{na.rm =} \ConstantTok{TRUE}\NormalTok{)}
\end{Highlighting}
\end{Shaded}

\begin{verbatim}
[1] 200.9152
\end{verbatim}

\begin{Shaded}
\begin{Highlighting}[]
\CommentTok{\# Median}
\FunctionTok{median}\NormalTok{(penguins}\SpecialCharTok{$}\NormalTok{flipper\_length\_mm, }\AttributeTok{na.rm =} \ConstantTok{TRUE}\NormalTok{)}
\end{Highlighting}
\end{Shaded}

\begin{verbatim}
[1] 197
\end{verbatim}

In this case we see that the mean and median is quite similar and both
are a good estimation of the center of the distribution. Usually the
mean is a good statistic to use, but there are two cases when the median
is a better choice. If there are some observations that are very
different to the majority of the observations, called
\href{https://en.wikipedia.org/wiki/Outlier}{\emph{outliers}}, the
median may be the better choice and if the distribution is very
\href{https://en.wikipedia.org/wiki/Skewness}{\emph{skewed}}, where the
distribution has a lot of very high or low values.

We can visualize the {mean} and {median} in our histogram by drawing
vertical lines. In this visualization we have also added the individual
observations as points below the histogram. Where there are many points,
i.e.~where there is less opacity, the count is higher.

\begin{Shaded}
\begin{Highlighting}[]
\NormalTok{mean\_flipper\_length }\OtherTok{\textless{}{-}} \FunctionTok{mean}\NormalTok{(penguins}\SpecialCharTok{$}\NormalTok{flipper\_length\_mm, }\AttributeTok{na.rm =} \ConstantTok{TRUE}\NormalTok{)}
\NormalTok{median\_flipper\_length }\OtherTok{\textless{}{-}} \FunctionTok{median}\NormalTok{(penguins}\SpecialCharTok{$}\NormalTok{flipper\_length\_mm, }\AttributeTok{na.rm =} \ConstantTok{TRUE}\NormalTok{)}
\NormalTok{p1 }\OtherTok{\textless{}{-}}\NormalTok{ penguins }\SpecialCharTok{|\textgreater{}}
  \FunctionTok{ggplot}\NormalTok{(}\FunctionTok{aes}\NormalTok{(}\AttributeTok{x =}\NormalTok{ flipper\_length\_mm, }\AttributeTok{y =} \DecValTok{0}\NormalTok{)) }\SpecialCharTok{+}
  \FunctionTok{geom\_jitter}\NormalTok{(}
    \AttributeTok{alpha =} \FloatTok{0.1}\NormalTok{, }
    \AttributeTok{size =} \DecValTok{5}\NormalTok{,}
    \AttributeTok{height =} \DecValTok{0}\NormalTok{) }\SpecialCharTok{+}
  \FunctionTok{geom\_vline}\NormalTok{(}
    \AttributeTok{xintercept =}\NormalTok{ mean\_flipper\_length,}
    \AttributeTok{col =} \StringTok{"\#f8766d"}\NormalTok{,}
    \AttributeTok{linewidth =} \DecValTok{1}\NormalTok{) }\SpecialCharTok{+}
  \FunctionTok{geom\_vline}\NormalTok{(}
    \AttributeTok{xintercept =}\NormalTok{ median\_flipper\_length,}
    \AttributeTok{col =} \StringTok{"\#00ba38"}\NormalTok{,}
    \AttributeTok{linewidth =} \DecValTok{1}\NormalTok{) }\SpecialCharTok{+}
  \FunctionTok{coord\_cartesian}\NormalTok{(}
    \AttributeTok{xlim =} \FunctionTok{c}\NormalTok{(}\DecValTok{170}\NormalTok{, }\DecValTok{232}\NormalTok{),}
    \AttributeTok{ylim =} \FunctionTok{c}\NormalTok{(}\SpecialCharTok{{-}}\DecValTok{3}\NormalTok{, }\DecValTok{3}\NormalTok{)) }\SpecialCharTok{+}
  \FunctionTok{scale\_x\_continuous}\NormalTok{(}
    \AttributeTok{breaks =} \FunctionTok{c}\NormalTok{(}
      \FunctionTok{round}\NormalTok{(mean\_flipper\_length),}
      \FunctionTok{round}\NormalTok{(median\_flipper\_length))) }\SpecialCharTok{+}
  \FunctionTok{labs}\NormalTok{(}\AttributeTok{x =} \StringTok{"Flipper length (mm)"}\NormalTok{) }\SpecialCharTok{+}
  \FunctionTok{guides}\NormalTok{(}\AttributeTok{col =} \ConstantTok{FALSE}\NormalTok{) }\SpecialCharTok{+}
  \FunctionTok{theme}\NormalTok{(}
    \AttributeTok{panel.grid.major =} \FunctionTok{element\_blank}\NormalTok{(),}
    \AttributeTok{panel.grid.minor =} \FunctionTok{element\_blank}\NormalTok{(),}
    \AttributeTok{axis.title.y =} \FunctionTok{element\_blank}\NormalTok{(),}
    \AttributeTok{axis.text.y =} \FunctionTok{element\_blank}\NormalTok{(),}
    \AttributeTok{axis.ticks.y =} \FunctionTok{element\_blank}\NormalTok{())}

\NormalTok{p2 }\OtherTok{\textless{}{-}}\NormalTok{ penguins }\SpecialCharTok{|\textgreater{}} 
  \FunctionTok{ggplot}\NormalTok{(}\FunctionTok{aes}\NormalTok{(}\AttributeTok{x =}\NormalTok{ flipper\_length\_mm)) }\SpecialCharTok{+} 
  \FunctionTok{geom\_histogram}\NormalTok{(}\AttributeTok{col =} \StringTok{"white"}\NormalTok{) }\SpecialCharTok{+}
  \FunctionTok{geom\_vline}\NormalTok{(}
    \AttributeTok{xintercept =}\NormalTok{ mean\_flipper\_length,}
    \AttributeTok{col =} \StringTok{"\#f8766d"}\NormalTok{,}
    \AttributeTok{linewidth =} \DecValTok{1}\NormalTok{) }\SpecialCharTok{+}
  \FunctionTok{geom\_vline}\NormalTok{(}
    \AttributeTok{xintercept =}\NormalTok{ median\_flipper\_length,}
    \AttributeTok{col =} \StringTok{"\#00ba38"}\NormalTok{,}
    \AttributeTok{linewidth =} \DecValTok{1}\NormalTok{) }\SpecialCharTok{+}
  \FunctionTok{coord\_cartesian}\NormalTok{(}\AttributeTok{xlim =} \FunctionTok{c}\NormalTok{(}\DecValTok{170}\NormalTok{, }\DecValTok{232}\NormalTok{)) }\SpecialCharTok{+}
  \FunctionTok{labs}\NormalTok{(}
    \AttributeTok{x =} \StringTok{"Flipper length (mm)"}\NormalTok{,}
    \AttributeTok{y =} \StringTok{"Count"}\NormalTok{) }\SpecialCharTok{+}
  \FunctionTok{theme}\NormalTok{(}
    \AttributeTok{axis.title.x =} \FunctionTok{element\_blank}\NormalTok{(), }
    \AttributeTok{axis.text.x =} \FunctionTok{element\_blank}\NormalTok{(),}
    \AttributeTok{axis.ticks.x =} \FunctionTok{element\_blank}\NormalTok{())}


\NormalTok{p2 }\SpecialCharTok{/}\NormalTok{ p1 }\SpecialCharTok{+} \FunctionTok{plot\_layout}\NormalTok{(}\AttributeTok{height =} \FunctionTok{c}\NormalTok{(}\DecValTok{20}\NormalTok{, }\DecValTok{1}\NormalTok{))}
\end{Highlighting}
\end{Shaded}

\includegraphics{basic_statistics_files/figure-pdf/unnamed-chunk-6-1.pdf}

\subsection{Quartiles and percentiles}\label{quartiles-and-percentiles}

Percentiles and quartiles are a statistical measures used to describe
the relative standing of a value within a data set. For example, if we
were to order the penguins in our data based on flipper length, the 90th
percentile is the flipper length of the penguin who has a longer
flippers than 90\% of penguins and 10th percentile is the flipper length
of the penguin who has longer flippers than only 10\% of penguins. We
can visualize the 10th and 90th percentile in a histogram using blue,
vertical lines where the 10th percentile is dashed.

\begin{Shaded}
\begin{Highlighting}[]
\NormalTok{percentile\_10 }\OtherTok{\textless{}{-}} \FunctionTok{quantile}\NormalTok{(penguins}\SpecialCharTok{$}\NormalTok{flipper\_length\_mm, }\AttributeTok{prob =} \FloatTok{0.1}\NormalTok{, }\AttributeTok{na.rm =} \ConstantTok{TRUE}\NormalTok{)}
\NormalTok{percentile\_90 }\OtherTok{\textless{}{-}} \FunctionTok{quantile}\NormalTok{(penguins}\SpecialCharTok{$}\NormalTok{flipper\_length\_mm, }\AttributeTok{prob =} \FloatTok{0.9}\NormalTok{, }\AttributeTok{na.rm =} \ConstantTok{TRUE}\NormalTok{)}
\NormalTok{p1 }\OtherTok{\textless{}{-}}\NormalTok{ penguins }\SpecialCharTok{|\textgreater{}}
  \FunctionTok{ggplot}\NormalTok{(}\FunctionTok{aes}\NormalTok{(}\AttributeTok{x =}\NormalTok{ flipper\_length\_mm, }\AttributeTok{y =} \DecValTok{0}\NormalTok{)) }\SpecialCharTok{+}
  \FunctionTok{geom\_jitter}\NormalTok{(}
    \AttributeTok{alpha =} \FloatTok{0.1}\NormalTok{, }
    \AttributeTok{size =} \DecValTok{5}\NormalTok{,}
    \AttributeTok{height =} \DecValTok{0}\NormalTok{) }\SpecialCharTok{+}
  \FunctionTok{geom\_vline}\NormalTok{(}
    \AttributeTok{xintercept =}\NormalTok{ percentile\_10,}
    \AttributeTok{col =} \StringTok{"\#619cff"}\NormalTok{,}
    \AttributeTok{linewidth =} \DecValTok{1}\NormalTok{) }\SpecialCharTok{+}
  \FunctionTok{geom\_vline}\NormalTok{(}
    \AttributeTok{xintercept =}\NormalTok{ percentile\_90,}
    \AttributeTok{col =} \StringTok{"\#619cff"}\NormalTok{,}
    \AttributeTok{linewidth =} \DecValTok{1}\NormalTok{) }\SpecialCharTok{+}
  \FunctionTok{scale\_x\_continuous}\NormalTok{(}
    \AttributeTok{breaks =} \FunctionTok{c}\NormalTok{(}
      \FunctionTok{round}\NormalTok{(percentile\_10[[}\DecValTok{1}\NormalTok{]], }\AttributeTok{digit =} \DecValTok{1}\NormalTok{),}
      \FunctionTok{round}\NormalTok{(percentile\_90[[}\DecValTok{1}\NormalTok{]], }\AttributeTok{digit =} \DecValTok{1}\NormalTok{))) }\SpecialCharTok{+}
  \FunctionTok{coord\_cartesian}\NormalTok{(}
    \AttributeTok{xlim =} \FunctionTok{c}\NormalTok{(}\DecValTok{170}\NormalTok{, }\DecValTok{232}\NormalTok{),}
    \AttributeTok{ylim =} \FunctionTok{c}\NormalTok{(}\SpecialCharTok{{-}}\DecValTok{3}\NormalTok{, }\DecValTok{3}\NormalTok{)) }\SpecialCharTok{+}
  \FunctionTok{labs}\NormalTok{(}\AttributeTok{x =} \StringTok{"Flipper length (mm)"}\NormalTok{) }\SpecialCharTok{+}
  \FunctionTok{guides}\NormalTok{(}\AttributeTok{col =} \ConstantTok{FALSE}\NormalTok{) }\SpecialCharTok{+}
  \FunctionTok{theme}\NormalTok{(}
    \AttributeTok{panel.grid.major =} \FunctionTok{element\_blank}\NormalTok{(),}
    \AttributeTok{panel.grid.minor =} \FunctionTok{element\_blank}\NormalTok{(),}
    \AttributeTok{axis.title.y =} \FunctionTok{element\_blank}\NormalTok{(),}
    \AttributeTok{axis.text.y =} \FunctionTok{element\_blank}\NormalTok{(),}
    \AttributeTok{axis.ticks.y =} \FunctionTok{element\_blank}\NormalTok{())}

\NormalTok{p2 }\OtherTok{\textless{}{-}}\NormalTok{ penguins }\SpecialCharTok{|\textgreater{}} 
  \FunctionTok{ggplot}\NormalTok{(}\FunctionTok{aes}\NormalTok{(}\AttributeTok{x =}\NormalTok{ flipper\_length\_mm)) }\SpecialCharTok{+} 
  \FunctionTok{geom\_histogram}\NormalTok{(}\AttributeTok{col =} \StringTok{"white"}\NormalTok{) }\SpecialCharTok{+}
  \FunctionTok{geom\_vline}\NormalTok{(}
    \AttributeTok{xintercept =}\NormalTok{ percentile\_10,}
    \AttributeTok{col =} \StringTok{"\#619cff"}\NormalTok{,}
    \AttributeTok{linetype =} \DecValTok{2}\NormalTok{,}
    \AttributeTok{linewidth =} \DecValTok{1}\NormalTok{) }\SpecialCharTok{+}
  \FunctionTok{geom\_vline}\NormalTok{(}
    \AttributeTok{xintercept =}\NormalTok{ percentile\_90,}
    \AttributeTok{col =} \StringTok{"\#619cff"}\NormalTok{,}
    \AttributeTok{linewidth =} \DecValTok{1}\NormalTok{) }\SpecialCharTok{+}
  \FunctionTok{coord\_cartesian}\NormalTok{(}\AttributeTok{xlim =} \FunctionTok{c}\NormalTok{(}\DecValTok{170}\NormalTok{, }\DecValTok{232}\NormalTok{)) }\SpecialCharTok{+}

    \FunctionTok{labs}\NormalTok{(}
    \AttributeTok{x =} \StringTok{"Flipper length (mm)"}\NormalTok{,}
    \AttributeTok{y =} \StringTok{"Count"}\NormalTok{) }\SpecialCharTok{+}
  \FunctionTok{theme}\NormalTok{(}
    \AttributeTok{axis.title.x =} \FunctionTok{element\_blank}\NormalTok{(), }
    \AttributeTok{axis.text.x =} \FunctionTok{element\_blank}\NormalTok{(),}
    \AttributeTok{axis.ticks.x =} \FunctionTok{element\_blank}\NormalTok{())}


\NormalTok{p2 }\SpecialCharTok{/}\NormalTok{ p1 }\SpecialCharTok{+} \FunctionTok{plot\_layout}\NormalTok{(}\AttributeTok{height =} \FunctionTok{c}\NormalTok{(}\DecValTok{20}\NormalTok{, }\DecValTok{1}\NormalTok{))}
\end{Highlighting}
\end{Shaded}

\includegraphics{basic_statistics_files/figure-pdf/unnamed-chunk-7-1.pdf}

We see that the 10th percentile is 185, i.e.~that 10\% of the penguins
have flippers shorter than 185 mm, and likewise for the 90th percentile,
90\% of the penguins have flippers shorter than 221 mm. We can calculate
percentiles using the \texttt{quantile} function in R, specifying the
percentile as a proportion, percentile / 100.

\begin{Shaded}
\begin{Highlighting}[]
\NormalTok{flipper\_length\_90th\_percentile }\OtherTok{\textless{}{-}} \FunctionTok{quantile}\NormalTok{(penguins}\SpecialCharTok{$}\NormalTok{flipper\_length\_mm, }\AttributeTok{prob =} \FloatTok{0.9}\NormalTok{, }\AttributeTok{na.rm =} \ConstantTok{TRUE}\NormalTok{)}
\NormalTok{flipper\_length\_90th\_percentile}
\end{Highlighting}
\end{Shaded}

\begin{verbatim}
  90% 
220.9 
\end{verbatim}

We see that the result agrees with what we observed in the histogram.

Quartiles are like percentiles, but they divide the data into 4 equally
sized parts instead of a 100 parts (one for each percent). The first
quartile is larger than 25\% of values, the second is larger than 50\%
of the values and the third is larger than 75\% of the values. Below we
visualize the quartiles in a histogram, where the first and third
quartiles are blue and the second quartile is green.

\begin{Shaded}
\begin{Highlighting}[]
\NormalTok{first\_quartile }\OtherTok{\textless{}{-}} \FunctionTok{quantile}\NormalTok{(penguins}\SpecialCharTok{$}\NormalTok{flipper\_length\_mm, }\AttributeTok{prob =} \FloatTok{0.25}\NormalTok{, }\AttributeTok{na.rm =} \ConstantTok{TRUE}\NormalTok{)}
\NormalTok{third\_quartile }\OtherTok{\textless{}{-}} \FunctionTok{quantile}\NormalTok{(penguins}\SpecialCharTok{$}\NormalTok{flipper\_length\_mm, }\AttributeTok{prob =} \FloatTok{0.75}\NormalTok{, }\AttributeTok{na.rm =} \ConstantTok{TRUE}\NormalTok{)}
\NormalTok{p1 }\OtherTok{\textless{}{-}}\NormalTok{ penguins }\SpecialCharTok{|\textgreater{}}
  \FunctionTok{ggplot}\NormalTok{(}\FunctionTok{aes}\NormalTok{(}\AttributeTok{x =}\NormalTok{ flipper\_length\_mm, }\AttributeTok{y =} \DecValTok{0}\NormalTok{)) }\SpecialCharTok{+}
  \FunctionTok{geom\_jitter}\NormalTok{(}
    \AttributeTok{alpha =} \FloatTok{0.1}\NormalTok{, }
    \AttributeTok{size =} \DecValTok{5}\NormalTok{,}
    \AttributeTok{height =} \DecValTok{0}\NormalTok{) }\SpecialCharTok{+}
  \FunctionTok{geom\_vline}\NormalTok{(}
    \AttributeTok{xintercept =}\NormalTok{ median\_flipper\_length,}
    \AttributeTok{col =} \StringTok{"\#00ba38"}\NormalTok{,}
    \AttributeTok{linewidth =} \DecValTok{1}\NormalTok{) }\SpecialCharTok{+}
  \FunctionTok{geom\_vline}\NormalTok{(}
    \AttributeTok{xintercept =}\NormalTok{ first\_quartile,}
    \AttributeTok{col =} \StringTok{"\#619cff"}\NormalTok{,}
    \AttributeTok{linewidth =} \DecValTok{1}\NormalTok{) }\SpecialCharTok{+}
  \FunctionTok{geom\_vline}\NormalTok{(}
    \AttributeTok{xintercept =}\NormalTok{ third\_quartile,}
    \AttributeTok{col =} \StringTok{"\#619cff"}\NormalTok{,}
    \AttributeTok{linewidth =} \DecValTok{1}\NormalTok{) }\SpecialCharTok{+}
  \FunctionTok{scale\_x\_continuous}\NormalTok{(}
    \AttributeTok{breaks =} \FunctionTok{c}\NormalTok{(}
      \FunctionTok{round}\NormalTok{(first\_quartile[[}\DecValTok{1}\NormalTok{]]),}
      \FunctionTok{round}\NormalTok{(median\_flipper\_length[[}\DecValTok{1}\NormalTok{]]),}
      \FunctionTok{round}\NormalTok{(third\_quartile[[}\DecValTok{1}\NormalTok{]]))) }\SpecialCharTok{+}
  \FunctionTok{coord\_cartesian}\NormalTok{(}
    \AttributeTok{xlim =} \FunctionTok{c}\NormalTok{(}\DecValTok{170}\NormalTok{, }\DecValTok{232}\NormalTok{),}
    \AttributeTok{ylim =} \FunctionTok{c}\NormalTok{(}\SpecialCharTok{{-}}\DecValTok{3}\NormalTok{, }\DecValTok{3}\NormalTok{)) }\SpecialCharTok{+}
  \FunctionTok{labs}\NormalTok{(}\AttributeTok{x =} \StringTok{"Flipper length (mm)"}\NormalTok{) }\SpecialCharTok{+}
  \FunctionTok{guides}\NormalTok{(}\AttributeTok{col =} \ConstantTok{FALSE}\NormalTok{) }\SpecialCharTok{+}
  \FunctionTok{theme}\NormalTok{(}
    \AttributeTok{panel.grid.major =} \FunctionTok{element\_blank}\NormalTok{(),}
    \AttributeTok{panel.grid.minor =} \FunctionTok{element\_blank}\NormalTok{(),}
    \AttributeTok{axis.title.y =} \FunctionTok{element\_blank}\NormalTok{(),}
    \AttributeTok{axis.text.y =} \FunctionTok{element\_blank}\NormalTok{(),}
    \AttributeTok{axis.ticks.y =} \FunctionTok{element\_blank}\NormalTok{())}

\NormalTok{p2 }\OtherTok{\textless{}{-}}\NormalTok{ penguins }\SpecialCharTok{|\textgreater{}} 
  \FunctionTok{ggplot}\NormalTok{(}\FunctionTok{aes}\NormalTok{(}\AttributeTok{x =}\NormalTok{ flipper\_length\_mm)) }\SpecialCharTok{+} 
  \FunctionTok{geom\_histogram}\NormalTok{(}\AttributeTok{col =} \StringTok{"white"}\NormalTok{) }\SpecialCharTok{+}
  \FunctionTok{geom\_vline}\NormalTok{(}
    \AttributeTok{xintercept =}\NormalTok{ median\_flipper\_length,}
    \AttributeTok{col =} \StringTok{"\#00ba38"}\NormalTok{,}
    \AttributeTok{linewidth =} \DecValTok{1}\NormalTok{) }\SpecialCharTok{+}
  \FunctionTok{geom\_vline}\NormalTok{(}
    \AttributeTok{xintercept =}\NormalTok{ first\_quartile,}
    \AttributeTok{col =} \StringTok{"\#619cff"}\NormalTok{,}
    \AttributeTok{linewidth =} \DecValTok{1}\NormalTok{) }\SpecialCharTok{+}
  \FunctionTok{geom\_vline}\NormalTok{(}
    \AttributeTok{xintercept =}\NormalTok{ third\_quartile,}
    \AttributeTok{col =} \StringTok{"\#619cff"}\NormalTok{,}
    \AttributeTok{linewidth =} \DecValTok{1}\NormalTok{) }\SpecialCharTok{+}
  \FunctionTok{coord\_cartesian}\NormalTok{(}\AttributeTok{xlim =} \FunctionTok{c}\NormalTok{(}\DecValTok{170}\NormalTok{, }\DecValTok{232}\NormalTok{)) }\SpecialCharTok{+}
  \FunctionTok{labs}\NormalTok{(}
    \AttributeTok{x =} \StringTok{"Flipper length (mm)"}\NormalTok{,}
    \AttributeTok{y =} \StringTok{"Count"}\NormalTok{) }\SpecialCharTok{+}
  \FunctionTok{theme}\NormalTok{(}
    \AttributeTok{axis.title.x =} \FunctionTok{element\_blank}\NormalTok{(), }
    \AttributeTok{axis.text.x =} \FunctionTok{element\_blank}\NormalTok{(),}
    \AttributeTok{axis.ticks.x =} \FunctionTok{element\_blank}\NormalTok{())}


\NormalTok{p2 }\SpecialCharTok{/}\NormalTok{ p1 }\SpecialCharTok{+} \FunctionTok{plot\_layout}\NormalTok{(}\AttributeTok{height =} \FunctionTok{c}\NormalTok{(}\DecValTok{20}\NormalTok{, }\DecValTok{1}\NormalTok{))}
\end{Highlighting}
\end{Shaded}

\includegraphics{basic_statistics_files/figure-pdf/unnamed-chunk-9-1.pdf}

We visualize the second quartile in green because it is the same as the
median, it is the value that is larger than half of all the values.

\subsection{\texorpdfstring{The \texttt{summary}
function}{The summary function}}\label{the-summary-function}

The \texttt{summary} function is a very useful function to quickly get
an overview of all the variables, their type, their descriptive
statistics, and how many missing values there are, all with just one
command. We take a look at the summary of the \texttt{penguins} dataset.

\begin{Shaded}
\begin{Highlighting}[]
\FunctionTok{summary}\NormalTok{(penguins)}
\end{Highlighting}
\end{Shaded}

\begin{verbatim}
      species          island    bill_length_mm  bill_depth_mm  
 Adelie   :152   Biscoe   :168   Min.   :32.10   Min.   :13.10  
 Chinstrap: 68   Dream    :124   1st Qu.:39.23   1st Qu.:15.60  
 Gentoo   :124   Torgersen: 52   Median :44.45   Median :17.30  
                                 Mean   :43.92   Mean   :17.15  
                                 3rd Qu.:48.50   3rd Qu.:18.70  
                                 Max.   :59.60   Max.   :21.50  
                                 NA's   :2       NA's   :2      
 flipper_length_mm  body_mass_g       sex           year     
 Min.   :172.0     Min.   :2700   female:165   Min.   :2007  
 1st Qu.:190.0     1st Qu.:3550   male  :168   1st Qu.:2007  
 Median :197.0     Median :4050   NA's  : 11   Median :2008  
 Mean   :200.9     Mean   :4202                Mean   :2008  
 3rd Qu.:213.0     3rd Qu.:4750                3rd Qu.:2009  
 Max.   :231.0     Max.   :6300                Max.   :2009  
 NA's   :2         NA's   :2                                 
\end{verbatim}

We see that we get the counts for the categorical variables (discrete
distributions) and the range (``Min.'', ``Max.''), the first and third
quartiles (``1st Qu.'', ``3rd Qu.''), the mean and median, and the
number of missing values (``NA's'') fir the numerical variables
(continuous distributions). The \texttt{summary} function is very good
to use to get an idea about all the variables in a dataset.

\subsection{Exercises}\label{exercises-6}

\begin{enumerate}
\def\labelenumi{\arabic{enumi}.}
\tightlist
\item
  What happens if you run the \texttt{mean} function without
  \texttt{na.rm\ =\ TRUE}?
\item
  Try calculating the mean and median of some other variables in the
  \texttt{penguins} dataset. What happens if you try to calculate the
  mean of a factor?
\item
  Use the \texttt{quantile} function with the parameter
  \texttt{prob\ =\ 0.5} for the variable \texttt{body\_mass\_g} to
  calculate the 50th percentile of the penguin body masses. Compare it
  with the output of the \texttt{median} function for the same variable.
\item
  Run the code below to show the summary of the variables in the
  \texttt{diamonds} dataset. Which variables are categorical
  vs.~numerical? Are there any missing values?
\end{enumerate}

\begin{Shaded}
\begin{Highlighting}[]
\FunctionTok{library}\NormalTok{(tidyverse)}
\FunctionTok{summary}\NormalTok{(diamonds)}
\end{Highlighting}
\end{Shaded}

\section{Relationships between
variables}\label{sec-basicstat-correlation}

So far we have only worked with one variable at a time and we have seen
how to study the variation of one variable, but often it is interesting
to know about the \emph{correlation} of two variables. Correlation tells
us about how much two variables \emph{linearly} correlate. If two
variables are positively (negatively) correlated, large values of one
variable often correspond to large (small) values of the other variable.
Correlation is always a number between -1 (negative correlation) and 1
(positive correlation). 0 indicates that the variables are not
correlated. The figure below shows the relationship between two
variables and the correlation between them.

\begin{Shaded}
\begin{Highlighting}[]
\FunctionTok{library}\NormalTok{(patchwork)}
\FunctionTok{set.seed}\NormalTok{(}\DecValTok{42}\NormalTok{) }\CommentTok{\# Setting a seed to get the same output every time.}

\CommentTok{\# Generating x and y variables based on different correlation scenarios}
\NormalTok{N }\OtherTok{\textless{}{-}} \DecValTok{1000}
\NormalTok{x }\OtherTok{\textless{}{-}} \FunctionTok{rnorm}\NormalTok{(N)}
\NormalTok{y1 }\OtherTok{\textless{}{-}} \SpecialCharTok{{-}}\NormalTok{x }\SpecialCharTok{+} \FunctionTok{rnorm}\NormalTok{(N, }\AttributeTok{sd =} \FloatTok{0.01}\NormalTok{)  }\CommentTok{\# Correlation \textasciitilde{} {-}1}
\NormalTok{y2 }\OtherTok{\textless{}{-}} \SpecialCharTok{{-}}\FloatTok{0.5} \SpecialCharTok{*}\NormalTok{ x }\SpecialCharTok{+} \FunctionTok{rnorm}\NormalTok{(N, }\AttributeTok{sd =} \FunctionTok{sqrt}\NormalTok{(}\FloatTok{0.75}\NormalTok{))  }\CommentTok{\# Correlation \textasciitilde{} {-}0.5}
\NormalTok{y3 }\OtherTok{\textless{}{-}} \FunctionTok{rnorm}\NormalTok{(N)  }\CommentTok{\# Correlation \textasciitilde{} 0}
\NormalTok{y4 }\OtherTok{\textless{}{-}} \FloatTok{0.5} \SpecialCharTok{*}\NormalTok{ x }\SpecialCharTok{+} \FunctionTok{rnorm}\NormalTok{(N, }\AttributeTok{sd =} \FunctionTok{sqrt}\NormalTok{(}\FloatTok{0.75}\NormalTok{))  }\CommentTok{\# Correlation \textasciitilde{} 0.5}
\NormalTok{y5 }\OtherTok{\textless{}{-}}\NormalTok{ x }\SpecialCharTok{+} \FunctionTok{rnorm}\NormalTok{(N, }\AttributeTok{sd =} \FloatTok{0.01}\NormalTok{)  }\CommentTok{\# Correlation \textasciitilde{} 1}

\CommentTok{\# Creating a data frame}
\NormalTok{corr\_df }\OtherTok{\textless{}{-}} \FunctionTok{data.frame}\NormalTok{(x, y1, y2, y3, y4, y5)}

\CommentTok{\# Creating individual plots with titles and modified themes}
\NormalTok{plot1 }\OtherTok{\textless{}{-}} \FunctionTok{ggplot}\NormalTok{(}\AttributeTok{data =}\NormalTok{ corr\_df, }\FunctionTok{aes}\NormalTok{(}\AttributeTok{x =}\NormalTok{ x, }\AttributeTok{y =}\NormalTok{ y1)) }\SpecialCharTok{+}
  \FunctionTok{geom\_point}\NormalTok{(}\AttributeTok{alpha =} \FloatTok{0.5}\NormalTok{) }\SpecialCharTok{+}
  \FunctionTok{geom\_smooth}\NormalTok{(}\AttributeTok{method =} \StringTok{"lm"}\NormalTok{, }\AttributeTok{se =}\NormalTok{ F, }\AttributeTok{col =} \StringTok{"purple"}\NormalTok{) }\SpecialCharTok{+}
  \FunctionTok{theme\_minimal}\NormalTok{() }\SpecialCharTok{+}
  \FunctionTok{labs}\NormalTok{(}
    \AttributeTok{y =} \StringTok{"y"}\NormalTok{, }
    \AttributeTok{subtitle =} \FunctionTok{round}\NormalTok{(}\FunctionTok{cor}\NormalTok{(corr\_df}\SpecialCharTok{$}\NormalTok{x, corr\_df}\SpecialCharTok{$}\NormalTok{y1), }\AttributeTok{digits =} \DecValTok{2}\NormalTok{) ) }\SpecialCharTok{+}
  \FunctionTok{theme}\NormalTok{(}
    \AttributeTok{aspect.ratio =} \DecValTok{1}\NormalTok{,}
    \AttributeTok{axis.text =} \FunctionTok{element\_blank}\NormalTok{(),}
    \AttributeTok{axis.ticks =} \FunctionTok{element\_blank}\NormalTok{(),}
    \AttributeTok{plot.title =} \FunctionTok{element\_text}\NormalTok{(}\AttributeTok{hjust =} \FloatTok{0.5}\NormalTok{)  }\CommentTok{\# Center the title horizontally}
\NormalTok{  )}

\NormalTok{plot2 }\OtherTok{\textless{}{-}} \FunctionTok{ggplot}\NormalTok{(}\AttributeTok{data =}\NormalTok{ corr\_df, }\FunctionTok{aes}\NormalTok{(}\AttributeTok{x =}\NormalTok{ x, }\AttributeTok{y =}\NormalTok{ y2)) }\SpecialCharTok{+}
  \FunctionTok{geom\_point}\NormalTok{(}\AttributeTok{alpha =} \FloatTok{0.5}\NormalTok{) }\SpecialCharTok{+}
  \FunctionTok{geom\_smooth}\NormalTok{(}\AttributeTok{method =} \StringTok{"lm"}\NormalTok{, }\AttributeTok{se =}\NormalTok{ F) }\SpecialCharTok{+}
  \FunctionTok{labs}\NormalTok{(}\AttributeTok{subtitle =} \FunctionTok{round}\NormalTok{(}\FunctionTok{cor}\NormalTok{(corr\_df}\SpecialCharTok{$}\NormalTok{x, corr\_df}\SpecialCharTok{$}\NormalTok{y2), }\AttributeTok{digits =} \DecValTok{2}\NormalTok{)) }\SpecialCharTok{+}
  \FunctionTok{theme\_minimal}\NormalTok{() }\SpecialCharTok{+}
  \FunctionTok{theme}\NormalTok{(}
    \AttributeTok{aspect.ratio =} \DecValTok{1}\NormalTok{,}
    \AttributeTok{axis.text =} \FunctionTok{element\_blank}\NormalTok{(),}
    \AttributeTok{axis.ticks =} \FunctionTok{element\_blank}\NormalTok{(),}
    \AttributeTok{axis.title.y =} \FunctionTok{element\_blank}\NormalTok{(),}
    \AttributeTok{plot.title =} \FunctionTok{element\_text}\NormalTok{(}\AttributeTok{hjust =} \FloatTok{0.5}\NormalTok{)  }\CommentTok{\# Center the title horizontally}
\NormalTok{  )}

\NormalTok{plot3 }\OtherTok{\textless{}{-}} \FunctionTok{ggplot}\NormalTok{(}\AttributeTok{data =}\NormalTok{ corr\_df, }\FunctionTok{aes}\NormalTok{(}\AttributeTok{x =}\NormalTok{ x, }\AttributeTok{y =}\NormalTok{ y3)) }\SpecialCharTok{+}
  \FunctionTok{geom\_point}\NormalTok{(}\AttributeTok{alpha =} \FloatTok{0.5}\NormalTok{) }\SpecialCharTok{+}
  \FunctionTok{geom\_smooth}\NormalTok{(}\AttributeTok{method =} \StringTok{"lm"}\NormalTok{, }\AttributeTok{se =}\NormalTok{ F) }\SpecialCharTok{+}
  \FunctionTok{labs}\NormalTok{(}
    \AttributeTok{title =} \StringTok{"Correlation"}\NormalTok{,}
    \AttributeTok{subtitle =} \FunctionTok{round}\NormalTok{(}\FunctionTok{cor}\NormalTok{(corr\_df}\SpecialCharTok{$}\NormalTok{x, corr\_df}\SpecialCharTok{$}\NormalTok{y3), }\AttributeTok{digits =} \DecValTok{2}\NormalTok{)) }\SpecialCharTok{+}
  \FunctionTok{theme\_minimal}\NormalTok{() }\SpecialCharTok{+}
  \FunctionTok{theme}\NormalTok{(}
    \AttributeTok{aspect.ratio =} \DecValTok{1}\NormalTok{,}
    \AttributeTok{axis.text =} \FunctionTok{element\_blank}\NormalTok{(),}
    \AttributeTok{axis.ticks =} \FunctionTok{element\_blank}\NormalTok{(),}
    \AttributeTok{axis.title.y =} \FunctionTok{element\_blank}\NormalTok{(),}
    \AttributeTok{plot.title =} \FunctionTok{element\_text}\NormalTok{(}\AttributeTok{hjust =} \FloatTok{0.5}\NormalTok{)  }\CommentTok{\# Center the title horizontally}
\NormalTok{  )}

\NormalTok{plot4 }\OtherTok{\textless{}{-}} \FunctionTok{ggplot}\NormalTok{(}\AttributeTok{data =}\NormalTok{ corr\_df, }\FunctionTok{aes}\NormalTok{(}\AttributeTok{x =}\NormalTok{ x, }\AttributeTok{y =}\NormalTok{ y4)) }\SpecialCharTok{+}
  \FunctionTok{geom\_point}\NormalTok{(}\AttributeTok{alpha =} \FloatTok{0.5}\NormalTok{) }\SpecialCharTok{+}
  \FunctionTok{geom\_smooth}\NormalTok{(}\AttributeTok{method =} \StringTok{"lm"}\NormalTok{, }\AttributeTok{se =}\NormalTok{ F) }\SpecialCharTok{+}
  \FunctionTok{labs}\NormalTok{(}\AttributeTok{subtitle =} \FunctionTok{round}\NormalTok{(}\FunctionTok{cor}\NormalTok{(corr\_df}\SpecialCharTok{$}\NormalTok{x, corr\_df}\SpecialCharTok{$}\NormalTok{y4), }\AttributeTok{digits =} \DecValTok{2}\NormalTok{)) }\SpecialCharTok{+}
  \FunctionTok{theme\_minimal}\NormalTok{() }\SpecialCharTok{+}
  \FunctionTok{theme}\NormalTok{(}
    \AttributeTok{aspect.ratio =} \DecValTok{1}\NormalTok{,}
    \AttributeTok{axis.text =} \FunctionTok{element\_blank}\NormalTok{(),}
    \AttributeTok{axis.ticks =} \FunctionTok{element\_blank}\NormalTok{(),}
    \AttributeTok{axis.title.y =} \FunctionTok{element\_blank}\NormalTok{(),}
    \AttributeTok{plot.title =} \FunctionTok{element\_text}\NormalTok{(}\AttributeTok{hjust =} \FloatTok{0.5}\NormalTok{)  }\CommentTok{\# Center the title horizontally}
\NormalTok{  )}

\NormalTok{plot5 }\OtherTok{\textless{}{-}} \FunctionTok{ggplot}\NormalTok{(}\AttributeTok{data =}\NormalTok{ corr\_df, }\FunctionTok{aes}\NormalTok{(}\AttributeTok{x =}\NormalTok{ x, }\AttributeTok{y =}\NormalTok{ y5)) }\SpecialCharTok{+}
  \FunctionTok{geom\_point}\NormalTok{(}\AttributeTok{alpha =} \FloatTok{0.5}\NormalTok{) }\SpecialCharTok{+}
  \FunctionTok{geom\_smooth}\NormalTok{(}\AttributeTok{method =} \StringTok{"lm"}\NormalTok{, }\AttributeTok{se =}\NormalTok{ F) }\SpecialCharTok{+}
  \FunctionTok{labs}\NormalTok{(}\AttributeTok{subtitle =} \FunctionTok{round}\NormalTok{(}\FunctionTok{cor}\NormalTok{(corr\_df}\SpecialCharTok{$}\NormalTok{x, corr\_df}\SpecialCharTok{$}\NormalTok{y5), }\AttributeTok{digits =} \DecValTok{1}\NormalTok{)) }\SpecialCharTok{+}
  \FunctionTok{theme\_minimal}\NormalTok{() }\SpecialCharTok{+}
  \FunctionTok{theme}\NormalTok{(}
    \AttributeTok{aspect.ratio =} \DecValTok{1}\NormalTok{,}
    \AttributeTok{axis.text =} \FunctionTok{element\_blank}\NormalTok{(),}
    \AttributeTok{axis.ticks =} \FunctionTok{element\_blank}\NormalTok{(),}
    \AttributeTok{axis.title.y =} \FunctionTok{element\_blank}\NormalTok{(),}
    \AttributeTok{plot.title =} \FunctionTok{element\_text}\NormalTok{(}\AttributeTok{hjust =} \FloatTok{0.5}\NormalTok{)  }\CommentTok{\# Center the title horizontally}
\NormalTok{  )}

\CommentTok{\# Using patchwork to arrange plots side by side}
\NormalTok{(plot1 }\SpecialCharTok{|}\NormalTok{ plot2 }\SpecialCharTok{|}\NormalTok{ plot3 }\SpecialCharTok{|}\NormalTok{ plot4 }\SpecialCharTok{|}\NormalTok{ plot5) }\SpecialCharTok{+} \FunctionTok{plot\_layout}\NormalTok{(}\AttributeTok{ncol =} \DecValTok{5}\NormalTok{)}
\end{Highlighting}
\end{Shaded}

\includegraphics{basic_statistics_files/figure-pdf/unnamed-chunk-12-1.pdf}

We see that the more the points form a line, the closer to -1 or 1 the
correlation is. We will come back to correlation when we deal with
multiple linear regression later in the course.

In the \texttt{penguins} dataset we would expect that bigger penguins
have bigger flippers, that is we would expect
\texttt{flipper\_length\_mm} to be positively correlated with
\texttt{body\_mass\_g}. The easiest way to investigate this is through a
\emph{scatter plot} and a line.

\begin{Shaded}
\begin{Highlighting}[]
\NormalTok{penguins }\SpecialCharTok{|\textgreater{}}
  \FunctionTok{ggplot}\NormalTok{(}\FunctionTok{aes}\NormalTok{(}\AttributeTok{x =}\NormalTok{ body\_mass\_g, }\AttributeTok{y =}\NormalTok{ flipper\_length\_mm)) }\SpecialCharTok{+}
  \FunctionTok{geom\_point}\NormalTok{() }\SpecialCharTok{+}
  \FunctionTok{geom\_smooth}\NormalTok{(}\AttributeTok{method =} \StringTok{"lm"}\NormalTok{, }\AttributeTok{se =} \ConstantTok{FALSE}\NormalTok{) }\SpecialCharTok{+}
  \FunctionTok{labs}\NormalTok{(}
    \AttributeTok{x =} \StringTok{"Body Mass (g)"}\NormalTok{,}
    \AttributeTok{y =} \StringTok{"Flipper Length (mm)"}\NormalTok{)}
\end{Highlighting}
\end{Shaded}

\includegraphics{basic_statistics_files/figure-pdf/unnamed-chunk-13-1.pdf}

We see a \emph{trend}: when body mass increases, flipper length
increases as well, on average. We can also calculate correlation as a
number using the \texttt{cor} function. Here we use
\texttt{use\ =\ "complete.obs"} to make sure we don't include missing
values. This parameter fills the same function as \texttt{na.rm} in the
\texttt{mean} function, but unfortunately the parameter name is
different.

\begin{Shaded}
\begin{Highlighting}[]
\FunctionTok{cor}\NormalTok{(penguins}\SpecialCharTok{$}\NormalTok{flipper\_length\_mm, penguins}\SpecialCharTok{$}\NormalTok{body\_mass\_g, }\AttributeTok{use =} \StringTok{"complete.obs"}\NormalTok{)}
\end{Highlighting}
\end{Shaded}

\begin{verbatim}
[1] 0.8712018
\end{verbatim}

We see that body mass and flipper length are highly positively
correlated, 0.87. Finding correlated variables is an important part of
data analysis as it \emph{may} indicate causal effects, such as the
correlation between smoking and lung cancer once did. However,
correlation is \emph{no guarantee} for causation, as the famous saying
goes

\begin{quote}
Correlation does not imply causation.
\end{quote}

To assert causation other assumptions or experiments are needed.

In this section you have learned about correlation between two numerical
variables. In the upcoming EDA part of the course you will learn how to
make visualizations to study how continuous variables vary with
categorical variables, and how categorical variables vary with each
other.

\subsection{Exercises}\label{exercises-7}

\begin{enumerate}
\def\labelenumi{\arabic{enumi}.}
\tightlist
\item
  Try to understand the output of the following code that computes the
  correlation of all numeric variables and creates a \emph{correlation}
  plot. Which variables are highly correlated? Does -1 mean high or low
  correlation? Note: Here we use the pipe operator,
  \texttt{\textbar{}\textgreater{}}, which allows you to pass arguments
  to function in what is called piping. You will learn more about this
  in the R4DS book.
\end{enumerate}

\begin{Shaded}
\begin{Highlighting}[]
\CommentTok{\#install.packages("ggcorrplot")}
\FunctionTok{library}\NormalTok{(ggcorrplot)}
\CommentTok{\# Correlation matrix}
\NormalTok{penguins }\SpecialCharTok{|\textgreater{}}
  \FunctionTok{select}\NormalTok{(}\FunctionTok{where}\NormalTok{(is.numeric)) }\SpecialCharTok{|\textgreater{}}
  \FunctionTok{cor}\NormalTok{(}\AttributeTok{use =} \StringTok{"complete.obs"}\NormalTok{) }
\end{Highlighting}
\end{Shaded}

\begin{verbatim}
                  bill_length_mm bill_depth_mm flipper_length_mm body_mass_g
bill_length_mm        1.00000000   -0.23505287         0.6561813  0.59510982
bill_depth_mm        -0.23505287    1.00000000        -0.5838512 -0.47191562
flipper_length_mm     0.65618134   -0.58385122         1.0000000  0.87120177
body_mass_g           0.59510982   -0.47191562         0.8712018  1.00000000
year                  0.05454458   -0.06035364         0.1696751  0.04220939
                         year
bill_length_mm     0.05454458
bill_depth_mm     -0.06035364
flipper_length_mm  0.16967511
body_mass_g        0.04220939
year               1.00000000
\end{verbatim}

\begin{Shaded}
\begin{Highlighting}[]
\CommentTok{\# Correlation plot}
\NormalTok{penguins }\SpecialCharTok{|\textgreater{}}
  \FunctionTok{select}\NormalTok{(}\FunctionTok{where}\NormalTok{(is.numeric)) }\SpecialCharTok{|\textgreater{}}
  \FunctionTok{cor}\NormalTok{(}\AttributeTok{use =} \StringTok{"complete.obs"}\NormalTok{) }\SpecialCharTok{|\textgreater{}}
  \FunctionTok{ggcorrplot}\NormalTok{()}
\end{Highlighting}
\end{Shaded}

\includegraphics{basic_statistics_files/figure-pdf/unnamed-chunk-15-1.pdf}

\section{Summary}\label{summary}

In this chapter we have covered many basic statistical concepts: from
variables, their distributions, to descriptive statistics. All of these
concepts are useful in data analysis and the more you work with data,
the more sense they make. For now it is enough that you have an idea of
these concepts and know where you can find more information about them
when you need it (in this chapter or on the web!). You are now ready to
dive into the core of the course: visualizing, transforming, tidying and
\emph{understanding} data using the course literature,
\href{https://r4ds.hadley.nz/}{R4DS}. This may help you understand some
of the concepts we have covered so far in a more visual way, so if it
feels a bit hard right now, hang in there and you'll see that things
will become more clear.

\part{Exploratory data analysis with R4DS}

Now that we have covered basic R programming and basic concepts in
statistics, we are ready to dive into Exploratory Data Analysis (EDA).
EDA is used to gain an understanding of a dataset and is an essential
for anyone working with data. In this part of the course we will be
using the course literature, \href{https://r4ds.hadley.nz/}{R for Data
Science}, and the \texttt{tidyverse} package. The book is written by one
of the authors of the \texttt{tidyverse} package, Hadley Wickham, along
with Mine Çetinkaya-Rundel, and Garrett Grolemund, so we are certain you
are in good hands.

\chapter{Reading instructions}\label{reading-instructions}

\section*{Introduction}\label{introduction}
\addcontentsline{toc}{section}{Introduction}

\markright{Introduction}

Recommended

\section*{Whole game}\label{whole-game}
\addcontentsline{toc}{section}{Whole game}

\markright{Whole game}

This chapter covers everything needed for Exploratory Data Analysis
(EDA).

\begin{itemize}
\tightlist
\item
  Chapters ``Data visualization'' and ``Data transformation'' are
  recommended.
\item
  The chapter is ``Data tidying'' is recommended as tidy data and
  pivoting are important concepts. However, there are no recommended
  exercises on this chapter, and it won't covered in the assignments.
\item
  The ``Workflow:'' chapters can be considered optional. If you are new
  to coding or R they provide useful tips to how to improve your
  workflow.
\item
  The chapter ``Data import'' can be skipped if you know how to use the
  \texttt{read\_csv} function.
\end{itemize}

\section*{Visualize}\label{visualize}
\addcontentsline{toc}{section}{Visualize}

\markright{Visualize}

This chapter covers how to visualize data and EDA, which is an important
part of data analysis.

\begin{itemize}
\tightlist
\item
  Chapter ``Layers'' gives a deeper understanding for visualizing with
  \texttt{ggplot}. It can be considered optional, but chapter
  ``Exploratory data analysis'' and ``Communication'' builds on some of
  the content in this chapter.
\item
  Chapter ``Exploratory data analysis'' is essential to the course.
\item
  Chapter ``Communication'' is much like chapter ``Layers'', it helps
  you create better visualizations. Make sure you read 11.2 and 11.4,
  the rest can be considered optional.
\end{itemize}

\chapter{Recommended exercises}\label{recommended-exercises}

These are the recommended exercises in R4DS with some comments. We
recommend reading through these recommendations before starting the
exercises as they are meant to help you learn more efficiently. Some of
the exercises in the book are rather straight forward, but some are more
difficult (as we also struggled with them when writing the solutions).
So if you feel that you don't know where to start with an exercise or
that you are stuck, consult the suggested solutions.

\subsection*{Data visualization}\label{data-visualization}
\addcontentsline{toc}{subsection}{Data visualization}

All exercises.

\begin{itemize}
\tightlist
\item
  1.2.4 Exercise 4: try using \texttt{geom\_density} as well and try
  using \texttt{geom\_density} and coloring by \texttt{cut}.
\end{itemize}

\subsection*{Data transformation}\label{data-transformation}
\addcontentsline{toc}{subsection}{Data transformation}

\textbf{3.2 Rows}

All exercises

Read 3.3 first and use the following functions: * \texttt{relocate} or
\texttt{select} to move columns of interest to the front of the tibble,
otherwise it is hard to see the result. * Use \texttt{mutate} in
Exercise 3 to create a new variable for speed and use that variable for
sorting, i.e.~\texttt{arrange}-ing. If you are having trouble figuring
out the mathematical expression, remember from high school physics the
relationship between time, distance, and speed. In Swedish
\href{https://sv.wikipedia.org/wiki/SVT-triangel}{``SVT''-triangeln}. If
you just want the answer, look at 3.3.1.

\textbf{3.3 Columns}

All exercises

\begin{itemize}
\tightlist
\item
  In exercise 1, try to think about the relationship. It is not
  necessary to successfully create a new variable that is the same as
  \texttt{dep\_delay}. See solutions for a suggestion of how, even
  though it may be beyond the scope of the course.
\item
  In exercise 4 for, use \texttt{?any\_of} and look at the function from
  the \texttt{tidyselect} package.
\end{itemize}

\textbf{3.5 Groups}

All exercises in 3.5

\begin{itemize}
\tightlist
\item
  In exercise 1, remember to use \texttt{na.rm} when calculating the
  average.
\end{itemize}

\subsection*{Exploratory data analysis}\label{exploratory-data-analysis}
\addcontentsline{toc}{subsection}{Exploratory data analysis}

Try to do all exercises but some are quite difficult. If you get stuck
consult the suggested solutions.

The exercises requires some basic knowledge about
\href{https://en.wikipedia.org/wiki/Diamond}{diamonds}: Measurements

\begin{itemize}
\tightlist
\item
  x, y, z: The dimensions of a diamond (width, length, height)
\end{itemize}

The four Cs of diamond grading *
\href{https://en.wikipedia.org/wiki/Carat_(mass)}{Carat}: A unit of mass
(like grams) * \href{https://en.wikipedia.org/wiki/Diamond_cut}{Cut}:
Style or design, but rather symmetry, proportioning, and polishing than
shape. A well cut diamond is more *luminous*. *
\href{https://en.wikipedia.org/wiki/Diamond_clarity}{Clarity}: The
existence and visual appearance of internal characteristics, called
inclusions, and surface defects, called blemishes. *
\href{https://en.wikipedia.org/wiki/Diamond_color}{Color}: A perfect
diamond has no hue or color.

\subsection*{Communication}\label{communication}
\addcontentsline{toc}{subsection}{Communication}

No recommended exercises, instead practice plotting as you work with
chapter 10 and assignments.

\part{Statistical inference and prediction}

Up to this point, we've been looking into datasets and learning about
the variables and relationships in them. To take our analysis beyond
datasets to knowledge about the world, we need to use \emph{statistical
inference}. Statistical inference enables us to learn about entire
\emph{populations} based on what we have found in our \emph{sample},
i.e.~to make inferences from data. Statistical inference is the first
step towards asking ``why?''-questions to our data and understanding our
world better. Prediction on the other hand lets us use previously
observed data to make powerful predictions about future observations
that can help us make better decisions or automate processes.

We will start by introducing hypothesis testing, the backbone of
statistical inference, and then linear and logistic regression which can
be used both for statistical inference and for prediction. Once you are
familiar with these tools you will be able to answer questions about
populations based on a sample, the relationship between both continuous
and categorical variables, and making predictions based on observed
variables.

{In this course we aim to introduce these concepts and showing the
possibilities with these methods without going into too much detail.
This will allow you to understand when to use which tool, but to be sure
you are using the tool correctly you have to go deeper into the theory
behind them, or consult an expert. You can think of it as learning about
the function of a car, a train, or an airplane. You can know when to use
which mode of transportation, even if you don't know how to operate each
one of them and may have to rely on experts to do that.}

\chapter{Hypothesis testing}\label{sec-inference}

Hypothesis testing is a systematic way of answering questions about a
population using data from a random sample. It is the backbone of
statistical inference and essential to have an understanding for to be
able to use more powerful techniques, like linear and logistic
regression. In this chapter we will be working with data presented by
the Swedish national public television broadcaster, SVT, on the current
election poll (19/9-2023)
\href{https://www.svt.se/special/valjarbarometern/}{Väljarbarometern}.

For a more technical/mathematical introduction to hypothesis testing,
please consult the chapter on multiple testing (Chapter 13, 2nd edition)
in \href{https://www.statlearning.com/}{An Introduction to Statistical
Learning}.

\section{Sample and population}\label{sample-and-population}

We will be using the \emph{sample} in Väljarbarometern (visited 19
September 2023) collected by Kantar Public by interviewing 3072 people
between August 21 and September 3, 2023. We call this a sample since it
is a small part of the whole \emph{population} of all citizens of
Sweden.

{The data comes in aggregated form, i.e.~in percentages, of the 3072
people but we have recreated the original sample below. Each row is the
answer of one participant in the poll and the column \texttt{party}
shows what the person said they would vote for. Below we show the code
for recreating the sample from aggregated data. You don't need to
understand the code, but if you want to recreate the sample yourself you
can use the function \texttt{recreate\_valjarbarometern}.}

\begin{Shaded}
\begin{Highlighting}[]
\FunctionTok{library}\NormalTok{(tidyverse)}

\NormalTok{recreate\_valjarbarometern }\OtherTok{\textless{}{-}} \ControlFlowTok{function}\NormalTok{(valjarbarometern\_aggregate, n)}
\NormalTok{\{}
  \CommentTok{\# Calculate number of observations of each party}
\NormalTok{  obs\_per\_party }\OtherTok{\textless{}{-}} \FunctionTok{round}\NormalTok{(n }\SpecialCharTok{*}\NormalTok{ valjarbarometern\_aggregate}\SpecialCharTok{$}\NormalTok{percentage }\SpecialCharTok{/} \DecValTok{100}\NormalTok{)}
  
  \CommentTok{\# Create a list of party observations}
\NormalTok{  party\_observations }\OtherTok{\textless{}{-}} \FunctionTok{mapply}\NormalTok{(rep, valjarbarometern\_aggregate}\SpecialCharTok{$}\NormalTok{party, obs\_per\_party)}
  
  \CommentTok{\# Flatten the list into a single vector}
\NormalTok{  sample\_data }\OtherTok{\textless{}{-}} \FunctionTok{unlist}\NormalTok{(party\_observations)}
  
  \CommentTok{\# Create a data frame with the samples}
\NormalTok{  valjarbarometern }\OtherTok{\textless{}{-}} \FunctionTok{tibble}\NormalTok{(}\AttributeTok{party =}\NormalTok{ sample\_data)}
  
  \FunctionTok{return}\NormalTok{(valjarbarometern)}
\NormalTok{\}}
\CommentTok{\# Number of people in survey}
\NormalTok{n }\OtherTok{=} \DecValTok{3072}

\NormalTok{parties }\OtherTok{\textless{}{-}} \FunctionTok{c}\NormalTok{(}\StringTok{"V"}\NormalTok{, }\StringTok{"S"}\NormalTok{, }\StringTok{"MP"}\NormalTok{, }\StringTok{"C"}\NormalTok{, }\StringTok{"L"}\NormalTok{, }\StringTok{"KD"}\NormalTok{, }\StringTok{"M"}\NormalTok{, }\StringTok{"SD"}\NormalTok{, }\StringTok{"Other"}\NormalTok{)}
\CommentTok{\# Aggregated sample}
\NormalTok{valjarbarometern\_aggregate }\OtherTok{\textless{}{-}} \FunctionTok{tibble}\NormalTok{(}
  \AttributeTok{party =} \FunctionTok{factor}\NormalTok{(parties, }\AttributeTok{levels =}\NormalTok{ parties), }\CommentTok{\# Using levels to preserve order}
  \AttributeTok{percentage =} \FunctionTok{c}\NormalTok{(}\FloatTok{7.6}\NormalTok{, }\FloatTok{37.6}\NormalTok{, }\FloatTok{4.3}\NormalTok{, }\FloatTok{3.9}\NormalTok{, }\DecValTok{3}\NormalTok{, }\FloatTok{3.1}\NormalTok{, }\DecValTok{19}\NormalTok{, }\FloatTok{19.7}\NormalTok{, }\FloatTok{1.7}\NormalTok{))}

\NormalTok{valjarbarometern }\OtherTok{\textless{}{-}} \FunctionTok{recreate\_valjarbarometern}\NormalTok{(}
\NormalTok{  valjarbarometern\_aggregate, n)}
\end{Highlighting}
\end{Shaded}

We can take a look at the summary of our sample

\begin{Shaded}
\begin{Highlighting}[]
\FunctionTok{summary}\NormalTok{(valjarbarometern)}
\end{Highlighting}
\end{Shaded}

\begin{verbatim}
     party     
 S      :1155  
 SD     : 605  
 M      : 584  
 V      : 233  
 MP     : 132  
 C      : 120  
 (Other): 239  
\end{verbatim}

The summary tells us how many people would vote for each party if the
election were today, in our sample. A more common way to look at this is
using a bar chart and showing the y-axis in percentage. The dashed line
indicates the
\href{https://www.informationsverige.se/en/om-sverige/att-paverka-i-sverige/det-svenska-valsystemet.html}{limit}
for being included in the Swedish parliament (Sveriges riksdag), 4\%.

\begin{Shaded}
\begin{Highlighting}[]
\CommentTok{\# Define the colors in the order V, S, MP, C, L, KD, M, SD, Other. Make sure party column follows the same order!}
\NormalTok{party\_colors }\OtherTok{\textless{}{-}} \FunctionTok{c}\NormalTok{(}\StringTok{"\#800000"}\NormalTok{, }\StringTok{"\#FF0000"}\NormalTok{, }\StringTok{"\#90EE90"}\NormalTok{, }\StringTok{"\#008000"}\NormalTok{, }\StringTok{"\#0054A6"}\NormalTok{, }\StringTok{"\#00008B"}\NormalTok{, }\StringTok{"\#00AEEF"}\NormalTok{, }\StringTok{"\#FFFF00"}\NormalTok{, }\StringTok{"\#808080"}\NormalTok{)}
\NormalTok{valjarbarometern }\SpecialCharTok{|\textgreater{}}
  \FunctionTok{ggplot}\NormalTok{(}\FunctionTok{aes}\NormalTok{(}\AttributeTok{x =}\NormalTok{ party, }\AttributeTok{fill =}\NormalTok{ party)) }\SpecialCharTok{+} 
  \FunctionTok{geom\_bar}\NormalTok{(}\FunctionTok{aes}\NormalTok{(}\AttributeTok{y =} \FunctionTok{after\_stat}\NormalTok{(count}\SpecialCharTok{/}\FunctionTok{sum}\NormalTok{(count) }\SpecialCharTok{*} \DecValTok{100}\NormalTok{))) }\SpecialCharTok{+} \CommentTok{\# Calculating percentages}
  \FunctionTok{geom\_hline}\NormalTok{(}
    \AttributeTok{yintercept =} \DecValTok{4}\NormalTok{,}
    \AttributeTok{col =} \StringTok{"black"}\NormalTok{,}
    \AttributeTok{linetype =} \DecValTok{2}\NormalTok{) }\SpecialCharTok{+}
  \FunctionTok{labs}\NormalTok{(}
    \AttributeTok{x =} \StringTok{"Party"}\NormalTok{,}
    \AttributeTok{y =} \StringTok{"Percentage"}\NormalTok{) }\SpecialCharTok{+}
  \FunctionTok{guides}\NormalTok{(}\AttributeTok{fill =} \StringTok{"none"}\NormalTok{) }\SpecialCharTok{+} \CommentTok{\# Remove fill legend}
  \FunctionTok{scale\_fill\_manual}\NormalTok{(}\AttributeTok{values =}\NormalTok{ party\_colors) }\CommentTok{\# Set colors to party colors}
\end{Highlighting}
\end{Shaded}

\includegraphics{hypothesis_testing_files/figure-pdf/unnamed-chunk-3-1.pdf}

In our bar chart we see that several parties seems to be at risk of
ending up below 4\% if there was an election today: MP (Miljöpartiet)
seems to be just above the 4\% limit and C (Centern) just below, while L
(Liberalerna) and KD (Kristdemokraterna) seem worst off. Let's take a
closer look at the percentages

\begin{Shaded}
\begin{Highlighting}[]
\NormalTok{n\_participants }\OtherTok{\textless{}{-}} \DecValTok{3072} \CommentTok{\# Number of participants in survey}
\NormalTok{valjarbarometern }\SpecialCharTok{|\textgreater{}}
  \FunctionTok{summarize}\NormalTok{(}
    \AttributeTok{percentage =} \FunctionTok{round}\NormalTok{(}\DecValTok{100} \SpecialCharTok{*} \FunctionTok{n}\NormalTok{() }\SpecialCharTok{/}\NormalTok{ n, }\AttributeTok{digits =} \DecValTok{1}\NormalTok{), }\CommentTok{\# Calculating percentage}
    \AttributeTok{.by =}\NormalTok{ party) }\CommentTok{\# Grouping by party}
\end{Highlighting}
\end{Shaded}

\begin{verbatim}
# A tibble: 9 x 2
  party percentage
  <fct>      <dbl>
1 V            7.6
2 S           37.6
3 MP           4.3
4 C            3.9
5 L            3  
6 KD           3.1
7 M           19  
8 SD          19.7
9 Other        1.7
\end{verbatim}

The summary confirms what we could tell from the bar chart, MP is just
above the 4\% limit with 4.3\% and C is just below the limit with 3.9\%.
But this result comes from a sample of 3072, what does that actually
tell us about how the whole \emph{population} of Sweden would vote? We
would like to know what the value would be if the whole population
voted. To learn about this from our sample, we need to perform
statistical inference, and in particular hypothesis testing.

\subsection{Exercises}\label{exercises-8}

\begin{enumerate}
\def\labelenumi{\arabic{enumi}.}
\tightlist
\item
  Previously we worked with the \texttt{palmerpenguins} dataset. What is
  the population in that dataset? Use \texttt{?penguins} after loading
  the package (\texttt{library(palmerpenguins)}) to read more about the
  data. Where and when was the data collected?
\end{enumerate}

\section{Hypothesis testing
(proportions)}\label{hypothesis-testing-proportions}

Hypothesis testing is method that allows us to draw inferences about a
population using a sample. In more simpler terms, hypothesis testing
let's us use a sample to answer questions about the population with some
uncertainty. The procedures follows four steps:

\textbf{1. Formulate hypotheses}

Formulate the hypotheses as a null hypothesis, \(H_0\), and an
alternative hypothesis, \(H_a\) (sometimes \(H_1\)). The hypotheses
should reflect the question that we have: in our case, for MP, we want
to know if the percentage in the population is larger than 4\%. We can
formulate the hypotheses for this question as

\begin{quote}
\textbf{\(H_0\)}: The percentage of people who would vote for MP is
\emph{smaller than or equal} to 4\%.
\end{quote}

\begin{quote}
\textbf{\(H_a\)}: The percentage of people who would vote for MP is
\emph{larger} than 4\%.
\end{quote}

\(H_0\) and \(H_a\) play different roles in hypothesis testing, \(H_0\)
is our baseline and \(H_a\) represents the hypothesis we want to prove.
The hypotheses should be mutually exclusive and exhaustive, that is they
should cover every possible outcome. In our example the percentage must
be a number between 0 and 100, which must be either smaller than or
equal to 4 or larger than 4, i.e.~our hypotheses cover all
possibilities.

{\textbf{2. Calculate test statistic}}

{Calculate a test statistic that summarizes the \emph{evidence} in our
data \emph{against} the null hypothesis. This part is were most of
mathematical and statistical theory is needed and \textbf{in this course
we will leave this to a test function in R.} Instead we will run use a
test funciton in R.}

\textbf{2. Use a test function in R}

Run a function in R that performs a statistical hypothesis test that
summarizes the \emph{evidence} in our data \emph{against} the null
hypothesis.

\textbf{3. \(p\)-value}

Calculate a \(p\)-value based on the evidence, which is done by the test
function. The \(p\)-value is of great importance in hypothesis testing
and it is important to have a basic understanding of it. In our case,
the \(p\)-value is the probability of observing 4.3\% or larger (further
from the null hypothesis) in a sample of 3072 people, \emph{if the null
hypothesis was true}. \(p\)-values can always be formulated in terms of
the null hypothesis and the observed parameter, in this case the
percentage voting for MP, as: ``The probability of observing the
\emph{observed value} or more extreme, i.e.~4.3\% or larger, in a sample
given that the null hypothesis is true''.

\textbf{4. Decision}

Decide to reject or not reject the null hypothesis based on the
\(p\)-value. This is a decision rule that is decided beforehand, based
on a \emph{significance level}, usually denoted \(\alpha\). If the
\(p\)-value is smaller than the significance level, we reject the null
hypothesis at significance level \(\alpha\). The most common
significance level is 0.05, and we will use it throughout this course.
If the \(p\)-value is \emph{smaller} than 0.05 we say that the evidence
in the data is \emph{statistically significant} at \emph{significance
level}, \(\alpha\), 0.05 to reject the null hypothesis. This decision
rule is dichotomous, i.e.~it is either true or false. This is commonly
used, but the \(p\)-value is valuable on it's own as a \(p\)-value of
0.06 is still a very interesting result, while a \(p\)-value of 0.9 is
not. A low \(p\)-value may indicate that ``there may be something
there'', even if it is not significant at significance level 0.05.

In practice we will not spend too much time on hypothesis testing alone
in this course, but instead focus more on regression analysis where
hypothesis testing is an important component. In fact, as you will see,
regression analysis can be used for hypothesis testing.

\subsection{Testing for significance}\label{testing-for-significance}

So how about MP and C? Is the percentage of people saying they would
vote for MP over 4\% with statistical significance? Let's perform a
hypothesis test. Note that we will be using \emph{proportions} and not
percentages, but it is a simple conversion, proportion = percentage /
100. Formulated in proportions our hypotheses are

\begin{quote}
\textbf{\(H_0\)}: The proportion of people who would vote for MP is
\emph{smaller} than 0.04.
\end{quote}

\begin{quote}
\textbf{\(H_a\)}: The proportion of people who would vote for MP is
\emph{larger} than 0.04.
\end{quote}

\begin{Shaded}
\begin{Highlighting}[]
\NormalTok{n\_participants }\OtherTok{\textless{}{-}} \DecValTok{3072} \CommentTok{\# Number of participants in survey}
\NormalTok{n\_mp\_voters }\OtherTok{\textless{}{-}}\NormalTok{ valjarbarometern }\SpecialCharTok{|\textgreater{}} \CommentTok{\# Number of participants who answered MP}
  \FunctionTok{filter}\NormalTok{(party }\SpecialCharTok{==} \StringTok{"MP"}\NormalTok{) }\SpecialCharTok{|\textgreater{}}
  \FunctionTok{summarize}\NormalTok{(}\AttributeTok{n =} \FunctionTok{n}\NormalTok{()) }\SpecialCharTok{|\textgreater{}}
  \FunctionTok{pull}\NormalTok{(n)}
\NormalTok{p\_h0 }\OtherTok{\textless{}{-}} \FloatTok{0.04} \CommentTok{\# Our null hypothesis}
\CommentTok{\# Conducting hypothesis test of proportions. Observe that we are testing if the proportion is *greater* than.}
\NormalTok{result }\OtherTok{\textless{}{-}} \FunctionTok{prop.test}\NormalTok{(n\_mp\_voters, n\_participants, p\_h0, }\AttributeTok{alternative =} \StringTok{"greater"}\NormalTok{)}
\NormalTok{result}
\end{Highlighting}
\end{Shaded}

\begin{verbatim}

    1-sample proportions test with continuity correction

data:  n_mp_voters out of n_participants, null probability p_h0
X-squared = 0.62989, df = 1, p-value = 0.2137
alternative hypothesis: true p is greater than 0.04
95 percent confidence interval:
 0.03719032 1.00000000
sample estimates:
         p 
0.04296875 
\end{verbatim}

The result of our test contains a lot of information, but we only need
to focus on the part that says ``p-value'' on the second line and read
its value, rounded 0.21. Note that we should \emph{not} read the p at
the second to last line, as it is the proportion we observed in our
sample. We compare the \(p\)-value with our significance level, 0.05 and
conclude that 0.21 is \emph{larger} than 0.05, and thus we \emph{can
not} reject our null hypothesis. This means we \emph{do not} have
evidence in the data that the percentage (proportion) of Sweden's
population that would vote for MP is larger than 4\% (0.04), i.e.~we
cannot reject the null hypothesis. So the conclusion of our test is that
if there was an election today MP would not be ``statistically
significantly'' safe.

We can also access the \(p\)-value directly by accessing the
\texttt{p.value} using the dollar operator, \texttt{\$}, from our
result.

\begin{Shaded}
\begin{Highlighting}[]
\NormalTok{result}\SpecialCharTok{$}\NormalTok{p.value}
\end{Highlighting}
\end{Shaded}

\begin{verbatim}
[1] 0.2136985
\end{verbatim}

\subsection{Exercises}\label{exercises-9}

\begin{enumerate}
\def\labelenumi{\arabic{enumi}.}
\tightlist
\item
  How would the hypotheses change if we were testing if the percentage
  that would vote on the C party is less than 4\%?
\item
  Perform a hypothesis test to see if the percentage that would vote on
  C is statistically significant \emph{less} than 4\%. Reuse the code
  from the chapter!
\item
  Find a party that is below the 4\% limit with statistical significance
  at significance level 0.05. Try yourself or expand the code below to
  see the answer. Reuse the code from the chapter!
\end{enumerate}

\section{Comments on hypothesis
testing}\label{comments-on-hypothesis-testing}

A few important things to note in general about hypothesis testing:

\begin{itemize}
\tightlist
\item
  We look for evidence \emph{against} the null hypothesis \(H_0\), but
  failing to reject the null hypothesis doesn't prove the null
  hypothesis is true. It only suggests that there isn't enough
  statistical evidence against it.
\item
  A significantly small p-value doesn't mean the difference between the
  mean under \(H_0\) and the actual population mean is large, only that
  it's likely not due to chance.
\item
  The significance level should be set \emph{beforehand}, not changed
  when the \(p\)-value is observed. Usually 0.05 is used in scientific
  analysis.
\item
  The significance level 0.05 is a strong indicator to reject the null
  hypothesis, but sometimes it may be too black and white. A \(p\)-value
  of 0.06 or 0.1 still indicates an interesting result, while a
  \(p\)-value of 0.9 does not. The \(p\)-value offers the opportunity
  for the reader to decide by themselves if they find the result
  interesting and worth further studies.
\item
  In practical applications, the outcome of hypothesis testing should be
  combined with domain knowledge and other information to make informed
  decisions or answer research questions.
\end{itemize}

When performing hypothesis testing two different errors can be committed
and if you want to further study hypothesis testing and/or statistics it
is important to know them.

\begin{itemize}
\tightlist
\item
  Type I Error: When you incorrectly reject a true null hypothesis. The
  probability of committing this error is \(\alpha\), the significance
  level.
\item
  Type II Error: When you fail to reject a false null hypothesis. The
  probability of committing this error is denoted by \(\beta\). The
  power of a test, \(1 - \beta\), is the probability of rejecting a
  false null hypothesis.
\end{itemize}

\section{Hypothesis test (mean)}\label{hypothesis-test-mean}

Up until now we have studied hypothesis test of proportions. A perhaps
more common test is hypothesis test of the mean and difference in means.
For this, we can look the \texttt{penguins} dataset again and
specifically the variable \texttt{body\_mass\_g}, the weight of the
penguins. Before we start looking at the variable we should think about
our sample and population. The data was collected during the years
2007-2009 in the Palmer Archipelago, Antarctica. If we believe that the
dataset is still valid for penguins in the Palmer archipelago today, we
can say that our population is all penguins of the species Adelie,
Chinstrap, and Gentoo in the Palmer archipelago. Now when we have
identified the population we are studying, we can start by looking at
the variable. We will start by visualizing the distribution over body
mass and the mean body mass marked by a dashed line.

\begin{Shaded}
\begin{Highlighting}[]
\FunctionTok{library}\NormalTok{(palmerpenguins)}
\NormalTok{penguins }\SpecialCharTok{|\textgreater{}} 
  \FunctionTok{ggplot}\NormalTok{(}\FunctionTok{aes}\NormalTok{(}\AttributeTok{x =}\NormalTok{ body\_mass\_g)) }\SpecialCharTok{+} 
  \FunctionTok{geom\_density}\NormalTok{() }\SpecialCharTok{+} 
  \FunctionTok{geom\_vline}\NormalTok{(}
    \AttributeTok{xintercept =} \FunctionTok{mean}\NormalTok{(penguins}\SpecialCharTok{$}\NormalTok{body\_mass\_g, }\AttributeTok{na.rm =} \ConstantTok{TRUE}\NormalTok{),}
    \AttributeTok{linetype =} \DecValTok{2}\NormalTok{) }\SpecialCharTok{+}
  \FunctionTok{labs}\NormalTok{(}
    \AttributeTok{x =} \StringTok{"Body mass (g)"}\NormalTok{,}
    \AttributeTok{y =} \StringTok{""}\NormalTok{) }\SpecialCharTok{+} 
  \FunctionTok{theme}\NormalTok{(}
    \AttributeTok{axis.text.y =} \FunctionTok{element\_blank}\NormalTok{(),}
    \AttributeTok{axis.ticks.y =} \FunctionTok{element\_blank}\NormalTok{())}
\end{Highlighting}
\end{Shaded}

\includegraphics{hypothesis_testing_files/figure-pdf/unnamed-chunk-7-1.pdf}

We see that there are quite a lot of observations with a smaller body
mass than the mean. Can we be sure that the population mean, i.e.~the
mean of all penguins body weights, including those outside our sample,
is above 4000 g? To test this we use a t-test. The t-test can be used we
we want to ask questions about the population mean based on our sample.
We formulate our hypothesis

\begin{quote}
\textbf{\(H_0\)}: The average body mass of all penguins in the Palmer
Archipelago is 4000 g.
\end{quote}

\begin{quote}
\textbf{\(H_a\)}: The average body mass of all penguins in the Palmer
Archipelago is \emph{more} than 4000 g.
\end{quote}

We can perform this test in R using the \texttt{t.test} function.

\begin{Shaded}
\begin{Highlighting}[]
\FunctionTok{t.test}\NormalTok{(penguins}\SpecialCharTok{$}\NormalTok{body\_mass\_g, }\AttributeTok{mu =} \DecValTok{4000}\NormalTok{, }\AttributeTok{alternative =} \StringTok{"greater"}\NormalTok{)}
\end{Highlighting}
\end{Shaded}

\begin{verbatim}

    One Sample t-test

data:  penguins$body_mass_g
t = 4.6525, df = 341, p-value = 2.35e-06
alternative hypothesis: true mean is greater than 4000
95 percent confidence interval:
 4130.231      Inf
sample estimates:
mean of x 
 4201.754 
\end{verbatim}

Again we are only interested in the \(p\)-value, and if it is smaller
than 0.05. We see that ``p-value = 2.35e-06''. Numbers with an ``e'' is
represented using \emph{scientific notation}, and it means
\(2.35 \cdot 10^{-06} = 0.000000235\), which is a very, \emph{very}
small number. In general the number after e can be understood as the
number of zeros before the first significant digit, e.g.~1.2e-02 =
0.0012. So we can conclude that at significance level 0.05 we have
evidence that the mean penguin weight is more than 4000 since our
\(p\)-value is smaller than 0.05.

\section{Difference in means}\label{difference-in-means}

Using your skills from EDA (or memory of the Data visualization chapter
in R4DS), you may have noticed that the density plot looked a bit
strange: it looks like there are several densities in our sample,
noticeable by the humps at around 3600, 4700, and 5500 grams. Let's
explore this further, and a good place to start is one of the
categorical variables that may indicate different groups. We can
visualize the density again, this time coloring the fill after the
different species.

\begin{Shaded}
\begin{Highlighting}[]
\CommentTok{\# warning: false}
\NormalTok{penguins }\SpecialCharTok{|\textgreater{}}
  \FunctionTok{drop\_na}\NormalTok{() }\SpecialCharTok{|\textgreater{}}
  \FunctionTok{ggplot}\NormalTok{(}\FunctionTok{aes}\NormalTok{(}\AttributeTok{x =}\NormalTok{ body\_mass\_g, }\AttributeTok{fill =}\NormalTok{ species)) }\SpecialCharTok{+}
  \FunctionTok{geom\_density}\NormalTok{(}\AttributeTok{alpha =} \FloatTok{0.3}\NormalTok{) }\SpecialCharTok{+}
  \FunctionTok{labs}\NormalTok{(}
    \AttributeTok{x =} \StringTok{"Body mass (g)"}\NormalTok{,}
    \AttributeTok{y =} \StringTok{""}\NormalTok{,}
    \AttributeTok{fill =} \StringTok{"Species"}\NormalTok{) }\SpecialCharTok{+} 
  \FunctionTok{theme}\NormalTok{(}
    \AttributeTok{axis.text.y =} \FunctionTok{element\_blank}\NormalTok{(),}
    \AttributeTok{axis.ticks.y =} \FunctionTok{element\_blank}\NormalTok{())}
\end{Highlighting}
\end{Shaded}

\includegraphics{hypothesis_testing_files/figure-pdf/unnamed-chunk-9-1.pdf}

That was revealing! But still there seems to be something going
on\ldots{} If we look at the distribution over Gentoo penguins, we see
two distinct humps again. Another categorical variable that we may
suspect will affect the body mass is sex. We take a closer look at only
the Gentoo penguins and color the fill based on sex instead.

\begin{Shaded}
\begin{Highlighting}[]
\FunctionTok{library}\NormalTok{(palmerpenguins)}
\NormalTok{penguins }\SpecialCharTok{|\textgreater{}}
  \FunctionTok{drop\_na}\NormalTok{() }\SpecialCharTok{|\textgreater{}}
  \FunctionTok{filter}\NormalTok{(species }\SpecialCharTok{==} \StringTok{"Gentoo"}\NormalTok{) }\SpecialCharTok{|\textgreater{}}
  \FunctionTok{ggplot}\NormalTok{(}\FunctionTok{aes}\NormalTok{(}\AttributeTok{x =}\NormalTok{ body\_mass\_g, }\AttributeTok{fill =}\NormalTok{ sex)) }\SpecialCharTok{+}
  \FunctionTok{geom\_density}\NormalTok{(}\AttributeTok{alpha =} \FloatTok{0.3}\NormalTok{) }\SpecialCharTok{+}
  \FunctionTok{labs}\NormalTok{(}
    \AttributeTok{x =} \StringTok{"Body mass (g)"}\NormalTok{,}
    \AttributeTok{y =} \StringTok{""}\NormalTok{,}
    \AttributeTok{fill =} \StringTok{"Sex"}\NormalTok{) }\SpecialCharTok{+} 
  \FunctionTok{theme}\NormalTok{(}
    \AttributeTok{axis.text.y =} \FunctionTok{element\_blank}\NormalTok{(),}
    \AttributeTok{axis.ticks.y =} \FunctionTok{element\_blank}\NormalTok{())}
\end{Highlighting}
\end{Shaded}

\includegraphics{hypothesis_testing_files/figure-pdf/unnamed-chunk-10-1.pdf}

We see a clear difference, but will this difference hold in the
population of Gentoo penguins? How big is our sample of Gentoo penguins?

\begin{Shaded}
\begin{Highlighting}[]
\NormalTok{penguins }\SpecialCharTok{|\textgreater{}}
  \FunctionTok{filter}\NormalTok{(species }\SpecialCharTok{==} \StringTok{"Gentoo"}\NormalTok{) }\SpecialCharTok{|\textgreater{}}
  \FunctionTok{drop\_na}\NormalTok{() }\SpecialCharTok{|\textgreater{}}
  \FunctionTok{summarize}\NormalTok{(}
    \AttributeTok{count =} \FunctionTok{n}\NormalTok{(),}
    \AttributeTok{.by =}\NormalTok{ sex)}
\end{Highlighting}
\end{Shaded}

\begin{verbatim}
# A tibble: 2 x 2
  sex    count
  <fct>  <int>
1 female    58
2 male      61
\end{verbatim}

We have 58 female Gentoos and 61 males, it is not that few, but
according to
\href{https://en.wikipedia.org/wiki/Gentoo_penguin}{Wikipedia} the
population of Gentoos is more than 600,000 birds! Is this difference
statistically significant? Again we can use a t-test to compare the mean
of two groups. We formulate our hypothesis

\begin{quote}
\textbf{\(H_0\)}: Female and male Gentoos have the same body mass on
average.
\end{quote}

\begin{quote}
\textbf{\(H_a\)}: Female and male Gentoos \emph{do not} have the same
body mass on average.
\end{quote}

We use \texttt{t.test} again to test for statistical significance.

\begin{Shaded}
\begin{Highlighting}[]
\NormalTok{female\_gentoo\_body\_mass }\OtherTok{\textless{}{-}}\NormalTok{ penguins }\SpecialCharTok{|\textgreater{}}
  \FunctionTok{filter}\NormalTok{(species }\SpecialCharTok{==} \StringTok{"Gentoo"} \SpecialCharTok{\&}\NormalTok{ sex }\SpecialCharTok{==} \StringTok{"female"}\NormalTok{) }\SpecialCharTok{|\textgreater{}}
  \FunctionTok{select}\NormalTok{(body\_mass\_g) }\SpecialCharTok{|\textgreater{}}
  \FunctionTok{drop\_na}\NormalTok{()}
\NormalTok{male\_gentoo\_body\_mass }\OtherTok{\textless{}{-}}\NormalTok{ penguins }\SpecialCharTok{|\textgreater{}}
  \FunctionTok{filter}\NormalTok{(species }\SpecialCharTok{==} \StringTok{"Gentoo"} \SpecialCharTok{\&}\NormalTok{ sex }\SpecialCharTok{==} \StringTok{"male"}\NormalTok{) }\SpecialCharTok{|\textgreater{}}
  \FunctionTok{select}\NormalTok{(body\_mass\_g) }\SpecialCharTok{|\textgreater{}}
  \FunctionTok{drop\_na}\NormalTok{()}

\FunctionTok{t.test}\NormalTok{(male\_gentoo\_body\_mass, female\_gentoo\_body\_mass)}
\end{Highlighting}
\end{Shaded}

\begin{verbatim}

    Welch Two Sample t-test

data:  male_gentoo_body_mass and female_gentoo_body_mass
t = 14.761, df = 116.64, p-value < 2.2e-16
alternative hypothesis: true difference in means is not equal to 0
95 percent confidence interval:
 697.0763 913.1130
sample estimates:
mean of x mean of y 
 5484.836  4679.741 
\end{verbatim}

We see that the \(p\)-value is very small, \(2.2 \cdot 10^{-16}\). We
conclude that we have evidence that the difference in average body
weight of female and male Gentoo penguins is statistically significant
at significance level 0.05.

Now we have looked at two different ways of performing hypothesis
testing about different variables, either proportions or means. In the
next chapter we will look at \emph{linear regression}, a method that we
can use to understand if there is an effect of one variable on another
and to make powerful predictions.

\subsection{Exercises}\label{exercises-10}

\begin{enumerate}
\def\labelenumi{\arabic{enumi}.}
\tightlist
\item
  Test if the mean flipper length is more than 180 mm. Try to write down
  the hypotheses before you perform your test. Is the result
  statistically significant?
\item
  Test if the mean flipper length differs between species. Formulate the
  hypothesis and draw conclusions from the tests. Based on your result,
  is flipper length a good way to distinguish between different penguins
  species? You can use the code below to get started, or try yourself
  first.
\end{enumerate}

\begin{Shaded}
\begin{Highlighting}[]
\NormalTok{gentoo\_flipper\_length }\OtherTok{\textless{}{-}}\NormalTok{ penguins }\SpecialCharTok{|\textgreater{}}
  \FunctionTok{filter}\NormalTok{(species }\SpecialCharTok{==} \StringTok{"Gentoo"}\NormalTok{) }\SpecialCharTok{|\textgreater{}}
  \FunctionTok{select}\NormalTok{(flipper\_length\_mm) }\SpecialCharTok{|\textgreater{}}
  \FunctionTok{drop\_na}\NormalTok{()}
\NormalTok{chinstrap\_flipper\_length }\OtherTok{\textless{}{-}}\NormalTok{ penguins }\SpecialCharTok{|\textgreater{}}
  \FunctionTok{filter}\NormalTok{(species }\SpecialCharTok{==} \StringTok{"Chinstrap"}\NormalTok{) }\SpecialCharTok{|\textgreater{}}
  \FunctionTok{select}\NormalTok{(flipper\_length\_mm) }\SpecialCharTok{|\textgreater{}}
  \FunctionTok{drop\_na}\NormalTok{()}
\NormalTok{adelie\_flipper\_length }\OtherTok{\textless{}{-}}\NormalTok{ penguins }\SpecialCharTok{|\textgreater{}}
  \FunctionTok{filter}\NormalTok{(species }\SpecialCharTok{==} \StringTok{"Adelie"}\NormalTok{) }\SpecialCharTok{|\textgreater{}}
  \FunctionTok{select}\NormalTok{(flipper\_length\_mm) }\SpecialCharTok{|\textgreater{}}
  \FunctionTok{drop\_na}\NormalTok{()}
\end{Highlighting}
\end{Shaded}

\begin{enumerate}
\def\labelenumi{\arabic{enumi}.}
\setcounter{enumi}{2}
\tightlist
\item
  Try to figure out how to perform a proportion test to study if there
  are more Chinstrap penguins than Adelie on Dream island. Start by
  visualizing the distribution using some visualization of Island and
  Species. Then you will need to count the number of Chinstrap penguins
  on Dream island as well as the total number of penguins on Dream
  island. Do you see any problem with the sample we are using for this
  test? You can use the code below as a hint.
\end{enumerate}

\begin{Shaded}
\begin{Highlighting}[]
\CommentTok{\# Visualization}
\NormalTok{penguins }\SpecialCharTok{|\textgreater{}} 
  \FunctionTok{ggplot}\NormalTok{(}\FunctionTok{aes}\NormalTok{(}\AttributeTok{x =}\NormalTok{ island, }\AttributeTok{fill =}\NormalTok{ species)) }\SpecialCharTok{+} 
  \FunctionTok{geom\_bar}\NormalTok{()}
\end{Highlighting}
\end{Shaded}

\includegraphics{hypothesis_testing_files/figure-pdf/unnamed-chunk-14-1.pdf}

\begin{Shaded}
\begin{Highlighting}[]
\CommentTok{\# Total number of penguins on Dream island}
\NormalTok{n\_total }\OtherTok{\textless{}{-}}\NormalTok{ penguins }\SpecialCharTok{|\textgreater{}} 
  \FunctionTok{filter}\NormalTok{(island }\SpecialCharTok{==} \StringTok{"Dream"}\NormalTok{) }\SpecialCharTok{|\textgreater{}}
  \FunctionTok{summarize}\NormalTok{(}\AttributeTok{n =} \FunctionTok{n}\NormalTok{()) }\SpecialCharTok{|\textgreater{}} 
  \FunctionTok{pull}\NormalTok{()}
\CommentTok{\# Number of chinstrap penguins on Dream island}
\NormalTok{n\_chinstrap }\OtherTok{\textless{}{-}}\NormalTok{ penguins }\SpecialCharTok{|\textgreater{}} 
  \FunctionTok{filter}\NormalTok{(species }\SpecialCharTok{==} \StringTok{"Chinstrap"} \SpecialCharTok{\&}\NormalTok{ island }\SpecialCharTok{==} \StringTok{"Dream"}\NormalTok{) }\SpecialCharTok{|\textgreater{}}
  \FunctionTok{summarize}\NormalTok{(}\AttributeTok{n =} \FunctionTok{n}\NormalTok{()) }\SpecialCharTok{|\textgreater{}} 
  \FunctionTok{pull}\NormalTok{()}
\end{Highlighting}
\end{Shaded}

\chapter{Linear regression}\label{sec-simple-linear-regression}

In the previous chapter we studied hypothesis testing and learned how we
can answer questions about the population using our data, at a
significance level. This is very useful to understand different
quantities in the population, but what if we want to find out how one
variable depends on another? In our previous example we saw that we
could understand the distribution of Antarctic penguins body masses by
first grouping by species and then, within each species, grouping by
sex. This is possible because species and sex are \emph{categorical}
variables, but we cannot use the same method if we want to know the
effect of \emph{numerical} variables. For this, we can use linear
regression, a statistical method used to model the (linear) relationship
between two continuous variables: a single predictor (independent)
variable and a response (dependent) variable. When we use linear
regression we model the relationship between two variables using a line,
characterized by the linear equation \(y = k \cdot x + m\). Independent
variables and dependent variables have several names in different fields
(and within fields!), and we will be using a few of them. For
independent variables we may say explanatory variables, predictors,
covariates, or features and for dependent variable we may say response
or outcome. It may cause some confusion in the beginning, but once you
understand that they are all the same it will help you understand data
analysis in different fields and from different sources.

For a more technical/mathematical introduction to linear regression,
please consult the chapter on linear regression (Chapter 3, 2nd edition)
in \href{https://www.statlearning.com/}{An Introduction to Statistical
Learning}.

\section{A linear relationship}\label{a-linear-relationship}

In this chapter we will continue to look at \texttt{body\_mass\_g} from
the \texttt{penguins} dataset. In the previous chapter we studied the
body weight with respect to species and sex, but our analysis was
limited to categorical variables. This time we will expand our analysis
to the relationship between body mass and other \emph{numerical}
variables. This may be useful if we can \emph{predict} a penguins weight
based on other body measures. It may be easier measuring the bill and
flippers on a penguin than getting it to stand on a scale(?).

We start by visualizing the relationship between flipper length and body
mass using a scatter plot (\texttt{geom\_point}).

\begin{Shaded}
\begin{Highlighting}[]
\FunctionTok{library}\NormalTok{(tidyverse)}
\FunctionTok{library}\NormalTok{(palmerpenguins)}

\NormalTok{penguins }\SpecialCharTok{|\textgreater{}} 
  \FunctionTok{ggplot}\NormalTok{(}\FunctionTok{aes}\NormalTok{(}\AttributeTok{x =}\NormalTok{ flipper\_length\_mm, }\AttributeTok{y =}\NormalTok{ body\_mass\_g)) }\SpecialCharTok{+} 
  \FunctionTok{geom\_point}\NormalTok{() }\SpecialCharTok{+}
  \FunctionTok{labs}\NormalTok{(}
    \AttributeTok{x =} \StringTok{"Flipper length (mm)"}\NormalTok{, }
    \AttributeTok{y =} \StringTok{"Body mass (g)"}\NormalTok{)}
\end{Highlighting}
\end{Shaded}

\includegraphics{linear_regression_files/figure-pdf/unnamed-chunk-1-1.pdf}

We can clearly see that the body mass increases with flipper length:
penguins with longer flippers have larger body mass, \emph{on average}.
It is important to note that this relationship holds on average and not
for every penguin, as then every point as we take a step to the right
would be higher up, but this is not the case. We illustrate what this
means below as we draw two vertical lines, creating three bins, then
calculate the mean body weight within the interval and plot it as red,
horizontal lines.

\begin{Shaded}
\begin{Highlighting}[]
\NormalTok{mu\_190 }\OtherTok{\textless{}{-}} \FunctionTok{mean}\NormalTok{(penguins }\SpecialCharTok{|\textgreater{}} \FunctionTok{filter}\NormalTok{(flipper\_length\_mm }\SpecialCharTok{\textless{}} \DecValTok{190}\NormalTok{) }\SpecialCharTok{|\textgreater{}} \FunctionTok{pull}\NormalTok{(body\_mass\_g))}
\NormalTok{mu\_210 }\OtherTok{\textless{}{-}} \FunctionTok{mean}\NormalTok{(penguins }\SpecialCharTok{|\textgreater{}} \FunctionTok{filter}\NormalTok{(flipper\_length\_mm }\SpecialCharTok{\textgreater{}} \DecValTok{190} \SpecialCharTok{\&}\NormalTok{ flipper\_length\_mm }\SpecialCharTok{\textless{}} \DecValTok{210}\NormalTok{) }\SpecialCharTok{|\textgreater{}} \FunctionTok{pull}\NormalTok{(body\_mass\_g))}
\NormalTok{mu\_230 }\OtherTok{\textless{}{-}} \FunctionTok{mean}\NormalTok{(penguins }\SpecialCharTok{|\textgreater{}} \FunctionTok{filter}\NormalTok{(flipper\_length\_mm }\SpecialCharTok{\textgreater{}} \DecValTok{210}\NormalTok{) }\SpecialCharTok{|\textgreater{}} \FunctionTok{pull}\NormalTok{(body\_mass\_g))}

\NormalTok{penguins }\SpecialCharTok{|\textgreater{}} 
  \FunctionTok{ggplot}\NormalTok{(}\FunctionTok{aes}\NormalTok{(}\AttributeTok{x =}\NormalTok{ flipper\_length\_mm, }\AttributeTok{y =}\NormalTok{ body\_mass\_g)) }\SpecialCharTok{+} 
  \FunctionTok{geom\_point}\NormalTok{() }\SpecialCharTok{+} 
  \FunctionTok{geom\_vline}\NormalTok{(}\AttributeTok{xintercept =} \DecValTok{190}\NormalTok{, }\AttributeTok{col =} \StringTok{"black"}\NormalTok{, }\AttributeTok{linetype =} \DecValTok{3}\NormalTok{) }\SpecialCharTok{+}
  \FunctionTok{geom\_vline}\NormalTok{(}\AttributeTok{xintercept =} \DecValTok{210}\NormalTok{, }\AttributeTok{col =} \StringTok{"black"}\NormalTok{, }\AttributeTok{linetype =} \DecValTok{3}\NormalTok{) }\SpecialCharTok{+}
  \FunctionTok{geom\_segment}\NormalTok{(}\AttributeTok{x =} \DecValTok{170}\NormalTok{, }\AttributeTok{xend =} \DecValTok{190}\NormalTok{,  }\AttributeTok{y =}\NormalTok{ mu\_190, }\AttributeTok{yend =}\NormalTok{ mu\_190, }\AttributeTok{color =} \StringTok{"red"}\NormalTok{) }\SpecialCharTok{+}
  \FunctionTok{geom\_segment}\NormalTok{(}\AttributeTok{x =} \DecValTok{190}\NormalTok{, }\AttributeTok{xend =} \DecValTok{210}\NormalTok{,  }\AttributeTok{y =}\NormalTok{ mu\_210, }\AttributeTok{yend =}\NormalTok{ mu\_210, }\AttributeTok{color =} \StringTok{"red"}\NormalTok{) }\SpecialCharTok{+}
  \FunctionTok{geom\_segment}\NormalTok{(}\AttributeTok{x =} \DecValTok{210}\NormalTok{, }\AttributeTok{xend =} \DecValTok{230}\NormalTok{,  }\AttributeTok{y =}\NormalTok{ mu\_230, }\AttributeTok{yend =}\NormalTok{ mu\_230, }\AttributeTok{color =} \StringTok{"red"}\NormalTok{) }\SpecialCharTok{+}
  \FunctionTok{labs}\NormalTok{(}
    \AttributeTok{x =} \StringTok{"Flipper length (mm)"}\NormalTok{, }
    \AttributeTok{y =} \StringTok{"Body mass (g)"}\NormalTok{)}
\end{Highlighting}
\end{Shaded}

\includegraphics{linear_regression_files/figure-pdf/unnamed-chunk-2-1.pdf}

Here we can clearly see that the average of our dependent variable, body
mass, increases as our independent variable, flipper length, increases.
We could follow the approach of the previous chapter and divide the
penguins into three groups and perform t-tests between the groups to see
if the difference is statistically significant.

\begin{Shaded}
\begin{Highlighting}[]
\NormalTok{flippers\_190 }\OtherTok{\textless{}{-}}\NormalTok{ penguins }\SpecialCharTok{|\textgreater{}} 
  \FunctionTok{filter}\NormalTok{(flipper\_length\_mm }\SpecialCharTok{\textless{}} \DecValTok{190}\NormalTok{) }\SpecialCharTok{|\textgreater{}} 
  \FunctionTok{pull}\NormalTok{(body\_mass\_g)}
\NormalTok{flippers\_210 }\OtherTok{\textless{}{-}}\NormalTok{ penguins }\SpecialCharTok{|\textgreater{}} 
  \FunctionTok{filter}\NormalTok{(flipper\_length\_mm }\SpecialCharTok{\textgreater{}} \DecValTok{190} \SpecialCharTok{\&}\NormalTok{ flipper\_length\_mm }\SpecialCharTok{\textless{}} \DecValTok{210}\NormalTok{) }\SpecialCharTok{|\textgreater{}} 
  \FunctionTok{pull}\NormalTok{(body\_mass\_g)}
\NormalTok{flippers\_230 }\OtherTok{\textless{}{-}}\NormalTok{ penguins }\SpecialCharTok{|\textgreater{}} 
  \FunctionTok{filter}\NormalTok{(flipper\_length\_mm }\SpecialCharTok{\textgreater{}} \DecValTok{210}\NormalTok{) }\SpecialCharTok{|\textgreater{}} 
  \FunctionTok{pull}\NormalTok{(body\_mass\_g)}

\FunctionTok{t.test}\NormalTok{(flippers\_190, flippers\_210)}\SpecialCharTok{$}\NormalTok{p.value}
\end{Highlighting}
\end{Shaded}

\begin{verbatim}
[1] 8.653392e-13
\end{verbatim}

\begin{Shaded}
\begin{Highlighting}[]
\FunctionTok{t.test}\NormalTok{(flippers\_210, flippers\_230)}\SpecialCharTok{$}\NormalTok{p.value}
\end{Highlighting}
\end{Shaded}

\begin{verbatim}
[1] 4.958129e-52
\end{verbatim}

We see that both \(p\)-values are significant at significance level
0.05. This is very useful information, as if we know that a penguin
falls into one of these bins we can assume it's body mass will be close
to average within that bin. If we make more bins we can make our
predictions more accurate, and maybe we can even say something about
\emph{how much} the body mass increases, on average, with the flipper
length. This is exactly the goal of linear regression, where we create a
line that represents one bin for every possible value for
\texttt{flipper\_length\_mm}. You have already seen linear regressions
before, graphically, using the \texttt{geom\_smooth} function with
parameter \texttt{method\ =\ lm}. \texttt{lm} stands for ``Linear
Model'', and is the name of the R function used to perform linear
regression. First, we visualize the line, then we will calculate our
linear model.

\begin{Shaded}
\begin{Highlighting}[]
\NormalTok{penguins }\SpecialCharTok{|\textgreater{}} 
  \FunctionTok{ggplot}\NormalTok{(}\FunctionTok{aes}\NormalTok{(}\AttributeTok{x =}\NormalTok{ flipper\_length\_mm, }\AttributeTok{y =}\NormalTok{ body\_mass\_g)) }\SpecialCharTok{+} 
  \FunctionTok{geom\_point}\NormalTok{() }\SpecialCharTok{+} 
  \FunctionTok{geom\_smooth}\NormalTok{(}\AttributeTok{method =} \StringTok{"lm"}\NormalTok{, }\AttributeTok{se =} \ConstantTok{FALSE}\NormalTok{) }\SpecialCharTok{+}
  \FunctionTok{labs}\NormalTok{(}
    \AttributeTok{x =} \StringTok{"Flipper length (mm)"}\NormalTok{, }
    \AttributeTok{y =} \StringTok{"Body mass (g)"}\NormalTok{)}
\end{Highlighting}
\end{Shaded}

\includegraphics{linear_regression_files/figure-pdf/unnamed-chunk-4-1.pdf}

The line you see describes how \texttt{body\_mass\_g} increases with
\texttt{flipper\_length\_mm}, on average. If you want to know the
average body mass for a specific flipper length you just have to follow
a vertical line from that value on the x axis to the line, as in the
example below where we use the line to get the average body mass of a
penguin with flipper length 200 mm.

\begin{Shaded}
\begin{Highlighting}[]
\NormalTok{x }\OtherTok{\textless{}{-}} \DecValTok{200}
\NormalTok{model }\OtherTok{\textless{}{-}}\NormalTok{ penguins }\SpecialCharTok{|\textgreater{}} 
  \FunctionTok{lm}\NormalTok{(}\AttributeTok{formula =}\NormalTok{ body\_mass\_g }\SpecialCharTok{\textasciitilde{}}\NormalTok{ flipper\_length\_mm)}
\NormalTok{prediction }\OtherTok{\textless{}{-}} \FunctionTok{data.frame}\NormalTok{(}
  \AttributeTok{flipper\_length\_mm =}\NormalTok{ x,}
  \AttributeTok{body\_mass\_g =}\NormalTok{ model }\SpecialCharTok{|\textgreater{}} \FunctionTok{predict}\NormalTok{(}\FunctionTok{data.frame}\NormalTok{(}\AttributeTok{flipper\_length\_mm =}\NormalTok{ x)))}

\NormalTok{penguins }\SpecialCharTok{|\textgreater{}} 
  \FunctionTok{ggplot}\NormalTok{(}\FunctionTok{aes}\NormalTok{(}\AttributeTok{x =}\NormalTok{ flipper\_length\_mm, }\AttributeTok{y =}\NormalTok{ body\_mass\_g)) }\SpecialCharTok{+} 
  \FunctionTok{geom\_point}\NormalTok{() }\SpecialCharTok{+} 
  \FunctionTok{geom\_smooth}\NormalTok{(}\AttributeTok{method =} \StringTok{"lm"}\NormalTok{, }\AttributeTok{se =} \ConstantTok{FALSE}\NormalTok{) }\SpecialCharTok{+} 
  \FunctionTok{geom\_segment}\NormalTok{(}\AttributeTok{x =}\NormalTok{ prediction}\SpecialCharTok{$}\NormalTok{flipper\_length\_mm, }\AttributeTok{xend =}\NormalTok{ prediction}\SpecialCharTok{$}\NormalTok{flipper\_length\_mm,  }\AttributeTok{y =} \DecValTok{0}\NormalTok{, }\AttributeTok{yend =}\NormalTok{ prediction}\SpecialCharTok{$}\NormalTok{body\_mass\_g, }\AttributeTok{color =} \StringTok{"black"}\NormalTok{, }\AttributeTok{linetype =} \DecValTok{2}\NormalTok{, }\AttributeTok{linewidth =} \FloatTok{0.2}\NormalTok{) }\SpecialCharTok{+}
  \FunctionTok{geom\_segment}\NormalTok{(}\AttributeTok{x =} \DecValTok{0}\NormalTok{, }\AttributeTok{xend =}\NormalTok{ prediction}\SpecialCharTok{$}\NormalTok{flipper\_length\_mm,  }\AttributeTok{y =}\NormalTok{ prediction}\SpecialCharTok{$}\NormalTok{body\_mass\_g, }\AttributeTok{yend =}\NormalTok{ prediction}\SpecialCharTok{$}\NormalTok{body\_mass\_g, }\AttributeTok{color =} \StringTok{"black"}\NormalTok{, }\AttributeTok{linetype =} \DecValTok{2}\NormalTok{, }\AttributeTok{linewidth =} \FloatTok{0.2}\NormalTok{) }\SpecialCharTok{+}
  \FunctionTok{scale\_x\_continuous}\NormalTok{(}\AttributeTok{breaks =} \FunctionTok{c}\NormalTok{(prediction}\SpecialCharTok{$}\NormalTok{flipper\_length\_mm)) }\SpecialCharTok{+}
  \FunctionTok{scale\_y\_continuous}\NormalTok{(}\AttributeTok{breaks =} \FunctionTok{round}\NormalTok{(}\FunctionTok{c}\NormalTok{(prediction}\SpecialCharTok{$}\NormalTok{body\_mass\_g))) }\SpecialCharTok{+}
  \FunctionTok{geom\_point}\NormalTok{(}\AttributeTok{data =}\NormalTok{ prediction, }\AttributeTok{col =} \StringTok{"red"}\NormalTok{, }\AttributeTok{size =} \DecValTok{3}\NormalTok{) }\SpecialCharTok{+} 
  \FunctionTok{labs}\NormalTok{(}
    \AttributeTok{x =} \StringTok{"Flipper length (mm)"}\NormalTok{, }
    \AttributeTok{y =} \StringTok{"Body mass (g)"}\NormalTok{)}
\end{Highlighting}
\end{Shaded}

\includegraphics{linear_regression_files/figure-pdf/unnamed-chunk-5-1.pdf}

\subsection{Exercises}\label{exercises-11}

\begin{enumerate}
\def\labelenumi{\arabic{enumi}.}
\tightlist
\item
  Copy the code below (the same as for the plot above) and change the
  value for \texttt{x} to 201. How much did the body mass increase? Note
  that we use the function \texttt{round} to round the value of body
  mass. This provides nicer output in the plot, but may induce a slight
  error (see next section where the exact numbers are used).
\end{enumerate}

\begin{Shaded}
\begin{Highlighting}[]
\NormalTok{x }\OtherTok{\textless{}{-}} \DecValTok{200}
\NormalTok{model }\OtherTok{\textless{}{-}}\NormalTok{ penguins }\SpecialCharTok{|\textgreater{}} 
  \FunctionTok{lm}\NormalTok{(}\AttributeTok{formula =}\NormalTok{ body\_mass\_g }\SpecialCharTok{\textasciitilde{}}\NormalTok{ flipper\_length\_mm)}
\NormalTok{prediction }\OtherTok{\textless{}{-}} \FunctionTok{data.frame}\NormalTok{(}
  \AttributeTok{flipper\_length\_mm =}\NormalTok{ x,}
  \AttributeTok{body\_mass\_g =}\NormalTok{ model }\SpecialCharTok{|\textgreater{}} \FunctionTok{predict}\NormalTok{(}\FunctionTok{data.frame}\NormalTok{(}\AttributeTok{flipper\_length\_mm =}\NormalTok{ x)))}

\NormalTok{penguins }\SpecialCharTok{|\textgreater{}} 
  \FunctionTok{ggplot}\NormalTok{(}\FunctionTok{aes}\NormalTok{(}\AttributeTok{x =}\NormalTok{ flipper\_length\_mm, }\AttributeTok{y =}\NormalTok{ body\_mass\_g)) }\SpecialCharTok{+} 
  \FunctionTok{geom\_point}\NormalTok{() }\SpecialCharTok{+} 
  \FunctionTok{geom\_smooth}\NormalTok{(}\AttributeTok{method =} \StringTok{"lm"}\NormalTok{, }\AttributeTok{se =} \ConstantTok{FALSE}\NormalTok{) }\SpecialCharTok{+} 
  \FunctionTok{geom\_segment}\NormalTok{(}\AttributeTok{x =}\NormalTok{ prediction}\SpecialCharTok{$}\NormalTok{flipper\_length\_mm, }\AttributeTok{xend =}\NormalTok{ prediction}\SpecialCharTok{$}\NormalTok{flipper\_length\_mm,  }\AttributeTok{y =} \DecValTok{0}\NormalTok{, }\AttributeTok{yend =}\NormalTok{ prediction}\SpecialCharTok{$}\NormalTok{body\_mass\_g, }\AttributeTok{color =} \StringTok{"black"}\NormalTok{, }\AttributeTok{linetype =} \DecValTok{2}\NormalTok{, }\AttributeTok{linewidth =} \FloatTok{0.2}\NormalTok{) }\SpecialCharTok{+}
  \FunctionTok{geom\_segment}\NormalTok{(}\AttributeTok{x =} \DecValTok{0}\NormalTok{, }\AttributeTok{xend =}\NormalTok{ prediction}\SpecialCharTok{$}\NormalTok{flipper\_length\_mm,  }\AttributeTok{y =}\NormalTok{ prediction}\SpecialCharTok{$}\NormalTok{body\_mass\_g, }\AttributeTok{yend =}\NormalTok{ prediction}\SpecialCharTok{$}\NormalTok{body\_mass\_g, }\AttributeTok{color =} \StringTok{"black"}\NormalTok{, }\AttributeTok{linetype =} \DecValTok{2}\NormalTok{, }\AttributeTok{linewidth =} \FloatTok{0.2}\NormalTok{) }\SpecialCharTok{+}
  \FunctionTok{scale\_x\_continuous}\NormalTok{(}\AttributeTok{breaks =} \FunctionTok{c}\NormalTok{(prediction}\SpecialCharTok{$}\NormalTok{flipper\_length\_mm)) }\SpecialCharTok{+}
  \FunctionTok{scale\_y\_continuous}\NormalTok{(}\AttributeTok{breaks =} \FunctionTok{round}\NormalTok{(}\FunctionTok{c}\NormalTok{(prediction}\SpecialCharTok{$}\NormalTok{body\_mass\_g))) }\SpecialCharTok{+}
  \FunctionTok{geom\_point}\NormalTok{(}\AttributeTok{data =}\NormalTok{ prediction, }\AttributeTok{col =} \StringTok{"red"}\NormalTok{, }\AttributeTok{size =} \DecValTok{3}\NormalTok{) }\SpecialCharTok{+} 
  \FunctionTok{labs}\NormalTok{(}
    \AttributeTok{x =} \StringTok{"Flipper length (mm)"}\NormalTok{, }
    \AttributeTok{y =} \StringTok{"Body mass (g)"}\NormalTok{)}
\end{Highlighting}
\end{Shaded}

\begin{enumerate}
\def\labelenumi{\arabic{enumi}.}
\setcounter{enumi}{1}
\tightlist
\item
  Change the \texttt{x}-value to 100. What is the predicted body mass?
  Is this a reasonable result? Why is this happening?
\end{enumerate}

\section{Fitting a linear model}\label{fitting-a-linear-model}

Now when we have seen what a linear model looks like we are ready to
understand the capabilities of this model. Just as before we can ask
ourselves if the linear relationship is statistically significant and
how much the average body weight increases with flipper length. We use
the \texttt{lm} function in R to estimate a linear model. When
estimating a model we need to specify the dependent variable, the
variable we are interested in as our outcome, and the independent
variable, the variable we expect will explain the outcome. We do so by
specifying a \texttt{formula} that follows the syntax
\texttt{formula\ =\ dependent\ \textasciitilde{}\ independent}, so in
our case we use
\texttt{formula\ =\ body\_mass\_g\ \textasciitilde{}\ flipper\_length\_mm}.

\begin{Shaded}
\begin{Highlighting}[]
\NormalTok{model }\OtherTok{\textless{}{-}}\NormalTok{ penguins }\SpecialCharTok{|\textgreater{}} 
  \FunctionTok{lm}\NormalTok{(}\AttributeTok{formula =}\NormalTok{ body\_mass\_g }\SpecialCharTok{\textasciitilde{}}\NormalTok{ flipper\_length\_mm)}
\NormalTok{model}
\end{Highlighting}
\end{Shaded}

\begin{verbatim}

Call:
lm(formula = body_mass_g ~ flipper_length_mm, data = penguins)

Coefficients:
      (Intercept)  flipper_length_mm  
         -5780.83              49.69  
\end{verbatim}

From the output we can see the \emph{coefficients} of the model. The
coefficients in this model are the \(m\) and \(k\) values in our linear
equation, \(y = k \cdot x + m\)\footnote{The letter \(k\) and \(m\) are
  typically used for the equation of the line in Sweden, apologies for
  any confusion for international students :).}, where the coefficient
under ``(Intercept)'' is the \(m\) and under ``flipper\_length\_mm'' is
the \(k\), the slope of the line. The slope tells us how much body mass
increases on average with every mm the flipper length increases. From
our example before we saw that a penguin with flipper length 200 mm on
average weigh 4156 g. According to our model a penguin with flipper
length 201 mm would weigh \(4156.282 + 49.69 \approx 4206\) g (using
exact values in the addition), increasing by one slope, i.e.~the value
under ``Coefficients: flipper\_length\_mm''. The coefficient of the
independent variable is of special interest, as it can tell us about the
effect one variable has on the other. The intercept, \(m\), in our case
represents the body mass of a penguin with flipper length 0, which
doesn't make much sense. The intercept may be interesting in some cases,
but you can safely ignore it in cases when it clearly doesn't make any
sense. It is important to remember the data used to fit the model, as
there are no penguins with flippers shorter than 172 mm or longer than
231 mm, and that our model is not fitted to predict values outside of
this. \textbf{Always be mindful when using your model outside the range
of your sample.}

The output from the model is not very detailed, instead we are more
interested the \texttt{summary} of the model.

\begin{Shaded}
\begin{Highlighting}[]
\NormalTok{model\_summary }\OtherTok{\textless{}{-}} \FunctionTok{summary}\NormalTok{(model)}
\NormalTok{model\_summary}
\end{Highlighting}
\end{Shaded}

\begin{verbatim}

Call:
lm(formula = body_mass_g ~ flipper_length_mm, data = penguins)

Residuals:
     Min       1Q   Median       3Q      Max 
-1058.80  -259.27   -26.88   247.33  1288.69 

Coefficients:
                   Estimate Std. Error t value Pr(>|t|)    
(Intercept)       -5780.831    305.815  -18.90   <2e-16 ***
flipper_length_mm    49.686      1.518   32.72   <2e-16 ***
---
Signif. codes:  0 '***' 0.001 '**' 0.01 '*' 0.05 '.' 0.1 ' ' 1

Residual standard error: 394.3 on 340 degrees of freedom
  (2 observations deleted due to missingness)
Multiple R-squared:  0.759, Adjusted R-squared:  0.7583 
F-statistic:  1071 on 1 and 340 DF,  p-value: < 2.2e-16
\end{verbatim}

This provides a much richer output. We can again see the coefficients,
but now with a lot more detail. We still find the estimated values of
the coefficients under the column ``Estimate'', but equally (if not
more) important is the value under
``Pr(\textgreater\textbar t\textbar)''. In this column we find a
\(p\)-value that tells how likely it is to get the linear relationship
we have in our data if there is in fact no relationship in the
underlying population of penguins. Specifically, the \(p\)-value tests
the hypothesis:

\begin{quote}
\textbf{\(H_0\)}: There is \emph{no} linear relationship between flipper
length and body mass.
\end{quote}

\begin{quote}
\textbf{\(H_a\)}: There \emph{is} a linear relationship between flipper
length and body mass.
\end{quote}

In other words, this means that the coefficient for
\texttt{flipper\_length\_mm} is either zero or non-zero:

\begin{quote}
\textbf{\(H_0\)}: The coefficient for flipper length \emph{is} 0.
\end{quote}

\begin{quote}
\textbf{\(H_a\)}: The coefficient for flipper length is \emph{not} 0.
\end{quote}

If the \(p\)-value is smaller than 0.05 we say that there is a
statistically significant linear relationship between our variables at
significance level 0.05. We also have a \(p\)-value for the intercept,
testing if the intercept is significantly different from 0, which is not
interesting in our case. We can also output the results more concisely
by accessing the \texttt{coefficients} using the dollar operator.

\begin{Shaded}
\begin{Highlighting}[]
\NormalTok{model\_summary}\SpecialCharTok{$}\NormalTok{coefficients}
\end{Highlighting}
\end{Shaded}

\begin{verbatim}
                     Estimate Std. Error   t value      Pr(>|t|)
(Intercept)       -5780.83136 305.814504 -18.90306  5.587301e-55
flipper_length_mm    49.68557   1.518404  32.72223 4.370681e-107
\end{verbatim}

Hypothesis testing gives linear regression a powerful interpretation: if
the \(p\)-value for the coefficient for the dependent variable we have
evidence of a linear relationship between the variables. Sometimes this
can also be interpreted as a \emph{causal effect}, that the predictor is
a cause of the outcome, however this requires more assumptions and we
should be careful drawing causal conclusions based on correlations. For
example, in medical studies this is of particular importance: we ideally
want a drug to have a causal effect on the outcome of a disease.

\subsection{Exercises}\label{exercises-12}

\begin{enumerate}
\def\labelenumi{\arabic{enumi}.}
\tightlist
\item
  Visualize and fit a linear model with flipper length as your dependent
  variable and body mass as your independent variable,
  \texttt{flipper\_length\_mm\ \textasciitilde{}\ body\_mass\_g}. How
  does the flipper length change on average when the body mass increases
  with one gram? Is it possible to distinguish any causal effect here?
\item
  The code below simulates data for ice cream sales. Copy the code and
  visualize the relationship (using \texttt{geom\_point} and
  \texttt{geom\_smooth} with \texttt{method\ =\ "lm"}) and fit a linear
  regression model of \texttt{sales\ \textasciitilde{}\ month}. Is the
  result as you expect? Do we have a linear relationship?
\end{enumerate}

\begin{Shaded}
\begin{Highlighting}[]
\CommentTok{\# Simulating data}
\NormalTok{day }\OtherTok{\textless{}{-}} \FunctionTok{seq}\NormalTok{(}\DecValTok{0}\NormalTok{, }\DecValTok{12}\NormalTok{, }\AttributeTok{length.out =} \DecValTok{356}\NormalTok{)}
\NormalTok{sales }\OtherTok{\textless{}{-}} \DecValTok{400} \SpecialCharTok{*} \FunctionTok{sin}\NormalTok{(pi }\SpecialCharTok{*}\NormalTok{ day }\SpecialCharTok{/} \DecValTok{12}\NormalTok{)}
\NormalTok{ice\_cream }\OtherTok{\textless{}{-}} \FunctionTok{tibble}\NormalTok{(}\AttributeTok{day =}\NormalTok{ day, }\AttributeTok{sales =}\NormalTok{ sales)}
\end{Highlighting}
\end{Shaded}

\begin{enumerate}
\def\labelenumi{\arabic{enumi}.}
\setcounter{enumi}{2}
\tightlist
\item
  Fit a linear model according to the formula
  \texttt{body\_mass\_g\ \textasciitilde{}\ species}. What do you expect
  will happen? This model will be discussed in the next section.
\item
  Run the code below that performs the t-test of female and male body
  weight in Gentoo penguins (from the previous chapter) and a fits a
  linear model \texttt{body\_mass\_g\ \textasciitilde{}\ species} for
  Gentoo penguins only. Compare the \(p\)-values.
\end{enumerate}

\begin{Shaded}
\begin{Highlighting}[]
\CommentTok{\# t{-}test}
\NormalTok{female\_gentoo\_body\_mass }\OtherTok{\textless{}{-}}\NormalTok{ penguins }\SpecialCharTok{|\textgreater{}}
  \FunctionTok{filter}\NormalTok{(species }\SpecialCharTok{==} \StringTok{"Gentoo"} \SpecialCharTok{\&}\NormalTok{ sex }\SpecialCharTok{==} \StringTok{"female"}\NormalTok{) }\SpecialCharTok{|\textgreater{}}
  \FunctionTok{select}\NormalTok{(body\_mass\_g) }\SpecialCharTok{|\textgreater{}}
  \FunctionTok{drop\_na}\NormalTok{()}
\NormalTok{male\_gentoo\_body\_mass }\OtherTok{\textless{}{-}}\NormalTok{ penguins }\SpecialCharTok{|\textgreater{}}
  \FunctionTok{filter}\NormalTok{(species }\SpecialCharTok{==} \StringTok{"Gentoo"} \SpecialCharTok{\&}\NormalTok{ sex }\SpecialCharTok{==} \StringTok{"male"}\NormalTok{) }\SpecialCharTok{|\textgreater{}}
  \FunctionTok{select}\NormalTok{(body\_mass\_g) }\SpecialCharTok{|\textgreater{}}
  \FunctionTok{drop\_na}\NormalTok{()}

\FunctionTok{t.test}\NormalTok{(male\_gentoo\_body\_mass, female\_gentoo\_body\_mass)}

\CommentTok{\# Linear regression}
\NormalTok{penguins }\SpecialCharTok{|\textgreater{}} 
  \FunctionTok{filter}\NormalTok{(species }\SpecialCharTok{==} \StringTok{"Gentoo"}\NormalTok{) }\SpecialCharTok{|\textgreater{}}
  \FunctionTok{lm}\NormalTok{(}\AttributeTok{formula =}\NormalTok{ body\_mass\_g }\SpecialCharTok{\textasciitilde{}}\NormalTok{ sex) }\SpecialCharTok{|\textgreater{}}
  \FunctionTok{summary}\NormalTok{()}
\end{Highlighting}
\end{Shaded}

\section{Categorical predictors}\label{categorical-predictors}

In the previous exercises we saw that we can fit a linear model with a
categorical variable as the independent variable, similarly to the
t-test in the previous chapter. We will take a look at the summary
output to better understand how this is done. We start by analyzing
\texttt{body\_weight\_g\ \textasciitilde{}\ species}. We can visualize
the distributions using a density plot with the mean values highlighted
as vertical lines.

\begin{Shaded}
\begin{Highlighting}[]
\NormalTok{mean\_bws }\OtherTok{\textless{}{-}}\NormalTok{ penguins }\SpecialCharTok{|\textgreater{}} 
  \FunctionTok{summarize}\NormalTok{(}
    \AttributeTok{body\_mass\_g =} \FunctionTok{mean}\NormalTok{(body\_mass\_g, }\AttributeTok{na.rm =} \ConstantTok{TRUE}\NormalTok{),}
    \AttributeTok{.by =}\NormalTok{ species)}
\NormalTok{penguins }\SpecialCharTok{|\textgreater{}} 
  \FunctionTok{ggplot}\NormalTok{(}\FunctionTok{aes}\NormalTok{(}\AttributeTok{x =}\NormalTok{ body\_mass\_g, }\AttributeTok{fill =}\NormalTok{ species, }\AttributeTok{col =}\NormalTok{ species)) }\SpecialCharTok{+} 
  \FunctionTok{geom\_density}\NormalTok{(}\AttributeTok{alpha =} \FloatTok{0.3}\NormalTok{) }\SpecialCharTok{+}
  \FunctionTok{geom\_vline}\NormalTok{(}
    \FunctionTok{aes}\NormalTok{(}
      \AttributeTok{xintercept =}\NormalTok{ body\_mass\_g,}
      \AttributeTok{col =}\NormalTok{ species),}
    \AttributeTok{data =}\NormalTok{ mean\_bws) }\SpecialCharTok{+}
  \FunctionTok{scale\_x\_continuous}\NormalTok{(}
    \AttributeTok{breaks =}\NormalTok{ mean\_bws}\SpecialCharTok{$}\NormalTok{body\_mass\_g,}
    \AttributeTok{labels =} \FunctionTok{round}\NormalTok{(mean\_bws}\SpecialCharTok{$}\NormalTok{body\_mass\_g),}
    \AttributeTok{guide =} \FunctionTok{guide\_axis}\NormalTok{(}\AttributeTok{n.dodge=}\DecValTok{3}\NormalTok{)) }\SpecialCharTok{+}
  \FunctionTok{labs}\NormalTok{(}
    \AttributeTok{x =} \StringTok{"Body mass (g)"}\NormalTok{,}
    \AttributeTok{y =} \StringTok{""}\NormalTok{,}
    \AttributeTok{col =} \StringTok{"Species"}\NormalTok{,}
    \AttributeTok{fill =} \StringTok{"Species"}\NormalTok{) }\SpecialCharTok{+}
  \FunctionTok{theme}\NormalTok{(}\AttributeTok{axis.text.y =} \FunctionTok{element\_blank}\NormalTok{(),}
        \AttributeTok{axis.ticks.y =} \FunctionTok{element\_blank}\NormalTok{())}
\end{Highlighting}
\end{Shaded}

\includegraphics{linear_regression_files/figure-pdf/unnamed-chunk-12-1.pdf}

We see that Gentoo penguins are in general larger, with most of the
density situated above the other two species and the mean body mass more
than one kg heavier than the other two species. Chinstrap and Adelie
penguins on the other hand have very similar body masses, with the
Chinstraps being a little heavier on average. A natural question is if
this difference is significant, is it only present in our sample or can
we say that there is a difference between the populations? Before we
used t-test, but now we will fit a linear model with \texttt{species} as
our independent variable. We immediately store the summary output and
look only at the coefficients.

\begin{Shaded}
\begin{Highlighting}[]
\NormalTok{model\_summary }\OtherTok{\textless{}{-}}\NormalTok{ penguins }\SpecialCharTok{|\textgreater{}} 
  \FunctionTok{lm}\NormalTok{(}\AttributeTok{formula =}\NormalTok{ body\_mass\_g }\SpecialCharTok{\textasciitilde{}}\NormalTok{ species) }\SpecialCharTok{|\textgreater{}} 
  \FunctionTok{summary}\NormalTok{()}
\NormalTok{model\_summary}\SpecialCharTok{$}\NormalTok{coefficients}
\end{Highlighting}
\end{Shaded}

\begin{verbatim}
                   Estimate Std. Error    t value      Pr(>|t|)
(Intercept)      3700.66225   37.61935 98.3712321 2.488024e-251
speciesChinstrap   32.42598   67.51168  0.4803018  6.313226e-01
speciesGentoo    1375.35401   56.14797 24.4951686  5.420612e-77
\end{verbatim}

We see that we have three coefficients, the intercept, one for
``speciesChinstrap'' and one for ``speciesGentoo''. Whenever we include
a categorical variable as a predictor the coefficients for that variable
represents the difference of that \emph{level} to a \emph{reference}, in
this case Adelie. Here we can actually interpret the intercept as
something useful, as it is the mean body weight of the Adelie penguins.
The coefficients for Chinstrap and Gentoo respectively is the difference
between the Adelie mean body weight and the Chinstrap and Gentoo mean
body weight respectively. From the output we can conclude that on
average Chinstrap penguins weigh \(\sim 32\) g more than Adelie penguins
and Gentoo penguins weigh \(\sim 1.4\) kg more than Adelie penguins. The
\(p\)-values, from the ``Pr(\textgreater\textbar t\textbar)'' column,
tells us if the mean of the Adelie species is different from the mean of
the Chinstrap and Gentoo respectively. \texttt{lm} selects the baseline
by itself. We can formulate the hypotheses for the tests generating the
\(p\)-values of ``speciesChinstrap'' and ``speciesGentoo''

\textbf{speciesChinstrap}

\begin{quote}
\textbf{\(H_0\)}: There is \emph{no} difference in the mean body weight
of Adelie and Chinstrap penguins.
\end{quote}

\begin{quote}
\textbf{\(H_a\)}: There \emph{is} a difference in the mean body weight
of Adelie and Chinstrap penguins.
\end{quote}

and likewise for

\textbf{speciesGentoo}

\begin{quote}
\textbf{\(H_0\)}: There is \emph{no} difference in the mean body weight
of Adelie and Gentoo penguins.
\end{quote}

\begin{quote}
\textbf{\(H_a\)}: There \emph{is} a difference in the mean body weight
of Adelie and Gentoo penguins.
\end{quote}

In this case the \(p\)-value for ``speciesGentoo'' is significant but
the \(p\)-value for ``speciesChinstrap'' is not. We can say that we
don't have evidence in the data that Adelie and Chinstrap penguins have
different body mass on average, but that Adelie and Gentoo have. We have
now used linear regression to perform a similar test to a t-test, but
testing multiple hypotheses at the same time, showing the versatility of
linear regression for hypotheses testing.

\subsection{Exercises}\label{exercises-13}

\begin{enumerate}
\def\labelenumi{\arabic{enumi}.}
\tightlist
\item
  Repeat the example above by fitting a linear model and making a
  density plot, but using \texttt{island} instead of species. Note which
  island \texttt{lm} chooses as the reference and think about how that
  impacts the result.
\item
  The code below sets Torgersen as the reference island. Copy the code
  and re-run the linear regression using the new tibble,
  i.e.~\texttt{penguins\_re}.
\end{enumerate}

\begin{Shaded}
\begin{Highlighting}[]
\NormalTok{penguins\_re }\OtherTok{\textless{}{-}}\NormalTok{ penguins }\SpecialCharTok{|\textgreater{}} \FunctionTok{mutate}\NormalTok{(}\AttributeTok{island =} \FunctionTok{relevel}\NormalTok{(island, }\AttributeTok{ref =} \StringTok{"Torgersen"}\NormalTok{))}
\end{Highlighting}
\end{Shaded}

\begin{enumerate}
\def\labelenumi{\arabic{enumi}.}
\setcounter{enumi}{2}
\tightlist
\item
  Make a bar chart that shows the distribution of the different species
  over the different islands. What do you think has the biggest effect
  on body mass, species or island?
\item
  Fit a linear model \texttt{body\_mass\_g\ \textasciitilde{}\ island}
  for only the Adelie penguins. Does island have an affect on body mass,
  i.e.~is one of the coefficients statistically significant at
  significance level 0.05?
\end{enumerate}

\section{A surprising result (Simpson's
paradox)}\label{a-surprising-result-simpsons-paradox}

After seeing the versatility of linear regression it may be tempting to
go ahead and fit linear models for any possible relationship in our data
without much hesitation. However, it is often a good idea to perform EDA
before any regression analysis to understand what relationships we
can/should model and if the result from the regression makes sense
according to our intuition. Sometimes a relationship will appear strange
or surprising when performing linear regression. We fit a linear model
on \texttt{body\_mass\_g\ \textasciitilde{}\ bill\_depth\_mm}, expecting
that bigger penguins have bigger bills, i.e.~bill depth.

\begin{Shaded}
\begin{Highlighting}[]
\NormalTok{model\_summary }\OtherTok{\textless{}{-}}\NormalTok{ penguins }\SpecialCharTok{|\textgreater{}} 
  \FunctionTok{lm}\NormalTok{(}\AttributeTok{formula =}\NormalTok{ body\_mass\_g }\SpecialCharTok{\textasciitilde{}}\NormalTok{ bill\_depth\_mm) }\SpecialCharTok{|\textgreater{}} 
  \FunctionTok{summary}\NormalTok{()}
\NormalTok{model\_summary}\SpecialCharTok{$}\NormalTok{coefficients}
\end{Highlighting}
\end{Shaded}

\begin{verbatim}
               Estimate Std. Error   t value     Pr(>|t|)
(Intercept)   7488.6524  335.21782 22.339661 1.133457e-68
bill_depth_mm -191.6428   19.41698 -9.869853 2.276941e-20
\end{verbatim}

We see that the coefficient for \texttt{bill\_depth\_mm} is negative,
saying that body mass decreases on average by \(\sim191\) g for every mm
the bill depth increases! This seems counter intuitive, so we can
further investigate by visualizing the relationship.

\begin{Shaded}
\begin{Highlighting}[]
\NormalTok{penguins }\SpecialCharTok{|\textgreater{}} 
  \FunctionTok{ggplot}\NormalTok{(}\FunctionTok{aes}\NormalTok{(}\AttributeTok{x =}\NormalTok{ bill\_depth\_mm, }\AttributeTok{y =}\NormalTok{ body\_mass\_g)) }\SpecialCharTok{+} 
  \FunctionTok{geom\_point}\NormalTok{() }\SpecialCharTok{+}
  \FunctionTok{geom\_smooth}\NormalTok{(}\AttributeTok{method =} \StringTok{"lm"}\NormalTok{, }\AttributeTok{se =} \ConstantTok{FALSE}\NormalTok{) }\SpecialCharTok{+}
  \FunctionTok{labs}\NormalTok{(}
    \AttributeTok{x =} \StringTok{"Bill depth (mm)"}\NormalTok{, }
    \AttributeTok{y =} \StringTok{"Body mass (g)"}\NormalTok{)}
\end{Highlighting}
\end{Shaded}

\includegraphics{linear_regression_files/figure-pdf/unnamed-chunk-16-1.pdf}

Our plot confirms our result, but something seems ``fishy'': there are
two clear ``point clouds'' in our plot and the line goes through both of
them. We can suspect the clouds represents different groups or
\emph{categories}. We have learned before that the body mass is very
different between the different species, so we can color the points
based on the species and fit individual lines for each species and see
if we can disentangle this strange relationship.

\begin{Shaded}
\begin{Highlighting}[]
\NormalTok{penguins }\SpecialCharTok{|\textgreater{}} 
  \FunctionTok{ggplot}\NormalTok{(}\FunctionTok{aes}\NormalTok{(}\AttributeTok{x =}\NormalTok{ bill\_depth\_mm, }\AttributeTok{y =}\NormalTok{ body\_mass\_g)) }\SpecialCharTok{+} 
  \FunctionTok{geom\_point}\NormalTok{(}\FunctionTok{aes}\NormalTok{(}\AttributeTok{col =}\NormalTok{ species)) }\SpecialCharTok{+}
  \FunctionTok{geom\_smooth}\NormalTok{(}\FunctionTok{aes}\NormalTok{(}\AttributeTok{col =}\NormalTok{ species), }\AttributeTok{method =} \StringTok{"lm"}\NormalTok{, }\AttributeTok{se =} \ConstantTok{FALSE}\NormalTok{) }\SpecialCharTok{+}
  \FunctionTok{labs}\NormalTok{(}
    \AttributeTok{x =} \StringTok{"Bill depth (mm)"}\NormalTok{, }
    \AttributeTok{y =} \StringTok{"Body mass (g)"}\NormalTok{,}
    \AttributeTok{col =} \StringTok{"Species"}\NormalTok{)}
\end{Highlighting}
\end{Shaded}

\includegraphics{linear_regression_files/figure-pdf/unnamed-chunk-17-1.pdf}

We see that the different species of penguins are responsible for the
negative slope, and within each species it seems that the slope is
indeed positive. This is an example of
\href{https://en.wikipedia.org/wiki/Simpson's_paradox}{Simpson's
paradox}, where a relationship is different within subgroups. We can
include \texttt{species} in our linear model using a + sign in the
formula,
\texttt{body\_mass\_g\ \textasciitilde{}\ bill\_depth\_mm\ +\ species}
to account for the difference.

\begin{Shaded}
\begin{Highlighting}[]
\NormalTok{model\_summary }\OtherTok{\textless{}{-}}\NormalTok{ penguins }\SpecialCharTok{|\textgreater{}} 
  \FunctionTok{lm}\NormalTok{(}\AttributeTok{formula =}\NormalTok{ body\_mass\_g }\SpecialCharTok{\textasciitilde{}}\NormalTok{ bill\_depth\_mm }\SpecialCharTok{+}\NormalTok{ species) }\SpecialCharTok{|\textgreater{}} 
  \FunctionTok{summary}\NormalTok{()}
\NormalTok{model\_summary}\SpecialCharTok{$}\NormalTok{coefficients}
\end{Highlighting}
\end{Shaded}

\begin{verbatim}
                    Estimate Std. Error    t value     Pr(>|t|)
(Intercept)      -1007.28112  323.56097 -3.1131107 2.009294e-03
bill_depth_mm      256.61461   17.56282 14.6112380 8.100884e-38
speciesChinstrap    13.37732   52.94712  0.2526544 8.006889e-01
speciesGentoo     2238.66811   73.68183 30.3829071 1.157366e-98
\end{verbatim}

In our new model we see that we have a positive relationship between
bill depth and body mass since we now take species into account, and the
model now says that body mass increases on average with \(\sim 256\) g
per mm the bill depth increases. In the next chapter you will learn more
about \emph{multiple linear regression}.

\subsection{Exercises}\label{exercises-14}

\begin{enumerate}
\def\labelenumi{\arabic{enumi}.}
\tightlist
\item
  Revisit \href{https://r4ds.hadley.nz/data-visualize}{Chapter 2 in
  R4DS} (see plot below) and think about if the coefficient for
  \texttt{species} would be significant when fitting the model
  \texttt{body\_mass\_g\ \textasciitilde{}\ flipper\_length\_mm\ +\ species}.
  Note that the line is calculated based on \emph{all} species,
  i.e.~\texttt{body\_mass\_g\ \textasciitilde{}\ flipper\_length\_mm}.
\end{enumerate}

\begin{Shaded}
\begin{Highlighting}[]
\NormalTok{penguins }\SpecialCharTok{|\textgreater{}} 
  \FunctionTok{ggplot}\NormalTok{(}\FunctionTok{aes}\NormalTok{(}\AttributeTok{x =}\NormalTok{ flipper\_length\_mm, }\AttributeTok{y =}\NormalTok{ body\_mass\_g)) }\SpecialCharTok{+}
  \FunctionTok{geom\_point}\NormalTok{(}\FunctionTok{aes}\NormalTok{(}\AttributeTok{col =}\NormalTok{ species)) }\SpecialCharTok{+}
  \FunctionTok{geom\_smooth}\NormalTok{(}\AttributeTok{method =} \StringTok{"lm"}\NormalTok{, }\AttributeTok{se =} \ConstantTok{FALSE}\NormalTok{, }\AttributeTok{col =} \StringTok{"black"}\NormalTok{) }\SpecialCharTok{+}
  \FunctionTok{labs}\NormalTok{(}
    \AttributeTok{x =} \StringTok{"Flipper length (mm)"}\NormalTok{,}
    \AttributeTok{y =} \StringTok{"Body mass (g)"}\NormalTok{,}
    \AttributeTok{col =} \StringTok{"Species"}\NormalTok{)}
\end{Highlighting}
\end{Shaded}

\includegraphics{linear_regression_files/figure-pdf/unnamed-chunk-19-1.pdf}

\begin{enumerate}
\def\labelenumi{\arabic{enumi}.}
\setcounter{enumi}{1}
\tightlist
\item
  Fit the model
  \texttt{body\_mass\_g\ \textasciitilde{}\ flipper\_length\_mm\ +\ species}
  to find out!
\end{enumerate}

\section{Prediction and prediction
interval}\label{prediction-and-prediction-interval}

We already predicted body mass based on flipper length using the fitted
line and those who looked at the code may have spotted that we used the
\texttt{predict} function. With \texttt{predict} we can use our model to
predict body weight based on flipper lengths that we did not observe in
our data. As before, we need to be careful when predicting outside the
range of our sample: flippers shorter than 172 mm or longer than 231 mm.
To use the \texttt{predict} function we need to supply a model and a
data frame (or a tibble) specifying the independent variables (the
predictors) that we want to use for prediction. Before we predicted the
body weight of a penguin with flipper length 200 mm, now we will do it
again using the \texttt{predict} function.

\begin{Shaded}
\begin{Highlighting}[]
\NormalTok{model }\OtherTok{\textless{}{-}}\NormalTok{ penguins }\SpecialCharTok{|\textgreater{}} 
  \FunctionTok{lm}\NormalTok{(}\AttributeTok{formula =}\NormalTok{ body\_mass\_g }\SpecialCharTok{\textasciitilde{}}\NormalTok{ flipper\_length\_mm)}
\NormalTok{model }\SpecialCharTok{|\textgreater{}} \FunctionTok{predict}\NormalTok{(}\FunctionTok{data.frame}\NormalTok{(}\AttributeTok{flipper\_length\_mm =} \DecValTok{200}\NormalTok{))}
\end{Highlighting}
\end{Shaded}

\begin{verbatim}
       1 
4156.282 
\end{verbatim}

We see that we get the same value as before (remember that we rounded
the value in the plot). The line represents the average body weight for
a given flipper length, but there is no guarantee that any one penguin
will fall precisely on the line. We can calculate a \emph{prediction
interval} using the \texttt{predict} function to say something about the
uncertainty in the prediction. To add a prediction interval we use the
argument \texttt{interval\ =\ "prediction"} and \emph{confidence level}
\texttt{level\ =\ 0.95}.

\begin{Shaded}
\begin{Highlighting}[]
\NormalTok{model }\SpecialCharTok{|\textgreater{}} \FunctionTok{predict}\NormalTok{(}\FunctionTok{data.frame}\NormalTok{(}\AttributeTok{flipper\_length\_mm =} \DecValTok{200}\NormalTok{), }\AttributeTok{interval =} \StringTok{"prediction"}\NormalTok{, }\AttributeTok{level =} \FloatTok{0.95}\NormalTok{)}
\end{Highlighting}
\end{Shaded}

\begin{verbatim}
       fit      lwr      upr
1 4156.282 3379.612 4932.951
\end{verbatim}

We now get two more outputs, the lower (lwr) and upper (upr) limits of
our prediction. We see that the interval is quite big, and there is a
lot of uncertainty regarding \emph{one} penguin. The prediction interval
can then be interpreted as ``we are 95\% \emph{confident} that a new
prediction will fall inside this interval'', i.e.~that if we observe a
previously unseen penguin with flippers 200 mm long, we are 95\%
confident its body mass will be within the interval. We can visualize
the prediction interval for all the flipper lengths observed in our
sample.

\begin{Shaded}
\begin{Highlighting}[]
\NormalTok{model }\OtherTok{\textless{}{-}}\NormalTok{ penguins }\SpecialCharTok{|\textgreater{}} 
  \FunctionTok{lm}\NormalTok{(}\AttributeTok{formula =}\NormalTok{ body\_mass\_g }\SpecialCharTok{\textasciitilde{}}\NormalTok{ flipper\_length\_mm)}
\NormalTok{predictions }\OtherTok{\textless{}{-}}\NormalTok{ model }\SpecialCharTok{|\textgreater{}} \FunctionTok{predict}\NormalTok{(penguins }\SpecialCharTok{|\textgreater{}} \FunctionTok{select}\NormalTok{(flipper\_length\_mm), }\AttributeTok{interval =} \StringTok{"prediction"}\NormalTok{, }\AttributeTok{level =} \FloatTok{0.95}\NormalTok{) }\SpecialCharTok{|\textgreater{}} \FunctionTok{as.data.frame}\NormalTok{()}
\NormalTok{penguins\_pred }\OtherTok{\textless{}{-}}\NormalTok{ penguins }\SpecialCharTok{|\textgreater{}} 
  \FunctionTok{mutate}\NormalTok{(}
    \AttributeTok{pred =}\NormalTok{ predictions}\SpecialCharTok{$}\NormalTok{fit,}
    \AttributeTok{lwr =}\NormalTok{ predictions}\SpecialCharTok{$}\NormalTok{lwr,}
    \AttributeTok{upr =}\NormalTok{ predictions}\SpecialCharTok{$}\NormalTok{upr)}

\NormalTok{penguins\_pred }\SpecialCharTok{|\textgreater{}} 
  \FunctionTok{ggplot}\NormalTok{(}\FunctionTok{aes}\NormalTok{(}\AttributeTok{x =}\NormalTok{ flipper\_length\_mm, }\AttributeTok{y =}\NormalTok{ body\_mass\_g)) }\SpecialCharTok{+} 
  \FunctionTok{geom\_ribbon}\NormalTok{(}\FunctionTok{aes}\NormalTok{(}\AttributeTok{ymin =}\NormalTok{ lwr, }\AttributeTok{ymax =}\NormalTok{ upr), }\AttributeTok{col =} \StringTok{"gray"}\NormalTok{, }\AttributeTok{alpha =} \FloatTok{0.2}\NormalTok{) }\SpecialCharTok{+}
  \FunctionTok{geom\_point}\NormalTok{() }\SpecialCharTok{+} 
  \FunctionTok{geom\_smooth}\NormalTok{(}\AttributeTok{method =} \StringTok{"lm"}\NormalTok{, }\AttributeTok{se =} \ConstantTok{FALSE}\NormalTok{) }\SpecialCharTok{+} 
  \FunctionTok{labs}\NormalTok{(}
    \AttributeTok{x =} \StringTok{"Flipper length (mm)"}\NormalTok{, }
    \AttributeTok{y =} \StringTok{"Body mass (g)"}\NormalTok{)}
\end{Highlighting}
\end{Shaded}

\includegraphics{linear_regression_files/figure-pdf/unnamed-chunk-22-1.pdf}

We see that most of our observations fall within the prediction
interval, but on some occasions outside. With more data, the interval
becomes tighter as the model is more ``sure'' about its predictions.

{Another common interval is the \emph{confidence interval}. It informs
about the uncertainty of the line itself, i.e.~the \emph{average} value
of body mass for a give flipper length. It is always tighter than the
prediction interval and you can visualize it by setting
\texttt{se\ =\ TRUE} in \texttt{geom\_smooth}. In this course we will
focus on prediction intervals.}

\subsection{Exercises}\label{exercises-15}

\begin{enumerate}
\def\labelenumi{\arabic{enumi}.}
\tightlist
\item
  Redo the predictions for flipper lengths 201 mm and 100 mm, this time
  using \texttt{predict} with
  \texttt{interval\ =\ "prediction",\ level\ =\ 0.95}.
\item
  What happens with the interval if you set it to 0.5? Does it become
  bigger or smaller? Try it out!
\end{enumerate}

\section{Assumptions of linear
regression}\label{assumptions-of-linear-regression}

So far we have only talked about the assumption that there is a linear
relationship between outcome and predictor. Linear regression comes with
several additional assumptions that are important for the prediction
interval to be valid. \textbf{We will not go into detail in these
assumptions or the theory behind them and you will not be tested on
them}, but we will list them here so you at least have seen them.

\begin{itemize}
\tightlist
\item
  \textbf{Linearity}: The relationship between the independent and
  dependent variables is linear. We cannot model non-linear
  relationships with a linear model as the model won't reflect reality
  and predictions will be bad.
\item
  \textbf{Homoscedasticity}: The variance of the residuals (the
  unexplained variance, see next chapter) is constant across all levels
  of the independent variables. This is needed for model inference (such
  as the prediction intervals and t-tests for the coefficients) to be
  accurate.
\item
  \textbf{Normality}: The residuals of the model are normally
  distributed. This is also needed for inference.
\item
  \textbf{Independence}: The observations are independent of each other
  (values of one observation do not influence the values of another). If
  they are not independent, the model may not generalize well to new
  data not used to fit the model.
\end{itemize}

When performing linear regression outside of this course you should
understand these assumptions in order to make valid analyses.

In this chapter we have introduced simple linear regression with one
predictor, except for when we included species along with bill depth. In
the next chapter we will explore using multiple predictors and selecting
between different models based on different variables in multiple linear
regression.

\chapter{Multiple linear regression}\label{sec-multlinreg}

In the last chapter we got a taste of multiple linear regression when we
included \texttt{species} to fit the relationship between
\texttt{body\_mass\_g} and \texttt{bill\_depth\_mm}. In this chapter we
will continue fitting regression models with multiple predictors,
capable of explaining more of the variation in the response (dependent
variable). Multiple linear regression is a powerful tool, but it
requires some care in order to create a model that is both good at
prediction and \emph{interpretable}, i.e.~a model that we can easily
understand with coefficients that make sense and that the model doesn't
get too big to understand. Sometimes we have to/should sacrifice
explanatory power in benefit of a simpler model.

To study multiple linear regression we will use an anonymized dataset on
the birth weight of children, given to us from the medical faculty at
Umeå University. Our goal is to analyze what affects the birth weight of
children given the variables in the dataset, which is available on
Canvas if you want to do the exercises or reproduce any analysis made in
this chapter.

\section{Birth weight dataset}\label{birth-weight-dataset}

The dataset was collected at a hospital in Sweden and consists of
observations of 230 children. We can take a \texttt{glimpse} at the
dataset to see what variables it contains. First we load the dataset
using \texttt{read\_csv}.

\begin{Shaded}
\begin{Highlighting}[]
\FunctionTok{library}\NormalTok{(tidyverse)}
\NormalTok{birth }\OtherTok{\textless{}{-}} \FunctionTok{read\_csv}\NormalTok{(}\StringTok{"datasets/birth.csv"}\NormalTok{) }\CommentTok{\# Remember to set the correct path}
\FunctionTok{glimpse}\NormalTok{(birth)}
\end{Highlighting}
\end{Shaded}

\begin{verbatim}
Rows: 230
Columns: 7
$ birth_weight_g         <dbl> 4073, 4060, 3071, 3546, 3222, 3420, 3106, 3274,~
$ pregnancy_duration_wks <dbl> 41, 40, 39, 39, 39, 41, 40, 40, 40, 41, 41, 38,~
$ sex                    <dbl> 1, 2, 2, 1, 2, 1, 2, 1, 1, 1, 2, 1, 1, 2, 1, 1,~
$ age                    <dbl> 42, 27, 37, 39, 31, 33, 38, 26, 30, 35, 29, 29,~
$ length_cm              <dbl> 176, 180, 171, 167, 158, 166, 169, 164, 167, 16~
$ weight_kg              <dbl> 65, 81, 70, 80, 54, 50, 50, 65, 80, 55, 57, 63,~
$ bmi                    <dbl> 20.98399, 25.00000, 23.93899, 28.68514, 21.6311~
\end{verbatim}

Our response variable is birth weight in grams and the predictors are:
pregnancy duration in weeks, the sex of the child, and age, length,
weight, and BMI of the mother. We can see that \texttt{sex} is coded as
a numerical variable (\texttt{double}) and we should start by recoding
it as a factor. In the dataset males are encoded as 1 and females as 2.

\begin{Shaded}
\begin{Highlighting}[]
\CommentTok{\# Recoding sex as factor}
\NormalTok{birth }\OtherTok{\textless{}{-}}\NormalTok{ birth }\SpecialCharTok{|\textgreater{}}  \FunctionTok{mutate}\NormalTok{(}\AttributeTok{sex =} \FunctionTok{as.factor}\NormalTok{(}\FunctionTok{ifelse}\NormalTok{(sex }\SpecialCharTok{==} \DecValTok{1}\NormalTok{, }\StringTok{"Male"}\NormalTok{, }\StringTok{"Female"}\NormalTok{)))}
\end{Highlighting}
\end{Shaded}

Lets explore our data using EDA. Since we are interested in linear
regression, we can start by plotting the response against the
predictors. We use scatter plots for the numerical variables and a box
plot for sex.

\begin{Shaded}
\begin{Highlighting}[]
\FunctionTok{library}\NormalTok{(patchwork)}

\NormalTok{p0 }\OtherTok{\textless{}{-}}\NormalTok{ birth }\SpecialCharTok{|\textgreater{}} 
  \FunctionTok{ggplot}\NormalTok{(}\FunctionTok{aes}\NormalTok{(}\AttributeTok{x =}\NormalTok{ sex, }\AttributeTok{y =}\NormalTok{ birth\_weight\_g, }\AttributeTok{fill =}\NormalTok{ sex)) }\SpecialCharTok{+} 
  \FunctionTok{geom\_boxplot}\NormalTok{() }\SpecialCharTok{+}
  \FunctionTok{guides}\NormalTok{(}\AttributeTok{fill =} \ConstantTok{FALSE}\NormalTok{)}

\NormalTok{p1 }\OtherTok{\textless{}{-}}\NormalTok{ birth }\SpecialCharTok{|\textgreater{}} 
  \FunctionTok{ggplot}\NormalTok{(}\FunctionTok{aes}\NormalTok{(}
    \AttributeTok{x =}\NormalTok{ pregnancy\_duration\_wks, }
    \AttributeTok{y =}\NormalTok{ birth\_weight\_g)) }\SpecialCharTok{+} 
  \FunctionTok{geom\_point}\NormalTok{() }\SpecialCharTok{+} 
  \FunctionTok{geom\_smooth}\NormalTok{(}\AttributeTok{method =} \StringTok{"lm"}\NormalTok{, }\AttributeTok{se =} \ConstantTok{FALSE}\NormalTok{) }

\NormalTok{p2 }\OtherTok{\textless{}{-}}\NormalTok{ birth }\SpecialCharTok{|\textgreater{}} 
  \FunctionTok{ggplot}\NormalTok{(}\FunctionTok{aes}\NormalTok{(}
    \AttributeTok{x =}\NormalTok{ age, }
    \AttributeTok{y =}\NormalTok{ birth\_weight\_g)) }\SpecialCharTok{+} 
  \FunctionTok{geom\_point}\NormalTok{() }\SpecialCharTok{+} 
  \FunctionTok{geom\_smooth}\NormalTok{(}\AttributeTok{method =} \StringTok{"lm"}\NormalTok{, }\AttributeTok{se =} \ConstantTok{FALSE}\NormalTok{)}

\NormalTok{p3 }\OtherTok{\textless{}{-}}\NormalTok{ birth }\SpecialCharTok{|\textgreater{}} 
  \FunctionTok{ggplot}\NormalTok{(}\FunctionTok{aes}\NormalTok{(}
    \AttributeTok{x =}\NormalTok{ length\_cm, }
    \AttributeTok{y =}\NormalTok{ birth\_weight\_g)) }\SpecialCharTok{+} 
  \FunctionTok{geom\_point}\NormalTok{() }\SpecialCharTok{+} 
  \FunctionTok{geom\_smooth}\NormalTok{(}\AttributeTok{method =} \StringTok{"lm"}\NormalTok{, }\AttributeTok{se =} \ConstantTok{FALSE}\NormalTok{)}

\NormalTok{p4 }\OtherTok{\textless{}{-}}\NormalTok{ birth }\SpecialCharTok{|\textgreater{}} 
  \FunctionTok{ggplot}\NormalTok{(}\FunctionTok{aes}\NormalTok{(}
    \AttributeTok{x =}\NormalTok{ weight\_kg, }
    \AttributeTok{y =}\NormalTok{ birth\_weight\_g)) }\SpecialCharTok{+} 
  \FunctionTok{geom\_point}\NormalTok{() }\SpecialCharTok{+} 
  \FunctionTok{geom\_smooth}\NormalTok{(}\AttributeTok{method =} \StringTok{"lm"}\NormalTok{, }\AttributeTok{se =} \ConstantTok{FALSE}\NormalTok{)}

\NormalTok{p5 }\OtherTok{\textless{}{-}}\NormalTok{ birth }\SpecialCharTok{|\textgreater{}} 
  \FunctionTok{ggplot}\NormalTok{(}\FunctionTok{aes}\NormalTok{(}
    \AttributeTok{x =}\NormalTok{ bmi, }
    \AttributeTok{y =}\NormalTok{ birth\_weight\_g)) }\SpecialCharTok{+} 
  \FunctionTok{geom\_point}\NormalTok{() }\SpecialCharTok{+} 
  \FunctionTok{geom\_smooth}\NormalTok{(}\AttributeTok{method =} \StringTok{"lm"}\NormalTok{, }\AttributeTok{se =} \ConstantTok{FALSE}\NormalTok{)}

\NormalTok{(p0 }\SpecialCharTok{+}\NormalTok{ p1 }\SpecialCharTok{+}\NormalTok{ p2) }\SpecialCharTok{/}\NormalTok{ (p3 }\SpecialCharTok{+}\NormalTok{ p4 }\SpecialCharTok{+}\NormalTok{ p5)}
\end{Highlighting}
\end{Shaded}

\includegraphics{linear_regression_multiple_files/figure-pdf/unnamed-chunk-3-1.pdf}

It seems like there is a linear relationship between birth weight and
pregnancy duration and mother length, weight, and possibly BMI. The
distribution of birth weight over the two sexes also seems to be
different, indicating that the mean birth weight may differ between
females and males. Another way of quickly assessing the strength of the
linear relationships between variables is a correlation plot. First we
compute the pairwise correlations of the numeric variables using
\texttt{cor}.

\begin{Shaded}
\begin{Highlighting}[]
\FunctionTok{library}\NormalTok{(ggcorrplot)}
\CommentTok{\# You can also use }
\CommentTok{\#library(corrplot)}
\NormalTok{birth }\SpecialCharTok{|\textgreater{}} 
  \FunctionTok{select}\NormalTok{(}\FunctionTok{where}\NormalTok{(is.numeric)) }\SpecialCharTok{|\textgreater{}} 
  \FunctionTok{cor}\NormalTok{() }\SpecialCharTok{|\textgreater{}} 
  \FunctionTok{ggcorrplot}\NormalTok{()}
\end{Highlighting}
\end{Shaded}

\includegraphics{linear_regression_multiple_files/figure-pdf/unnamed-chunk-4-1.pdf}

Looking at the row or column for \texttt{birth\_weight\_g} lets us
quickly see that all other variables are positively correlated with
\texttt{birth\_weight\_g}, as red indicates positive correlation (see
the legend ``Corr''), confirming what we saw in our the scatter plots
before. We will revisit correlation from a the perspective of linear
regression in a moment, but first we will ask ourselves if these results
are statistically significant, i.e.~if the result holds in the
population or is just the result of chance, by performing simple linear
regression.

\subsection{Exercises}\label{exercises-16}

\begin{enumerate}
\def\labelenumi{\arabic{enumi}.}
\tightlist
\item
  What is the population that generated this sample? Would you like to
  know more about the sample to better specify the population? Are there
  more variables that could be interesting for predicting birth weight?
\item
  Perform simple linear regression of
  \texttt{birth\_weight\_g\ \textasciitilde{}\ age} to see if there is a
  significant linear relationship. State the null and alternative
  hypotheses and analyze the result using significance level 0.05.
\item
  Perform simple linear regression of
  \texttt{birth\_weight\_g\ \textasciitilde{}\ length\_cm},
  \texttt{birth\_weight\_g\ \textasciitilde{}\ weight\_kg}, and
  \texttt{birth\_weight\_g\ \textasciitilde{}\ bmi} to see if there is a
  significant linear relationship between the variables.
\end{enumerate}

\section{Comparing models}\label{comparing-models}

Before we can start performing multiple linear regression we must learn
how to \emph{compare} models and what constitutes a \emph{good} model.

In the previous exercises we fitted several linear models to predict the
birth weight of the child. Two of them were significant, the models
using \texttt{length\_cm} and \texttt{weight\_kg} as predictors, while
\texttt{age} and \texttt{bmi} did not show a significant linear
relationship with \texttt{birth\_weight\_g}. A natural question to ask
is which of the two models
\texttt{birth\_weight\_g\ \textasciitilde{}\ length\_cm} and
\texttt{birth\_weight\_g\ \textasciitilde{}\ weight\_kg} is the best,
but how do we determine if a model is better than another?

\subsection{Correlation revisited}\label{correlation-revisited}

In Section~\ref{sec-basicstat-correlation} we briefly touched upon
correlation and looked at this figure, where strong correlation is
indicated by values close to -1 and 1 and values close to 0 indicate
weak or no correlation.

\begin{Shaded}
\begin{Highlighting}[]
\FunctionTok{library}\NormalTok{(patchwork)}
\FunctionTok{set.seed}\NormalTok{(}\DecValTok{42}\NormalTok{) }\CommentTok{\# Setting a seed to get the same output every time.}

\CommentTok{\# Generating x and y variables based on different correlation scenarios}
\NormalTok{N }\OtherTok{\textless{}{-}} \DecValTok{1000}
\NormalTok{x }\OtherTok{\textless{}{-}} \FunctionTok{rnorm}\NormalTok{(N)}
\NormalTok{y1 }\OtherTok{\textless{}{-}} \SpecialCharTok{{-}}\NormalTok{x }\SpecialCharTok{+} \FunctionTok{rnorm}\NormalTok{(N, }\AttributeTok{sd =} \FloatTok{0.01}\NormalTok{)  }\CommentTok{\# Correlation \textasciitilde{} {-}1}
\NormalTok{y2 }\OtherTok{\textless{}{-}} \SpecialCharTok{{-}}\FloatTok{0.5} \SpecialCharTok{*}\NormalTok{ x }\SpecialCharTok{+} \FunctionTok{rnorm}\NormalTok{(N, }\AttributeTok{sd =} \FunctionTok{sqrt}\NormalTok{(}\FloatTok{0.75}\NormalTok{))  }\CommentTok{\# Correlation \textasciitilde{} {-}0.5}
\NormalTok{y3 }\OtherTok{\textless{}{-}} \FunctionTok{rnorm}\NormalTok{(N)  }\CommentTok{\# Correlation \textasciitilde{} 0}
\NormalTok{y4 }\OtherTok{\textless{}{-}} \FloatTok{0.5} \SpecialCharTok{*}\NormalTok{ x }\SpecialCharTok{+} \FunctionTok{rnorm}\NormalTok{(N, }\AttributeTok{sd =} \FunctionTok{sqrt}\NormalTok{(}\FloatTok{0.75}\NormalTok{))  }\CommentTok{\# Correlation \textasciitilde{} 0.5}
\NormalTok{y5 }\OtherTok{\textless{}{-}}\NormalTok{ x }\SpecialCharTok{+} \FunctionTok{rnorm}\NormalTok{(N, }\AttributeTok{sd =} \FloatTok{0.01}\NormalTok{)  }\CommentTok{\# Correlation \textasciitilde{} 1}

\CommentTok{\# Creating a data frame}
\NormalTok{corr\_df }\OtherTok{\textless{}{-}} \FunctionTok{data.frame}\NormalTok{(x, y1, y2, y3, y4, y5)}

\CommentTok{\# Creating individual plots with titles and modified themes}
\NormalTok{plot1 }\OtherTok{\textless{}{-}} \FunctionTok{ggplot}\NormalTok{(}\AttributeTok{data =}\NormalTok{ corr\_df, }\FunctionTok{aes}\NormalTok{(}\AttributeTok{x =}\NormalTok{ x, }\AttributeTok{y =}\NormalTok{ y1)) }\SpecialCharTok{+}
  \FunctionTok{geom\_point}\NormalTok{(}\AttributeTok{alpha =} \FloatTok{0.5}\NormalTok{) }\SpecialCharTok{+}
  \FunctionTok{geom\_smooth}\NormalTok{(}\AttributeTok{method =} \StringTok{"lm"}\NormalTok{, }\AttributeTok{se =}\NormalTok{ F, }\AttributeTok{col =} \StringTok{"purple"}\NormalTok{) }\SpecialCharTok{+}
  \FunctionTok{theme\_minimal}\NormalTok{() }\SpecialCharTok{+}
  \FunctionTok{labs}\NormalTok{(}
    \AttributeTok{y =} \StringTok{"y"}\NormalTok{, }
    \AttributeTok{subtitle =} \FunctionTok{round}\NormalTok{(}\FunctionTok{cor}\NormalTok{(corr\_df}\SpecialCharTok{$}\NormalTok{x, corr\_df}\SpecialCharTok{$}\NormalTok{y1), }\AttributeTok{digits =} \DecValTok{2}\NormalTok{) ) }\SpecialCharTok{+}
  \FunctionTok{theme}\NormalTok{(}
    \AttributeTok{aspect.ratio =} \DecValTok{1}\NormalTok{,}
    \AttributeTok{axis.text =} \FunctionTok{element\_blank}\NormalTok{(),}
    \AttributeTok{axis.ticks =} \FunctionTok{element\_blank}\NormalTok{(),}
    \AttributeTok{plot.title =} \FunctionTok{element\_text}\NormalTok{(}\AttributeTok{hjust =} \FloatTok{0.5}\NormalTok{)  }\CommentTok{\# Center the title horizontally}
\NormalTok{  )}

\NormalTok{plot2 }\OtherTok{\textless{}{-}} \FunctionTok{ggplot}\NormalTok{(}\AttributeTok{data =}\NormalTok{ corr\_df, }\FunctionTok{aes}\NormalTok{(}\AttributeTok{x =}\NormalTok{ x, }\AttributeTok{y =}\NormalTok{ y2)) }\SpecialCharTok{+}
  \FunctionTok{geom\_point}\NormalTok{(}\AttributeTok{alpha =} \FloatTok{0.5}\NormalTok{) }\SpecialCharTok{+}
  \FunctionTok{geom\_smooth}\NormalTok{(}\AttributeTok{method =} \StringTok{"lm"}\NormalTok{, }\AttributeTok{se =}\NormalTok{ F) }\SpecialCharTok{+}
  \FunctionTok{labs}\NormalTok{(}\AttributeTok{subtitle =} \FunctionTok{round}\NormalTok{(}\FunctionTok{cor}\NormalTok{(corr\_df}\SpecialCharTok{$}\NormalTok{x, corr\_df}\SpecialCharTok{$}\NormalTok{y2), }\AttributeTok{digits =} \DecValTok{2}\NormalTok{)) }\SpecialCharTok{+}
  \FunctionTok{theme\_minimal}\NormalTok{() }\SpecialCharTok{+}
  \FunctionTok{theme}\NormalTok{(}
    \AttributeTok{aspect.ratio =} \DecValTok{1}\NormalTok{,}
    \AttributeTok{axis.text =} \FunctionTok{element\_blank}\NormalTok{(),}
    \AttributeTok{axis.ticks =} \FunctionTok{element\_blank}\NormalTok{(),}
    \AttributeTok{axis.title.y =} \FunctionTok{element\_blank}\NormalTok{(),}
    \AttributeTok{plot.title =} \FunctionTok{element\_text}\NormalTok{(}\AttributeTok{hjust =} \FloatTok{0.5}\NormalTok{)  }\CommentTok{\# Center the title horizontally}
\NormalTok{  )}

\NormalTok{plot3 }\OtherTok{\textless{}{-}} \FunctionTok{ggplot}\NormalTok{(}\AttributeTok{data =}\NormalTok{ corr\_df, }\FunctionTok{aes}\NormalTok{(}\AttributeTok{x =}\NormalTok{ x, }\AttributeTok{y =}\NormalTok{ y3)) }\SpecialCharTok{+}
  \FunctionTok{geom\_point}\NormalTok{(}\AttributeTok{alpha =} \FloatTok{0.5}\NormalTok{) }\SpecialCharTok{+}
  \FunctionTok{geom\_smooth}\NormalTok{(}\AttributeTok{method =} \StringTok{"lm"}\NormalTok{, }\AttributeTok{se =}\NormalTok{ F) }\SpecialCharTok{+}
  \FunctionTok{labs}\NormalTok{(}
    \AttributeTok{title =} \StringTok{"Correlation"}\NormalTok{,}
    \AttributeTok{subtitle =} \FunctionTok{round}\NormalTok{(}\FunctionTok{cor}\NormalTok{(corr\_df}\SpecialCharTok{$}\NormalTok{x, corr\_df}\SpecialCharTok{$}\NormalTok{y3), }\AttributeTok{digits =} \DecValTok{2}\NormalTok{)) }\SpecialCharTok{+}
  \FunctionTok{theme\_minimal}\NormalTok{() }\SpecialCharTok{+}
  \FunctionTok{theme}\NormalTok{(}
    \AttributeTok{aspect.ratio =} \DecValTok{1}\NormalTok{,}
    \AttributeTok{axis.text =} \FunctionTok{element\_blank}\NormalTok{(),}
    \AttributeTok{axis.ticks =} \FunctionTok{element\_blank}\NormalTok{(),}
    \AttributeTok{axis.title.y =} \FunctionTok{element\_blank}\NormalTok{(),}
    \AttributeTok{plot.title =} \FunctionTok{element\_text}\NormalTok{(}\AttributeTok{hjust =} \FloatTok{0.5}\NormalTok{)  }\CommentTok{\# Center the title horizontally}
\NormalTok{  )}

\NormalTok{plot4 }\OtherTok{\textless{}{-}} \FunctionTok{ggplot}\NormalTok{(}\AttributeTok{data =}\NormalTok{ corr\_df, }\FunctionTok{aes}\NormalTok{(}\AttributeTok{x =}\NormalTok{ x, }\AttributeTok{y =}\NormalTok{ y4)) }\SpecialCharTok{+}
  \FunctionTok{geom\_point}\NormalTok{(}\AttributeTok{alpha =} \FloatTok{0.5}\NormalTok{) }\SpecialCharTok{+}
  \FunctionTok{geom\_smooth}\NormalTok{(}\AttributeTok{method =} \StringTok{"lm"}\NormalTok{, }\AttributeTok{se =}\NormalTok{ F) }\SpecialCharTok{+}
  \FunctionTok{labs}\NormalTok{(}\AttributeTok{subtitle =} \FunctionTok{round}\NormalTok{(}\FunctionTok{cor}\NormalTok{(corr\_df}\SpecialCharTok{$}\NormalTok{x, corr\_df}\SpecialCharTok{$}\NormalTok{y4), }\AttributeTok{digits =} \DecValTok{2}\NormalTok{)) }\SpecialCharTok{+}
  \FunctionTok{theme\_minimal}\NormalTok{() }\SpecialCharTok{+}
  \FunctionTok{theme}\NormalTok{(}
    \AttributeTok{aspect.ratio =} \DecValTok{1}\NormalTok{,}
    \AttributeTok{axis.text =} \FunctionTok{element\_blank}\NormalTok{(),}
    \AttributeTok{axis.ticks =} \FunctionTok{element\_blank}\NormalTok{(),}
    \AttributeTok{axis.title.y =} \FunctionTok{element\_blank}\NormalTok{(),}
    \AttributeTok{plot.title =} \FunctionTok{element\_text}\NormalTok{(}\AttributeTok{hjust =} \FloatTok{0.5}\NormalTok{)  }\CommentTok{\# Center the title horizontally}
\NormalTok{  )}

\NormalTok{plot5 }\OtherTok{\textless{}{-}} \FunctionTok{ggplot}\NormalTok{(}\AttributeTok{data =}\NormalTok{ corr\_df, }\FunctionTok{aes}\NormalTok{(}\AttributeTok{x =}\NormalTok{ x, }\AttributeTok{y =}\NormalTok{ y5)) }\SpecialCharTok{+}
  \FunctionTok{geom\_point}\NormalTok{(}\AttributeTok{alpha =} \FloatTok{0.5}\NormalTok{) }\SpecialCharTok{+}
  \FunctionTok{geom\_smooth}\NormalTok{(}\AttributeTok{method =} \StringTok{"lm"}\NormalTok{, }\AttributeTok{se =}\NormalTok{ F) }\SpecialCharTok{+}
  \FunctionTok{labs}\NormalTok{(}\AttributeTok{subtitle =} \FunctionTok{round}\NormalTok{(}\FunctionTok{cor}\NormalTok{(corr\_df}\SpecialCharTok{$}\NormalTok{x, corr\_df}\SpecialCharTok{$}\NormalTok{y5), }\AttributeTok{digits =} \DecValTok{1}\NormalTok{)) }\SpecialCharTok{+}
  \FunctionTok{theme\_minimal}\NormalTok{() }\SpecialCharTok{+}
  \FunctionTok{theme}\NormalTok{(}
    \AttributeTok{aspect.ratio =} \DecValTok{1}\NormalTok{,}
    \AttributeTok{axis.text =} \FunctionTok{element\_blank}\NormalTok{(),}
    \AttributeTok{axis.ticks =} \FunctionTok{element\_blank}\NormalTok{(),}
    \AttributeTok{axis.title.y =} \FunctionTok{element\_blank}\NormalTok{(),}
    \AttributeTok{plot.title =} \FunctionTok{element\_text}\NormalTok{(}\AttributeTok{hjust =} \FloatTok{0.5}\NormalTok{)  }\CommentTok{\# Center the title horizontally}
\NormalTok{  )}

\CommentTok{\# Using patchwork to arrange plots side by side}
\NormalTok{(plot1 }\SpecialCharTok{|}\NormalTok{ plot2 }\SpecialCharTok{|}\NormalTok{ plot3 }\SpecialCharTok{|}\NormalTok{ plot4 }\SpecialCharTok{|}\NormalTok{ plot5) }\SpecialCharTok{+} \FunctionTok{plot\_layout}\NormalTok{(}\AttributeTok{ncol =} \DecValTok{5}\NormalTok{)}
\end{Highlighting}
\end{Shaded}

\includegraphics{linear_regression_multiple_files/figure-pdf/unnamed-chunk-5-1.pdf}

Now that we know linear regression we are better equipped to understand
what correlation really means: two variables are highly correlated if
they \emph{linearly describe} each other well. They correlate perfectly
if all the points fall in a line, as for correlation -1 and 1. In this
case, estimating a linear regression model will yield a very small
\(p\)-value, while fitting a regression model where the correlation is
closer to zero is more likely to yield an non-significant linear
relationship. It is important to remember that the correlation does not
tell us anything about the strength of the relationship, i.e.~the size
of the coefficient, but only how well the points fall on the regression
line and if the coefficient is positive or negative.

Correlation sounds like a good candidate to assert which one of our two
models are the best at describing our outcome, and indeed correlation
can be used for this in simple linear regression, but for multiple
linear regression there is no one correlation for the whole model.
Instead, the \emph{coefficient of determination}, or R\(^2\) pronounced
``R-squared'' is used. In the case of simple linear regression it has a
simple relation ship with correlation, as R\(^2\) = Correlation\(^2\).
In both simple and multiple linear regression the interpretation of
R\(^2\) is very intuitive, which makes it a very good measure of how
well a model describes the outcome.

\subsection{\texorpdfstring{Coefficient of determination,
R\(^2\)}{Coefficient of determination, R\^{}2}}\label{coefficient-of-determination-r2}

R\(^2\) can be interpreted as ``the percentage of the variation in the
outcome the model describes''. Specifically, it is the percentage of the
variation around the mean of the outcome that the model describes. Let's
visualize this for a subset of the points for
\texttt{birth\_weight\_g\ \textasciitilde{}\ length\_cm}. Here we
highlight three points, the mean value by a horizontal dotted line, and
their total variation from the mean with vertical dashed lines.

\begin{Shaded}
\begin{Highlighting}[]
\CommentTok{\# Limits}
\NormalTok{ylims }\OtherTok{\textless{}{-}} \FunctionTok{c}\NormalTok{(}\DecValTok{3000}\NormalTok{, }\DecValTok{4000}\NormalTok{)}
\CommentTok{\# Predictions}
\NormalTok{model }\OtherTok{\textless{}{-}}\NormalTok{ birth }\SpecialCharTok{|\textgreater{}} \FunctionTok{lm}\NormalTok{(}\AttributeTok{formula =}\NormalTok{ birth\_weight\_g }\SpecialCharTok{\textasciitilde{}}\NormalTok{ length\_cm)}
\CommentTok{\# Subset points to show}
\FunctionTok{set.seed}\NormalTok{(}\DecValTok{1337}\NormalTok{)}
\NormalTok{subset\_birth }\OtherTok{\textless{}{-}}\NormalTok{ birth }\SpecialCharTok{|\textgreater{}} 
  \FunctionTok{mutate}\NormalTok{(}
    \AttributeTok{pred =} \FunctionTok{predict}\NormalTok{(model),}
    \AttributeTok{resid =} \FunctionTok{resid}\NormalTok{(model)) }\SpecialCharTok{|\textgreater{}} 
  \FunctionTok{filter}\NormalTok{(}
\NormalTok{    (length\_cm }\SpecialCharTok{\textless{}} \FunctionTok{mean}\NormalTok{(birth}\SpecialCharTok{$}\NormalTok{length\_cm) }\SpecialCharTok{{-}} \DecValTok{5} \SpecialCharTok{\&}\NormalTok{ birth\_weight\_g }\SpecialCharTok{\textless{}}\NormalTok{ pred }\SpecialCharTok{{-}} \DecValTok{100}\NormalTok{) }\SpecialCharTok{|}
\NormalTok{    (length\_cm }\SpecialCharTok{\textgreater{}} \FunctionTok{mean}\NormalTok{(birth}\SpecialCharTok{$}\NormalTok{length\_cm) }\SpecialCharTok{+} \DecValTok{5} \SpecialCharTok{\&}\NormalTok{ birth\_weight\_g }\SpecialCharTok{\textgreater{}}\NormalTok{ pred }\SpecialCharTok{+} \DecValTok{100}\NormalTok{)) }\SpecialCharTok{|\textgreater{}}  
  \FunctionTok{filter}\NormalTok{(birth\_weight\_g }\SpecialCharTok{\textless{}} \DecValTok{4000} \SpecialCharTok{\&}\NormalTok{ birth\_weight\_g }\SpecialCharTok{\textgreater{}} \DecValTok{3000}\NormalTok{) }\SpecialCharTok{|\textgreater{}} 
  \FunctionTok{sample\_n}\NormalTok{(}\DecValTok{3}\NormalTok{)}
\CommentTok{\# Calculate mean of birth weight}
\NormalTok{mean\_bw }\OtherTok{\textless{}{-}}\NormalTok{ birth }\SpecialCharTok{|\textgreater{}} \FunctionTok{pull}\NormalTok{(birth\_weight\_g) }\SpecialCharTok{|\textgreater{}} \FunctionTok{mean}\NormalTok{()}

\NormalTok{variation\_plot }\OtherTok{\textless{}{-}}\NormalTok{ birth }\SpecialCharTok{|\textgreater{}} 
  \FunctionTok{ggplot}\NormalTok{(}\FunctionTok{aes}\NormalTok{(}
    \AttributeTok{x =}\NormalTok{ length\_cm, }
    \AttributeTok{y =}\NormalTok{ birth\_weight\_g)) }\SpecialCharTok{+} 
  \FunctionTok{geom\_hline}\NormalTok{(}
    \AttributeTok{yintercept =}\NormalTok{ mean\_bw,}
    \AttributeTok{linetype =} \DecValTok{3}\NormalTok{,}
    \AttributeTok{linewidth =} \DecValTok{1}\NormalTok{) }\SpecialCharTok{+}
  \FunctionTok{geom\_point}\NormalTok{(}\AttributeTok{alpha =} \FloatTok{0.1}\NormalTok{) }\SpecialCharTok{+} 
  \FunctionTok{geom\_smooth}\NormalTok{(}\AttributeTok{method =} \StringTok{"lm"}\NormalTok{, }\AttributeTok{se =} \ConstantTok{FALSE}\NormalTok{) }\SpecialCharTok{+}
  \FunctionTok{geom\_segment}\NormalTok{( }\CommentTok{\# Total variation}
    \AttributeTok{data =}\NormalTok{ subset\_birth,}
    \FunctionTok{aes}\NormalTok{(}
      \AttributeTok{x =}\NormalTok{ length\_cm }\SpecialCharTok{+} \FloatTok{0.1}\NormalTok{,}
      \AttributeTok{xend =}\NormalTok{ length\_cm }\SpecialCharTok{+} \FloatTok{0.1}\NormalTok{,}
      \AttributeTok{y =}\NormalTok{  mean\_bw,}
      \AttributeTok{yend =}\NormalTok{ birth\_weight\_g,}
      \AttributeTok{col =} \StringTok{"Total"}\NormalTok{),}
    \AttributeTok{linetype =} \DecValTok{2}\NormalTok{,}
    \AttributeTok{linewidth =} \DecValTok{1}\NormalTok{,}
    \AttributeTok{alpha =} \DecValTok{1}\NormalTok{) }\SpecialCharTok{+}
  \FunctionTok{geom\_point}\NormalTok{(}\AttributeTok{data =}\NormalTok{ subset\_birth, }\AttributeTok{size =} \DecValTok{4}\NormalTok{) }\SpecialCharTok{+} 
  \FunctionTok{labs}\NormalTok{(}
    \AttributeTok{x =} \StringTok{"Length (cm)"}\NormalTok{,}
    \AttributeTok{y =} \StringTok{"Birth weight (g)"}\NormalTok{,}
    \AttributeTok{col =} \StringTok{"Variation"}\NormalTok{) }\SpecialCharTok{+}
  \FunctionTok{scale\_color\_manual}\NormalTok{(}
    \AttributeTok{values =} \FunctionTok{c}\NormalTok{(}
      \StringTok{"Total"} \OtherTok{=} \StringTok{"black"}\NormalTok{,}
      \StringTok{"Explained"} \OtherTok{=} \StringTok{"darkgreen"}\NormalTok{, }
      \StringTok{"Unexplained"} \OtherTok{=} \StringTok{"darkorange"}\NormalTok{),}
    \AttributeTok{breaks =} \FunctionTok{c}\NormalTok{(}
      \StringTok{"Total"}\NormalTok{, }\StringTok{"Explained"}\NormalTok{, }\StringTok{"Unexplained"}\NormalTok{)) }\SpecialCharTok{+}
  \FunctionTok{coord\_cartesian}\NormalTok{(}\AttributeTok{ylim =}\NormalTok{ ylims)}
\NormalTok{variation\_plot}
\end{Highlighting}
\end{Shaded}

\includegraphics{linear_regression_multiple_files/figure-pdf/unnamed-chunk-6-1.pdf}

We see the total variation as the distance from the mean value to the
points in the y-direction. If all points align on the mean value, there
is no variation to explain, but as long as there is some variation, it
is possible that it can be explained by some other variable, in this
case \texttt{length\_cm}. If the regression line (in blue) would have
been completely flat, \texttt{length\_cm} wouldn't have explained
anything more that then mean of \texttt{birth\_weight\_g}. But how much
of the variation does the regression line explain? We can illustrate
this for the higlighted points by visualizing the distance from the mean
to the regression line (explained variation) and from the prediction
line to the observed value, i.e.~the point (un-explained variation).

\begin{Shaded}
\begin{Highlighting}[]
\NormalTok{variation\_plot }\SpecialCharTok{+} 
  \FunctionTok{geom\_segment}\NormalTok{( }\CommentTok{\# Unexplained variation}
    \AttributeTok{data =}\NormalTok{ subset\_birth,}
    \FunctionTok{aes}\NormalTok{(}
      \AttributeTok{x =}\NormalTok{ length\_cm,}
      \AttributeTok{xend =}\NormalTok{ length\_cm,}
      \AttributeTok{y =}\NormalTok{  pred,}
      \AttributeTok{yend =}\NormalTok{ birth\_weight\_g,}
      \AttributeTok{col =} \StringTok{"Unexplained"}\NormalTok{),}
    \AttributeTok{linewidth =} \FloatTok{1.5}\NormalTok{,}
    \AttributeTok{alpha =} \FloatTok{0.7}
\NormalTok{    ) }\SpecialCharTok{+}
  \FunctionTok{geom\_segment}\NormalTok{( }\CommentTok{\# Explained variation}
    \AttributeTok{data =}\NormalTok{ subset\_birth,}
    \FunctionTok{aes}\NormalTok{(}
      \AttributeTok{x =}\NormalTok{ length\_cm,}
      \AttributeTok{xend =}\NormalTok{ length\_cm,}
      \AttributeTok{y =}\NormalTok{  mean\_bw,}
      \AttributeTok{yend =}\NormalTok{ pred,}
      \AttributeTok{col =} \StringTok{"Explained"}\NormalTok{),}
    \AttributeTok{linewidth =} \FloatTok{1.5}\NormalTok{,}
    \AttributeTok{alpha =} \FloatTok{0.7}
\NormalTok{    ) }\SpecialCharTok{+}
  \FunctionTok{geom\_point}\NormalTok{(}\AttributeTok{data =}\NormalTok{ subset\_birth, }\AttributeTok{size =} \DecValTok{4}\NormalTok{)}
\end{Highlighting}
\end{Shaded}

\includegraphics{linear_regression_multiple_files/figure-pdf/unnamed-chunk-7-1.pdf}

In this figure we see that part of the variation is explained by the
line (green), but not all (orange). R\(^2\) is calculated by summing
together the squared explained variation (green) and dividing by the
squared total variation (dashed, black) for all points, mathematically
R\(^2\) = \(\sum_{i}\) Explained\(_i^2\) / Total\(_i^2\), where \(i\)
indicates each observation, i.e.~each point. If the explained variation
is the same as the total variation, R\(^2\) is 1 and the predictor
explains 100\% of the variation in the response (compare with the
correlation). If the predictor doesn't explain anything, i.e.~the line
is flat, R\(^2\) is 0. You are not expected to be able to calculate
R\(^2\), as we will let \texttt{lm} do that for us, but you are expected
to be able to interpret the value of R\(^2\).

For a more detailed explanation of R\(^2\),
\href{https://en.wikipedia.org/wiki/Coefficient_of_determination}{Wikipedia}
offers a good alternative with both visualizations and mathematical
formulas.

\subsection{Selecting between two simple linear regression
models}\label{selecting-between-two-simple-linear-regression-models}

Armed with R\(^2\) as a measure of how well a model describes the
variation in the data, we can now compare the two models using either
\texttt{length\_cm} or \texttt{weight\_kg} to explain the variation in
\texttt{birth\_weight\_g}. It turns out that we have been calculating
R\(^2\)-values with every \texttt{lm} call we have made, and we can see
it by calling the \texttt{summary} function. We start by looking at the
model of \texttt{birth\_weight\_g\ \textasciitilde{}\ weight\_kg}.

\begin{Shaded}
\begin{Highlighting}[]
\CommentTok{\# Mother weight}
\NormalTok{birth }\SpecialCharTok{|\textgreater{}} 
  \FunctionTok{lm}\NormalTok{(}\AttributeTok{formula =}\NormalTok{ birth\_weight\_g }\SpecialCharTok{\textasciitilde{}}\NormalTok{ weight\_kg) }\SpecialCharTok{|\textgreater{}} 
  \FunctionTok{summary}\NormalTok{()}
\end{Highlighting}
\end{Shaded}

\begin{verbatim}

Call:
lm(formula = birth_weight_g ~ weight_kg, data = birth)

Residuals:
     Min       1Q   Median       3Q      Max 
-1259.97  -270.84     2.46   296.60  1164.35 

Coefficients:
            Estimate Std. Error t value Pr(>|t|)    
(Intercept) 2974.058    179.757  16.545  < 2e-16 ***
weight_kg      8.032      2.748   2.923  0.00382 ** 
---
Signif. codes:  0 '***' 0.001 '**' 0.01 '*' 0.05 '.' 0.1 ' ' 1

Residual standard error: 458.4 on 228 degrees of freedom
Multiple R-squared:  0.03612,   Adjusted R-squared:  0.03189 
F-statistic: 8.543 on 1 and 228 DF,  p-value: 0.003818
\end{verbatim}

In the second line from the bottom we find ``Multiple R-squared'' and
find that it is 0.03612. We can interpret this as ``The weight of the
mother explains 3.6\% of the variation in birth weight of infants''.
What about \texttt{length\_cm}?

\begin{Shaded}
\begin{Highlighting}[]
\CommentTok{\# Mother length}
\NormalTok{birth }\SpecialCharTok{|\textgreater{}} 
  \FunctionTok{lm}\NormalTok{(}\AttributeTok{formula =}\NormalTok{ birth\_weight\_g }\SpecialCharTok{\textasciitilde{}}\NormalTok{ length\_cm) }\SpecialCharTok{|\textgreater{}} 
  \FunctionTok{summary}\NormalTok{()}
\end{Highlighting}
\end{Shaded}

\begin{verbatim}

Call:
lm(formula = birth_weight_g ~ length_cm, data = birth)

Residuals:
     Min       1Q   Median       3Q      Max 
-1253.71  -291.45    39.43   274.12  1291.05 

Coefficients:
            Estimate Std. Error t value Pr(>|t|)    
(Intercept)  108.034    767.698   0.141    0.888    
length_cm     20.251      4.591   4.411 1.59e-05 ***
---
Signif. codes:  0 '***' 0.001 '**' 0.01 '*' 0.05 '.' 0.1 ' ' 1

Residual standard error: 448.1 on 228 degrees of freedom
Multiple R-squared:  0.07863,   Adjusted R-squared:  0.07459 
F-statistic: 19.46 on 1 and 228 DF,  p-value: 1.585e-05
\end{verbatim}

The mothers length explains a little bit more of the variation in infant
birth weight, about 7.8\%. For convenience, we may want to output only
the R\(^2\) value from our regression. Again we can use the dollar
operator.

\begin{Shaded}
\begin{Highlighting}[]
\CommentTok{\# Mother weight}
\NormalTok{weight\_model\_summary }\OtherTok{\textless{}{-}}\NormalTok{ birth }\SpecialCharTok{|\textgreater{}} 
  \FunctionTok{lm}\NormalTok{(}\AttributeTok{formula =}\NormalTok{ birth\_weight\_g }\SpecialCharTok{\textasciitilde{}}\NormalTok{ weight\_kg) }\SpecialCharTok{|\textgreater{}} 
  \FunctionTok{summary}\NormalTok{()}
\NormalTok{weight\_model\_summary}\SpecialCharTok{$}\NormalTok{r.squared}
\end{Highlighting}
\end{Shaded}

\begin{verbatim}
[1] 0.03611514
\end{verbatim}

\begin{Shaded}
\begin{Highlighting}[]
\CommentTok{\# Mother length}
\NormalTok{length\_model\_summary }\OtherTok{\textless{}{-}}\NormalTok{ birth }\SpecialCharTok{|\textgreater{}} 
  \FunctionTok{lm}\NormalTok{(}\AttributeTok{formula =}\NormalTok{ birth\_weight\_g }\SpecialCharTok{\textasciitilde{}}\NormalTok{ length\_cm) }\SpecialCharTok{|\textgreater{}} 
  \FunctionTok{summary}\NormalTok{()}
\NormalTok{length\_model\_summary}\SpecialCharTok{$}\NormalTok{r.squared}
\end{Highlighting}
\end{Shaded}

\begin{verbatim}
[1] 0.07863332
\end{verbatim}

These results may be underwhelming, as the predictors only explain a
small percentage of the total variation. This motivates us to use
multiple linear regression, capable of including multiple predictors at
the same time to create models with more explanatory power than simple
linear regression.

\subsection{Exercises}\label{exercises-17}

\begin{enumerate}
\def\labelenumi{\arabic{enumi}.}
\tightlist
\item
  What R\(^2\) value do you expect for the
  \texttt{birth\_weight\_g\ \textasciitilde{}\ age} model? Is it higher
  or lower than \texttt{weight\_kg} or \texttt{length\_cm}? Use
  \texttt{summary} to find out and confirm what you see by a scatter
  plot with a line using \texttt{geom\_smooth}.
\item
  Fit a multiple linear regression model using
  \texttt{formula\ =\ birth\_weight\_g\ \textasciitilde{}\ weight\_kg\ +\ length\_cm}.
  Did R\(^2\) increase?
\item
  Run the following code (requires the \texttt{plotly} package) to
  create an interactive 3D visualization of
  \texttt{birth\_weight\_g\ \textasciitilde{}\ weight\_kg\ +\ weight\_cm}.
  Un-comment the three different ``VIEW:'' sections, one at a time, to
  view the plot from different angles. Try to understand how the surface
  represents multiple linear regression. Can you visualize more
  dimensions? Note: You don't need to understand the code, rather try to
  understand the content. Hint: You may zoom in and out by scrolling
  while hoovering over the plot.
\end{enumerate}

\begin{Shaded}
\begin{Highlighting}[]
\CommentTok{\#install.packages("plotly")}
\FunctionTok{library}\NormalTok{(plotly)}

\CommentTok{\# Fit the multiple linear regression model}
\NormalTok{model }\OtherTok{\textless{}{-}}\NormalTok{ birth }\SpecialCharTok{|\textgreater{}} \FunctionTok{lm}\NormalTok{(}\AttributeTok{formula =}\NormalTok{ birth\_weight\_g }\SpecialCharTok{\textasciitilde{}}\NormalTok{ weight\_kg }\SpecialCharTok{+}\NormalTok{ length\_cm)}

\CommentTok{\# Create a grid of values for weight\_kg and length\_cm}
\NormalTok{weight\_kg\_range }\OtherTok{\textless{}{-}} \FunctionTok{seq}\NormalTok{(}\FunctionTok{min}\NormalTok{(birth}\SpecialCharTok{$}\NormalTok{weight\_kg), }\FunctionTok{max}\NormalTok{(birth}\SpecialCharTok{$}\NormalTok{weight\_kg), }\AttributeTok{length.out =} \DecValTok{50}\NormalTok{)}
\NormalTok{length\_cm\_range }\OtherTok{\textless{}{-}} \FunctionTok{seq}\NormalTok{(}\FunctionTok{min}\NormalTok{(birth}\SpecialCharTok{$}\NormalTok{length\_cm), }\FunctionTok{max}\NormalTok{(birth}\SpecialCharTok{$}\NormalTok{length\_cm), }\AttributeTok{length.out =} \DecValTok{50}\NormalTok{)}
\NormalTok{grid }\OtherTok{\textless{}{-}} \FunctionTok{expand.grid}\NormalTok{(}\AttributeTok{weight\_kg =}\NormalTok{ weight\_kg\_range, }\AttributeTok{length\_cm =}\NormalTok{ length\_cm\_range)}

\CommentTok{\# Predict birth\_weight\_g for the grid}
\NormalTok{grid}\SpecialCharTok{$}\NormalTok{birth\_weight\_g }\OtherTok{\textless{}{-}} \FunctionTok{predict}\NormalTok{(model, }\AttributeTok{newdata =}\NormalTok{ grid)}

\CommentTok{\# Create 3D scatter plot}
\NormalTok{p }\OtherTok{\textless{}{-}} \FunctionTok{plot\_ly}\NormalTok{(}
\NormalTok{  birth, }
  \AttributeTok{x =} \SpecialCharTok{\textasciitilde{}}\NormalTok{length\_cm, }
  \AttributeTok{y =} \SpecialCharTok{\textasciitilde{}}\NormalTok{weight\_kg, }
  \AttributeTok{z =} \SpecialCharTok{\textasciitilde{}}\NormalTok{birth\_weight\_g, }
  \AttributeTok{type =} \StringTok{\textquotesingle{}scatter3d\textquotesingle{}}\NormalTok{, }
  \AttributeTok{mode =} \StringTok{\textquotesingle{}markers\textquotesingle{}}\NormalTok{, }
  \AttributeTok{size =} \DecValTok{1}\NormalTok{)}

\CommentTok{\# Add regression surface}
\NormalTok{p }\OtherTok{\textless{}{-}}\NormalTok{ p }\SpecialCharTok{|\textgreater{}} 
  \FunctionTok{add\_surface}\NormalTok{(}
    \AttributeTok{x =} \SpecialCharTok{\textasciitilde{}}\NormalTok{length\_cm\_range, }
    \AttributeTok{y =} \SpecialCharTok{\textasciitilde{}}\NormalTok{weight\_kg\_range, }
    \AttributeTok{z =} \SpecialCharTok{\textasciitilde{}}\FunctionTok{matrix}\NormalTok{(grid}\SpecialCharTok{$}\NormalTok{birth\_weight\_g, }\AttributeTok{nrow =} \DecValTok{50}\NormalTok{), }
    \AttributeTok{opacity =} \FloatTok{0.5}\NormalTok{)}

\NormalTok{p }\OtherTok{\textless{}{-}}\NormalTok{ p }\SpecialCharTok{|\textgreater{}} 
  \FunctionTok{colorbar}\NormalTok{(}\AttributeTok{title =} \StringTok{"\textless{}b\textgreater{}Birth weight (g)\textless{}/b\textgreater{}"}\NormalTok{) }\SpecialCharTok{|\textgreater{}}
  \FunctionTok{layout}\NormalTok{(}
    \AttributeTok{scene =} \FunctionTok{list}\NormalTok{(}
      \AttributeTok{xaxis =} \FunctionTok{list}\NormalTok{(}\AttributeTok{title=}\StringTok{"Weight (kg)"}\NormalTok{), }
      \AttributeTok{yaxis =} \FunctionTok{list}\NormalTok{(}\AttributeTok{title=}\StringTok{"Length (cm)"}\NormalTok{),}
      \AttributeTok{zaxis =} \FunctionTok{list}\NormalTok{(}\AttributeTok{title=}\StringTok{"Birth weight (g)"}\NormalTok{)))}

\CommentTok{\# VIEW: birth\_weight\_g \textasciitilde{} length\_cm perspective}
\CommentTok{\#p \textless{}{-} p |\textgreater{} }
\CommentTok{\#  layout(scene = list(camera = list(eye = list(x = {-}2.25, y = 0, z = 0))))}

\CommentTok{\# VIEW: birth\_weight\_g \textasciitilde{} weight\_kg perspective}
\CommentTok{\#p \textless{}{-} p |\textgreater{} }
\CommentTok{\#  layout(scene = list(camera = list(eye = list(x = 0, y = 2.25, z = 0))))}

\CommentTok{\# VIEW: weight\_kg \textasciitilde{} length\_cm perspective}
\CommentTok{\#p \textless{}{-} p |\textgreater{} }
\CommentTok{\#  layout(scene = list(camera = list(eye = list(x = 0, y = 0, z = 2.25))))}

\CommentTok{\# Show the plot}
\NormalTok{p}
\end{Highlighting}
\end{Shaded}

\section{Fitting multiple linear regression
models}\label{fitting-multiple-linear-regression-models}

Now that we can compare models we are ready to jump into multiple linear
regression. Fitting multiple linear regression models in R is hardly any
more difficult than fitting simple linear models, just add more
variables using the \texttt{+} sign in the \texttt{formula}. We start by
adding all the variables to our dataset, including the categorical
'\texttt{sex}.

\subsection{Naive approach}\label{naive-approach}

More predictors will explain more of the variation, so we start by
naively using all possible predictors in our dataset. To include all
variables except the outcome we can use dot, \texttt{.}, in the formula.
It stands for ``all variables in the \texttt{data.frame} except for the
outcome'', which is perfect when we want to include all possible
predictors.

\begin{Shaded}
\begin{Highlighting}[]
\NormalTok{model }\OtherTok{\textless{}{-}}\NormalTok{ birth }\SpecialCharTok{|\textgreater{}} 
  \FunctionTok{lm}\NormalTok{(}\AttributeTok{formula =}\NormalTok{ birth\_weight\_g }\SpecialCharTok{\textasciitilde{}}\NormalTok{ .)}
\NormalTok{model }\SpecialCharTok{|\textgreater{}} \FunctionTok{summary}\NormalTok{()}
\end{Highlighting}
\end{Shaded}

\begin{verbatim}

Call:
lm(formula = birth_weight_g ~ ., data = birth)

Residuals:
    Min      1Q  Median      3Q     Max 
-986.78 -245.33    3.51  247.62 1101.82 

Coefficients:
                        Estimate Std. Error t value Pr(>|t|)    
(Intercept)            -3507.946   4552.540  -0.771  0.44179    
pregnancy_duration_wks   175.974     20.442   8.609 1.35e-15 ***
sexMale                  152.066     50.903   2.987  0.00313 ** 
age                        2.034      5.282   0.385  0.70050    
length_cm                 -2.782     26.820  -0.104  0.91749    
weight_kg                 29.015     35.086   0.827  0.40914    
bmi                      -66.988     96.855  -0.692  0.48989    
---
Signif. codes:  0 '***' 0.001 '**' 0.01 '*' 0.05 '.' 0.1 ' ' 1

Residual standard error: 378.1 on 223 degrees of freedom
Multiple R-squared:  0.3583,    Adjusted R-squared:  0.3411 
F-statistic: 20.76 on 6 and 223 DF,  p-value: < 2.2e-16
\end{verbatim}

We start by looking at R\(^2\) and see that it has improved a lot, now
35.8\% percent of the variation in \texttt{birth\_weight\_g} is
explained. But looking at the coefficient estimates and \(p\)-values
(``Pr(\textgreater\textbar t\textbar)'') something seems off. The
coefficients for \texttt{length\_cm} and \texttt{bmi} are now negative,
despite both being positive in their respective simple linear
regressions, and \texttt{length\_cm} and \texttt{weight\_kg} are no
longer significant. We have improved our R\(^2\), but sacrificed the
\emph{interpretability} of our model, we can no longer make any sense of
it. Why does this happen? To answer that we need to better understand
how adding more predictors to the model affects each predictor in the
model.

\subsection{Interpreting multiple linear
regression}\label{interpreting-multiple-linear-regression}

When we perform multiple linear regression the significance and
coefficient estimate of each predictor is \emph{dependent} on the other
variables included in the model. This may change the coefficient
estimates and \(p\)-values for a predictor when using it in a multiple
linear regression model compared to a simple regression model. The
hypotheses in multiple linear regression highlights this. For example,
the hypotheses for the coefficient of \texttt{length\_cm} in the model
including all potential predictors are

\begin{quote}
\textbf{\(H_0\)}: There is \emph{no} linear relationship between the
mothers length and the birth weight of the child, given that pregnancy
duration, sex, mother age, weight, and BMI is also part of the model.
\end{quote}

\begin{quote}
\textbf{\(H_a\)}: There \emph{is} a linear relationship between the
mothers length and the birth weight of the child, given that pregnancy
duration, \ldots, and BMI are also part of the model.
\end{quote}

Both the values of the coefficients and the \(p\)-values depend on the
other variables included in the model. The interpretation of the
coefficients follows the hypotheses in depending on the other variables
in the model:

\begin{quote}
When the mothers length increases by 1 cm the birth weight of the child
decreases by 2.782 grams on average, \emph{given that the other
variables in the model are held fixed}.
\end{quote}

Looking at our model, this statement becomes an issue: BMI is a function
of both weight and length, BMI = Weight / Length\(^2\). The
interpretation of the coefficient for BMI is then

\begin{quote}
When the mothers BMI increases by 1 the birth weight decreases by
\$\sim\$67 g, given that the other variables in the model,
\emph{including weight and length}, are held fixed.
\end{quote}

But this cannot be, as if length and weight are held fixed, BMI cannot
change! The interpretation of the model breaks down when all three of
these variables are included in the model, causing the un-intuitive
\(p\) values and coefficient values. When the predictors are (strongly)
correlated, as in this case with \texttt{bmi}, \texttt{weight\_kg}, and
\texttt{length\_cm}, the model suffers from \emph{multicollinearity},
invalidating the inference of the model, i.e.~interpretation of
coefficients and \(p\)-values.

\subsubsection{Exercises}\label{exercises-18}

\begin{enumerate}
\def\labelenumi{\arabic{enumi}.}
\tightlist
\item
  The code below simulates data for ice cream sales and snake bites
  during the summer months. Copy the code and visualize the relationship
  between \texttt{ice\_cream\_sales} and \texttt{snake\_bites} (using
  \texttt{geom\_point} and \texttt{geom\_smooth} with
  \texttt{method\ =\ "lm"}) and fit a linear regression model of
  \texttt{ice\_cream\_sales\ \textasciitilde{}\ snake\_bites}. Do we
  have a linear relationship? Is it significant?
\end{enumerate}

\begin{Shaded}
\begin{Highlighting}[]
\CommentTok{\# Simulating data}
\FunctionTok{set.seed}\NormalTok{(}\DecValTok{42}\NormalTok{) }\CommentTok{\# Important to include this!}
\NormalTok{n\_days }\OtherTok{\textless{}{-}} \DecValTok{153} 
\NormalTok{n\_years }\OtherTok{\textless{}{-}} \DecValTok{5}
\CommentTok{\# Daily observations stored as a decimal number of the month}
\NormalTok{month }\OtherTok{\textless{}{-}} \FunctionTok{rep}\NormalTok{(}\FunctionTok{seq}\NormalTok{(}\DecValTok{4}\NormalTok{, }\DecValTok{8}\NormalTok{, }\AttributeTok{length.out =}\NormalTok{ n\_days), n\_years) }
\NormalTok{ice\_cream\_sales }\OtherTok{\textless{}{-}} \DecValTok{2000} \SpecialCharTok{+} \DecValTok{3000} \SpecialCharTok{*} \FunctionTok{sin}\NormalTok{(pi }\SpecialCharTok{*}\NormalTok{ month }\SpecialCharTok{/} \DecValTok{14}\NormalTok{) }\SpecialCharTok{+} \FunctionTok{rnorm}\NormalTok{(}\FunctionTok{length}\NormalTok{(month), }\DecValTok{0}\NormalTok{, }\DecValTok{200}\NormalTok{)}
\NormalTok{snake\_bites }\OtherTok{\textless{}{-}} \DecValTok{20} \SpecialCharTok{*} \FunctionTok{sin}\NormalTok{(pi }\SpecialCharTok{*}\NormalTok{ month }\SpecialCharTok{/} \DecValTok{14}\NormalTok{) }\SpecialCharTok{+} \FunctionTok{rnorm}\NormalTok{(}\FunctionTok{length}\NormalTok{(month), }\DecValTok{0}\NormalTok{, }\DecValTok{8}\NormalTok{)}

\NormalTok{ice\_cream }\OtherTok{\textless{}{-}} \FunctionTok{tibble}\NormalTok{(}\AttributeTok{month =}\NormalTok{ month, }\AttributeTok{ice\_cream\_sales =}\NormalTok{ ice\_cream\_sales, }\AttributeTok{snake\_bites =}\NormalTok{ snake\_bites)}
\end{Highlighting}
\end{Shaded}

\begin{enumerate}
\def\labelenumi{\arabic{enumi}.}
\setcounter{enumi}{1}
\tightlist
\item
  Would you recommend an ice cream company to increase snake bites in
  order to increase ice cream sales? In other words, do you think that
  the effect of snake bites on ice cream sales is a causal effect?
  Causal effect in this case means that if we increase the number of
  snake bites we will also increase the number of ice cream sales.
\item
  Now fit a model including \texttt{month} as well,
  i.e.~\texttt{ice\_cream\_sales\ \textasciitilde{}\ .} and state the
  hypotheses for the coefficient for \texttt{snake\_bites}. What
  changed? Is the coefficient for \texttt{snake\_bites} still
  significantly different from 0 with significance level 0.05?
\end{enumerate}

\subsection{Multicollinearity}\label{multicollinearity}

If the coefficients and \(p\)-values change drastically from simple
linear regression models or from visualizations in scatter plots or
correlation plots we should suspect multicollinearity in our model. It
is also possible to look at the potential predictors in the dataset and
their correlation in a scatter plot or in the correlation plot to try to
anticipate multicollinearity.

\begin{Shaded}
\begin{Highlighting}[]
\NormalTok{p1 }\OtherTok{\textless{}{-}}\NormalTok{ birth }\SpecialCharTok{|\textgreater{}} 
  \FunctionTok{ggplot}\NormalTok{(}\FunctionTok{aes}\NormalTok{(}\AttributeTok{x =}\NormalTok{ weight\_kg, }\AttributeTok{y =}\NormalTok{ bmi)) }\SpecialCharTok{+}
  \FunctionTok{geom\_point}\NormalTok{() }\SpecialCharTok{+}
  \FunctionTok{geom\_smooth}\NormalTok{(}\AttributeTok{method =} \StringTok{"lm"}\NormalTok{, }\AttributeTok{se =} \ConstantTok{FALSE}\NormalTok{)}
\NormalTok{p2 }\OtherTok{\textless{}{-}}\NormalTok{ birth }\SpecialCharTok{|\textgreater{}} 
  \FunctionTok{select}\NormalTok{(bmi, weight\_kg, length\_cm) }\SpecialCharTok{|\textgreater{}} 
  \FunctionTok{select}\NormalTok{(}\FunctionTok{where}\NormalTok{(is.numeric)) }\SpecialCharTok{|\textgreater{}} 
  \FunctionTok{cor}\NormalTok{() }\SpecialCharTok{|\textgreater{}} 
  \FunctionTok{ggcorrplot}\NormalTok{()}

\NormalTok{p1 }\SpecialCharTok{+}\NormalTok{ p2}
\end{Highlighting}
\end{Shaded}

\begin{verbatim}
`geom_smooth()` using formula = 'y ~ x'
\end{verbatim}

\includegraphics{linear_regression_multiple_files/figure-pdf/unnamed-chunk-14-1.pdf}

We see a strong correlation between \texttt{bmi} and \texttt{weight\_kg}
and moderate correlation between \texttt{weight\_kg} and
\texttt{length\_cm}. To assess multicollinearity in a multiple
regression model we can use the Variance Inflation Factor (VIF) which
can be calculated using the \texttt{vif} function from the \texttt{car}
package.

\begin{Shaded}
\begin{Highlighting}[]
\FunctionTok{library}\NormalTok{(car)}
\FunctionTok{vif}\NormalTok{(model)}
\end{Highlighting}
\end{Shaded}

\begin{verbatim}
pregnancy_duration_wks                    sex                    age 
              1.045718               1.039959               1.058718 
             length_cm              weight_kg                    bmi 
             47.932159             239.468621             208.562969 
\end{verbatim}

VIF-values around 5 should alert you to potential issues with
multicollinearity, and values above 10 suggest that the model suffers
from multicollinearity. As we can see our model has VIF-values over 200
for \texttt{weight\_kg} and \texttt{bmi}! In this scenario we have
already understood why, because of the relationship between BMI and
weight and length.

Fitting a model with all possible predictors improved our R\(^2\), but
at the cost of interpretability. In addition to loosing
interpretablility, a more complex model is more prone to
\href{https://en.wikipedia.org/wiki/Overfitting}{\emph{overfitting}} to
the data in the sample, leading to worse predictions about the
population. Using our model to predict new, unobserved data points,
i.e.~birth weights in this example, is what we want to use the model for
if we are interested in prediction. So even if we don't care about the
inference and the interpretability of the model, it may be a bad idea to
include all variables. Instead we should do \emph{model selection},
looking at several models with different predictors and different number
of predictors to find a model that is both interpretable and offers good
predictive power, i.e.~high R\(^2\).

\section{Model selection}\label{model-selection}

To select a model that is both interpretable and useful for prediction
is a difficult task, and a good quote to remember by statistician
\href{https://en.wikipedia.org/wiki/All_models_are_wrong}{George Box} is
``All models are wrong, but some are useful''. When fitting models it is
up to you to argue why the model is useful. Many times it is a good idea
to use expert (domain) knowledge to motivate decisions in model
selection. If you have access to knowledge about the data, either by
your own reasoning, by knowledge about how the data was collected, or
from an expert in the domain, you can include or excluding variables
that you know are not useful. For example in the \texttt{penguins}
dataset you may assume that the year the penguin was observed will not
have an effect on the birth weight of the penguin. A reasonable
assumption if we believe that the conditions haven't changed from year
to year. If you don't have any background knowledge accessible there are
systematic methods that allow us to select which variables to include.

\subsection{Using domain knowledge}\label{using-domain-knowledge}

Returning to our birth weight model with all variables in the dataset as
predictors, we will now try to simplify our model to increase
interpretability. We learned before that the relationship between the
outcome and a predictor may become insignificant when we include another
variable in the model and that our model suffers from multicollinearity.
We identified the problem as being the relationship between BMI, weight,
and length. Since we know that a model including all of these variables
cannot be interpreted, we may choose to create a model removing one of
the variables. BMI is a good candidate, as it is strongly correlated
with weight. To remove a variable from the model we can use dot again in
combination with a minus sign, \texttt{-} to use all variables
\emph{except} \texttt{bmi}.

\begin{Shaded}
\begin{Highlighting}[]
\CommentTok{\# Domain expertise (de) model}
\NormalTok{model\_de }\OtherTok{\textless{}{-}}\NormalTok{ birth }\SpecialCharTok{|\textgreater{}} 
  \FunctionTok{lm}\NormalTok{(}\AttributeTok{formula =}\NormalTok{ birth\_weight\_g }\SpecialCharTok{\textasciitilde{}}\NormalTok{ . }\SpecialCharTok{{-}}\NormalTok{ bmi)}
\NormalTok{model\_summary }\OtherTok{\textless{}{-}}\NormalTok{ model\_de }\SpecialCharTok{|\textgreater{}} \FunctionTok{summary}\NormalTok{()}
\NormalTok{model\_summary}
\end{Highlighting}
\end{Shaded}

\begin{verbatim}

Call:
lm(formula = birth_weight_g ~ . - bmi, data = birth)

Residuals:
    Min      1Q  Median      3Q     Max 
-979.93 -247.04    1.17  244.13 1105.81 

Coefficients:
                        Estimate Std. Error t value Pr(>|t|)    
(Intercept)            -6573.071   1040.665  -6.316 1.43e-09 ***
pregnancy_duration_wks   176.335     20.411   8.639 1.08e-15 ***
sexMale                  148.713     50.613   2.938 0.003646 ** 
age                        2.241      5.267   0.425 0.670904    
length_cm                 15.535      4.232   3.671 0.000302 ***
weight_kg                  4.807      2.435   1.974 0.049634 *  
---
Signif. codes:  0 '***' 0.001 '**' 0.01 '*' 0.05 '.' 0.1 ' ' 1

Residual standard error: 377.7 on 224 degrees of freedom
Multiple R-squared:  0.357, Adjusted R-squared:  0.3426 
F-statistic: 24.87 on 5 and 224 DF,  p-value: < 2.2e-16
\end{verbatim}

\begin{Shaded}
\begin{Highlighting}[]
\FunctionTok{vif}\NormalTok{(model\_de)}
\end{Highlighting}
\end{Shaded}

\begin{verbatim}
pregnancy_duration_wks                    sex                    age 
              1.045037               1.030526               1.055327 
             length_cm              weight_kg 
              1.196283               1.156474 
\end{verbatim}

Now we see that there is no problem with multicollinearity (lowest
possible VIF is 1), all the coefficient estimates have the expected
sign, and all relationships are statistically significant except for
\texttt{age}. This means that \texttt{age} has no significant linear
relationship with the outcome, given the other variables in the model,
which is a good reason to remove it from the model.

\begin{Shaded}
\begin{Highlighting}[]
\NormalTok{model\_de }\OtherTok{\textless{}{-}}\NormalTok{ birth }\SpecialCharTok{|\textgreater{}} 
  \FunctionTok{lm}\NormalTok{(}\AttributeTok{formula =}\NormalTok{ birth\_weight\_g }\SpecialCharTok{\textasciitilde{}}\NormalTok{ . }\SpecialCharTok{{-}}\NormalTok{ bmi }\SpecialCharTok{{-}}\NormalTok{ age)}
\NormalTok{model\_summary }\OtherTok{\textless{}{-}}\NormalTok{ model\_de }\SpecialCharTok{|\textgreater{}} \FunctionTok{summary}\NormalTok{()}
\NormalTok{model\_summary}
\end{Highlighting}
\end{Shaded}

\begin{verbatim}

Call:
lm(formula = birth_weight_g ~ . - bmi - age, data = birth)

Residuals:
     Min       1Q   Median       3Q      Max 
-1000.23  -249.24    -6.36   249.70  1102.63 

Coefficients:
                        Estimate Std. Error t value Pr(>|t|)    
(Intercept)            -6504.850   1026.365  -6.338 1.26e-09 ***
pregnancy_duration_wks   175.099     20.167   8.683 7.97e-16 ***
sexMale                  149.730     50.465   2.967 0.003332 ** 
length_cm                 15.842      4.163   3.806 0.000182 ***
weight_kg                  4.800      2.431   1.975 0.049545 *  
---
Signif. codes:  0 '***' 0.001 '**' 0.01 '*' 0.05 '.' 0.1 ' ' 1

Residual standard error: 377 on 225 degrees of freedom
Multiple R-squared:  0.3564,    Adjusted R-squared:  0.345 
F-statistic: 31.15 on 4 and 225 DF,  p-value: < 2.2e-16
\end{verbatim}

After removing \texttt{age} we see that our R\(^2\) decreased, which
will always be the case when removing a variable from the model, as all
variables explain \emph{some} of the variation in the outcome, but if
the change in R\(^2\) is small there is a good chance that the variation
the variable explains is only due to chance and may possibly lead to
overfitting. In this case we can use ``Adjusted R-squared'', adjusted
R\(^2\) or R\(^2\)-adj, which we find next to ``Multiple R-squared'' in
the \texttt{summary} output. It is based on R\(^2\), but gives a penalty
for including more variables in the model in an attempt to promote
smaller, more easily interpretable models. Comparing with the previous
model the adjusted R\(^2\) has increased after removing \texttt{age},
indicating that \texttt{age} didn't explain enough of the variation in
\texttt{birth\_weight\_g} to warrant inclusion in the model. Adjusted
R\(^2\) can be used to remove variables even if the variable has a
significant relationship with the outcome, as this may happen when
sample sizes are large, even if the predictor explains only a small part
of the variation. In general, smaller models are easier to understand
and should be preferred when not sacrificing much explanatory power,
i.e.~R\(^2\).

In our latest model all relationships are statistically significant and
our final value of R-squared is 0.3564, compared to 0.3583 when
including all variables in the dataset. In this case we almost didn't
lose any explanatory power while making our model more interpretable by
reducing the number of variables, making sure the model doesn't suffer
from multicollinearity, and by only including statistically significant
variables. Note that when reporting the final model we use to R\(^2\)
and \emph{not} the adjusted R\(^2\).

In this example we could use our knowledge of the relationship between
BMI and weight to design a good model. If we don't have knowledge about
the variables we have to resort to systematic model selection
strategies.

\subsection{\texorpdfstring{Using \texttt{.} in
formula}{Using . in formula}}\label{sec-multlinreg-usingdot}

When using \texttt{.} in the formula for fitting a regression model,
\texttt{lm} implicitly includes all variables, including those removed
by \texttt{-}. This has two implications: the model is trained (fitted)
only on the observations where \emph{all} the variables are complete,
i.e.~all rows with no missing values and, when the model is used, it
expects a value for the removed variables as well, even if it doesn't
use them during the prediction (see below for prediction). This is good
when comparing models, as it ensures that the same data is used for all
possible models, but when fitting the final model it is better to
\emph{not} use \texttt{.} in your formula, as it ensures the model is
fitted on as much data as possible.

This example illustrates how \texttt{lm} handles missing values when
using the \texttt{.} and not. We use \texttt{palmerpenguins} because it
contains missing values and output only the R\(^2\) from the
\texttt{summary} to illustrate the difference between the models. Both
models are modeling
\texttt{body\_mass\_g\ \textasciitilde{}\ flipper\_length\_mm}, but use
different data for fitting the model.

\begin{Shaded}
\begin{Highlighting}[]
\FunctionTok{library}\NormalTok{(palmerpenguins)}
\CommentTok{\# Model fitted by using the dot operator in formula}
\NormalTok{model\_dot }\OtherTok{\textless{}{-}}\NormalTok{ penguins }\SpecialCharTok{|\textgreater{}} \FunctionTok{lm}\NormalTok{(}\AttributeTok{formula =}\NormalTok{ body\_mass\_g }\SpecialCharTok{\textasciitilde{}}\NormalTok{ . }\SpecialCharTok{{-}}\NormalTok{ species }\SpecialCharTok{{-}}\NormalTok{ island }\SpecialCharTok{{-}}\NormalTok{ bill\_length\_mm }\SpecialCharTok{{-}}\NormalTok{ bill\_depth\_mm }\SpecialCharTok{{-}}\NormalTok{ sex }\SpecialCharTok{{-}}\NormalTok{ year)}
\NormalTok{model\_dot\_summary }\OtherTok{\textless{}{-}}\NormalTok{ model\_dot }\SpecialCharTok{|\textgreater{}} \FunctionTok{summary}\NormalTok{()}
\NormalTok{model\_dot\_summary}\SpecialCharTok{$}\NormalTok{r.squared}
\end{Highlighting}
\end{Shaded}

\begin{verbatim}
[1] 0.7620922
\end{verbatim}

\begin{Shaded}
\begin{Highlighting}[]
\CommentTok{\# Model fitted by explicitly listing the predictor}
\NormalTok{model\_no\_dot }\OtherTok{\textless{}{-}}\NormalTok{ penguins }\SpecialCharTok{|\textgreater{}} \FunctionTok{lm}\NormalTok{(}\AttributeTok{formula =}\NormalTok{ body\_mass\_g }\SpecialCharTok{\textasciitilde{}}\NormalTok{ flipper\_length\_mm)}
\NormalTok{model\_no\_dot\_summary }\OtherTok{\textless{}{-}}\NormalTok{ model\_no\_dot }\SpecialCharTok{|\textgreater{}} \FunctionTok{summary}\NormalTok{()}
\NormalTok{model\_no\_dot\_summary}\SpecialCharTok{$}\NormalTok{r.squared}
\end{Highlighting}
\end{Shaded}

\begin{verbatim}
[1] 0.7589925
\end{verbatim}

We can recreate the R\(^2\) of \texttt{model\_dot} using only
\texttt{flipper\_length\_mm} as a predictor by dropping all the
observations (rows) with missing values before fitting
\texttt{model\_no\_dot}. The model is then fitted using the same data as
for \texttt{model\_dot}.

\begin{Shaded}
\begin{Highlighting}[]
\CommentTok{\# Model fitted by explicitly listing the predictor, dropping missing values first}
\NormalTok{model\_no\_dot }\OtherTok{\textless{}{-}}\NormalTok{ penguins }\SpecialCharTok{|\textgreater{}} 
  \FunctionTok{drop\_na}\NormalTok{() }\SpecialCharTok{|\textgreater{}} 
  \FunctionTok{lm}\NormalTok{(}\AttributeTok{formula =}\NormalTok{ body\_mass\_g }\SpecialCharTok{\textasciitilde{}}\NormalTok{ flipper\_length\_mm)}
\NormalTok{model\_no\_dot\_summary }\OtherTok{\textless{}{-}}\NormalTok{ model\_no\_dot }\SpecialCharTok{|\textgreater{}} \FunctionTok{summary}\NormalTok{()}
\NormalTok{model\_no\_dot\_summary}\SpecialCharTok{$}\NormalTok{r.squared}
\end{Highlighting}
\end{Shaded}

\begin{verbatim}
[1] 0.7620922
\end{verbatim}

When comparing models handling missing values explicitly or using the
dot is good. When fitting your final model you should avoid using
\texttt{.}. In these lecture notes we do not consider this as we are not
intending to use our models for practical applications, but it is good
to be aware of when applying \texttt{lm}.

\subsection{Exercises}\label{exercises-19}

\begin{enumerate}
\def\labelenumi{\arabic{enumi}.}
\tightlist
\item
  Visualize and fit a regression model of
  \texttt{bmi\ \textasciitilde{}\ weight\_kg} and confirm the functional
  relationship, i.e.~the linear relationship.
\item
  Run the code below. Why does the regression model fail? How can we fit
  a model with formula \texttt{y\ \textasciitilde{}\ x}?
\end{enumerate}

\begin{Shaded}
\begin{Highlighting}[]
\FunctionTok{set.seed}\NormalTok{(}\DecValTok{42}\NormalTok{) }\CommentTok{\# Setting seed for reproducibility}
\NormalTok{x }\OtherTok{\textless{}{-}} \FunctionTok{seq}\NormalTok{(}\DecValTok{1}\NormalTok{, }\DecValTok{100}\NormalTok{) }\SpecialCharTok{+} \FunctionTok{rnorm}\NormalTok{(}\DecValTok{100}\NormalTok{, }\DecValTok{0}\NormalTok{, }\DecValTok{2}\NormalTok{)}
\NormalTok{y }\OtherTok{\textless{}{-}} \DecValTok{3} \SpecialCharTok{*}\NormalTok{ x }\SpecialCharTok{+} \FunctionTok{rnorm}\NormalTok{(}\DecValTok{100}\NormalTok{, }\DecValTok{0}\NormalTok{, }\DecValTok{2}\NormalTok{)}
\NormalTok{z }\OtherTok{\textless{}{-}} \FunctionTok{rep}\NormalTok{(}\ConstantTok{NA}\NormalTok{, }\DecValTok{100}\NormalTok{)}

\NormalTok{df }\OtherTok{\textless{}{-}} \FunctionTok{data.frame}\NormalTok{(}\AttributeTok{x =}\NormalTok{ x, }\AttributeTok{y =}\NormalTok{ y, }\AttributeTok{z =}\NormalTok{ z)}

\NormalTok{df }\SpecialCharTok{|\textgreater{}} \FunctionTok{summary}\NormalTok{()}

\NormalTok{model }\OtherTok{\textless{}{-}}\NormalTok{ df }\SpecialCharTok{|\textgreater{}} \FunctionTok{lm}\NormalTok{(}\AttributeTok{formula =}\NormalTok{ y }\SpecialCharTok{\textasciitilde{}}\NormalTok{ . }\SpecialCharTok{{-}}\NormalTok{ z)}
\end{Highlighting}
\end{Shaded}

\section{Model selection
strategies}\label{sec-multlinreg-modelselectionstrategies}

When selecting models without any previous knowledge about the variables
we can resort to several different strategies. Below we describe the
different strategies along with domain (expert) knowledge to give an
overview of the different options.

\begin{itemize}
\tightlist
\item
  \textbf{Domain (expert) knowledge}: If there is domain knowledge
  available about the variables it is a good idea to use it. Sometimes
  this is possible to deduce by reading about the dataset and how it was
  produced, or from previous experience, and in other cases it may be
  necessary to involve a real expert.
\item
  \textbf{Forward selection}: If domain knowledge is not accessible it
  is possible to perform simple linear regression with each of the
  variables in the dataset and then start by adding the one which
  results in the highest R\(^2\) (or some other measure of model fit).
  In the next step one more variable is added, one by one, and the
  variable that increases R\(^2\) the most is added. The process is
  continued until some stopping rule is satisfied. A good stopping rule
  is to add variables until R\(^2\) adjusted is no longer increased.
\item
  \textbf{Backward selection}: All variables in the dataset are used as
  predictors, then the one with highest \(p\)-value is removed and a new
  model is fitted. Then the process is repeated until all variables have
  sufficiently low \(p\)-values (below 0.05 or some other threshold).
\item
  \textbf{Mixed selection}: Starts with forward selection, but in each
  step variables with high \(p\)-values are removed according to some
  threshold (0.05).
\end{itemize}

In our previous example we started like in backward selection, fitting a
model with all the variables (except the response) as predictors, then
removed \texttt{bmi} based on a VIF analysis and domain knowledge, and
then \texttt{age} based on the high \(p\)-value. There is no one way to
arrive at a good model and you may combine the different approaches as
long as you are \emph{transparent with how the model was selected}. If
you don't have any domain knowledge then it is a good idea to resort to
one of the above mentioned strategies. We will now redo the model
selection using backward selection without using any domain knowledge
and ignoring the high VIF values when selecting what predictors to
remove.

\subsection{Backward selection in the birth weight
dataset}\label{backward-selection-in-the-birth-weight-dataset}

The first step of backward selection is fitting all variables (except
outcome) as predictors, just as before.

\begin{Shaded}
\begin{Highlighting}[]
\CommentTok{\# Backward selection (bs) model}
\NormalTok{model\_bs }\OtherTok{\textless{}{-}}\NormalTok{ birth }\SpecialCharTok{|\textgreater{}} \FunctionTok{lm}\NormalTok{(}\AttributeTok{formula =}\NormalTok{ birth\_weight\_g }\SpecialCharTok{\textasciitilde{}}\NormalTok{ .)}
\NormalTok{model\_summary }\OtherTok{\textless{}{-}}\NormalTok{ model\_bs }\SpecialCharTok{|\textgreater{}} \FunctionTok{summary}\NormalTok{()}
\NormalTok{model\_summary}\SpecialCharTok{$}\NormalTok{coefficients}
\end{Highlighting}
\end{Shaded}

\begin{verbatim}
                           Estimate Std. Error    t value     Pr(>|t|)
(Intercept)            -3507.945867 4552.54025 -0.7705469 4.417911e-01
pregnancy_duration_wks   175.974265   20.44187  8.6085219 1.347298e-15
sexMale                  152.066373   50.90347  2.9873480 3.128843e-03
age                        2.034181    5.28159  0.3851456 7.004966e-01
length_cm                 -2.781630   26.82002 -0.1037147 9.174890e-01
weight_kg                 29.014634   35.08557  0.8269678 4.091393e-01
bmi                      -66.988384   96.85528 -0.6916338 4.898868e-01
\end{verbatim}

We proceed and remove \texttt{length\_cm} as it has the highest
\(p\)-value.

\begin{Shaded}
\begin{Highlighting}[]
\NormalTok{model\_bs }\OtherTok{\textless{}{-}}\NormalTok{ birth }\SpecialCharTok{|\textgreater{}} \FunctionTok{lm}\NormalTok{(}\AttributeTok{formula =}\NormalTok{ birth\_weight\_g }\SpecialCharTok{\textasciitilde{}}\NormalTok{ . }\SpecialCharTok{{-}}\NormalTok{ length\_cm)}
\NormalTok{model\_summary }\OtherTok{\textless{}{-}}\NormalTok{ model\_bs }\SpecialCharTok{|\textgreater{}} \FunctionTok{summary}\NormalTok{()}
\NormalTok{model\_summary}\SpecialCharTok{$}\NormalTok{coefficients}
\end{Highlighting}
\end{Shaded}

\begin{verbatim}
                           Estimate Std. Error    t value     Pr(>|t|)
(Intercept)            -3971.487861 864.180000 -4.5956720 7.203834e-06
pregnancy_duration_wks   176.015166  20.392883  8.6312056 1.138982e-15
sexMale                  151.522005  50.520219  2.9992349 3.012494e-03
age                        2.050085   5.267693  0.3891809 6.975121e-01
weight_kg                 25.415620   5.168363  4.9175379 1.694973e-06
bmi                      -57.069217  15.267424 -3.7379730 2.355433e-04
\end{verbatim}

Continuing we remove \texttt{age} as it has the only insignificant
\(p\)-value.

\begin{Shaded}
\begin{Highlighting}[]
\NormalTok{model\_bs }\OtherTok{\textless{}{-}}\NormalTok{ birth }\SpecialCharTok{|\textgreater{}} \FunctionTok{lm}\NormalTok{(}\AttributeTok{formula =}\NormalTok{ birth\_weight\_g }\SpecialCharTok{\textasciitilde{}}\NormalTok{ . }\SpecialCharTok{{-}}\NormalTok{ length\_cm }\SpecialCharTok{{-}}\NormalTok{ age)}
\NormalTok{model\_bs }\SpecialCharTok{|\textgreater{}} \FunctionTok{summary}\NormalTok{()}
\end{Highlighting}
\end{Shaded}

\begin{verbatim}

Call:
lm(formula = birth_weight_g ~ . - length_cm - age, data = birth)

Residuals:
     Min       1Q   Median       3Q      Max 
-1005.24  -242.74    -3.17   254.24  1099.70 

Coefficients:
                       Estimate Std. Error t value Pr(>|t|)    
(Intercept)            -3862.19     815.73  -4.735 3.88e-06 ***
pregnancy_duration_wks   174.88      20.14   8.681 8.05e-16 ***
sexMale                  152.50      50.36   3.028  0.00275 ** 
weight_kg                 25.79       5.07   5.086 7.68e-07 ***
bmi                      -58.12      15.00  -3.875  0.00014 ***
---
Signif. codes:  0 '***' 0.001 '**' 0.01 '*' 0.05 '.' 0.1 ' ' 1

Residual standard error: 376.6 on 225 degrees of freedom
Multiple R-squared:  0.3579,    Adjusted R-squared:  0.3465 
F-statistic: 31.35 on 4 and 225 DF,  p-value: < 2.2e-16
\end{verbatim}

We end up with a model with a significant relationship with all the
predictors, but including both \texttt{bmi} and \texttt{weight\_kg}
makes the model harder to interpret as BMI depends on weight. The model
also has slight multicollinearity issues, even if it is not as
catastrophic as before, making the inference about the model less
reliable.

\begin{Shaded}
\begin{Highlighting}[]
\FunctionTok{vif}\NormalTok{(model\_bs)}
\end{Highlighting}
\end{Shaded}

\begin{verbatim}
pregnancy_duration_wks                    sex              weight_kg 
              1.023938               1.026370               5.041159 
                   bmi 
              5.042461 
\end{verbatim}

Comparing with our previous model that achieved R\(^2 = 0.3564\), the
backward selection model is slightly better with R\(^2 = 0.3579\). In
this case we may prefer the first model as it is more interpretable at
the cost of only a very small amount of explained variance.

\subsection{Exercises}\label{exercises-20}

\begin{enumerate}
\def\labelenumi{\arabic{enumi}.}
\tightlist
\item
  Perform forward selection on the \texttt{birth} dataset, stop when
  R\(^2\) adjusted is not increased any more. Which variable has the
  lowest \(p\)-value in the single regression models? Is this
  surprising? What is the final model? Feeling lazy? Just look at the
  solution.
\item
  Based on the results of exercise 1, would the result using mixed
  selection been any different?
\item
  Perform backward selection on the \texttt{penguins} dataset to find a
  model for predicting \texttt{body\_mass\_g}. Start by excluding
  \texttt{year} as we do not believe it to be relevant. Take
  multicollinearity into account and use reasoning and R\(^2\) adjusted
  to select between models. Hint 1: When using \texttt{vif} with
  categorical variables with more than three levels the output changes.
  Look under ``GVIF'' for continuous variables and categorical variables
  with \emph{only two} levels, i.e.~\texttt{sex}, and ignore
  multicollinearity for \texttt{species} and \texttt{island}. Hint 2:
  This is a difficult task and many models can be considered to be
  correct.
\end{enumerate}

\section{Prediction}\label{prediction}

Predicting using multiple linear regression models follows the same
structure as in simple linear regression. The only difference is that we
now need to provide a value for all the predictors in the model. We give
an example using the model we found using backward selection in the
previous section, predicting the birth weight of a female child born in
week 40 with a mother that is 170 cm long, weighs 68 kg, and has a BMI
of 24.

\textbf{Note:} The \texttt{lm} function in R is sometimes troublesome to
work with. If we fit our model with the formula
\texttt{birth\_weight\_g\ \textasciitilde{}\ .\ -\ length\_cm\ -\ age}
we still need to specify \texttt{age} and \texttt{length\_cm} for our
prediction. Instead fitting the model by \emph{adding} all the included
predictors allows us to make predictions using only the included
variables. So when using our model for prediction, explicitly state the
variables that are \emph{included} in the model is better. Another
option is to set the variables of \texttt{length\_cm} and \texttt{age}
to some arbitrary, but valid number.

\begin{Shaded}
\begin{Highlighting}[]
\CommentTok{\# Setting length\_cm and age to 0 and using previously fitted model}
\NormalTok{prediction }\OtherTok{\textless{}{-}} \FunctionTok{predict}\NormalTok{(}
\NormalTok{  model\_bs, }
  \FunctionTok{data.frame}\NormalTok{(}
    \AttributeTok{sex =} \StringTok{"Female"}\NormalTok{, }
    \AttributeTok{pregnancy\_duration\_wks =} \DecValTok{40}\NormalTok{, }
    \AttributeTok{weight\_kg =} \DecValTok{68}\NormalTok{, }
    \AttributeTok{bmi =} \DecValTok{24}\NormalTok{, }
    \AttributeTok{length\_cm =} \DecValTok{170}\NormalTok{, }
    \AttributeTok{age =} \DecValTok{0}\NormalTok{))}
\NormalTok{prediction}
\end{Highlighting}
\end{Shaded}

\begin{verbatim}
       1 
3491.656 
\end{verbatim}

\begin{Shaded}
\begin{Highlighting}[]
\CommentTok{\# Refitting model with only included predictors in the formula}
\NormalTok{model\_bs2 }\OtherTok{\textless{}{-}}\NormalTok{ birth }\SpecialCharTok{|\textgreater{}} \FunctionTok{lm}\NormalTok{(}\AttributeTok{formula =}\NormalTok{ birth\_weight\_g }\SpecialCharTok{\textasciitilde{}}\NormalTok{ pregnancy\_duration\_wks }\SpecialCharTok{+}\NormalTok{ sex }\SpecialCharTok{+}\NormalTok{ weight\_kg }\SpecialCharTok{+}\NormalTok{ bmi)}
\NormalTok{prediction2 }\OtherTok{\textless{}{-}} \FunctionTok{predict}\NormalTok{(}
\NormalTok{  model\_bs2, }
  \FunctionTok{data.frame}\NormalTok{(}
    \AttributeTok{sex =} \StringTok{"Female"}\NormalTok{, }
    \AttributeTok{pregnancy\_duration\_wks =} \DecValTok{40}\NormalTok{, }
    \AttributeTok{weight\_kg =} \DecValTok{68}\NormalTok{, }
    \AttributeTok{bmi =} \DecValTok{24}\NormalTok{))}
\NormalTok{prediction2}
\end{Highlighting}
\end{Shaded}

\begin{verbatim}
       1 
3491.656 
\end{verbatim}

We see that the prediction turns out the same either way.

\subsection{Exercises}\label{exercises-21}

\begin{enumerate}
\def\labelenumi{\arabic{enumi}.}
\tightlist
\item
  Use the model we found using domain knowledge or forward selection and
  compare the prediction with the result of the backward selection
  model.
\item
  Make a prediction interval for the backward selection model and the
  model you compared to before. Is the predictions different if we take
  uncertainty into account? That is, does the intervals overlap? Hint:
  Add the argument \texttt{interval\ =\ "prediction"} to the predict
  call.
\end{enumerate}

\section{Summary}\label{summary-1}

In this chapter we have looked at how to compare simple and multiple
linear regression models using R\(^2\), the explained variance in the
outcome. We have learned how to deal with multicollinearity and several
strategies for systematic model selection, which we can use when no
domain knowledge is available. Finally we looked at how to use the
\texttt{predict} function for multiple linear regression models.

This chapter is quite substantial, but the most important takeaway is:

\begin{quote}
Building models is hard and there are many good models but even more bad
models. Motivate your choice of model based on domain knowledge if
possible and be transparent in how the model was selected.
\end{quote}

As long as the reader is able to follow how you made your decisions when
selecting your model you have done a good job.

\chapter{Logistic regression}\label{logistic-regression}

When our outcome (dependent) is categorical, linear regression is not
the best option because there is no order and/or distance between the
values the variable can take on, see
Section~\ref{sec-basicstat-categorical}. Modeling the relationship
between continuous (or categorical) variables and a categorical variable
is sometimes called \emph{classification}, because we want to classify
an observation in a certain class, in a certain category. This could be
any qualitative variable, such as eye color or customer satisfaction.
For binary variables, one of the most common methods to achieve this is
logistic regression. This allows us to ask yes and no questions.

For a more technical/mathematical introduction to logistic regression,
please consult the chapter on logistic regression in
\href{https://www.statlearning.com/}{An Introduction to Statistical
Learning}.

\section{The Pima Indians dataset}\label{the-pima-indians-dataset}

To demonstrate logistic regression we will use the
\texttt{PimaIndiansDiabetes2} dataset from the \texttt{mlbench} package.
The dataset is commonly used to test classification algorithms,
predicting the onset of diabetes. The subjects consists of women aged 21
or older of the \href{https://en.wikipedia.org/wiki/Akimel_O'odham}{Pima
Indian} population. Measurements of the potential predictors were made
every other year, but only one measurement per person is included in our
dataset. A subject was classified as diabetic, according to WHO's
criteria, if a subject had a
\href{https://www.ncbi.nlm.nih.gov/books/NBK532915/}{Glucose Tolerance
Test (GTT)} show more than 200 mg/dl 2 hours after ingesting 75 g of
carbohydrate solution within five years of the included measurement. To
better understand the case selection, look at section 2.3 in the
\href{https://www.ncbi.nlm.nih.gov/pmc/articles/PMC2245318/}{original
paper}. Below is a table of the potential predictors and the outcome,
\texttt{diabetes}, and their descriptions from running
\texttt{?PimaIndiansDiabetes2}.

\begin{longtable}[]{@{}ll@{}}
\toprule\noalign{}
Variable & Description \\
\midrule\noalign{}
\endhead
\bottomrule\noalign{}
\endlastfoot
pregnant & Number of times pregnant \\
glucose & Plasma glucose concentration (glucose tolerance test) \\
pressure & Diastolic blood pressure (mm Hg) \\
triceps & Triceps skin fold thickness (mm) \\
insulin & 2-Hour serum insulin (mu U/ml) \\
mass & Body mass index (weight in kg/(height in m)\^{}2) \\
pedigree & Diabetes pedigree function \\
age & Age (years) \\
diabetes & Class variable, pos/neg \\
\end{longtable}

We can see that the information doesn't tell us about how diabetes was
classified, risking confusion as \texttt{glucose} represents a GTT at
\emph{an earlier time} than when diabetes was classified. If it would
have been the at the same time \texttt{glucose} would have been a
perfect predictor for diabetes, as \texttt{glucose} \(> 200\) would
imply diabetic and \texttt{glucose} \(< 200\) not diabetic. This
highlights the importance of understanding the dataset and the context,
as presenting these results to a diabetes researcher without this
information would be embarrassing.

Now when we know how the dataset was created, let's take a
\texttt{glimpse} at it.

\begin{Shaded}
\begin{Highlighting}[]
\FunctionTok{library}\NormalTok{(tidyverse)}
\CommentTok{\#install.packages("mlbench") \# Install to get PimaIndiansDiabetes2 dataset}
\FunctionTok{library}\NormalTok{(mlbench)}
\FunctionTok{data}\NormalTok{(}\StringTok{"PimaIndiansDiabetes2"}\NormalTok{) }\CommentTok{\# Using 2 because it is a corrected version.}
\FunctionTok{glimpse}\NormalTok{(PimaIndiansDiabetes2)}
\end{Highlighting}
\end{Shaded}

\begin{verbatim}
Rows: 768
Columns: 9
$ pregnant <dbl> 6, 1, 8, 1, 0, 5, 3, 10, 2, 8, 4, 10, 10, 1, 5, 7, 0, 7, 1, 1~
$ glucose  <dbl> 148, 85, 183, 89, 137, 116, 78, 115, 197, 125, 110, 168, 139,~
$ pressure <dbl> 72, 66, 64, 66, 40, 74, 50, NA, 70, 96, 92, 74, 80, 60, 72, N~
$ triceps  <dbl> 35, 29, NA, 23, 35, NA, 32, NA, 45, NA, NA, NA, NA, 23, 19, N~
$ insulin  <dbl> NA, NA, NA, 94, 168, NA, 88, NA, 543, NA, NA, NA, NA, 846, 17~
$ mass     <dbl> 33.6, 26.6, 23.3, 28.1, 43.1, 25.6, 31.0, 35.3, 30.5, NA, 37.~
$ pedigree <dbl> 0.627, 0.351, 0.672, 0.167, 2.288, 0.201, 0.248, 0.134, 0.158~
$ age      <dbl> 50, 31, 32, 21, 33, 30, 26, 29, 53, 54, 30, 34, 57, 59, 51, 3~
$ diabetes <fct> pos, neg, pos, neg, pos, neg, pos, neg, pos, pos, neg, pos, n~
\end{verbatim}

We see that we have some missing values and that all potential
predictors are numeric, while our outcome is \emph{binary}. Knowing that
a GTT of over 200 is classified as diabetic, we can start by if
\texttt{glucose} is a good predictor for \texttt{diabetes}. It seems
reasonable, as a person developing diabetes may already have higher
glucose levels. We can visualize the relationship using boxplots to get
an idea of the distribution of glucose levels in the diabetic and
non-diabetic groups.

\begin{Shaded}
\begin{Highlighting}[]
\NormalTok{PimaIndiansDiabetes2 }\SpecialCharTok{|\textgreater{}}
  \FunctionTok{ggplot}\NormalTok{(}\FunctionTok{aes}\NormalTok{(}\AttributeTok{x =}\NormalTok{ diabetes, }\AttributeTok{y =}\NormalTok{ glucose)) }\SpecialCharTok{+} 
  \FunctionTok{geom\_boxplot}\NormalTok{() }\SpecialCharTok{+}
  \FunctionTok{labs}\NormalTok{(}
    \AttributeTok{x =} \StringTok{"Diabetes"}\NormalTok{,}
    \AttributeTok{y =} \StringTok{"Glucose level (mg/dl)"}\NormalTok{)}
\end{Highlighting}
\end{Shaded}

\includegraphics{logistic_regression_files/figure-pdf/unnamed-chunk-2-1.pdf}

We see that the distribution of glucose level is shifted more towards
higher values of the pre-diabetic women compared to the women not
diagnosed with diabetes at a later stage. This indicates that glucose
may be have a relationship with the onset of diabetes, which may not be
surprising, but is the relationship significant? Do we expect it to hold
in the population? Let's try to model the probability of having diabetes
for a given glucose level using a linear model.

\subsection{Exercises}\label{exercises-22}

\begin{enumerate}
\def\labelenumi{\arabic{enumi}.}
\tightlist
\item
  What is the population that we our sample, dataset, is drawn from?
\item
  Visualize the relationship between \texttt{pressure}, diastolic blood
  pressure, and \texttt{diabetes}. Do you suspect that \texttt{pressure}
  may be correlated with diabetes? In the next chapter we will
  investigate this in more detail.
\end{enumerate}

\section{Using linear regression}\label{using-linear-regression}

We may try to \texttt{mutate} diabetes to a numerical variable, where
having diabetes is indicated by 1 and not having diabetes as 0, and use
linear regression to study the effect of glucose level on the onset of
diabetes.

\begin{Shaded}
\begin{Highlighting}[]
\NormalTok{PimaIndiansDiabetes2\_numeric }\OtherTok{\textless{}{-}}\NormalTok{ PimaIndiansDiabetes2 }\SpecialCharTok{|\textgreater{}}
  \FunctionTok{mutate}\NormalTok{(}\AttributeTok{diabetes\_num =} \FunctionTok{ifelse}\NormalTok{(diabetes }\SpecialCharTok{==} \StringTok{"pos"}\NormalTok{, }\DecValTok{1}\NormalTok{, }\DecValTok{0}\NormalTok{))}
\NormalTok{PimaIndiansDiabetes2\_numeric }\SpecialCharTok{|\textgreater{}}  
  \FunctionTok{lm}\NormalTok{(}\AttributeTok{formula =}\NormalTok{ diabetes\_num }\SpecialCharTok{\textasciitilde{}}\NormalTok{ glucose) }\SpecialCharTok{|\textgreater{}} 
  \FunctionTok{summary}\NormalTok{()}
\end{Highlighting}
\end{Shaded}

\begin{verbatim}

Call:
lm(formula = diabetes_num ~ glucose, data = PimaIndiansDiabetes2_numeric)

Residuals:
    Min      1Q  Median      3Q     Max 
-0.9304 -0.2970 -0.1193  0.3322  0.9888 

Coefficients:
             Estimate Std. Error t value Pr(>|t|)    
(Intercept) -0.591346   0.061721  -9.581   <2e-16 ***
glucose      0.007724   0.000492  15.701   <2e-16 ***
---
Signif. codes:  0 '***' 0.001 '**' 0.01 '*' 0.05 '.' 0.1 ' ' 1

Residual standard error: 0.4147 on 761 degrees of freedom
  (5 observations deleted due to missingness)
Multiple R-squared:  0.2447,    Adjusted R-squared:  0.2437 
F-statistic: 246.5 on 1 and 761 DF,  p-value: < 2.2e-16
\end{verbatim}

We see that the relationship is positive and significant, but what does
this actually mean? Let's interpret the coefficient in the model. The
coefficient for \texttt{glucose} tells us that \texttt{diabetes}
increase with \(\sim0.007\) when \texttt{glucose} increases by one, but
what does this mean? Diabetes is either 0 or 1, it cannot even be 0.007,
so how can it increase by 0.007? What we want to model is rather the
\emph{probability} of \texttt{diabetes} to be 1 given a certain glucose
level. We can visualize our linear model to see if we maybe can do that.

\begin{Shaded}
\begin{Highlighting}[]
\NormalTok{PimaIndiansDiabetes2\_numeric }\SpecialCharTok{|\textgreater{}}
  \FunctionTok{ggplot}\NormalTok{(}\FunctionTok{aes}\NormalTok{(}\AttributeTok{x =}\NormalTok{ glucose, }\AttributeTok{y =}\NormalTok{ diabetes\_num)) }\SpecialCharTok{+}
  \FunctionTok{geom\_smooth}\NormalTok{(}\AttributeTok{method =} \StringTok{"lm"}\NormalTok{, }\AttributeTok{se =}\NormalTok{ F, }\AttributeTok{col =} \StringTok{"black"}\NormalTok{) }\SpecialCharTok{+}
  \FunctionTok{geom\_point}\NormalTok{(}\FunctionTok{aes}\NormalTok{(}\AttributeTok{col =}\NormalTok{ diabetes)) }\SpecialCharTok{+}
  \FunctionTok{labs}\NormalTok{(}
   \AttributeTok{x =} \StringTok{"Glucose level (mg/dl)"}\NormalTok{,}
   \AttributeTok{y =} \StringTok{"Probability of diabetes"}\NormalTok{,}
   \AttributeTok{col =} \StringTok{"Diabetes"}\NormalTok{)}
\end{Highlighting}
\end{Shaded}

\begin{verbatim}
`geom_smooth()` using formula = 'y ~ x'
\end{verbatim}

\begin{verbatim}
Warning: Removed 5 rows containing non-finite outside the scale range
(`stat_smooth()`).
\end{verbatim}

\begin{verbatim}
Warning: Removed 5 rows containing missing values or values outside the scale range
(`geom_point()`).
\end{verbatim}

\includegraphics{logistic_regression_files/figure-pdf/unnamed-chunk-4-1.pdf}

Some glucose levels seem to give a negative probability for having
diabetes, which is not possible, so the interpretation of the model as
assigning a probability of diabetes to a given glucose level doesn't
work either. Linear regression is not suitable for either categorical
outcomes or for estimating probabilities. Instead we should use
\emph{logistic regression}.

\subsection{Exercises}\label{exercises-23}

\begin{enumerate}
\def\labelenumi{\arabic{enumi}.}
\tightlist
\item
  Linear regression, as you know, is useful to investigate the
  relationship between two continuous variables. Perform a linear
  regression \texttt{pressure\ \textasciitilde{}\ mass}, where
  \texttt{mass} is body mass index (BMI), and draw a conclusion based on
  the \(p\)-value if there is a linear relationship between the
  variables. State the null and alternative hypotheses.
\item
  Formulate the hypotheses for the test in the linear regression. Do
  they make sense?
\end{enumerate}

\section{Using logistic regression}\label{using-logistic-regression}

Unlike linear regression, logistic regression models the probability of
having diabetes for a given glucose level. It models the relationships
between the explanatory variables and the outcome through the
\emph{logistic function}. In this course we will not go into details of
the function itself, but we can visualize it

{Technically logistic regression models the log-odds, but we can always
transform the log-odds to probabilities.}

\begin{Shaded}
\begin{Highlighting}[]
\NormalTok{PimaIndiansDiabetes2 }\SpecialCharTok{|\textgreater{}}
  \FunctionTok{mutate}\NormalTok{(}\AttributeTok{diabetes\_num =} \FunctionTok{ifelse}\NormalTok{(diabetes }\SpecialCharTok{==} \StringTok{"pos"}\NormalTok{, }\DecValTok{1}\NormalTok{, }\DecValTok{0}\NormalTok{)) }\SpecialCharTok{|\textgreater{}} \CommentTok{\# Necessary for the geom\_smooth}
  \FunctionTok{ggplot}\NormalTok{(}\FunctionTok{aes}\NormalTok{(}\AttributeTok{x =}\NormalTok{ glucose, }\AttributeTok{y =}\NormalTok{ diabetes\_num)) }\SpecialCharTok{+}
  \FunctionTok{geom\_smooth}\NormalTok{(}\AttributeTok{method =} \StringTok{"glm"}\NormalTok{, }\AttributeTok{method.args =} \FunctionTok{list}\NormalTok{(}\AttributeTok{family =} \StringTok{"binomial"}\NormalTok{), }\AttributeTok{se =} \ConstantTok{FALSE}\NormalTok{) }\SpecialCharTok{+}
  \FunctionTok{geom\_point}\NormalTok{(}\FunctionTok{aes}\NormalTok{(}\AttributeTok{col =}\NormalTok{ diabetes)) }\SpecialCharTok{+}
  \FunctionTok{labs}\NormalTok{(}
   \AttributeTok{x =} \StringTok{"Glucose level (mg/dl)"}\NormalTok{,}
   \AttributeTok{y =} \StringTok{"Probability of diabetes"}\NormalTok{,}
   \AttributeTok{col =} \StringTok{"Diabetes"}\NormalTok{)}
\end{Highlighting}
\end{Shaded}

\includegraphics{logistic_regression_files/figure-pdf/unnamed-chunk-5-1.pdf}

Using logistic regression, all probabilities fall nicely within 0 and 1,
and we can interpret the relationship in terms of probabilities. In this
course we will only be interested in two things from the logistic
regression: The sign of the regression coefficient coefficient, i.e.~if
it is positive or negative, and if it is statistically significant
(\(p < 0.05\)). A positive coefficient means that it increases the
probability of having diabetes, and vice versa for a negative
coefficient. To create our logistic regression model, we use the
\texttt{glm} function. GLM stands for ``Generalized Linear Model'', and
it can model more relationships than logistic, but it is the only
relationship we will work with in this course. To select the logistic
regression we must use the argument \texttt{family\ =\ "binomial"}.
before performing the regression we will state our hypotheses:

\begin{quote}
\textbf{\(H_0\)}: There is \emph{no} relationship between glucose level
and being diagnosed with diabetes within five years.
\end{quote}

\begin{quote}
\textbf{\(H_a\)}: There \emph{is} a relationship between glucose level
and being diagnosed with diabetes within five years.
\end{quote}

Now we are ready to fit our logistic model.

\begin{Shaded}
\begin{Highlighting}[]
\NormalTok{model }\OtherTok{\textless{}{-}} \FunctionTok{glm}\NormalTok{(diabetes }\SpecialCharTok{\textasciitilde{}}\NormalTok{ glucose, }\AttributeTok{data =}\NormalTok{ PimaIndiansDiabetes2, }\AttributeTok{family =} \StringTok{"binomial"}\NormalTok{)}
\FunctionTok{summary}\NormalTok{(model)}
\end{Highlighting}
\end{Shaded}

\begin{verbatim}

Call:
glm(formula = diabetes ~ glucose, family = "binomial", data = PimaIndiansDiabetes2)

Coefficients:
             Estimate Std. Error z value Pr(>|z|)    
(Intercept) -5.715088   0.438100  -13.04   <2e-16 ***
glucose      0.040634   0.003382   12.01   <2e-16 ***
---
Signif. codes:  0 '***' 0.001 '**' 0.01 '*' 0.05 '.' 0.1 ' ' 1

(Dispersion parameter for binomial family taken to be 1)

    Null deviance: 986.70  on 762  degrees of freedom
Residual deviance: 786.56  on 761  degrees of freedom
  (5 observations deleted due to missingness)
AIC: 790.56

Number of Fisher Scoring iterations: 4
\end{verbatim}

Under ``Coefficients'' we find ``glucose'' on row three, and the
estimate of the coefficient under ``estimate''. We can see that it is a
positive coefficient, so having a higher glucose level increases the
probability of having diabetes, again as expected. We can also find the
\(p\)-value in the fourth column,
``Pr(\textgreater\textbar z\textbar)''. We see that we have a positive
relationship between glucose level and the later onset of diabetes that
is statistically significant. We can conclude that having a higher
glucose level increases the risk of the onset of diabetes in the
population of women in the Pima Indians tribe, at significance level
0.05.

\subsection{Exercises}\label{exercises-24}

\begin{enumerate}
\def\labelenumi{\arabic{enumi}.}
\tightlist
\item
  Since glucose level is used to diagnose diabetes this analysis may not
  be very insightful. Instead it would be interesting to find other
  variables that are useful for predicting diabetes. Fit logistic
  regressions for (a) \texttt{pressure} and (b) \texttt{mass} and
  analyze the results. Remember to formulate the hypotheses. In the next
  chapter we will continue to study these variables and their
  relationship with diabetes using multiple logistic regression.
\end{enumerate}

\section{Prediction}\label{prediction-1}

Just as with linear regression we can use our logistic model to predict
the probability of having diabetes given a glucose level. Again we use
the \texttt{predict} function with the parameter
\texttt{type\ =\ "response"} to get the output in predicted
probabilities. We can try to predict the probability of having diabetes
if the glucose level is 150 mg/dl.

\begin{Shaded}
\begin{Highlighting}[]
\FunctionTok{predict}\NormalTok{(model, }\FunctionTok{data.frame}\NormalTok{(}\AttributeTok{glucose =} \DecValTok{150}\NormalTok{), }\AttributeTok{type =} \StringTok{"response"}\NormalTok{)}
\end{Highlighting}
\end{Shaded}

\begin{verbatim}
        1 
0.5938624 
\end{verbatim}

According to our model, the probability of having diabetes of a woman in
the Pima Indian tribe with glucose level 150 mg/dl is 0.59.

Now we have studied the relationship between one continuous variable and
the outcome. Just as with linear regression before, we can study the
relationship between multiple predictors and \texttt{diabetes} using
multiple logistic regression, which is the topic of the next chapter.

\subsection{Exercises}\label{exercises-25}

\begin{enumerate}
\def\labelenumi{\arabic{enumi}.}
\tightlist
\item
  Use your previously fitted models of
  \texttt{diabetes\ \textasciitilde{}\ pressure} and
  \texttt{diabetes\ \textasciitilde{}\ mass} to predict probabilities of
  getting diabetes given \texttt{pressure\ =\ 90} and
  \texttt{mass\ =\ 25}.
\item
  Based on the following plot, what glucose level would generate a
  prediction of above 0.9 probability of diabetes? Verify using
  \texttt{predict}.
\end{enumerate}

\begin{Shaded}
\begin{Highlighting}[]
\NormalTok{PimaIndiansDiabetes2 }\SpecialCharTok{|\textgreater{}}
  \FunctionTok{mutate}\NormalTok{(}\AttributeTok{diabetes\_num =} \FunctionTok{ifelse}\NormalTok{(diabetes }\SpecialCharTok{==} \StringTok{"pos"}\NormalTok{, }\DecValTok{1}\NormalTok{, }\DecValTok{0}\NormalTok{)) }\SpecialCharTok{|\textgreater{}}
  \FunctionTok{ggplot}\NormalTok{(}\FunctionTok{aes}\NormalTok{(}\AttributeTok{x =}\NormalTok{ glucose, }\AttributeTok{y =}\NormalTok{ diabetes\_num)) }\SpecialCharTok{+}
  \FunctionTok{geom\_smooth}\NormalTok{(}\AttributeTok{method =} \StringTok{"glm"}\NormalTok{, }\AttributeTok{method.args =} \FunctionTok{list}\NormalTok{(}\AttributeTok{family =} \StringTok{"binomial"}\NormalTok{), }\AttributeTok{se =} \ConstantTok{FALSE}\NormalTok{) }\SpecialCharTok{+}
  \FunctionTok{geom\_point}\NormalTok{(}\FunctionTok{aes}\NormalTok{(}\AttributeTok{col =}\NormalTok{ diabetes)) }\SpecialCharTok{+}
  \FunctionTok{geom\_hline}\NormalTok{(}
    \AttributeTok{yintercept =} \FloatTok{0.9}\NormalTok{,}
    \AttributeTok{col =} \StringTok{"black"}\NormalTok{,}
    \AttributeTok{linetype =} \DecValTok{2}\NormalTok{) }\SpecialCharTok{+}
  \FunctionTok{geom\_vline}\NormalTok{(}
    \AttributeTok{xintercept =} \DecValTok{195}\NormalTok{,}
    \AttributeTok{col =} \StringTok{"black"}\NormalTok{,}
    \AttributeTok{linetype =} \DecValTok{2}\NormalTok{) }\SpecialCharTok{+}
  \FunctionTok{labs}\NormalTok{(}
   \AttributeTok{x =} \StringTok{"Glucose level (mg/dl)"}\NormalTok{,}
   \AttributeTok{y =} \StringTok{"Probability of diabetes"}\NormalTok{,}
   \AttributeTok{col =} \StringTok{"Diabetes"}\NormalTok{) }\SpecialCharTok{+}
  \FunctionTok{scale\_x\_continuous}\NormalTok{(}\AttributeTok{breaks =} \FunctionTok{seq}\NormalTok{(}\DecValTok{40}\NormalTok{, }\DecValTok{200}\NormalTok{, }\AttributeTok{by =} \DecValTok{10}\NormalTok{)) }\SpecialCharTok{+}
  \FunctionTok{scale\_y\_continuous}\NormalTok{(}\AttributeTok{breaks =} \FunctionTok{seq}\NormalTok{(}\DecValTok{0}\NormalTok{, }\DecValTok{1}\NormalTok{, }\AttributeTok{by =} \FloatTok{0.1}\NormalTok{))}
\end{Highlighting}
\end{Shaded}

\includegraphics{logistic_regression_files/figure-pdf/unnamed-chunk-8-1.pdf}

\begin{enumerate}
\def\labelenumi{\arabic{enumi}.}
\setcounter{enumi}{2}
\tightlist
\item
  Predict the probability of the onset of diabetes given
  \texttt{glucose\ =\ 200}, i.e.~the threshold for being diagnosed with
  diabetes. Does your prediction make sense? Is the model trained for
  this?
\end{enumerate}

\chapter{Multiple logistic
regression}\label{multiple-logistic-regression}

Just as we expanded linear regression to multiple linear rgression, we
will now study multiple logistic regression. You will see that the step
from knowing multiple linear regression and logistic regression to
multiple logistic regression will be a very small step. When performing
multiple logistic regression the hypotheses and interpretations are
similar to that of multiple linear regression, they also depend on the
other variables included in the model. So any conclusions that we draw
based on multiple regression is \emph{always} conditioned on the other
variables in the model.

\section{Blood pressure and diabetes}\label{blood-pressure-and-diabetes}

We will again use the \texttt{PimaIndianDiabetes2} dataset, but this
time we will start by looking at \texttt{pressure}, the diastolic blood
pressure, as our predictor. We start by looking at the distribution over
\texttt{pressure}.

\begin{Shaded}
\begin{Highlighting}[]
\FunctionTok{library}\NormalTok{(tidyverse)}
\CommentTok{\#install.packages("mlbench") \# Install to get PimaIndiansDiabetes2 dataset}
\FunctionTok{library}\NormalTok{(mlbench)}
\FunctionTok{data}\NormalTok{(}\StringTok{"PimaIndiansDiabetes2"}\NormalTok{) }\CommentTok{\# Using 2 because it is a corrected version.}
\FunctionTok{library}\NormalTok{(patchwork)}
\NormalTok{p1 }\OtherTok{\textless{}{-}}\NormalTok{ PimaIndiansDiabetes2 }\SpecialCharTok{|\textgreater{}}
  \FunctionTok{ggplot}\NormalTok{(}\FunctionTok{aes}\NormalTok{(}\AttributeTok{y =}\NormalTok{ pressure)) }\SpecialCharTok{+} 
  \FunctionTok{facet\_grid}\NormalTok{(}\SpecialCharTok{\textasciitilde{}}\NormalTok{ diabetes) }\SpecialCharTok{+}
  \FunctionTok{geom\_boxplot}\NormalTok{() }\SpecialCharTok{+}
  \FunctionTok{theme}\NormalTok{(}
    \AttributeTok{axis.text.x =} \FunctionTok{element\_blank}\NormalTok{(),  }\CommentTok{\# Remove axis text}
    \AttributeTok{axis.ticks.x =} \FunctionTok{element\_blank}\NormalTok{()  }\CommentTok{\# Remove axis title}
\NormalTok{  ) }\SpecialCharTok{+}
  \FunctionTok{labs}\NormalTok{(}\AttributeTok{y =} \StringTok{"Diastolic blood pressure (mmHg)"}\NormalTok{)}
\NormalTok{p2 }\OtherTok{\textless{}{-}}\NormalTok{ PimaIndiansDiabetes2 }\SpecialCharTok{|\textgreater{}}
  \FunctionTok{ggplot}\NormalTok{(}\FunctionTok{aes}\NormalTok{(}\AttributeTok{x =}\NormalTok{ pressure, }\AttributeTok{fill =}\NormalTok{ diabetes)) }\SpecialCharTok{+} 
  \FunctionTok{geom\_density}\NormalTok{(}\AttributeTok{alpha =} \FloatTok{0.5}\NormalTok{) }\SpecialCharTok{+}
  \FunctionTok{theme}\NormalTok{(}
    \AttributeTok{legend.position =} \StringTok{"top"}\NormalTok{,}
    \AttributeTok{axis.text.y =} \FunctionTok{element\_blank}\NormalTok{(),  }\CommentTok{\# Remove axis text}
    \AttributeTok{axis.ticks.y =} \FunctionTok{element\_blank}\NormalTok{()  }\CommentTok{\# Remove axis title}
\NormalTok{  ) }\SpecialCharTok{+}
  \FunctionTok{labs}\NormalTok{(}
    \AttributeTok{x =} \StringTok{"Diastolic blood pressure (mmHg)"}\NormalTok{,}
    \AttributeTok{y =} \StringTok{""}\NormalTok{,}
    \AttributeTok{fill =} \StringTok{"Diabetes"}
\NormalTok{    )}
\NormalTok{p1 }\SpecialCharTok{+}\NormalTok{ p2}
\end{Highlighting}
\end{Shaded}

\includegraphics{logistic_regression_multiple_files/figure-pdf/unnamed-chunk-1-1.pdf}

We see that there seems to be a slight relationship between blood
pressure and diabetes, but the difference is not as clear as with
glucose. Let's see if the result is significant in a simple logistic
regression.

\begin{Shaded}
\begin{Highlighting}[]
\NormalTok{model }\OtherTok{\textless{}{-}} \FunctionTok{glm}\NormalTok{(diabetes }\SpecialCharTok{\textasciitilde{}}\NormalTok{ pressure, }\AttributeTok{data =}\NormalTok{ PimaIndiansDiabetes2, }\AttributeTok{family =} \StringTok{"binomial"}\NormalTok{)}
\NormalTok{model\_summary }\OtherTok{\textless{}{-}}\NormalTok{ model }\SpecialCharTok{|\textgreater{}} \FunctionTok{summary}\NormalTok{()}
\NormalTok{model\_summary}\SpecialCharTok{$}\NormalTok{coefficients}
\end{Highlighting}
\end{Shaded}

\begin{verbatim}
               Estimate  Std. Error   z value     Pr(>|z|)
(Intercept) -2.83025863 0.491750376 -5.755478 8.639681e-09
pressure     0.02988395 0.006587373  4.536550 5.718197e-06
\end{verbatim}

We see that there is a statistically significant, positive relationship
between blood pressure and having diabetes (\(p = 5.72 \cdot 10^{-6}\))
at significance level 0.05. What if we include more variables in our
model, can there be some other explanation for the relationship between
blood pressure and diabetes?

\section{Fitting multiple logistic regression
models}\label{fitting-multiple-logistic-regression-models}

Now we will create our multiple logistic regression model, including
both \texttt{pressure} and \texttt{mass}, BMI. In the previous chapter
we saw that \texttt{pressure} and \texttt{mass} have a positive linear
correlation, indicating that BMI may have an effect on blood pressure,
so how will it affect our logistic regression model? We state the
hypotheses for the coefficient of both \texttt{mass} and
\texttt{pressure}:

\begin{quote}
\textbf{\(H_0\)}: There is \emph{no} relationship between BMI and being
diagnosed with diabetes, given that \emph{blood pressure} is part of the
model.
\end{quote}

\begin{quote}
\textbf{\(H_a\)}: There \emph{is} a relationship between BMI level and
being diagnosed with diabetes, given that \emph{blood pressure} is part
of the model.
\end{quote}

and

\begin{quote}
\textbf{\(H_0\)}: There is \emph{no} relationship between blood pressure
and being diagnosed with diabetes, given that \emph{BMI} is part of the
model.
\end{quote}

\begin{quote}
\textbf{\(H_a\)}: There \emph{is} a relationship between blood pressure
level and being diagnosed with diabetes, given that \emph{BMI} is part
of the model.*
\end{quote}

Now we are ready to fit our multiple logistic regression model.

\begin{Shaded}
\begin{Highlighting}[]
\NormalTok{model }\OtherTok{\textless{}{-}}\NormalTok{ PimaIndiansDiabetes2 }\SpecialCharTok{|\textgreater{}} 
  \FunctionTok{glm}\NormalTok{(}\AttributeTok{formula =}\NormalTok{ diabetes }\SpecialCharTok{\textasciitilde{}}\NormalTok{ pressure }\SpecialCharTok{+}\NormalTok{ mass, }\AttributeTok{family =} \StringTok{"binomial"}\NormalTok{)}
\NormalTok{model }\SpecialCharTok{|\textgreater{}} \FunctionTok{summary}\NormalTok{()}
\end{Highlighting}
\end{Shaded}

\begin{verbatim}

Call:
glm(formula = diabetes ~ pressure + mass, family = "binomial", 
    data = PimaIndiansDiabetes2)

Coefficients:
             Estimate Std. Error z value Pr(>|z|)    
(Intercept) -4.877158   0.604581  -8.067 7.20e-16 ***
pressure     0.016488   0.006971   2.365    0.018 *  
mass         0.091729   0.013218   6.940 3.93e-12 ***
---
Signif. codes:  0 '***' 0.001 '**' 0.01 '*' 0.05 '.' 0.1 ' ' 1

(Dispersion parameter for binomial family taken to be 1)

    Null deviance: 938.74  on 728  degrees of freedom
Residual deviance: 864.32  on 726  degrees of freedom
  (39 observations deleted due to missingness)
AIC: 870.32

Number of Fisher Scoring iterations: 4
\end{verbatim}

We see that both \texttt{mass} and \texttt{pressure} ar statistically
significant when included in the model and that both increase the
probability of diabetes, i.e.~their coefficients are both positive.

\section{Model selection}\label{model-selection-1}

We may also choose to include the variable \texttt{triceps}, which
measures the triceps skin fold thickness that correlates with body fat.
Let's fit a model using \texttt{triceps} instead of \texttt{mass}.

\begin{Shaded}
\begin{Highlighting}[]
\NormalTok{model }\OtherTok{\textless{}{-}}\NormalTok{ PimaIndiansDiabetes2 }\SpecialCharTok{|\textgreater{}} 
  \FunctionTok{glm}\NormalTok{(}\AttributeTok{formula =}\NormalTok{ diabetes }\SpecialCharTok{\textasciitilde{}}\NormalTok{ pressure }\SpecialCharTok{+}\NormalTok{ triceps, }\AttributeTok{family =} \StringTok{"binomial"}\NormalTok{)}
\NormalTok{model }\SpecialCharTok{|\textgreater{}} \FunctionTok{summary}\NormalTok{()}
\end{Highlighting}
\end{Shaded}

\begin{verbatim}

Call:
glm(formula = diabetes ~ pressure + triceps, family = "binomial", 
    data = PimaIndiansDiabetes2)

Coefficients:
             Estimate Std. Error z value Pr(>|z|)    
(Intercept) -4.003582   0.628387  -6.371 1.88e-10 ***
pressure     0.024955   0.008167   3.055  0.00225 ** 
triceps      0.050017   0.009960   5.022 5.12e-07 ***
---
Signif. codes:  0 '***' 0.001 '**' 0.01 '*' 0.05 '.' 0.1 ' ' 1

(Dispersion parameter for binomial family taken to be 1)

    Null deviance: 685.23  on 538  degrees of freedom
Residual deviance: 638.43  on 536  degrees of freedom
  (229 observations deleted due to missingness)
AIC: 644.43

Number of Fisher Scoring iterations: 4
\end{verbatim}

Again we see that both variables are significant at significance level
0.05, but which model is the best? In linear regression we could use the
adjusted R\(^2\), but it doesn't apply to logistic regression. Instead
we are going to use the \emph{Akaike Information Criterion} (AIC).

To be able to tell which logistic regression model is better, we can
look at the summary output and second to last line: ``AIC''. AIC
estimates the prediction error of the model, where a lower AIC is
better. We can compare the two models that we fitted before, either by
looking at the summary output (see above), or using the \texttt{\$}
operator and \texttt{aic}.

\begin{Shaded}
\begin{Highlighting}[]
\CommentTok{\# Triceps model}
\NormalTok{model\_triceps }\OtherTok{\textless{}{-}}\NormalTok{ PimaIndiansDiabetes2 }\SpecialCharTok{|\textgreater{}} 
  \FunctionTok{glm}\NormalTok{(}\AttributeTok{formula =}\NormalTok{ diabetes }\SpecialCharTok{\textasciitilde{}}\NormalTok{ pressure }\SpecialCharTok{+}\NormalTok{ triceps, }\AttributeTok{family =} \StringTok{"binomial"}\NormalTok{)}
\NormalTok{model\_triceps\_summary }\OtherTok{\textless{}{-}}\NormalTok{ model\_triceps }\SpecialCharTok{|\textgreater{}} \FunctionTok{summary}\NormalTok{()}
\NormalTok{model\_triceps\_summary}\SpecialCharTok{$}\NormalTok{aic}
\end{Highlighting}
\end{Shaded}

\begin{verbatim}
[1] 644.4286
\end{verbatim}

\begin{Shaded}
\begin{Highlighting}[]
\CommentTok{\# Mass}
\NormalTok{model\_bmi }\OtherTok{\textless{}{-}}\NormalTok{ PimaIndiansDiabetes2 }\SpecialCharTok{|\textgreater{}} 
  \FunctionTok{glm}\NormalTok{(}\AttributeTok{formula =}\NormalTok{ diabetes }\SpecialCharTok{\textasciitilde{}}\NormalTok{ pressure }\SpecialCharTok{+}\NormalTok{ mass, }\AttributeTok{family =} \StringTok{"binomial"}\NormalTok{)}
\NormalTok{model\_bmi\_summary }\OtherTok{\textless{}{-}}\NormalTok{ model\_bmi }\SpecialCharTok{|\textgreater{}} \FunctionTok{summary}\NormalTok{()}
\NormalTok{model\_bmi\_summary}\SpecialCharTok{$}\NormalTok{aic}
\end{Highlighting}
\end{Shaded}

\begin{verbatim}
[1] 870.3233
\end{verbatim}

We see that the model including \texttt{triceps} and \texttt{pressure}
has a lower AIC and is thus a preferable model.

AIC is more similar to R\(^2\) adjusted than R\(^2\) in that it also
penalizes more complex models. There is a ``cost'' to adding a new
variable to the model and if it doesn't increase how well the model
describes the data, the AIC will increase.

\subsection{Exercises}\label{exercises-26}

\begin{enumerate}
\def\labelenumi{\arabic{enumi}.}
\tightlist
\item
  Which of the two models \texttt{diabetes\ \textasciitilde{}\ triceps}
  and \texttt{diabetes\ \textasciitilde{}\ pressure} has the lowest AIC?
\item
  Fit a logistic regression model for
  \texttt{diabetes\ \textasciitilde{}\ pressure\ +\ mass\ +\ triceps}.
  Is the AIC value lower? What about the VIF values? Remember you can
  calculate the VIF value using the \texttt{vif} function from the
  \texttt{car} library.
\end{enumerate}

\subsection{Model selection
strategies}\label{model-selection-strategies}

When creating a multiple logistic regression model you may use any model
selection strategy discussed in
Section~\ref{sec-multlinreg-modelselectionstrategies}. Whenever possible
it is good to use domain or expert knowledge (D/E knowledge), and
otherwise it is good to try either backward, forwards, or mixed
selection and motivate the final model either based on AIC and
significance or D/E knowledge.

Adding more variables tend to reduce the AIC, as long as they improve
the model more than the penalty of adding a variable. We saw this when
comparing the models of with \texttt{triceps}, \texttt{mass}, and
\texttt{pressure} alone vs the models including two or more variables.
We can add all the variables in the dataset using the dot, \texttt{.},
and build a model using backward selection.

\begin{Shaded}
\begin{Highlighting}[]
\NormalTok{model }\OtherTok{\textless{}{-}}\NormalTok{ PimaIndiansDiabetes2 }\SpecialCharTok{|\textgreater{}} 
  \FunctionTok{glm}\NormalTok{(}\AttributeTok{formula =}\NormalTok{ diabetes }\SpecialCharTok{\textasciitilde{}}\NormalTok{ ., }\AttributeTok{family =} \StringTok{"binomial"}\NormalTok{)}
\NormalTok{model }\SpecialCharTok{|\textgreater{}} \FunctionTok{summary}\NormalTok{()}
\end{Highlighting}
\end{Shaded}

\begin{verbatim}

Call:
glm(formula = diabetes ~ ., family = "binomial", data = PimaIndiansDiabetes2)

Coefficients:
              Estimate Std. Error z value Pr(>|z|)    
(Intercept) -1.004e+01  1.218e+00  -8.246  < 2e-16 ***
pregnant     8.216e-02  5.543e-02   1.482  0.13825    
glucose      3.827e-02  5.768e-03   6.635 3.24e-11 ***
pressure    -1.420e-03  1.183e-02  -0.120  0.90446    
triceps      1.122e-02  1.708e-02   0.657  0.51128    
insulin     -8.253e-04  1.306e-03  -0.632  0.52757    
mass         7.054e-02  2.734e-02   2.580  0.00989 ** 
pedigree     1.141e+00  4.274e-01   2.669  0.00760 ** 
age          3.395e-02  1.838e-02   1.847  0.06474 .  
---
Signif. codes:  0 '***' 0.001 '**' 0.01 '*' 0.05 '.' 0.1 ' ' 1

(Dispersion parameter for binomial family taken to be 1)

    Null deviance: 498.10  on 391  degrees of freedom
Residual deviance: 344.02  on 383  degrees of freedom
  (376 observations deleted due to missingness)
AIC: 362.02

Number of Fisher Scoring iterations: 5
\end{verbatim}

The AIC value reduced to 362.02, which is much lower than the previous
models. But we have several variables that are not significantly related
to \texttt{diabetes}, with \texttt{pressure} now having a \(p\)-value
larger than 0.9. We remove \texttt{pressure} from our model using the
formula \texttt{diabetes\ \textasciitilde{}\ .\ -\ pressure}.

\begin{Shaded}
\begin{Highlighting}[]
\NormalTok{model }\OtherTok{\textless{}{-}}\NormalTok{ PimaIndiansDiabetes2 }\SpecialCharTok{|\textgreater{}} 
  \FunctionTok{glm}\NormalTok{(}\AttributeTok{formula =}\NormalTok{ diabetes }\SpecialCharTok{\textasciitilde{}}\NormalTok{ . }\SpecialCharTok{{-}}\NormalTok{ pressure, }\AttributeTok{family =} \StringTok{"binomial"}\NormalTok{)}
\NormalTok{model\_summary }\OtherTok{\textless{}{-}}\NormalTok{ model }\SpecialCharTok{|\textgreater{}} \FunctionTok{summary}\NormalTok{()}
\NormalTok{model\_summary}\SpecialCharTok{$}\NormalTok{aic}
\end{Highlighting}
\end{Shaded}

\begin{verbatim}
[1] 360.0356
\end{verbatim}

We see that the AIC reduced even more by removing pressure, illustrating
the penalty in AIC for including variables that doesn't improve the
models ability to explain the data.

\textbf{Note:} using \texttt{.} in the formula in \texttt{glm} removes
all observations with missing values in the data before fitting the
model. When comparing models \texttt{.} is good, but when fitting the
final model intended for inference or prediction it is recommended to
either handle missing values explicitly or stating included predictors
explicitly in the formula using \texttt{+}. See
Section~\ref{sec-multlinreg-usingdot} for more details.

\subsection{Exercises}\label{exercises-27}

\begin{enumerate}
\def\labelenumi{\arabic{enumi}.}
\tightlist
\item
  Continue the backward selection process based on the model above. What
  is the lowest AIC you can achieve? Note that this exercise is
  contained in Assignment 2.
\end{enumerate}

\section{Prediction}\label{prediction-2}

Just as with simple logistic regression we can use our model to predict.
We need to provide a \texttt{data.frame} (or \texttt{tibble}) with
values for the predictors in the fitted model. We can use the model we
fitted before,
\texttt{diabetes\ \textasciitilde{}\ pressure\ +\ triceps}, to predict
the probability of the onset of diabetes within five years for a woman
with diastolic blood pressure 75 mm Hg and triceps skin fold thickness
30 mm. To get the output as a probability we also need to provide the
parameter \texttt{type\ =\ "response"}.

\begin{Shaded}
\begin{Highlighting}[]
\CommentTok{\# Triceps model}
\NormalTok{model\_triceps }\OtherTok{\textless{}{-}}\NormalTok{ PimaIndiansDiabetes2 }\SpecialCharTok{|\textgreater{}} 
  \FunctionTok{glm}\NormalTok{(}\AttributeTok{formula =}\NormalTok{ diabetes }\SpecialCharTok{\textasciitilde{}}\NormalTok{ pressure }\SpecialCharTok{+}\NormalTok{ triceps, }\AttributeTok{family =} \StringTok{"binomial"}\NormalTok{)}
\NormalTok{model\_triceps }\SpecialCharTok{|\textgreater{}} \FunctionTok{predict}\NormalTok{(}\FunctionTok{data.frame}\NormalTok{(}\AttributeTok{pressure =} \DecValTok{75}\NormalTok{, }\AttributeTok{triceps =} \DecValTok{30}\NormalTok{), }\AttributeTok{type =} \StringTok{"response"}\NormalTok{)}
\end{Highlighting}
\end{Shaded}

\begin{verbatim}
        1 
0.3471786 
\end{verbatim}

The probability of the onset of diabetes within five years for a women
of the Pima Indian population with diastolic blood pressure 75 mm Hg and
30 mm triceps skin fold thickness is \(\sim 0.35\) according to our
model. If we want our prediction to fall into the classes ``diabetes
positive'' or ``diabetes negative'', we may use 0.5 as a cutoff. Then
any predicted probability of more than 0.5 will be considered as
positive and below 0.5 as negative. We can use the \texttt{ifelse}
function to achieve this.

\begin{Shaded}
\begin{Highlighting}[]
\CommentTok{\# Triceps model}
\NormalTok{prediction\_prob }\OtherTok{\textless{}{-}}\NormalTok{ model\_triceps }\SpecialCharTok{|\textgreater{}} \FunctionTok{predict}\NormalTok{(}\FunctionTok{data.frame}\NormalTok{(}\AttributeTok{pressure =} \DecValTok{75}\NormalTok{, }\AttributeTok{triceps =} \DecValTok{30}\NormalTok{), }\AttributeTok{type =} \StringTok{"response"}\NormalTok{)}
\FunctionTok{ifelse}\NormalTok{(prediction\_prob }\SpecialCharTok{\textgreater{}} \FloatTok{0.5}\NormalTok{, }\StringTok{"pos"}\NormalTok{, }\StringTok{"neg"}\NormalTok{)}
\end{Highlighting}
\end{Shaded}

\begin{verbatim}
    1 
"neg" 
\end{verbatim}

In this case, the prediction is negative, since the probability was
below 0.5.

\section{Accuracy}\label{accuracy}

In addition to AIC, we can evaluate the performance of our model using
prediction by measuring the \emph{accuracy} of our model. Accuracy is
the ratio of correct predictions over the total number of predictions,
\(\text{Accuracy} = \#\text{Correct predictions} / \#\text{Total predictions}\).
We can calculate the accuracy of the model above on our dataset by
predicting the probability of diabetes for all (complete) observations
and then comparing to the observed values of \texttt{diabetes}. To
compare with \texttt{diabetes}, we set the predicted outcome to
positive, \texttt{pos}, if the predicted probability is larger than 0.5.

\begin{Shaded}
\begin{Highlighting}[]
\NormalTok{model }\OtherTok{\textless{}{-}}\NormalTok{ PimaIndiansDiabetes2 }\SpecialCharTok{|\textgreater{}} 
  \FunctionTok{glm}\NormalTok{(}\AttributeTok{formula =}\NormalTok{ diabetes }\SpecialCharTok{\textasciitilde{}}\NormalTok{ . }\SpecialCharTok{{-}}\NormalTok{ pressure, }\AttributeTok{family =} \StringTok{"binomial"}\NormalTok{)}
\CommentTok{\# Calculating the accuracy}
\NormalTok{accuracy }\OtherTok{\textless{}{-}}\NormalTok{ PimaIndiansDiabetes2 }\SpecialCharTok{|\textgreater{}} 
  \CommentTok{\# Calculated the predicted probabilities for all observations}
  \FunctionTok{mutate}\NormalTok{(}\AttributeTok{predicted\_prob =}\NormalTok{ model }\SpecialCharTok{|\textgreater{}} \FunctionTok{predict}\NormalTok{(PimaIndiansDiabetes2, }\AttributeTok{type =} \StringTok{"response"}\NormalTok{)) }\SpecialCharTok{|\textgreater{}} 
  \CommentTok{\# Predicting the outcome}
  \FunctionTok{mutate}\NormalTok{(}\AttributeTok{predicted\_outcome =} \FunctionTok{ifelse}\NormalTok{(predicted\_prob }\SpecialCharTok{\textgreater{}} \FloatTok{0.5}\NormalTok{, }\StringTok{"pos"}\NormalTok{, }\StringTok{"neg"}\NormalTok{)) }\SpecialCharTok{|\textgreater{}} 
  \CommentTok{\# Dropping predictions with missing values}
  \FunctionTok{filter}\NormalTok{(}\SpecialCharTok{!}\FunctionTok{is.na}\NormalTok{(predicted\_outcome)) }\SpecialCharTok{|\textgreater{}} 
  \CommentTok{\# Summarizing Accuracy = \# Correct predictions / \# Total predictions}
  \FunctionTok{summarize}\NormalTok{(}\AttributeTok{accuracy =} \FunctionTok{mean}\NormalTok{(predicted\_outcome }\SpecialCharTok{==}\NormalTok{ diabetes)) }\SpecialCharTok{|\textgreater{}}
  \CommentTok{\# Pull accuracy to get the value}
  \FunctionTok{pull}\NormalTok{(accuracy)}

\NormalTok{accuracy}
\end{Highlighting}
\end{Shaded}

\begin{verbatim}
[1] 0.7831633
\end{verbatim}

We see that our model has an accuracy of \(\sim 0.78\), predicting the
observed outcome in 78\% of the cases in the dataset. Accuracy may be
intuitive to measure the predictive performance of the model, but it may
be misleading. If we look at how well the model predicts within each
class, i.e.~diabetic positive and negative, we see that the model is
better att predicting negative outcomes than positive. Here we visualize
each observation in two columns, one for diabetes negative and one for
positive, and then color on wether the prediction is correct or not.

\begin{Shaded}
\begin{Highlighting}[]
\FunctionTok{library}\NormalTok{(scales)}
\NormalTok{hex }\OtherTok{\textless{}{-}} \FunctionTok{hue\_pal}\NormalTok{()(}\DecValTok{3}\NormalTok{)}
\FunctionTok{set.seed}\NormalTok{(}\DecValTok{42}\NormalTok{) }\CommentTok{\# For reproducibility}
\NormalTok{PimaIndiansDiabetes2 }\SpecialCharTok{|\textgreater{}} 
  \CommentTok{\# Calculated the predicted probabilities for all observations}
  \FunctionTok{mutate}\NormalTok{(}\AttributeTok{predicted\_prob =}\NormalTok{ model }\SpecialCharTok{|\textgreater{}} \FunctionTok{predict}\NormalTok{(PimaIndiansDiabetes2, }\AttributeTok{type =} \StringTok{"response"}\NormalTok{)) }\SpecialCharTok{|\textgreater{}} 
  \CommentTok{\# Dropping predictions with missing values}
  \FunctionTok{drop\_na}\NormalTok{() }\SpecialCharTok{|\textgreater{}} 
  \CommentTok{\# Predicting the outcome}
  \FunctionTok{mutate}\NormalTok{(}\AttributeTok{predicted\_outcome =} \FunctionTok{factor}\NormalTok{(}\FunctionTok{ifelse}\NormalTok{(predicted\_prob }\SpecialCharTok{\textgreater{}} \FloatTok{0.5}\NormalTok{, }\StringTok{"pos"}\NormalTok{, }\StringTok{"neg"}\NormalTok{))) }\SpecialCharTok{|\textgreater{}} 
  \FunctionTok{mutate}\NormalTok{(}\AttributeTok{correct\_prediction =} \FunctionTok{factor}\NormalTok{(}\FunctionTok{ifelse}\NormalTok{(predicted\_outcome }\SpecialCharTok{==}\NormalTok{ diabetes, }\StringTok{"yes"}\NormalTok{, }\StringTok{"no"}\NormalTok{))) }\SpecialCharTok{|\textgreater{}} 
  \FunctionTok{mutate}\NormalTok{(}\AttributeTok{index =} \FunctionTok{row\_number}\NormalTok{(), }\AttributeTok{.by =}\NormalTok{ diabetes) }\SpecialCharTok{|\textgreater{}} 
  \FunctionTok{arrange}\NormalTok{(index) }\SpecialCharTok{|\textgreater{}} 
  \FunctionTok{ggplot}\NormalTok{(}\FunctionTok{aes}\NormalTok{(}\AttributeTok{x =}\NormalTok{ diabetes, }\AttributeTok{y =}\NormalTok{ index, }\AttributeTok{col =}\NormalTok{ correct\_prediction)) }\SpecialCharTok{+}
  \FunctionTok{geom\_jitter}\NormalTok{(}\AttributeTok{alpha =} \FloatTok{0.7}\NormalTok{) }\SpecialCharTok{+}
  \FunctionTok{labs}\NormalTok{(}
    \AttributeTok{x =} \StringTok{"Diabetes"}\NormalTok{,}
    \AttributeTok{col =} \StringTok{"Correct prediction"}\NormalTok{) }\SpecialCharTok{+}
  \FunctionTok{scale\_color\_manual}\NormalTok{(}\AttributeTok{values =} \FunctionTok{c}\NormalTok{(hex[}\DecValTok{1}\NormalTok{], hex[}\DecValTok{2}\NormalTok{])) }\SpecialCharTok{+}
  \FunctionTok{theme}\NormalTok{(}
    \AttributeTok{legend.position =} \StringTok{"top"}\NormalTok{,}
    \AttributeTok{axis.title.y =} \FunctionTok{element\_blank}\NormalTok{(),}
    \AttributeTok{axis.text.y =} \FunctionTok{element\_blank}\NormalTok{(),}
    \AttributeTok{axis.ticks.y =} \FunctionTok{element\_blank}\NormalTok{(),}
\NormalTok{    )}
\end{Highlighting}
\end{Shaded}

\includegraphics{logistic_regression_multiple_files/figure-pdf/unnamed-chunk-11-1.pdf}

We see that our model is better at predicting negative outcomes than
positive outcomes. We may consider another cut-off to be more strict or
lenient in how many cases we classify as diabetic. We may want to lower
the cut-off to screen more women for diabetes rather than fewer, if that
is what the model is used for. In general it is good to also report
other metrics, such as \emph{precision}, \emph{recall}, and
\emph{F-score}, to give a better idea of the performance of the model.
In this course we will not be considering either of these, and simply
use 0.5.

{Usually when calculating accuracy the data is split in a training and
test set, where the model is trained on the training set and then the
accuracy is tested on the test set. This way the models overfitting to
the sample is evaluated as well.}

\subsection{Exercises}\label{exercises-28}

\begin{enumerate}
\def\labelenumi{\arabic{enumi}.}
\tightlist
\item
  Calculate the accuracy of the model including all potential
  predictors, i.e.~include \texttt{pressure}. Is the accuracy better?
\item
  Calculate the accuracy of the best model you found before (i.e.~lowest
  AIC). How does it compare to the model including all variables?
\item
  Run the code below to test the accuracy of the models with and without
  \texttt{pressure} using a training and test set. Try varying the seed.
  Do you think we should include pressure or not? Note: You don't have
  to use a training and test set in the assignments in this course.
\end{enumerate}

\begin{Shaded}
\begin{Highlighting}[]
\CommentTok{\# Set seed for reproducibility}
\FunctionTok{set.seed}\NormalTok{(}\DecValTok{42}\NormalTok{) }

\CommentTok{\# Create an index to keep track of subjects}
\NormalTok{pima }\OtherTok{\textless{}{-}}\NormalTok{ PimaIndiansDiabetes2 }\SpecialCharTok{|\textgreater{}} 
  \FunctionTok{drop\_na}\NormalTok{() }\SpecialCharTok{|\textgreater{}} \CommentTok{\# Using only complete observations}
  \FunctionTok{mutate}\NormalTok{(}\AttributeTok{index =} \DecValTok{1}\SpecialCharTok{:}\FunctionTok{n}\NormalTok{())}

\CommentTok{\# Create training and testing set}
\NormalTok{training\_set }\OtherTok{\textless{}{-}}\NormalTok{ pima }\SpecialCharTok{|\textgreater{}} 
  \FunctionTok{sample\_frac}\NormalTok{(}\FloatTok{0.8}\NormalTok{) }\CommentTok{\# Use 80\% of the data in training set}
\NormalTok{test\_set }\OtherTok{\textless{}{-}}\NormalTok{ pima }\SpecialCharTok{|\textgreater{}} 
  \FunctionTok{anti\_join}\NormalTok{(training\_set, }\AttributeTok{by =} \StringTok{"index"}\NormalTok{)}

\CommentTok{\# Train model on training set}
\NormalTok{model }\OtherTok{\textless{}{-}}\NormalTok{ training\_set }\SpecialCharTok{|\textgreater{}} 
  \FunctionTok{glm}\NormalTok{(}\AttributeTok{formula =}\NormalTok{ diabetes }\SpecialCharTok{\textasciitilde{}}\NormalTok{ ., }\AttributeTok{family =} \StringTok{"binomial"}\NormalTok{)}

\CommentTok{\# Evaluate model by accuracy on test set}
\NormalTok{predictions }\OtherTok{\textless{}{-}}\NormalTok{ model }\SpecialCharTok{|\textgreater{}} 
  \FunctionTok{predict}\NormalTok{(test\_set, }\AttributeTok{type =} \StringTok{"response"}\NormalTok{)}
\NormalTok{test\_set }\SpecialCharTok{|\textgreater{}} 
  \FunctionTok{mutate}\NormalTok{(}\AttributeTok{prediction =} \FunctionTok{ifelse}\NormalTok{(predictions }\SpecialCharTok{\textgreater{}} \FloatTok{0.5}\NormalTok{, }\StringTok{"pos"}\NormalTok{, }\StringTok{"neg"}\NormalTok{)) }\SpecialCharTok{|\textgreater{}} 
  \FunctionTok{summarize}\NormalTok{(}\AttributeTok{accuracy =} \FunctionTok{mean}\NormalTok{(prediction }\SpecialCharTok{==}\NormalTok{ diabetes))}
\end{Highlighting}
\end{Shaded}

\begin{enumerate}
\def\labelenumi{\arabic{enumi}.}
\setcounter{enumi}{3}
\tightlist
\item
  Consider a very unbalanced dataset of 100 observations with a binary
  outcome, ``yes'' or ``no''. In the dataset, 90\% of the observations
  have outcome ``yes''. What accuracy would a model that always predicts
  ``yes'' get? Hint: This exercise is better done with pen and paper,
  using the formula for accuracy.
\end{enumerate}

\section{Summary}\label{summary-2}

This concludes the part on logistic regression. We have seen that
logistic regression is suitable when our outcome is a \emph{binary}
categorical variable, how to compare logistic regression models, and how
to compute accuracy.

When faced with a modelling problem, it is important to understand if it
is a regression problem or classification problem and select an
appropriate method, e.g.~linear or logistic regression, where regression
problems deal with continuous outcomes and classification problems
categorical outcomes.

After these chapters you should be able to identify the type of problem
and apply the appropriate method.

Good luck!

\chapter{Missing values and imputation}\label{sec-missing}

Missing data is ubiquitous when working with real data. We typically
(and should) represent missing values in our data by \texttt{NA}, which
stands for ``Not Available'' or in other words, missing. The most
obvious issue with missing data is that most of the algorithms we use
rely on complete observations, observations without any missing values.
Our algorithms simply will not produce a result if the data contains
missing values. But this is not the only problem of missing data as,
depending on \emph{why} the data is missing, it may bias our inference
and creating predictive models that do not perform as expected.

In this chapter we give a brief introduction to missing data and how we
can handle it.

\section{Why is data missing?}\label{why-is-data-missing}

To illustrate how missing data can occur and in what ways we will rely
on an artificial example with the \texttt{palmerpenguins} dataset. We
first look at the \texttt{summary} of the dataset to remember the
variables and see if it has any missing values.

\begin{Shaded}
\begin{Highlighting}[]
\FunctionTok{library}\NormalTok{(tidyverse)}
\FunctionTok{library}\NormalTok{(palmerpenguins)}
\FunctionTok{summary}\NormalTok{(penguins)}
\end{Highlighting}
\end{Shaded}

\begin{verbatim}
      species          island    bill_length_mm  bill_depth_mm  
 Adelie   :152   Biscoe   :168   Min.   :32.10   Min.   :13.10  
 Chinstrap: 68   Dream    :124   1st Qu.:39.23   1st Qu.:15.60  
 Gentoo   :124   Torgersen: 52   Median :44.45   Median :17.30  
                                 Mean   :43.92   Mean   :17.15  
                                 3rd Qu.:48.50   3rd Qu.:18.70  
                                 Max.   :59.60   Max.   :21.50  
                                 NA's   :2       NA's   :2      
 flipper_length_mm  body_mass_g       sex           year     
 Min.   :172.0     Min.   :2700   female:165   Min.   :2007  
 1st Qu.:190.0     1st Qu.:3550   male  :168   1st Qu.:2007  
 Median :197.0     Median :4050   NA's  : 11   Median :2008  
 Mean   :200.9     Mean   :4202                Mean   :2008  
 3rd Qu.:213.0     3rd Qu.:4750                3rd Qu.:2009  
 Max.   :231.0     Max.   :6300                Max.   :2009  
 NA's   :2         NA's   :2                                 
\end{verbatim}

We note that there are a few missing values in the dataset and we will
now try to understand different mechanisms that can cause these missing
values, \texttt{NA}s. To understand how values become missing, we need
to understand how the data was recorded. From the
\href{https://journals.plos.org/plosone/article?id=10.1371/journal.pone.0090081}{original
reference} of the study, the \texttt{palmerpenguins} dataset was
collected by the following procedure

\begin{quote}
Each study season, Adélie penguin study nests (n = 30) were distributed
equally between the three study islands, with 10 nests located on each
island. Gentoo penguin study nests (n = 30) were all located on Biscoe
Island, while chinstrap penguin study nests (n = 15) were all located on
Dream Island (Fig. 1c). The reduced sample size for chinstraps was due
to the overall smaller number of individuals breeding at rookeries on
Dream Island. Each season, study nests, where pairs of adults were
present, were individually marked and chosen before the onset of
egg-laying, and consistently monitored. When study nests were found at
the one-egg stage, both adults were captured to obtain blood samples
used for molecular sexing and SI analyses, and measurements of
structural size and body mass. At the time of capture, each adult
penguin was quickly blood sampled (\textasciitilde1 ml) from the
brachial vein using a sterile 3 ml syringe and heparinized infusion
needle. Collected blood was stored in 1.5 ml micro-centrifuge tubes that
were kept cool. In the field, a small amount of whole blood was smeared
on clean filter paper stored in a 1.5 ml micro-centrifuge tube for
molecular sexing. Measurements of culmen length and depth (using dial
calipers \$\pm\$60.1 mm), right flipper (using a ruler 61 mm), and body
mass (using 5 kg \(\pm\) 625 g or 10 kg \(\pm\) 650 g Pesola spring
scales and a weigh bag) were obtained to quantify body size variation.
After handling, individuals at study nests were further monitored to
ensure the pair reached clutch completion, i.e., two eggs.
\end{quote}

Now that we have information about how the data was collected we can
start to think about reasons for missing values. But first we need to
understand more generally how values can become missing. Missing values
are usually categorized into three different patterns: Missing
Completely At Random (MCAR), Missing At Random (MAR), and Not Missing At
Random (NMAR). In this course we will only focus on the first form,
MCAR, and assume that it holds. For those who are interested there is a
short section at the end of this chapter that introduces the other
types.

\section{Handling missing values when
MCAR}\label{handling-missing-values-when-mcar}

When values are ``missing completely at random'' there is no systematic
explanation for why a specific value is missing. In our example this may
happen if for example the blood sample of a penguin was lost, resulting
in a missing value in the \texttt{sex} variable, or if one of the
researchers would spill coffee on the noted values, ruining the paper.
These events are ``completely random'' and does not depend on
characteristics of the penguins themselves. In this case we can use
intuitive methods to handle the missing values since the missing values
are not introducing any bias in our data.

\subsection{Complete case analysis}\label{complete-case-analysis}

If the missing values are believed to be missing completely at random
(MCAR) and there are only a few missing values we may use the
observations that are complete in our dataset. This is called
\emph{complete case analysis} and we can achieve it using the
\texttt{drop\_na} function

\begin{Shaded}
\begin{Highlighting}[]
\NormalTok{penguins\_complete }\OtherTok{\textless{}{-}}\NormalTok{ penguins }\SpecialCharTok{|\textgreater{}} \FunctionTok{drop\_na}\NormalTok{()}
\end{Highlighting}
\end{Shaded}

\subsection{Mean imputation}\label{mean-imputation}

If we have several missing values such that we have to drop a
substantial part of our dataset but we still believe that the values are
MCAR we may \emph{impute} the values. When we impute the values of a
variable we use information from our data to replace the missing value,
to impute it. A simple method of imputation is \emph{single value
imputation} where all missing values for one variable is replaced with
the same value. If the variable is continuous, a natural and common
value to use is the mean of the variable, called \emph{mean imputation}.
We can impute the missing flipper lengths in the penguins data using
mean imputation with the \texttt{mutate} and \texttt{replace\_na}
functions from \texttt{dplyr}.

\begin{Shaded}
\begin{Highlighting}[]
\NormalTok{mean\_flipper\_length }\OtherTok{\textless{}{-}} \FunctionTok{mean}\NormalTok{(penguins}\SpecialCharTok{$}\NormalTok{flipper\_length\_mm, }\AttributeTok{na.rm =} \ConstantTok{TRUE}\NormalTok{)}
\NormalTok{penguins\_imputed }\OtherTok{\textless{}{-}}\NormalTok{ penguins }\SpecialCharTok{|\textgreater{}} 
  \FunctionTok{mutate}\NormalTok{(}\AttributeTok{flipper\_length\_mm =} \FunctionTok{replace\_na}\NormalTok{(flipper\_length\_mm, mean\_flipper\_length))}
\NormalTok{penguins\_imputed }\SpecialCharTok{|\textgreater{}} \FunctionTok{summary}\NormalTok{()}
\end{Highlighting}
\end{Shaded}

We see that the flipper length variable now has no missing values, while
keeping the same mean value.

\subsection{Mode imputation}\label{mode-imputation}

If our variable is categorical we cannot use mean imputation. Instead,
we can use the mode, the most common value, to impute the missing
values. We can impute any missing values of the \texttt{sex} variable
using mode imputation using the \texttt{replace\_na} function from
\texttt{dplyr}.

\begin{Shaded}
\begin{Highlighting}[]
\NormalTok{mode\_sex }\OtherTok{\textless{}{-}}\NormalTok{ penguins }\SpecialCharTok{|\textgreater{}}
  \FunctionTok{count}\NormalTok{(sex) }\SpecialCharTok{|\textgreater{}}
  \FunctionTok{filter}\NormalTok{(n }\SpecialCharTok{==} \FunctionTok{max}\NormalTok{(n)) }\SpecialCharTok{|\textgreater{}}
  \FunctionTok{pull}\NormalTok{(sex)}
\NormalTok{penguins\_imputed }\OtherTok{\textless{}{-}}\NormalTok{ penguins\_imputed }\SpecialCharTok{|\textgreater{}} 
  \FunctionTok{mutate}\NormalTok{(}\AttributeTok{sex =} \FunctionTok{replace\_na}\NormalTok{(sex, mode\_sex))}
\NormalTok{penguins\_imputed }\SpecialCharTok{|\textgreater{}} \FunctionTok{summary}\NormalTok{()}
\end{Highlighting}
\end{Shaded}

We see that the \texttt{sex} variable now has no missing values. In this
case, there was a small difference between females and males in the sex
distribution and using the mode might not have been the best option. For
more advanced methods of imputation, see the advanced section at the end
of this chapter.

\section{Extra: Other types of missing values and advanced
imputation}\label{extra-other-types-of-missing-values-and-advanced-imputation}

In this section we introduce the other types of missing values. This is
beyond the scope of the course, but may be of interest to some students.

\subsection{Missing at random (MAR)}\label{missing-at-random-mar}

Values that are ``missing at random'' are missing because of some reason
that is recorded in the data. In our example, the penguins body weight
is measured by using a scale and a weigh bag. Maybe penguins with longer
flippers are harder to fit in the weigh bag, making the researchers more
prone to failing weighing penguins with long flippers. Or maybe female
penguins are harder to gather a blood sample from due to smaller
brachial veins. This results in a missingness patterns in the
\texttt{body\_weight\_g} and \texttt{sex} variables that can to some
extent be explained by the flipper length or sex of the penguin. Thus
the missingness pattern \emph{depends} on our data and we can use it to
better make educated guesses while imputing.

\subsection{Not missing at random
(MNAR)}\label{not-missing-at-random-mnar}

Not missing at random happens when values are systematically missing,
but we have not observed the reason for why the specific values are
missing. For example, if one researcher is responsible for taking
measurements of chinstrap penguins on Dream Island and this researcher
is not able to catch faster penguins then faster chinstrap penguins will
be systematically excluded from our dataset. There is no information in
the dataset that we can use to recover the information about the
missingness since we have not observed any fast chinstrap penguins at
all.

This is the hardest type of missing values to handle as we have no way
of repairing it using the available data and it may influence our
inference and prediction. If we suspect this type of missing values the
best option is to collect more data, compensating for the missingness in
the original data.

\subsection{Regression imputation}\label{regression-imputation}

If data is MAR we can use our dataset to make more informed imputations.
One way that may be appealing after learning about linear and logistic
regression is to use this methods to assign values instead, performing
\emph{regression imputation}. We illustrate a simple form of regression
imputation to make a more educated guess for the missing values of the
\texttt{sex} variable that we imputed by the mode earlier. First, we
make a logistic regression model to predict the sex of the penguins
based on the other variables in the dataset.

\begin{Shaded}
\begin{Highlighting}[]
\NormalTok{sex\_model }\OtherTok{\textless{}{-}}\NormalTok{ penguins }\SpecialCharTok{|\textgreater{}} \FunctionTok{glm}\NormalTok{(}\AttributeTok{formula =}\NormalTok{ sex }\SpecialCharTok{\textasciitilde{}}\NormalTok{ ., }\AttributeTok{family =}\NormalTok{ binomial)}
\end{Highlighting}
\end{Shaded}

Lets select only the significant variables to have a simpler model. This
can be useful if we have more missing values in other columns, which
would prevent us from using these columns for imputation. Once we have
our model with only significant variables we can use it to predict the
sex of the penguins with missing values.

\begin{Shaded}
\begin{Highlighting}[]
\CommentTok{\# Fit new model}
\NormalTok{sex\_model }\OtherTok{\textless{}{-}}\NormalTok{ penguins }\SpecialCharTok{|\textgreater{}} \FunctionTok{glm}\NormalTok{(}\AttributeTok{formula =}\NormalTok{ sex }\SpecialCharTok{\textasciitilde{}}\NormalTok{ species }\SpecialCharTok{+}\NormalTok{ bill\_length\_mm }\SpecialCharTok{+}\NormalTok{ bill\_depth\_mm }\SpecialCharTok{+}\NormalTok{ body\_mass\_g, }\AttributeTok{family =}\NormalTok{ binomial)}

\CommentTok{\# Select penguins with missing sex}
\NormalTok{penguins\_missing\_sex }\OtherTok{\textless{}{-}}\NormalTok{ penguins }\SpecialCharTok{|\textgreater{}}
  \FunctionTok{filter}\NormalTok{(}\FunctionTok{is.na}\NormalTok{(sex))}

\CommentTok{\# Predict sex probability}
\NormalTok{predicted\_sex\_probability }\OtherTok{\textless{}{-}} \FunctionTok{predict}\NormalTok{(sex\_model, }\AttributeTok{newdata =}\NormalTok{ penguins\_missing\_sex, }\AttributeTok{type =} \StringTok{"response"}\NormalTok{)}

\CommentTok{\# Convert probabilities to factor levels}
\NormalTok{predicted\_sex }\OtherTok{\textless{}{-}}\NormalTok{ (predicted\_sex\_probability }\SpecialCharTok{\textgreater{}} \FloatTok{0.5}\NormalTok{) }\SpecialCharTok{|\textgreater{}}
  \FunctionTok{ifelse}\NormalTok{(}\StringTok{"male"}\NormalTok{, }\StringTok{"female"}\NormalTok{) }\SpecialCharTok{|\textgreater{}}
  \FunctionTok{factor}\NormalTok{(}\AttributeTok{levels =} \FunctionTok{levels}\NormalTok{(penguins}\SpecialCharTok{$}\NormalTok{sex))}

\NormalTok{predicted\_sex}
\end{Highlighting}
\end{Shaded}

We see that one penguin still has a missing value. This observation has
missing values in all variables in our \texttt{sex\_model} except
species. Interestingly our model predict most penguins to be female. We
can be fairly confident in our predictions as body measurements and
species should make it quite easy to differentiate between the sexes of
the penguins, as illustrated by the following density plots of body
mass, also showing the missing values.

\begin{Shaded}
\begin{Highlighting}[]
\CommentTok{\# Plot the distributions of the body weight and color by sex and facet by species}
\NormalTok{penguins }\SpecialCharTok{|\textgreater{}} 
  \FunctionTok{ggplot}\NormalTok{(}\FunctionTok{aes}\NormalTok{(}\AttributeTok{x =}\NormalTok{ body\_mass\_g, }\AttributeTok{fill =}\NormalTok{ sex)) }\SpecialCharTok{+}
  \FunctionTok{geom\_density}\NormalTok{(}\AttributeTok{alpha =} \FloatTok{0.5}\NormalTok{) }\SpecialCharTok{+}
  \FunctionTok{facet\_wrap}\NormalTok{(}\SpecialCharTok{\textasciitilde{}}\NormalTok{species) }\SpecialCharTok{+}
  \FunctionTok{theme}\NormalTok{(}\AttributeTok{axis.text.y =} \FunctionTok{element\_blank}\NormalTok{(),  }\CommentTok{\# Remove y{-}axis text}
        \AttributeTok{axis.ticks.y =} \FunctionTok{element\_blank}\NormalTok{()) }\SpecialCharTok{+}\CommentTok{\# Remove y{-}axis ticks}
  \FunctionTok{labs}\NormalTok{(}
    \AttributeTok{x =} \StringTok{"Body mass (g)"}\NormalTok{,}
    \AttributeTok{y =} \StringTok{""}\NormalTok{,}
    \AttributeTok{fill =} \StringTok{"Sex"}\NormalTok{)  }\SpecialCharTok{+}
  \FunctionTok{theme}\NormalTok{(}\AttributeTok{legend.position =} \StringTok{"top"}\NormalTok{) }\SpecialCharTok{+}
  \FunctionTok{guides}\NormalTok{(}\AttributeTok{color =} \FunctionTok{guide\_legend}\NormalTok{(}\AttributeTok{nrow =} \DecValTok{1}\NormalTok{))}
\end{Highlighting}
\end{Shaded}

\includegraphics{missing_values_and_imputation_files/figure-pdf/unnamed-chunk-7-1.pdf}

Regression imputation is a useful tool that you already know how to do
if you have completed the regression chapters.

\subsection{Multiple imputation}\label{multiple-imputation}

Multiple imputation is the most popular imputation. It consists of
sampling from the distribution of observed values to create new dataset.
It is a conceptually more difficult method, but if you are interested
you can have a look at the
\href{https://cran.r-project.org/web/packages/mice/index.html}{mice
package} in R.

\part{Solutions to exercises}

Here are some \emph{suggested} solutions to the recommended exercises.
There are usually many ways to solve problems using code, some ways may
be better than these solutions, so if you have solved exercises in a
different way doesn't mean that it is wrong. Throughout the solutions we
have prioritized solving the problems in the way that we prefer to,
sometimes using code or visualizations that may be outside the scope of
this course. The idea is to show you the potential of R and
\texttt{tidyverse} and \texttt{ggplot} specifically. This means that you
are not expected to be able to solve every problem as it's solved in
these solutions.

If you have any questions, spot any errors, or have any suggestions for
improvement in the lecture notes or solutions, please let us know by
sending an email to the teachers of the course with the tag
{[}DAwR-LN{]} in the subject.

\chapter{Introduction to R}\label{introduction-to-r-1}

\section{Objects}\label{objects-1}

\subsection{Exercise 1}\label{exercise-1}

In R it is possible to use \texttt{=} and \texttt{\textless{}-} to
assign values to objects, but \texttt{\textless{}-} is the preferred
way. \texttt{=} is used to assign function parameters a value, which you
will encounter later.

\section{Functions}\label{functions-1}

\subsection{Exercise 1}\label{exercise-1-1}

\begin{Shaded}
\begin{Highlighting}[]
\NormalTok{multiply }\OtherTok{\textless{}{-}} \ControlFlowTok{function}\NormalTok{(x, y) \{}
\NormalTok{  x }\SpecialCharTok{*}\NormalTok{ y}
\NormalTok{\}}

\FunctionTok{multiply}\NormalTok{(}\DecValTok{36}\NormalTok{, }\DecValTok{52}\NormalTok{)}
\end{Highlighting}
\end{Shaded}

\begin{verbatim}
[1] 1872
\end{verbatim}

\subsection{Exercise 2}\label{exercise-2}

Make sure you know which variable appears on what sign of the minus sign
when using the function.

\begin{Shaded}
\begin{Highlighting}[]
\NormalTok{subtraction }\OtherTok{\textless{}{-}} \ControlFlowTok{function}\NormalTok{(x, y) \{}
\NormalTok{  x }\SpecialCharTok{{-}}\NormalTok{ y}
\NormalTok{\}}

\FunctionTok{subtraction}\NormalTok{(}\DecValTok{36}\NormalTok{, }\DecValTok{52}\NormalTok{)}
\end{Highlighting}
\end{Shaded}

\begin{verbatim}
[1] -16
\end{verbatim}

\subsection{Exercise 3}\label{exercise-3}

Here we see the usage of \texttt{=} in assigning function parameters
values.

\begin{Shaded}
\begin{Highlighting}[]
\FunctionTok{subtraction}\NormalTok{(}\AttributeTok{y =} \DecValTok{52}\NormalTok{, }\AttributeTok{x =} \DecValTok{36}\NormalTok{)}
\end{Highlighting}
\end{Shaded}

\begin{verbatim}
[1] -16
\end{verbatim}

\subsection{Exercise 4}\label{exercise-4}

\begin{Shaded}
\begin{Highlighting}[]
\FunctionTok{subtraction}\NormalTok{(y }\OtherTok{\textless{}{-}} \DecValTok{52}\NormalTok{, x }\OtherTok{\textless{}{-}} \DecValTok{36}\NormalTok{)}
\end{Highlighting}
\end{Shaded}

\begin{verbatim}
[1] 16
\end{verbatim}

Instead of passing x and y as parameters to the function,
\texttt{\textless{}-} sets the values of the objects \texttt{x} and
\texttt{y} to 36 and 52 respectively, which we can see by printing the
variables:

\begin{Shaded}
\begin{Highlighting}[]
\NormalTok{x}
\end{Highlighting}
\end{Shaded}

\begin{verbatim}
[1] 36
\end{verbatim}

\begin{Shaded}
\begin{Highlighting}[]
\NormalTok{y}
\end{Highlighting}
\end{Shaded}

\begin{verbatim}
[1] 52
\end{verbatim}

The output is the same as the following code

\begin{Shaded}
\begin{Highlighting}[]
\NormalTok{x }\OtherTok{\textless{}{-}} \DecValTok{36}
\NormalTok{y }\OtherTok{\textless{}{-}} \DecValTok{52}

\FunctionTok{subtraction}\NormalTok{(y, x)}
\end{Highlighting}
\end{Shaded}

\begin{verbatim}
[1] 16
\end{verbatim}

\textbf{Never use \texttt{\textless{}-} inside a function call.}

\section{Vectors}\label{vectors-1}

\subsection{Exercise 1}\label{exercise-1-2}

The operations, \(+, -, \cdot, /\), are performed \(elementwise\), that
is the first element of bot vectors gets
added/subtracted/multiplied/divided and so on.

\begin{Shaded}
\begin{Highlighting}[]
\NormalTok{vec1 }\OtherTok{\textless{}{-}} \FunctionTok{c}\NormalTok{(}\DecValTok{1}\NormalTok{, }\DecValTok{2}\NormalTok{)}
\NormalTok{vec2 }\OtherTok{\textless{}{-}} \FunctionTok{c}\NormalTok{(}\DecValTok{3}\NormalTok{, }\DecValTok{4}\NormalTok{)}

\NormalTok{vec1 }\SpecialCharTok{+}\NormalTok{ vec2}
\end{Highlighting}
\end{Shaded}

\begin{verbatim}
[1] 4 6
\end{verbatim}

\begin{Shaded}
\begin{Highlighting}[]
\NormalTok{vec1 }\SpecialCharTok{{-}}\NormalTok{ vec2}
\end{Highlighting}
\end{Shaded}

\begin{verbatim}
[1] -2 -2
\end{verbatim}

\begin{Shaded}
\begin{Highlighting}[]
\NormalTok{vec1 }\SpecialCharTok{*}\NormalTok{ vec2}
\end{Highlighting}
\end{Shaded}

\begin{verbatim}
[1] 3 8
\end{verbatim}

\begin{Shaded}
\begin{Highlighting}[]
\NormalTok{vec1 }\SpecialCharTok{/}\NormalTok{ vec2}
\end{Highlighting}
\end{Shaded}

\begin{verbatim}
[1] 0.3333333 0.5000000
\end{verbatim}

\subsection{Exercise 2}\label{exercise-2-1}

We try running the code:

\begin{Shaded}
\begin{Highlighting}[]
\NormalTok{vec1 }\OtherTok{\textless{}{-}} \FunctionTok{c}\NormalTok{(}\DecValTok{1}\NormalTok{, }\DecValTok{2}\NormalTok{)}
\NormalTok{vec2 }\OtherTok{\textless{}{-}} \FunctionTok{c}\NormalTok{(}\DecValTok{3}\NormalTok{, }\DecValTok{4}\NormalTok{)}
\NormalTok{vec3 }\OtherTok{\textless{}{-}} \FunctionTok{c}\NormalTok{(}\DecValTok{3}\NormalTok{, }\StringTok{"four"}\NormalTok{)}

\NormalTok{vec1 }\SpecialCharTok{+}\NormalTok{ vec2}
\end{Highlighting}
\end{Shaded}

\begin{verbatim}
[1] 4 6
\end{verbatim}

\begin{Shaded}
\begin{Highlighting}[]
\NormalTok{vec2 }\SpecialCharTok{+}\NormalTok{ vec3}
\end{Highlighting}
\end{Shaded}

\begin{verbatim}
Error in vec2 + vec3: non-numeric argument to binary operator
\end{verbatim}

We get an error for \texttt{vec2\ +\ vec3}. Vectors must always contain
the same data type, so in this case \texttt{vec3} becomes a vector of
strings:

\begin{Shaded}
\begin{Highlighting}[]
\NormalTok{vec3}
\end{Highlighting}
\end{Shaded}

\begin{verbatim}
[1] "3"    "four"
\end{verbatim}

We can see that 3 is now also a string by the quotes, \texttt{"3"}.

\section{Data frames}\label{data-frames-1}

\subsection{Exercise 1}\label{exercise-1-3}

We run the code and see that the output is similar to the \texttt{head}
function.

\begin{Shaded}
\begin{Highlighting}[]
\CommentTok{\# We comment the installation to not run it twice on our machine}
\CommentTok{\# install.packages("palmerpenguins")}
\FunctionTok{library}\NormalTok{(palmerpenguins)}
\NormalTok{penguins}
\end{Highlighting}
\end{Shaded}

\begin{verbatim}
# A tibble: 344 x 8
   species island    bill_length_mm bill_depth_mm flipper_length_mm body_mass_g
   <fct>   <fct>              <dbl>         <dbl>             <int>       <int>
 1 Adelie  Torgersen           39.1          18.7               181        3750
 2 Adelie  Torgersen           39.5          17.4               186        3800
 3 Adelie  Torgersen           40.3          18                 195        3250
 4 Adelie  Torgersen           NA            NA                  NA          NA
 5 Adelie  Torgersen           36.7          19.3               193        3450
 6 Adelie  Torgersen           39.3          20.6               190        3650
 7 Adelie  Torgersen           38.9          17.8               181        3625
 8 Adelie  Torgersen           39.2          19.6               195        4675
 9 Adelie  Torgersen           34.1          18.1               193        3475
10 Adelie  Torgersen           42            20.2               190        4250
# i 334 more rows
# i 2 more variables: sex <fct>, year <int>
\end{verbatim}

\subsection{Exercise 2}\label{exercise-2-2}

We type the dollar operator straight after the tibble and then the
column name to print the values in it.

\begin{Shaded}
\begin{Highlighting}[]
\NormalTok{penguins}\SpecialCharTok{$}\NormalTok{flipper\_length\_mm}
\end{Highlighting}
\end{Shaded}

\begin{verbatim}
  [1] 181 186 195  NA 193 190 181 195 193 190 186 180 182 191 198 185 195 197
 [19] 184 194 174 180 189 185 180 187 183 187 172 180 178 178 188 184 195 196
 [37] 190 180 181 184 182 195 186 196 185 190 182 179 190 191 186 188 190 200
 [55] 187 191 186 193 181 194 185 195 185 192 184 192 195 188 190 198 190 190
 [73] 196 197 190 195 191 184 187 195 189 196 187 193 191 194 190 189 189 190
 [91] 202 205 185 186 187 208 190 196 178 192 192 203 183 190 193 184 199 190
[109] 181 197 198 191 193 197 191 196 188 199 189 189 187 198 176 202 186 199
[127] 191 195 191 210 190 197 193 199 187 190 191 200 185 193 193 187 188 190
[145] 192 185 190 184 195 193 187 201 211 230 210 218 215 210 211 219 209 215
[163] 214 216 214 213 210 217 210 221 209 222 218 215 213 215 215 215 216 215
[181] 210 220 222 209 207 230 220 220 213 219 208 208 208 225 210 216 222 217
[199] 210 225 213 215 210 220 210 225 217 220 208 220 208 224 208 221 214 231
[217] 219 230 214 229 220 223 216 221 221 217 216 230 209 220 215 223 212 221
[235] 212 224 212 228 218 218 212 230 218 228 212 224 214 226 216 222 203 225
[253] 219 228 215 228 216 215 210 219 208 209 216 229 213 230 217 230 217 222
[271] 214  NA 215 222 212 213 192 196 193 188 197 198 178 197 195 198 193 194
[289] 185 201 190 201 197 181 190 195 181 191 187 193 195 197 200 200 191 205
[307] 187 201 187 203 195 199 195 210 192 205 210 187 196 196 196 201 190 212
[325] 187 198 199 201 193 203 187 197 191 203 202 194 206 189 195 207 202 193
[343] 210 198
\end{verbatim}

\subsection{Exercise 3}\label{exercise-3-1}

When we press tab we see all the columns in the tibble as suggestions.
In general you can use tab to autocomplete in R.

\subsection{Exercise 4}\label{exercise-4-1}

We see some output when loading \texttt{tidyverse} which you can ignore
for now. If you have questions about the functions, run
\texttt{?\textless{}function\textgreater{}}, e.g.~\texttt{?glimpse} in
the console.

\begin{Shaded}
\begin{Highlighting}[]
\FunctionTok{library}\NormalTok{(tidyverse)}
\end{Highlighting}
\end{Shaded}

\begin{verbatim}
-- Attaching core tidyverse packages ------------------------ tidyverse 2.0.0 --
v dplyr     1.1.4     v readr     2.1.5
v forcats   1.0.0     v stringr   1.5.1
v ggplot2   3.5.1     v tibble    3.2.1
v lubridate 1.9.3     v tidyr     1.3.1
v purrr     1.0.2     
-- Conflicts ------------------------------------------ tidyverse_conflicts() --
x dplyr::filter() masks stats::filter()
x dplyr::lag()    masks stats::lag()
i Use the conflicted package (<http://conflicted.r-lib.org/>) to force all conflicts to become errors
\end{verbatim}

\begin{Shaded}
\begin{Highlighting}[]
\FunctionTok{glimpse}\NormalTok{(penguins)}
\end{Highlighting}
\end{Shaded}

\begin{verbatim}
Rows: 344
Columns: 8
$ species           <fct> Adelie, Adelie, Adelie, Adelie, Adelie, Adelie, Adel~
$ island            <fct> Torgersen, Torgersen, Torgersen, Torgersen, Torgerse~
$ bill_length_mm    <dbl> 39.1, 39.5, 40.3, NA, 36.7, 39.3, 38.9, 39.2, 34.1, ~
$ bill_depth_mm     <dbl> 18.7, 17.4, 18.0, NA, 19.3, 20.6, 17.8, 19.6, 18.1, ~
$ flipper_length_mm <int> 181, 186, 195, NA, 193, 190, 181, 195, 193, 190, 186~
$ body_mass_g       <int> 3750, 3800, 3250, NA, 3450, 3650, 3625, 4675, 3475, ~
$ sex               <fct> male, female, female, NA, female, male, female, male~
$ year              <int> 2007, 2007, 2007, 2007, 2007, 2007, 2007, 2007, 2007~
\end{verbatim}

\begin{Shaded}
\begin{Highlighting}[]
\FunctionTok{summary}\NormalTok{(penguins)}
\end{Highlighting}
\end{Shaded}

\begin{verbatim}
      species          island    bill_length_mm  bill_depth_mm  
 Adelie   :152   Biscoe   :168   Min.   :32.10   Min.   :13.10  
 Chinstrap: 68   Dream    :124   1st Qu.:39.23   1st Qu.:15.60  
 Gentoo   :124   Torgersen: 52   Median :44.45   Median :17.30  
                                 Mean   :43.92   Mean   :17.15  
                                 3rd Qu.:48.50   3rd Qu.:18.70  
                                 Max.   :59.60   Max.   :21.50  
                                 NA's   :2       NA's   :2      
 flipper_length_mm  body_mass_g       sex           year     
 Min.   :172.0     Min.   :2700   female:165   Min.   :2007  
 1st Qu.:190.0     1st Qu.:3550   male  :168   1st Qu.:2007  
 Median :197.0     Median :4050   NA's  : 11   Median :2008  
 Mean   :200.9     Mean   :4202                Mean   :2008  
 3rd Qu.:213.0     3rd Qu.:4750                3rd Qu.:2009  
 Max.   :231.0     Max.   :6300                Max.   :2009  
 NA's   :2         NA's   :2                                 
\end{verbatim}

\begin{Shaded}
\begin{Highlighting}[]
\FunctionTok{head}\NormalTok{(penguins)}
\end{Highlighting}
\end{Shaded}

\begin{verbatim}
# A tibble: 6 x 8
  species island    bill_length_mm bill_depth_mm flipper_length_mm body_mass_g
  <fct>   <fct>              <dbl>         <dbl>             <int>       <int>
1 Adelie  Torgersen           39.1          18.7               181        3750
2 Adelie  Torgersen           39.5          17.4               186        3800
3 Adelie  Torgersen           40.3          18                 195        3250
4 Adelie  Torgersen           NA            NA                  NA          NA
5 Adelie  Torgersen           36.7          19.3               193        3450
6 Adelie  Torgersen           39.3          20.6               190        3650
# i 2 more variables: sex <fct>, year <int>
\end{verbatim}

\begin{Shaded}
\begin{Highlighting}[]
\FunctionTok{tail}\NormalTok{(penguins)}
\end{Highlighting}
\end{Shaded}

\begin{verbatim}
# A tibble: 6 x 8
  species   island bill_length_mm bill_depth_mm flipper_length_mm body_mass_g
  <fct>     <fct>           <dbl>         <dbl>             <int>       <int>
1 Chinstrap Dream            45.7          17                 195        3650
2 Chinstrap Dream            55.8          19.8               207        4000
3 Chinstrap Dream            43.5          18.1               202        3400
4 Chinstrap Dream            49.6          18.2               193        3775
5 Chinstrap Dream            50.8          19                 210        4100
6 Chinstrap Dream            50.2          18.7               198        3775
# i 2 more variables: sex <fct>, year <int>
\end{verbatim}

\begin{Shaded}
\begin{Highlighting}[]
\FunctionTok{nrow}\NormalTok{(penguins)}
\end{Highlighting}
\end{Shaded}

\begin{verbatim}
[1] 344
\end{verbatim}

\begin{Shaded}
\begin{Highlighting}[]
\FunctionTok{ncol}\NormalTok{(penguins)}
\end{Highlighting}
\end{Shaded}

\begin{verbatim}
[1] 8
\end{verbatim}

\begin{Shaded}
\begin{Highlighting}[]
\FunctionTok{names}\NormalTok{(penguins)}
\end{Highlighting}
\end{Shaded}

\begin{verbatim}
[1] "species"           "island"            "bill_length_mm"   
[4] "bill_depth_mm"     "flipper_length_mm" "body_mass_g"      
[7] "sex"               "year"             
\end{verbatim}

\chapter{Basic statistics}\label{basic-statistics}

\section{Variables}\label{variables-1}

\subsection{Exercise 1}\label{exercise-1-4}

We run the \texttt{summary} function on the penguins dataset after
loading \texttt{palmerpenguins}.

\begin{Shaded}
\begin{Highlighting}[]
\FunctionTok{library}\NormalTok{(palmerpenguins)}
\FunctionTok{summary}\NormalTok{(penguins)}
\end{Highlighting}
\end{Shaded}

\begin{verbatim}
      species          island    bill_length_mm  bill_depth_mm  
 Adelie   :152   Biscoe   :168   Min.   :32.10   Min.   :13.10  
 Chinstrap: 68   Dream    :124   1st Qu.:39.23   1st Qu.:15.60  
 Gentoo   :124   Torgersen: 52   Median :44.45   Median :17.30  
                                 Mean   :43.92   Mean   :17.15  
                                 3rd Qu.:48.50   3rd Qu.:18.70  
                                 Max.   :59.60   Max.   :21.50  
                                 NA's   :2       NA's   :2      
 flipper_length_mm  body_mass_g       sex           year     
 Min.   :172.0     Min.   :2700   female:165   Min.   :2007  
 1st Qu.:190.0     1st Qu.:3550   male  :168   1st Qu.:2007  
 Median :197.0     Median :4050   NA's  : 11   Median :2008  
 Mean   :200.9     Mean   :4202                Mean   :2008  
 3rd Qu.:213.0     3rd Qu.:4750                3rd Qu.:2009  
 Max.   :231.0     Max.   :6300                Max.   :2009  
 NA's   :2         NA's   :2                                 
\end{verbatim}

We see that the continuous variables show summary statistics and the
categorical shows counts of the different levels, like Adelie,
Chinstrap, and Gentoo in the species variable. You will learn more about
these concepts in the Descriptive statistics section.

\subsection{Exercise 2}\label{exercise-2-3}

We see that the variables have both changed to \texttt{double},
indicated by the \texttt{\textless{}dbl\textgreater{}} in the
\texttt{glimpse} output.

\begin{Shaded}
\begin{Highlighting}[]
\FunctionTok{library}\NormalTok{(tidyverse)}
\end{Highlighting}
\end{Shaded}

\begin{verbatim}
-- Attaching core tidyverse packages ------------------------ tidyverse 2.0.0 --
v dplyr     1.1.4     v readr     2.1.5
v forcats   1.0.0     v stringr   1.5.1
v ggplot2   3.5.1     v tibble    3.2.1
v lubridate 1.9.3     v tidyr     1.3.1
v purrr     1.0.2     
-- Conflicts ------------------------------------------ tidyverse_conflicts() --
x dplyr::filter() masks stats::filter()
x dplyr::lag()    masks stats::lag()
i Use the conflicted package (<http://conflicted.r-lib.org/>) to force all conflicts to become errors
\end{verbatim}

\begin{Shaded}
\begin{Highlighting}[]
\NormalTok{penguins}\SpecialCharTok{$}\NormalTok{flipper\_length\_mm }\OtherTok{\textless{}{-}} \FunctionTok{as.numeric}\NormalTok{(penguins}\SpecialCharTok{$}\NormalTok{flipper\_length\_mm)}
\NormalTok{penguins}\SpecialCharTok{$}\NormalTok{body\_mass\_g }\OtherTok{\textless{}{-}} \FunctionTok{as.numeric}\NormalTok{(penguins}\SpecialCharTok{$}\NormalTok{body\_mass\_g)}
\FunctionTok{glimpse}\NormalTok{(penguins)}
\end{Highlighting}
\end{Shaded}

\begin{verbatim}
Rows: 344
Columns: 8
$ species           <fct> Adelie, Adelie, Adelie, Adelie, Adelie, Adelie, Adel~
$ island            <fct> Torgersen, Torgersen, Torgersen, Torgersen, Torgerse~
$ bill_length_mm    <dbl> 39.1, 39.5, 40.3, NA, 36.7, 39.3, 38.9, 39.2, 34.1, ~
$ bill_depth_mm     <dbl> 18.7, 17.4, 18.0, NA, 19.3, 20.6, 17.8, 19.6, 18.1, ~
$ flipper_length_mm <dbl> 181, 186, 195, NA, 193, 190, 181, 195, 193, 190, 186~
$ body_mass_g       <dbl> 3750, 3800, 3250, NA, 3450, 3650, 3625, 4675, 3475, ~
$ sex               <fct> male, female, female, NA, female, male, female, male~
$ year              <int> 2007, 2007, 2007, 2007, 2007, 2007, 2007, 2007, 2007~
\end{verbatim}

\begin{Shaded}
\begin{Highlighting}[]
\FunctionTok{summary}\NormalTok{(penguins)}
\end{Highlighting}
\end{Shaded}

\begin{verbatim}
      species          island    bill_length_mm  bill_depth_mm  
 Adelie   :152   Biscoe   :168   Min.   :32.10   Min.   :13.10  
 Chinstrap: 68   Dream    :124   1st Qu.:39.23   1st Qu.:15.60  
 Gentoo   :124   Torgersen: 52   Median :44.45   Median :17.30  
                                 Mean   :43.92   Mean   :17.15  
                                 3rd Qu.:48.50   3rd Qu.:18.70  
                                 Max.   :59.60   Max.   :21.50  
                                 NA's   :2       NA's   :2      
 flipper_length_mm  body_mass_g       sex           year     
 Min.   :172.0     Min.   :2700   female:165   Min.   :2007  
 1st Qu.:190.0     1st Qu.:3550   male  :168   1st Qu.:2007  
 Median :197.0     Median :4050   NA's  : 11   Median :2008  
 Mean   :200.9     Mean   :4202                Mean   :2008  
 3rd Qu.:213.0     3rd Qu.:4750                3rd Qu.:2009  
 Max.   :231.0     Max.   :6300                Max.   :2009  
 NA's   :2         NA's   :2                                 
\end{verbatim}

\section{Descriptive statistics}\label{descriptive-statistics-1}

\subsection{Exercise 1}\label{exercise-1-5}

If we don't specify \texttt{na.rm\ =\ TRUE} the \texttt{mean} function
will include the missing values and the function will return a missing
value, \texttt{NA}.

\begin{Shaded}
\begin{Highlighting}[]
\CommentTok{\# Not specifying na.rm = TRUE}
\FunctionTok{mean}\NormalTok{(penguins}\SpecialCharTok{$}\NormalTok{flipper\_length\_mm)}
\end{Highlighting}
\end{Shaded}

\begin{verbatim}
[1] NA
\end{verbatim}

\subsection{Exercise 2}\label{exercise-2-4}

Calculating the mean of a factor returns a warning saying that the
argument is not numeric or logical and returns \texttt{NA}. What would
the mean of a categorical (\texttt{factor}) mean? Does it make sense?

\begin{Shaded}
\begin{Highlighting}[]
\FunctionTok{mean}\NormalTok{(penguins}\SpecialCharTok{$}\NormalTok{species, }\AttributeTok{na.rm =} \ConstantTok{TRUE}\NormalTok{)}
\end{Highlighting}
\end{Shaded}

\begin{verbatim}
Warning in mean.default(penguins$species, na.rm = TRUE): argument is not
numeric or logical: returning NA
\end{verbatim}

\begin{verbatim}
[1] NA
\end{verbatim}

\subsection{Exercise 3}\label{exercise-3-2}

The median and the 50th percentile (or second quartile) are the same
value! The median is the value that splits the data into two parts with
the same number of observations and the 50th percentile is the value
that puts half of the data below the value and half above. Do you see
that this is two ways of saying the same thing?

\begin{Shaded}
\begin{Highlighting}[]
\CommentTok{\# 50th percentile}
\NormalTok{penguins }\SpecialCharTok{|\textgreater{}} 
  \FunctionTok{pull}\NormalTok{(body\_mass\_g) }\SpecialCharTok{|\textgreater{}} 
  \FunctionTok{quantile}\NormalTok{(}\AttributeTok{prob =} \FloatTok{0.5}\NormalTok{, }\AttributeTok{na.rm =} \ConstantTok{TRUE}\NormalTok{)}
\end{Highlighting}
\end{Shaded}

\begin{verbatim}
 50% 
4050 
\end{verbatim}

\begin{Shaded}
\begin{Highlighting}[]
\CommentTok{\# Median}
\NormalTok{penguins }\SpecialCharTok{|\textgreater{}} 
  \FunctionTok{pull}\NormalTok{(body\_mass\_g) }\SpecialCharTok{|\textgreater{}} 
  \FunctionTok{median}\NormalTok{(}\AttributeTok{na.rm =} \ConstantTok{TRUE}\NormalTok{)}
\end{Highlighting}
\end{Shaded}

\begin{verbatim}
[1] 4050
\end{verbatim}

\subsection{Exercise 4}\label{exercise-4-2}

\texttt{cut}, \texttt{color}, \texttt{clarity} are categorical, the
other numerical. There are no missing values.

\begin{Shaded}
\begin{Highlighting}[]
\FunctionTok{library}\NormalTok{(tidyverse)}
\FunctionTok{summary}\NormalTok{(diamonds)}
\end{Highlighting}
\end{Shaded}

\section{Relationships between
variables}\label{relationships-between-variables}

\subsection{Exercise 1}\label{exercise-1-6}

We see that the values in the matrix correspond to the colors in the
correlation plot. High correlation is indicated by values close to
either 1 (positive correlation) or -1 (negative correlation).

\begin{Shaded}
\begin{Highlighting}[]
\FunctionTok{library}\NormalTok{(ggcorrplot)}
\NormalTok{penguins }\SpecialCharTok{|\textgreater{}}
  \FunctionTok{select}\NormalTok{(}\FunctionTok{where}\NormalTok{(is.numeric)) }\SpecialCharTok{|\textgreater{}}
  \FunctionTok{cor}\NormalTok{(}\AttributeTok{use =} \StringTok{"complete.obs"}\NormalTok{)}
\end{Highlighting}
\end{Shaded}

\begin{verbatim}
                  bill_length_mm bill_depth_mm flipper_length_mm body_mass_g
bill_length_mm        1.00000000   -0.23505287         0.6561813  0.59510982
bill_depth_mm        -0.23505287    1.00000000        -0.5838512 -0.47191562
flipper_length_mm     0.65618134   -0.58385122         1.0000000  0.87120177
body_mass_g           0.59510982   -0.47191562         0.8712018  1.00000000
year                  0.05454458   -0.06035364         0.1696751  0.04220939
                         year
bill_length_mm     0.05454458
bill_depth_mm     -0.06035364
flipper_length_mm  0.16967511
body_mass_g        0.04220939
year               1.00000000
\end{verbatim}

\begin{Shaded}
\begin{Highlighting}[]
\NormalTok{penguins }\SpecialCharTok{|\textgreater{}}
  \FunctionTok{select}\NormalTok{(}\FunctionTok{where}\NormalTok{(is.numeric)) }\SpecialCharTok{|\textgreater{}}
  \FunctionTok{cor}\NormalTok{(}\AttributeTok{use =} \StringTok{"complete.obs"}\NormalTok{) }\SpecialCharTok{|\textgreater{}}
  \FunctionTok{ggcorrplot}\NormalTok{()}
\end{Highlighting}
\end{Shaded}

\includegraphics{solutions_basic_statistics_files/figure-pdf/unnamed-chunk-7-1.pdf}

\chapter{R4DS: Data visualization}\label{r4ds-data-visualization}

In these solutions we use the piping operator,
\texttt{\textbar{}\textgreater{}} (also \texttt{\%\textgreater{}\%}).
You will encounter this later in the book, in Chapter 4.4. It is very
useful to write clear and understandable code. In these exercises the
only difference between using the piping operator or not is

\begin{Shaded}
\begin{Highlighting}[]
\NormalTok{penguins }\SpecialCharTok{|\textgreater{}}
  \FunctionTok{ggplot}\NormalTok{(}\FunctionTok{aes}\NormalTok{(...))}
\end{Highlighting}
\end{Shaded}

instead of

\begin{Shaded}
\begin{Highlighting}[]
\FunctionTok{ggplot}\NormalTok{(}\AttributeTok{data =}\NormalTok{ penguins, }\AttributeTok{mapping =} \FunctionTok{aes}\NormalTok{(...))}
\end{Highlighting}
\end{Shaded}

and it doesn't matter which one you choose to do for this exercise.
However, as you progress in the course, piping will be good practice and
therefore we use it in our solutions so you can use this as a reference
later in the course. We also leave out the explicit arguments
\texttt{data\ =} and \texttt{mapping\ =} as that is in line with the
later parts of the book.

\section{Setup}\label{setup}

First we need to load the necessary libraries.

\begin{Shaded}
\begin{Highlighting}[]
\FunctionTok{library}\NormalTok{(tidyverse)}
\end{Highlighting}
\end{Shaded}

\begin{verbatim}
-- Attaching core tidyverse packages ------------------------ tidyverse 2.0.0 --
v dplyr     1.1.4     v readr     2.1.5
v forcats   1.0.0     v stringr   1.5.1
v ggplot2   3.5.1     v tibble    3.2.1
v lubridate 1.9.3     v tidyr     1.3.1
v purrr     1.0.2     
-- Conflicts ------------------------------------------ tidyverse_conflicts() --
x dplyr::filter() masks stats::filter()
x dplyr::lag()    masks stats::lag()
i Use the conflicted package (<http://conflicted.r-lib.org/>) to force all conflicts to become errors
\end{verbatim}

\begin{Shaded}
\begin{Highlighting}[]
\FunctionTok{library}\NormalTok{(palmerpenguins)}
\end{Highlighting}
\end{Shaded}

It is also good practice to take a \texttt{glimpse} at the data before
we begin, even if we are not looking at something in particular.

\begin{Shaded}
\begin{Highlighting}[]
\FunctionTok{glimpse}\NormalTok{(penguins)}
\end{Highlighting}
\end{Shaded}

\begin{verbatim}
Rows: 344
Columns: 8
$ species           <fct> Adelie, Adelie, Adelie, Adelie, Adelie, Adelie, Adel~
$ island            <fct> Torgersen, Torgersen, Torgersen, Torgersen, Torgerse~
$ bill_length_mm    <dbl> 39.1, 39.5, 40.3, NA, 36.7, 39.3, 38.9, 39.2, 34.1, ~
$ bill_depth_mm     <dbl> 18.7, 17.4, 18.0, NA, 19.3, 20.6, 17.8, 19.6, 18.1, ~
$ flipper_length_mm <int> 181, 186, 195, NA, 193, 190, 181, 195, 193, 190, 186~
$ body_mass_g       <int> 3750, 3800, 3250, NA, 3450, 3650, 3625, 4675, 3475, ~
$ sex               <fct> male, female, female, NA, female, male, female, male~
$ year              <int> 2007, 2007, 2007, 2007, 2007, 2007, 2007, 2007, 2007~
\end{verbatim}

That's a number of penguins. Anyway\ldots{}

\section{2.2 First steps}\label{first-steps}

\subsection{Exercise 1}\label{exercise-1-7}

\begin{Shaded}
\begin{Highlighting}[]
\CommentTok{\# Using piping}
\NormalTok{penguins }\SpecialCharTok{|\textgreater{}} \FunctionTok{nrow}\NormalTok{()}
\end{Highlighting}
\end{Shaded}

\begin{verbatim}
[1] 344
\end{verbatim}

\begin{Shaded}
\begin{Highlighting}[]
\NormalTok{penguins }\SpecialCharTok{|\textgreater{}} \FunctionTok{ncol}\NormalTok{()}
\end{Highlighting}
\end{Shaded}

\begin{verbatim}
[1] 8
\end{verbatim}

\begin{Shaded}
\begin{Highlighting}[]
\CommentTok{\# Using functions with arguments}
\FunctionTok{nrow}\NormalTok{(penguins)}
\end{Highlighting}
\end{Shaded}

\begin{verbatim}
[1] 344
\end{verbatim}

\begin{Shaded}
\begin{Highlighting}[]
\FunctionTok{ncol}\NormalTok{(penguins)}
\end{Highlighting}
\end{Shaded}

\begin{verbatim}
[1] 8
\end{verbatim}

\subsection{Exercise 3}\label{exercise-3-3}

\begin{Shaded}
\begin{Highlighting}[]
\NormalTok{penguins }\SpecialCharTok{|\textgreater{}}
  \FunctionTok{ggplot}\NormalTok{(}\FunctionTok{aes}\NormalTok{(}\AttributeTok{x =}\NormalTok{ bill\_length\_mm, }\AttributeTok{y =}\NormalTok{ bill\_depth\_mm)) }\SpecialCharTok{+}
  \FunctionTok{geom\_point}\NormalTok{()}
\end{Highlighting}
\end{Shaded}

\begin{verbatim}
Warning: Removed 2 rows containing missing values or values outside the scale range
(`geom_point()`).
\end{verbatim}

\includegraphics{solutions_r4ds_data_visualization_files/figure-pdf/unnamed-chunk-4-1.pdf}

There seems to be multiple groups within the data (in this case, the
different species), as there seems to be a positive correlation when
bill length is 0-40 mm and another one when bill length is between 40-60
mm.

\subsection{Exercise 4}\label{exercise-4-3}

\begin{Shaded}
\begin{Highlighting}[]
\NormalTok{penguins }\SpecialCharTok{|\textgreater{}}
  \FunctionTok{ggplot}\NormalTok{(}\FunctionTok{aes}\NormalTok{(}\AttributeTok{x =}\NormalTok{ species, }\AttributeTok{y =}\NormalTok{ bill\_depth\_mm)) }\SpecialCharTok{+}
  \FunctionTok{geom\_boxplot}\NormalTok{()}
\end{Highlighting}
\end{Shaded}

\begin{verbatim}
Warning: Removed 2 rows containing non-finite outside the scale range
(`stat_boxplot()`).
\end{verbatim}

\includegraphics{solutions_r4ds_data_visualization_files/figure-pdf/unnamed-chunk-5-1.pdf}

As suspected, the bill depth is different between the different species.
Box plot is a good graphic to visualize continuous data for different
categories. Another option would be to do separate histograms or density
plots for the different species:

\begin{Shaded}
\begin{Highlighting}[]
\NormalTok{penguins }\SpecialCharTok{|\textgreater{}}
  \FunctionTok{ggplot}\NormalTok{(}\FunctionTok{aes}\NormalTok{(}\AttributeTok{fill =}\NormalTok{ species, }\AttributeTok{x =}\NormalTok{ bill\_depth\_mm)) }\SpecialCharTok{+}
  \FunctionTok{geom\_histogram}\NormalTok{(}\AttributeTok{alpha =} \FloatTok{0.5}\NormalTok{)}
\end{Highlighting}
\end{Shaded}

\begin{verbatim}
`stat_bin()` using `bins = 30`. Pick better value with `binwidth`.
\end{verbatim}

\begin{verbatim}
Warning: Removed 2 rows containing non-finite outside the scale range
(`stat_bin()`).
\end{verbatim}

\includegraphics{solutions_r4ds_data_visualization_files/figure-pdf/unnamed-chunk-6-1.pdf}

Note that we set \texttt{alpha\ =\ 0.5} to make the bins slightly
transparent. Which one to do you prefer?

\subsection{Exercise 5}\label{exercise-5}

We haven't defined any aesthetic mappings! Plots need at least an x or y
aestethic mapped.

\subsection{Exercise 6}\label{exercise-6}

\begin{Shaded}
\begin{Highlighting}[]
\NormalTok{penguins }\SpecialCharTok{|\textgreater{}}
  \FunctionTok{ggplot}\NormalTok{(}\FunctionTok{aes}\NormalTok{(}\AttributeTok{x =}\NormalTok{ bill\_length\_mm, }\AttributeTok{y =}\NormalTok{ bill\_depth\_mm)) }\SpecialCharTok{+}
  \FunctionTok{geom\_point}\NormalTok{(}\AttributeTok{na.rm =} \ConstantTok{TRUE}\NormalTok{)}
\end{Highlighting}
\end{Shaded}

\includegraphics{solutions_r4ds_data_visualization_files/figure-pdf/unnamed-chunk-7-1.pdf}

The \texttt{na.rm} argument specifies how the plot deals with
\emph{missing values}. The default behavior is to remove the NAs, but if
we specify it explicitly, like in this exercise, the plot does not
output any warning. Compare with exercise 3!

\subsection{Exercise 7}\label{exercise-7}

\begin{Shaded}
\begin{Highlighting}[]
\NormalTok{penguins }\SpecialCharTok{|\textgreater{}}
  \FunctionTok{ggplot}\NormalTok{(}\FunctionTok{aes}\NormalTok{(}\AttributeTok{x =}\NormalTok{ bill\_length\_mm, }\AttributeTok{y =}\NormalTok{ bill\_depth\_mm)) }\SpecialCharTok{+}
  \FunctionTok{geom\_point}\NormalTok{(}\AttributeTok{na.rm =} \ConstantTok{TRUE}\NormalTok{) }\SpecialCharTok{+}
  \FunctionTok{labs}\NormalTok{(}\AttributeTok{caption =} \StringTok{"Data come from the palmerpenguins package"}\NormalTok{)}
\end{Highlighting}
\end{Shaded}

\includegraphics{solutions_r4ds_data_visualization_files/figure-pdf/unnamed-chunk-8-1.pdf}

Use \texttt{?labs} to see the documentation for the \texttt{labs}
function.

\subsection{Exercise 8}\label{exercise-8}

\begin{Shaded}
\begin{Highlighting}[]
\NormalTok{penguins }\SpecialCharTok{|\textgreater{}}
  \FunctionTok{ggplot}\NormalTok{(}\FunctionTok{aes}\NormalTok{(}\AttributeTok{x =}\NormalTok{ flipper\_length\_mm, }\AttributeTok{y =}\NormalTok{ body\_mass\_g)) }\SpecialCharTok{+}
  \FunctionTok{geom\_point}\NormalTok{(}\FunctionTok{aes}\NormalTok{(}\AttributeTok{col =}\NormalTok{ bill\_depth\_mm), }\AttributeTok{na.rm =} \ConstantTok{TRUE}\NormalTok{) }\SpecialCharTok{+}
  \FunctionTok{geom\_smooth}\NormalTok{()}
\end{Highlighting}
\end{Shaded}

\begin{verbatim}
`geom_smooth()` using method = 'loess' and formula = 'y ~ x'
\end{verbatim}

\begin{verbatim}
Warning: Removed 2 rows containing non-finite outside the scale range
(`stat_smooth()`).
\end{verbatim}

\includegraphics{solutions_r4ds_data_visualization_files/figure-pdf/unnamed-chunk-9-1.pdf}

Here we defined \texttt{col} only for the \texttt{geom\_point}. Does it
make a difference if we define it globally in \texttt{ggplot}? Spoiler:
it doesn't! It only affects \texttt{geom\_point}, because
\texttt{geom\_smooth} creates a line based on many points, so it doesn't
make sense to color parts of the line.

\subsection{Exercise 9}\label{exercise-9}

\begin{Shaded}
\begin{Highlighting}[]
\FunctionTok{ggplot}\NormalTok{(}
  \AttributeTok{data =}\NormalTok{ penguins,}
  \AttributeTok{mapping =} \FunctionTok{aes}\NormalTok{(}\AttributeTok{x =}\NormalTok{ flipper\_length\_mm, }\AttributeTok{y =}\NormalTok{ body\_mass\_g, }\AttributeTok{color =}\NormalTok{ island)}
\NormalTok{) }\SpecialCharTok{+}
  \FunctionTok{geom\_point}\NormalTok{() }\SpecialCharTok{+}
  \FunctionTok{geom\_smooth}\NormalTok{(}\AttributeTok{se =} \ConstantTok{FALSE}\NormalTok{)}
\end{Highlighting}
\end{Shaded}

\begin{verbatim}
`geom_smooth()` using method = 'loess' and formula = 'y ~ x'
\end{verbatim}

\begin{verbatim}
Warning: Removed 2 rows containing non-finite outside the scale range
(`stat_smooth()`).
\end{verbatim}

\begin{verbatim}
Warning: Removed 2 rows containing missing values or values outside the scale range
(`geom_point()`).
\end{verbatim}

\includegraphics{solutions_r4ds_data_visualization_files/figure-pdf/unnamed-chunk-10-1.pdf}

We now color based on the categorical variable \emph{island} instead.
Now we see that our \texttt{geom\_smooth} indeed is affected, as it now
creates one line for each of the different islands.

We have also set \texttt{se\ =\ FALSE}, which means that we do not see
the shaded area around the curve. \emph{se} stands for Standard Error
and tells us about the uncertainty of the line at different points. We
do not expect you to fully understand or know this concept now, but it
can be useful to at least get an idea of whats going on.

\subsection{Exercise 10}\label{exercise-10}

\begin{Shaded}
\begin{Highlighting}[]
\CommentTok{\# Here we use patchwork to set the plots next to each other.}
\FunctionTok{library}\NormalTok{(patchwork)}

\NormalTok{p1 }\OtherTok{\textless{}{-}} \FunctionTok{ggplot}\NormalTok{(}
  \AttributeTok{data =}\NormalTok{ penguins,}
  \AttributeTok{mapping =} \FunctionTok{aes}\NormalTok{(}\AttributeTok{x =}\NormalTok{ flipper\_length\_mm, }\AttributeTok{y =}\NormalTok{ body\_mass\_g)}
\NormalTok{) }\SpecialCharTok{+}
  \FunctionTok{geom\_point}\NormalTok{() }\SpecialCharTok{+}
  \FunctionTok{geom\_smooth}\NormalTok{()}

\NormalTok{p2 }\OtherTok{\textless{}{-}} \FunctionTok{ggplot}\NormalTok{() }\SpecialCharTok{+}
  \FunctionTok{geom\_point}\NormalTok{(}
    \AttributeTok{data =}\NormalTok{ penguins,}
    \AttributeTok{mapping =} \FunctionTok{aes}\NormalTok{(}\AttributeTok{x =}\NormalTok{ flipper\_length\_mm, }\AttributeTok{y =}\NormalTok{ body\_mass\_g)}
\NormalTok{  ) }\SpecialCharTok{+}
  \FunctionTok{geom\_smooth}\NormalTok{(}
    \AttributeTok{data =}\NormalTok{ penguins,}
    \AttributeTok{mapping =} \FunctionTok{aes}\NormalTok{(}\AttributeTok{x =}\NormalTok{ flipper\_length\_mm, }\AttributeTok{y =}\NormalTok{ body\_mass\_g)}
\NormalTok{  )}

\NormalTok{p1 }\SpecialCharTok{+}\NormalTok{ p2}
\end{Highlighting}
\end{Shaded}

\begin{verbatim}
`geom_smooth()` using method = 'loess' and formula = 'y ~ x'
\end{verbatim}

\begin{verbatim}
Warning: Removed 2 rows containing non-finite outside the scale range
(`stat_smooth()`).
\end{verbatim}

\begin{verbatim}
Warning: Removed 2 rows containing missing values or values outside the scale range
(`geom_point()`).
\end{verbatim}

\begin{verbatim}
`geom_smooth()` using method = 'loess' and formula = 'y ~ x'
\end{verbatim}

\begin{verbatim}
Warning: Removed 2 rows containing non-finite outside the scale range (`stat_smooth()`).
Removed 2 rows containing missing values or values outside the scale range
(`geom_point()`).
\end{verbatim}

\includegraphics{solutions_r4ds_data_visualization_files/figure-pdf/unnamed-chunk-11-1.pdf}

We see that it doesn't make a difference. We can define aesthetics
globally in \texttt{ggplot} or locally in each \texttt{geom\_}. However,
it is better to define as much as possible in \texttt{ggplot}, as it
makes it clear for you writing the code and someone reading the code
what the plot will show.

Here we used \texttt{library(patchwork} to put the plots next to each
other. You can try it out, but you don't have to use it!

\section{2.4 Visualizing distributions}\label{visualizing-distributions}

\subsection{Exercise 1}\label{exercise-1-8}

\begin{Shaded}
\begin{Highlighting}[]
\NormalTok{penguins }\SpecialCharTok{|\textgreater{}}
  \FunctionTok{ggplot}\NormalTok{(}\FunctionTok{aes}\NormalTok{(}\AttributeTok{y =}\NormalTok{ species)) }\SpecialCharTok{+}
  \FunctionTok{geom\_bar}\NormalTok{()}
\end{Highlighting}
\end{Shaded}

\includegraphics{solutions_r4ds_data_visualization_files/figure-pdf/unnamed-chunk-12-1.pdf}

The bars now extend horizontally.

\subsection{Exercise 2}\label{exercise-2-5}

\begin{Shaded}
\begin{Highlighting}[]
\FunctionTok{library}\NormalTok{(patchwork)}
\NormalTok{p1 }\OtherTok{\textless{}{-}}\NormalTok{ penguins }\SpecialCharTok{|\textgreater{}}
  \FunctionTok{ggplot}\NormalTok{(}\FunctionTok{aes}\NormalTok{(}\AttributeTok{y =}\NormalTok{ species)) }\SpecialCharTok{+}
  \FunctionTok{geom\_bar}\NormalTok{(}\AttributeTok{col =} \StringTok{"red"}\NormalTok{)}
\NormalTok{p2 }\OtherTok{\textless{}{-}}\NormalTok{ penguins }\SpecialCharTok{|\textgreater{}}
  \FunctionTok{ggplot}\NormalTok{(}\FunctionTok{aes}\NormalTok{(}\AttributeTok{y =}\NormalTok{ species)) }\SpecialCharTok{+}
  \FunctionTok{geom\_bar}\NormalTok{(}\AttributeTok{fill =} \StringTok{"red"}\NormalTok{)}

\NormalTok{p1 }\SpecialCharTok{+}\NormalTok{ p2}
\end{Highlighting}
\end{Shaded}

\includegraphics{solutions_r4ds_data_visualization_files/figure-pdf/unnamed-chunk-13-1.pdf}

\texttt{fill} is usually the more interesting property to adjust when it
is applicable. Colors and fills are useful to play around with to make
it easier for the reader, and sometimes ourselves. Maybe a more useful
coloring is one color for each species and a colorscheme that is
suitable to colorblind people?

\begin{Shaded}
\begin{Highlighting}[]
\CommentTok{\# We import the ggthemes package to get the scale\_fill\_colorblind}
\FunctionTok{library}\NormalTok{(ggthemes)}
\NormalTok{penguins }\SpecialCharTok{|\textgreater{}}
  \FunctionTok{ggplot}\NormalTok{(}\FunctionTok{aes}\NormalTok{(}\AttributeTok{y =}\NormalTok{ species, }\AttributeTok{fill =}\NormalTok{ species)) }\SpecialCharTok{+}
  \FunctionTok{geom\_bar}\NormalTok{() }\SpecialCharTok{+} 
  \FunctionTok{scale\_fill\_colorblind}\NormalTok{()}
\end{Highlighting}
\end{Shaded}

\includegraphics{solutions_r4ds_data_visualization_files/figure-pdf/unnamed-chunk-14-1.pdf}

\subsection{Exercise 3}\label{exercise-3-4}

To find out, we can look at \texttt{?geom\_histogram}.

\begin{quote}
bins\\
Number of bins. Overridden by binwidth. Defaults to 30.
\end{quote}

Selecting number of bins can affect the perception of the histogram, and
of the data.

\begin{Shaded}
\begin{Highlighting}[]
\FunctionTok{library}\NormalTok{(patchwork)}
\NormalTok{p1 }\OtherTok{\textless{}{-}}\NormalTok{ penguins }\SpecialCharTok{|\textgreater{}}
  \FunctionTok{ggplot}\NormalTok{(}\FunctionTok{aes}\NormalTok{(}\AttributeTok{fill =}\NormalTok{ species, }\AttributeTok{x =}\NormalTok{ bill\_depth\_mm)) }\SpecialCharTok{+}
  \FunctionTok{geom\_histogram}\NormalTok{(}\AttributeTok{bins =} \DecValTok{50}\NormalTok{, }\AttributeTok{alpha =} \FloatTok{0.5}\NormalTok{)}
\NormalTok{p2 }\OtherTok{\textless{}{-}}\NormalTok{ penguins }\SpecialCharTok{|\textgreater{}}
  \FunctionTok{ggplot}\NormalTok{(}\FunctionTok{aes}\NormalTok{(}\AttributeTok{fill =}\NormalTok{ species, }\AttributeTok{x =}\NormalTok{ bill\_depth\_mm)) }\SpecialCharTok{+}
  \FunctionTok{geom\_histogram}\NormalTok{(}\AttributeTok{bins =} \DecValTok{10}\NormalTok{, }\AttributeTok{alpha =} \FloatTok{0.5}\NormalTok{)}

\NormalTok{p1 }\SpecialCharTok{+}\NormalTok{ p2}
\end{Highlighting}
\end{Shaded}

\begin{verbatim}
Warning: Removed 2 rows containing non-finite outside the scale range (`stat_bin()`).
Removed 2 rows containing non-finite outside the scale range (`stat_bin()`).
\end{verbatim}

\includegraphics{solutions_r4ds_data_visualization_files/figure-pdf/unnamed-chunk-15-1.pdf}

Which one is better depends on the data and what features you are
interested in. If the histogram doesn't look like you would expect,
changing number of bins may be a good idea. Try using \texttt{binwidth}
as well, as that may be more intuitive to derive from the data.

\subsection{Exercise 4}\label{exercise-4-4}

First it's a good idea to get a \texttt{glimpse} of the
\texttt{diamonds} data. \texttt{glimpse} is a function that gives a
quick overview, a \emph{glimpse}, of a data frame.

\begin{Shaded}
\begin{Highlighting}[]
\FunctionTok{glimpse}\NormalTok{(diamonds)}
\end{Highlighting}
\end{Shaded}

\begin{verbatim}
Rows: 53,940
Columns: 10
$ carat   <dbl> 0.23, 0.21, 0.23, 0.29, 0.31, 0.24, 0.24, 0.26, 0.22, 0.23, 0.~
$ cut     <ord> Ideal, Premium, Good, Premium, Good, Very Good, Very Good, Ver~
$ color   <ord> E, E, E, I, J, J, I, H, E, H, J, J, F, J, E, E, I, J, J, J, I,~
$ clarity <ord> SI2, SI1, VS1, VS2, SI2, VVS2, VVS1, SI1, VS2, VS1, SI1, VS1, ~
$ depth   <dbl> 61.5, 59.8, 56.9, 62.4, 63.3, 62.8, 62.3, 61.9, 65.1, 59.4, 64~
$ table   <dbl> 55, 61, 65, 58, 58, 57, 57, 55, 61, 61, 55, 56, 61, 54, 62, 58~
$ price   <int> 326, 326, 327, 334, 335, 336, 336, 337, 337, 338, 339, 340, 34~
$ x       <dbl> 3.95, 3.89, 4.05, 4.20, 4.34, 3.94, 3.95, 4.07, 3.87, 4.00, 4.~
$ y       <dbl> 3.98, 3.84, 4.07, 4.23, 4.35, 3.96, 3.98, 4.11, 3.78, 4.05, 4.~
$ z       <dbl> 2.43, 2.31, 2.31, 2.63, 2.75, 2.48, 2.47, 2.53, 2.49, 2.39, 2.~
\end{verbatim}

Then we can create the histogram over \texttt{carat}:

\begin{Shaded}
\begin{Highlighting}[]
\NormalTok{diamonds }\SpecialCharTok{|\textgreater{}}
  \FunctionTok{ggplot}\NormalTok{(}\FunctionTok{aes}\NormalTok{(}\AttributeTok{x =}\NormalTok{ carat)) }\SpecialCharTok{+}
  \FunctionTok{geom\_histogram}\NormalTok{()}
\end{Highlighting}
\end{Shaded}

\begin{verbatim}
`stat_bin()` using `bins = 30`. Pick better value with `binwidth`.
\end{verbatim}

\includegraphics{solutions_r4ds_data_visualization_files/figure-pdf/unnamed-chunk-17-1.pdf}

\subsection{Extra exercise}\label{extra-exercise}

\begin{Shaded}
\begin{Highlighting}[]
\NormalTok{diamonds }\SpecialCharTok{|\textgreater{}}
  \FunctionTok{ggplot}\NormalTok{(}\FunctionTok{aes}\NormalTok{(}\AttributeTok{x =}\NormalTok{ carat)) }\SpecialCharTok{+}
  \FunctionTok{geom\_density}\NormalTok{()}
\end{Highlighting}
\end{Shaded}

\includegraphics{solutions_r4ds_data_visualization_files/figure-pdf/unnamed-chunk-18-1.pdf}

Using \texttt{geom\_density} is another way of visualizing a
distribution. It is less detailed than the histogram but gives a good
overview of the data. It can also be useful to visualize many categories
at the same time, which can get messy when using a histogram. Note
however that we do not escape the problem of selecting a suitable number
of bins, as something similar exists in geom\_density, the parameter
\texttt{bw}. We do not expect you to understand this parameter or even
use it, but it is good to be aware of it.

An example of multiple densities, colored by \texttt{cut}.

\begin{Shaded}
\begin{Highlighting}[]
\NormalTok{diamonds }\SpecialCharTok{|\textgreater{}}
  \FunctionTok{ggplot}\NormalTok{(}\FunctionTok{aes}\NormalTok{(}\AttributeTok{x =}\NormalTok{ carat, }\AttributeTok{col =}\NormalTok{ cut)) }\SpecialCharTok{+}
  \FunctionTok{geom\_density}\NormalTok{() }
\end{Highlighting}
\end{Shaded}

\includegraphics{solutions_r4ds_data_visualization_files/figure-pdf/unnamed-chunk-19-1.pdf}

\section{2.5 Visualizing relationships}\label{visualizing-relationships}

\subsection{Exercise 1}\label{exercise-1-9}

Again we can use \texttt{glimpse} to get an overview, and get to know
which kinds of variables the data frame contains:

\begin{Shaded}
\begin{Highlighting}[]
\FunctionTok{glimpse}\NormalTok{(mpg)}
\end{Highlighting}
\end{Shaded}

\begin{verbatim}
Rows: 234
Columns: 11
$ manufacturer <chr> "audi", "audi", "audi", "audi", "audi", "audi", "audi", "~
$ model        <chr> "a4", "a4", "a4", "a4", "a4", "a4", "a4", "a4 quattro", "~
$ displ        <dbl> 1.8, 1.8, 2.0, 2.0, 2.8, 2.8, 3.1, 1.8, 1.8, 2.0, 2.0, 2.~
$ year         <int> 1999, 1999, 2008, 2008, 1999, 1999, 2008, 1999, 1999, 200~
$ cyl          <int> 4, 4, 4, 4, 6, 6, 6, 4, 4, 4, 4, 6, 6, 6, 6, 6, 6, 8, 8, ~
$ trans        <chr> "auto(l5)", "manual(m5)", "manual(m6)", "auto(av)", "auto~
$ drv          <chr> "f", "f", "f", "f", "f", "f", "f", "4", "4", "4", "4", "4~
$ cty          <int> 18, 21, 20, 21, 16, 18, 18, 18, 16, 20, 19, 15, 17, 17, 1~
$ hwy          <int> 29, 29, 31, 30, 26, 26, 27, 26, 25, 28, 27, 25, 25, 25, 2~
$ fl           <chr> "p", "p", "p", "p", "p", "p", "p", "p", "p", "p", "p", "p~
$ class        <chr> "compact", "compact", "compact", "compact", "compact", "c~
\end{verbatim}

\subsection{Exercise 2}\label{exercise-2-6}

One example using \texttt{color} and \texttt{size} at the same time.

\begin{Shaded}
\begin{Highlighting}[]
\NormalTok{mpg }\SpecialCharTok{|\textgreater{}}
  \FunctionTok{ggplot}\NormalTok{(}\FunctionTok{aes}\NormalTok{(}\AttributeTok{x =}\NormalTok{ hwy, }\AttributeTok{y =}\NormalTok{ displ, }\AttributeTok{col =}\NormalTok{ cty, }\AttributeTok{size =}\NormalTok{ cyl)) }\SpecialCharTok{+}
  \FunctionTok{geom\_point}\NormalTok{()}
\end{Highlighting}
\end{Shaded}

\includegraphics{solutions_r4ds_data_visualization_files/figure-pdf/unnamed-chunk-21-1.pdf}

Note that we cannot assign \texttt{shape} to a continuous variable! For
that we need the variable to be a \emph{factor}, as in this example

\begin{Shaded}
\begin{Highlighting}[]
\NormalTok{mpg }\SpecialCharTok{|\textgreater{}}
  \FunctionTok{ggplot}\NormalTok{(}\FunctionTok{aes}\NormalTok{(}\AttributeTok{x =}\NormalTok{ hwy, }\AttributeTok{y =}\NormalTok{ displ, }\AttributeTok{col =}\NormalTok{ cty, }\AttributeTok{shape =} \FunctionTok{as.factor}\NormalTok{(cyl))) }\SpecialCharTok{+}
  \FunctionTok{geom\_point}\NormalTok{()}
\end{Highlighting}
\end{Shaded}

\includegraphics{solutions_r4ds_data_visualization_files/figure-pdf/unnamed-chunk-22-1.pdf}

\subsection{Exercise 3}\label{exercise-3-5}

\texttt{linewidth} doesn't have any effect on a scatterplot,
\texttt{geom\_point}, so it will be ignored.

\subsection{Exercise 4}\label{exercise-4-5}

It is possible to map the same variable to multiple aesthetics and it
can even be useful sometimes, for example when we have a reason to use a
specific color theme, but also want our plot to be easily understood by
someone who is colorblind, see example below.

\begin{Shaded}
\begin{Highlighting}[]
\NormalTok{mpg }\SpecialCharTok{|\textgreater{}}
  \FunctionTok{ggplot}\NormalTok{(}\FunctionTok{aes}\NormalTok{(}\AttributeTok{x =}\NormalTok{ hwy, }\AttributeTok{y =}\NormalTok{ displ, }\AttributeTok{col =} \FunctionTok{as.factor}\NormalTok{(cyl), }\AttributeTok{shape =} \FunctionTok{as.factor}\NormalTok{(cyl))) }\SpecialCharTok{+}
  \FunctionTok{geom\_point}\NormalTok{()}
\end{Highlighting}
\end{Shaded}

\includegraphics{solutions_r4ds_data_visualization_files/figure-pdf/unnamed-chunk-23-1.pdf}

\subsection{Exercise 5}\label{exercise-5-1}

As we saw in 2.2.4, exercise 3, there seemed to be groups within the
scatter plot between \emph{bill\_length\_mm} and \emph{bill\_depth\_mm}.
Now we will investigate this by coloring the plot by species.

\begin{Shaded}
\begin{Highlighting}[]
\NormalTok{penguins }\SpecialCharTok{|\textgreater{}}
  \FunctionTok{ggplot}\NormalTok{(}\FunctionTok{aes}\NormalTok{(}\AttributeTok{x =}\NormalTok{ bill\_length\_mm, }\AttributeTok{y =}\NormalTok{ bill\_depth\_mm, }\AttributeTok{col =}\NormalTok{ species)) }\SpecialCharTok{+}
  \FunctionTok{geom\_point}\NormalTok{()}
\end{Highlighting}
\end{Shaded}

\begin{verbatim}
Warning: Removed 2 rows containing missing values or values outside the scale range
(`geom_point()`).
\end{verbatim}

\includegraphics{solutions_r4ds_data_visualization_files/figure-pdf/unnamed-chunk-24-1.pdf}

Now it becomes obvious that the different species in our data was
responsible for different patterns.

Interesting fact: This is an example of Simpson's paradox: when a
relationship between two variables doesn't show up for the whole
population (all penguins) but shows up in sub-populations (the different
species).

We will try faceting as well

\begin{Shaded}
\begin{Highlighting}[]
\NormalTok{penguins }\SpecialCharTok{|\textgreater{}}
  \FunctionTok{ggplot}\NormalTok{(}\FunctionTok{aes}\NormalTok{(}\AttributeTok{x =}\NormalTok{ bill\_length\_mm, }\AttributeTok{y =}\NormalTok{ bill\_depth\_mm, }\AttributeTok{col =}\NormalTok{ species)) }\SpecialCharTok{+}
  \FunctionTok{geom\_point}\NormalTok{() }\SpecialCharTok{+}
  \FunctionTok{facet\_wrap}\NormalTok{(}\SpecialCharTok{\textasciitilde{}}\NormalTok{species)}
\end{Highlighting}
\end{Shaded}

\begin{verbatim}
Warning: Removed 2 rows containing missing values or values outside the scale range
(`geom_point()`).
\end{verbatim}

\includegraphics{solutions_r4ds_data_visualization_files/figure-pdf/unnamed-chunk-25-1.pdf}

Which one do you think is most useful?

\subsection{Exercise 6}\label{exercise-6-1}

Legends take the name from the aesthetic mapping that creates them, in
this case \texttt{species}. When we manually rename a legend, as in
\texttt{labs(color\ =\ "Species")}, it only affects the specified
aesthetic, i.e.~\texttt{color}. We can remedy this by specifying the
same label for shape, as below.

\begin{Shaded}
\begin{Highlighting}[]
\NormalTok{penguins }\SpecialCharTok{|\textgreater{}}
  \FunctionTok{ggplot}\NormalTok{(}\FunctionTok{aes}\NormalTok{(}\AttributeTok{x =}\NormalTok{ bill\_length\_mm, }\AttributeTok{y =}\NormalTok{ bill\_depth\_mm, }\AttributeTok{col =}\NormalTok{ species, }\AttributeTok{shape =}\NormalTok{ species)) }\SpecialCharTok{+}
  \FunctionTok{geom\_point}\NormalTok{() }\SpecialCharTok{+}
  \FunctionTok{labs}\NormalTok{(}\AttributeTok{color =} \StringTok{"Species"}\NormalTok{,}
       \AttributeTok{shape =} \StringTok{"Species"}\NormalTok{)}
\end{Highlighting}
\end{Shaded}

\begin{verbatim}
Warning: Removed 2 rows containing missing values or values outside the scale range
(`geom_point()`).
\end{verbatim}

\includegraphics{solutions_r4ds_data_visualization_files/figure-pdf/unnamed-chunk-26-1.pdf}

\subsection{Exercise 7}\label{exercise-7-1}

\begin{Shaded}
\begin{Highlighting}[]
\FunctionTok{library}\NormalTok{(patchwork)}
\NormalTok{p1 }\OtherTok{\textless{}{-}}\NormalTok{ penguins }\SpecialCharTok{|\textgreater{}}
  \FunctionTok{ggplot}\NormalTok{(}\FunctionTok{aes}\NormalTok{(}\AttributeTok{x =}\NormalTok{ island, }\AttributeTok{fill =}\NormalTok{ species)) }\SpecialCharTok{+}
  \FunctionTok{geom\_bar}\NormalTok{(}\AttributeTok{position =} \StringTok{"fill"}\NormalTok{)}
\NormalTok{p2 }\OtherTok{\textless{}{-}}\NormalTok{ penguins }\SpecialCharTok{|\textgreater{}}
  \FunctionTok{ggplot}\NormalTok{(}\FunctionTok{aes}\NormalTok{(}\AttributeTok{x =}\NormalTok{ species, }\AttributeTok{fill =}\NormalTok{ island)) }\SpecialCharTok{+}
  \FunctionTok{geom\_bar}\NormalTok{(}\AttributeTok{position =} \StringTok{"fill"}\NormalTok{)}

\NormalTok{p1 }\SpecialCharTok{+}\NormalTok{ p2}
\end{Highlighting}
\end{Shaded}

\includegraphics{solutions_r4ds_data_visualization_files/figure-pdf/unnamed-chunk-27-1.pdf}

In the first plot, \texttt{p1}, we can answer what fraction of the total
number of penguins at each island is of a particular species. In
\texttt{p2} we can answer what fraction of each species resides on which
island.

\section{2.6 Saving your plots}\label{saving-your-plots}

Note that it is good practice to save you plots as PDF whenever it is
possible. PDF is a form of \emph{vector graphics}, this enables you to
zoom as much as you want without the image becoming ``pixely''. It is
also good practice to always reflect on and specify \texttt{width} and
\texttt{height} of your plots to make sure the format suits your needs.

\subsection{Exercise 1 \& 2}\label{exercise-1-2-1}

The last plot is always the one to be saved. If you have multiple plots
and want to save all of them, put a \texttt{ggsave} after each plot. To
save as pdf, specify \texttt{.pdf} in the filename.

\begin{Shaded}
\begin{Highlighting}[]
\NormalTok{penguins }\SpecialCharTok{|\textgreater{}}
  \FunctionTok{ggplot}\NormalTok{(}\FunctionTok{aes}\NormalTok{(}\AttributeTok{x =}\NormalTok{ island, }\AttributeTok{fill =}\NormalTok{ species)) }\SpecialCharTok{+}
  \FunctionTok{geom\_bar}\NormalTok{(}\AttributeTok{position =} \StringTok{"fill"}\NormalTok{)}
\FunctionTok{ggsave}\NormalTok{(}\StringTok{"species{-}by{-}island.pdf"}\NormalTok{, }\AttributeTok{width =} \DecValTok{10}\NormalTok{, }\AttributeTok{height =} \DecValTok{8}\NormalTok{)}
\NormalTok{penguins }\SpecialCharTok{|\textgreater{}}
  \FunctionTok{ggplot}\NormalTok{(}\FunctionTok{aes}\NormalTok{(}\AttributeTok{x =}\NormalTok{ species, }\AttributeTok{fill =}\NormalTok{ island)) }\SpecialCharTok{+}
  \FunctionTok{geom\_bar}\NormalTok{(}\AttributeTok{position =} \StringTok{"fill"}\NormalTok{)}
\FunctionTok{ggsave}\NormalTok{(}\StringTok{"island{-}by{-}species.pdf"}\NormalTok{, }\AttributeTok{width =} \DecValTok{10}\NormalTok{, }\AttributeTok{height =} \DecValTok{8}\NormalTok{)}
\end{Highlighting}
\end{Shaded}

\chapter{R4DS: Data transformation}\label{r4ds-data-transformation}

To run these solutions you need to load the necessary libraries.

\begin{Shaded}
\begin{Highlighting}[]
\FunctionTok{library}\NormalTok{(tidyverse)}
\end{Highlighting}
\end{Shaded}

\begin{verbatim}
-- Attaching core tidyverse packages ------------------------ tidyverse 2.0.0 --
v dplyr     1.1.4     v readr     2.1.5
v forcats   1.0.0     v stringr   1.5.1
v ggplot2   3.5.1     v tibble    3.2.1
v lubridate 1.9.3     v tidyr     1.3.1
v purrr     1.0.2     
-- Conflicts ------------------------------------------ tidyverse_conflicts() --
x dplyr::filter() masks stats::filter()
x dplyr::lag()    masks stats::lag()
i Use the conflicted package (<http://conflicted.r-lib.org/>) to force all conflicts to become errors
\end{verbatim}

\begin{Shaded}
\begin{Highlighting}[]
\FunctionTok{library}\NormalTok{(nycflights13)}
\end{Highlighting}
\end{Shaded}

As always, its good to take a \texttt{glimpse} at the data

\begin{Shaded}
\begin{Highlighting}[]
\FunctionTok{glimpse}\NormalTok{(flights)}
\end{Highlighting}
\end{Shaded}

\begin{verbatim}
Rows: 336,776
Columns: 19
$ year           <int> 2013, 2013, 2013, 2013, 2013, 2013, 2013, 2013, 2013, 2~
$ month          <int> 1, 1, 1, 1, 1, 1, 1, 1, 1, 1, 1, 1, 1, 1, 1, 1, 1, 1, 1~
$ day            <int> 1, 1, 1, 1, 1, 1, 1, 1, 1, 1, 1, 1, 1, 1, 1, 1, 1, 1, 1~
$ dep_time       <int> 517, 533, 542, 544, 554, 554, 555, 557, 557, 558, 558, ~
$ sched_dep_time <int> 515, 529, 540, 545, 600, 558, 600, 600, 600, 600, 600, ~
$ dep_delay      <dbl> 2, 4, 2, -1, -6, -4, -5, -3, -3, -2, -2, -2, -2, -2, -1~
$ arr_time       <int> 830, 850, 923, 1004, 812, 740, 913, 709, 838, 753, 849,~
$ sched_arr_time <int> 819, 830, 850, 1022, 837, 728, 854, 723, 846, 745, 851,~
$ arr_delay      <dbl> 11, 20, 33, -18, -25, 12, 19, -14, -8, 8, -2, -3, 7, -1~
$ carrier        <chr> "UA", "UA", "AA", "B6", "DL", "UA", "B6", "EV", "B6", "~
$ flight         <int> 1545, 1714, 1141, 725, 461, 1696, 507, 5708, 79, 301, 4~
$ tailnum        <chr> "N14228", "N24211", "N619AA", "N804JB", "N668DN", "N394~
$ origin         <chr> "EWR", "LGA", "JFK", "JFK", "LGA", "EWR", "EWR", "LGA",~
$ dest           <chr> "IAH", "IAH", "MIA", "BQN", "ATL", "ORD", "FLL", "IAD",~
$ air_time       <dbl> 227, 227, 160, 183, 116, 150, 158, 53, 140, 138, 149, 1~
$ distance       <dbl> 1400, 1416, 1089, 1576, 762, 719, 1065, 229, 944, 733, ~
$ hour           <dbl> 5, 5, 5, 5, 6, 5, 6, 6, 6, 6, 6, 6, 6, 6, 6, 5, 6, 6, 6~
$ minute         <dbl> 15, 29, 40, 45, 0, 58, 0, 0, 0, 0, 0, 0, 0, 0, 0, 59, 0~
$ time_hour      <dttm> 2013-01-01 05:00:00, 2013-01-01 05:00:00, 2013-01-01 0~
\end{verbatim}

That's a lot of flights! Anyway\ldots{}

\section{4.2 Rows}\label{rows}

\subsection{Exercise 1}\label{exercise-1-10}

TO avoid extensive outputs, we only provide the code for these
solutions.

\begin{Shaded}
\begin{Highlighting}[]
\CommentTok{\# Had an arrival delay of two or more hours}
\NormalTok{flights }\SpecialCharTok{|\textgreater{}}
  \FunctionTok{filter}\NormalTok{(arr\_delay }\SpecialCharTok{\textgreater{}} \DecValTok{120}\NormalTok{)}

\CommentTok{\# Flew to Houston (IAH or HOU)}
\NormalTok{flights }\SpecialCharTok{|\textgreater{}}
  \FunctionTok{filter}\NormalTok{(dest }\SpecialCharTok{==} \StringTok{"IAH"} \SpecialCharTok{|}\NormalTok{ dest }\SpecialCharTok{==} \StringTok{"HOU"}\NormalTok{ )}

\CommentTok{\# Were operated by United, American, or Delta}
\CommentTok{\# Here we use "|" for or.}
\NormalTok{flights }\SpecialCharTok{|\textgreater{}}
  \FunctionTok{filter}\NormalTok{(carrier }\SpecialCharTok{==} \StringTok{"UA"} \SpecialCharTok{|}\NormalTok{ carrier }\SpecialCharTok{==} \StringTok{"AA"} \SpecialCharTok{|}\NormalTok{ carrier }\SpecialCharTok{==} \StringTok{"DL"}\NormalTok{)}

\CommentTok{\# Here we use "\%in\%".}
\NormalTok{flights }\SpecialCharTok{|\textgreater{}}
  \FunctionTok{filter}\NormalTok{(carrier }\SpecialCharTok{\%in\%} \FunctionTok{c}\NormalTok{(}\StringTok{"UA"}\NormalTok{, }\StringTok{"AA"}\NormalTok{, }\StringTok{"DL"}\NormalTok{))}

\CommentTok{\# Departed in summer (July, August, and September)}
\NormalTok{flights }\SpecialCharTok{|\textgreater{}}
  \FunctionTok{filter}\NormalTok{(month }\SpecialCharTok{\%in\%} \FunctionTok{c}\NormalTok{(}\DecValTok{7}\NormalTok{, }\DecValTok{8}\NormalTok{, }\DecValTok{9}\NormalTok{)) }

\CommentTok{\# Arrived more than two hours late, but didn’t leave late}
\NormalTok{flights }\SpecialCharTok{|\textgreater{}}
  \FunctionTok{filter}\NormalTok{(arr\_delay }\SpecialCharTok{\textgreater{}} \DecValTok{120} \SpecialCharTok{\&}\NormalTok{ dep\_delay }\SpecialCharTok{\textless{}=} \DecValTok{0}\NormalTok{) }

\CommentTok{\# Were delayed by at least an hour, but made up over 30 minutes in flight}
\NormalTok{flights }\SpecialCharTok{|\textgreater{}}
  \FunctionTok{filter}\NormalTok{(dep\_delay }\SpecialCharTok{\textgreater{}} \DecValTok{60} \SpecialCharTok{\&}\NormalTok{ arr\_delay }\SpecialCharTok{\textless{}} \DecValTok{30}\NormalTok{) }
\end{Highlighting}
\end{Shaded}

In this solution we used both \texttt{\textbar{}} and \texttt{\%in\%}.
When you have multiple conditions the \texttt{\%in\%} operator can be
more concise and easier to understand. \texttt{\textbar{}} or more
specifically \texttt{\textbar{}\textbar{}} is a more traditional way of
indicating \emph{or} conditions, but this is outside of the scope of
this course.

\subsection{Exercise 2}\label{exercise-2-7}

Here we fist \texttt{arrange} by \texttt{dep\_delay} in
\texttt{desc}-ending order to find the flights with the longest delay.
We also use \texttt{relocate}, which is introduced in 4.2.6 to move the
column to the left in the tibble so it is visible in our output.

\begin{Shaded}
\begin{Highlighting}[]
\NormalTok{flights }\SpecialCharTok{|\textgreater{}}
  \FunctionTok{arrange}\NormalTok{(}\FunctionTok{desc}\NormalTok{(dep\_delay)) }\SpecialCharTok{|\textgreater{}}
  \FunctionTok{relocate}\NormalTok{(dep\_delay) }\SpecialCharTok{|\textgreater{}}
  \FunctionTok{head}\NormalTok{(}\DecValTok{3}\NormalTok{)}
\end{Highlighting}
\end{Shaded}

\begin{verbatim}
# A tibble: 3 x 19
  dep_delay  year month   day dep_time sched_dep_time arr_time sched_arr_time
      <dbl> <int> <int> <int>    <int>          <int>    <int>          <int>
1      1301  2013     1     9      641            900     1242           1530
2      1137  2013     6    15     1432           1935     1607           2120
3      1126  2013     1    10     1121           1635     1239           1810
# i 11 more variables: arr_delay <dbl>, carrier <chr>, flight <int>,
#   tailnum <chr>, origin <chr>, dest <chr>, air_time <dbl>, distance <dbl>,
#   hour <dbl>, minute <dbl>, time_hour <dttm>
\end{verbatim}

Then we \texttt{arrange} by \texttt{dep\_time} to find the flights that
left earliest in the morning. Here we use select, also introduced in
4.2.6 to select only the column that we arranged by.

\begin{Shaded}
\begin{Highlighting}[]
\NormalTok{flights }\SpecialCharTok{|\textgreater{}}
  \FunctionTok{arrange}\NormalTok{(dep\_time) }\SpecialCharTok{|\textgreater{}}
  \FunctionTok{select}\NormalTok{(dep\_time) }\SpecialCharTok{|\textgreater{}}
  \FunctionTok{head}\NormalTok{(}\DecValTok{3}\NormalTok{)}
\end{Highlighting}
\end{Shaded}

\begin{verbatim}
# A tibble: 3 x 1
  dep_time
     <int>
1        1
2        1
3        1
\end{verbatim}

Did you do it differently? I interpreted the question in this way, but
maybe it was actually supposed to be

\begin{Shaded}
\begin{Highlighting}[]
\NormalTok{flights }\SpecialCharTok{|\textgreater{}}
  \FunctionTok{arrange}\NormalTok{(}\FunctionTok{desc}\NormalTok{(dep\_delay), dep\_time) }\SpecialCharTok{|\textgreater{}}
  \FunctionTok{relocate}\NormalTok{(dep\_delay, dep\_time) }\SpecialCharTok{|\textgreater{}}
  \FunctionTok{head}\NormalTok{(}\DecValTok{3}\NormalTok{)}
\end{Highlighting}
\end{Shaded}

\begin{verbatim}
# A tibble: 3 x 19
  dep_delay dep_time  year month   day sched_dep_time arr_time sched_arr_time
      <dbl>    <int> <int> <int> <int>          <int>    <int>          <int>
1      1301      641  2013     1     9            900     1242           1530
2      1137     1432  2013     6    15           1935     1607           2120
3      1126     1121  2013     1    10           1635     1239           1810
# i 11 more variables: arr_delay <dbl>, carrier <chr>, flight <int>,
#   tailnum <chr>, origin <chr>, dest <chr>, air_time <dbl>, distance <dbl>,
#   hour <dbl>, minute <dbl>, time_hour <dttm>
\end{verbatim}

\subsection{Exercise 3}\label{exercise-3-6}

Here we remember our high school physics, speed = distance / time.
Remember to use \texttt{desc} to get the

\begin{Shaded}
\begin{Highlighting}[]
\NormalTok{flights }\SpecialCharTok{|\textgreater{}}
  \FunctionTok{arrange}\NormalTok{(}\FunctionTok{desc}\NormalTok{(distance }\SpecialCharTok{/}\NormalTok{ air\_time)) }\SpecialCharTok{|\textgreater{}}
  \FunctionTok{relocate}\NormalTok{(distance, air\_time) }\SpecialCharTok{|\textgreater{}}
  \FunctionTok{head}\NormalTok{(}\DecValTok{3}\NormalTok{)}
\end{Highlighting}
\end{Shaded}

\begin{verbatim}
# A tibble: 3 x 19
  distance air_time  year month   day dep_time sched_dep_time dep_delay arr_time
     <dbl>    <dbl> <int> <int> <int>    <int>          <int>     <dbl>    <int>
1      762       65  2013     5    25     1709           1700         9     1923
2     1008       93  2013     7     2     1558           1513        45     1745
3      594       55  2013     5    13     2040           2025        15     2225
# i 10 more variables: sched_arr_time <int>, arr_delay <dbl>, carrier <chr>,
#   flight <int>, tailnum <chr>, origin <chr>, dest <chr>, hour <dbl>,
#   minute <dbl>, time_hour <dttm>
\end{verbatim}

A better solution would be to use \texttt{mutate} to calculate a new
variable, \texttt{speed}:

\begin{Shaded}
\begin{Highlighting}[]
\NormalTok{flights }\SpecialCharTok{|\textgreater{}}
  \FunctionTok{mutate}\NormalTok{(}\AttributeTok{speed =}\NormalTok{ distance }\SpecialCharTok{/}\NormalTok{ air\_time) }\SpecialCharTok{|\textgreater{}}
  \FunctionTok{arrange}\NormalTok{(}\FunctionTok{desc}\NormalTok{(speed)) }\SpecialCharTok{|\textgreater{}}
  \FunctionTok{relocate}\NormalTok{(speed) }\SpecialCharTok{|\textgreater{}}
  \FunctionTok{head}\NormalTok{(}\DecValTok{3}\NormalTok{)}
\end{Highlighting}
\end{Shaded}

\begin{verbatim}
# A tibble: 3 x 20
  speed  year month   day dep_time sched_dep_time dep_delay arr_time
  <dbl> <int> <int> <int>    <int>          <int>     <dbl>    <int>
1  11.7  2013     5    25     1709           1700         9     1923
2  10.8  2013     7     2     1558           1513        45     1745
3  10.8  2013     5    13     2040           2025        15     2225
# i 12 more variables: sched_arr_time <int>, arr_delay <dbl>, carrier <chr>,
#   flight <int>, tailnum <chr>, origin <chr>, dest <chr>, air_time <dbl>,
#   distance <dbl>, hour <dbl>, minute <dbl>, time_hour <dttm>
\end{verbatim}

\subsection{Exercise 4}\label{exercise-4-6}

\begin{Shaded}
\begin{Highlighting}[]
\NormalTok{flights }\SpecialCharTok{|\textgreater{}}
  \FunctionTok{distinct}\NormalTok{(month, day) }\SpecialCharTok{|\textgreater{}}
  \FunctionTok{count}\NormalTok{()}
\end{Highlighting}
\end{Shaded}

\begin{verbatim}
# A tibble: 1 x 1
      n
  <int>
1   365
\end{verbatim}

We see that there are 365 distinct values and 2013 was not a leap year,
so there were flights everyday of 2013. It is not necessary to use
\texttt{count}, but the output is more concise.

\subsection{Exercise 5}\label{exercise-5-2}

We simply \texttt{arrange} with (longest flights) or without (shortest
flights) \texttt{desc}.

\begin{Shaded}
\begin{Highlighting}[]
\NormalTok{flights }\SpecialCharTok{|\textgreater{}}
  \FunctionTok{arrange}\NormalTok{(}\FunctionTok{desc}\NormalTok{(distance)) }\SpecialCharTok{|\textgreater{}}
  \FunctionTok{relocate}\NormalTok{(distance) }\SpecialCharTok{|\textgreater{}}
  \FunctionTok{head}\NormalTok{(}\DecValTok{3}\NormalTok{)}
\end{Highlighting}
\end{Shaded}

\begin{verbatim}
# A tibble: 3 x 19
  distance  year month   day dep_time sched_dep_time dep_delay arr_time
     <dbl> <int> <int> <int>    <int>          <int>     <dbl>    <int>
1     4983  2013     1     1      857            900        -3     1516
2     4983  2013     1     2      909            900         9     1525
3     4983  2013     1     3      914            900        14     1504
# i 11 more variables: sched_arr_time <int>, arr_delay <dbl>, carrier <chr>,
#   flight <int>, tailnum <chr>, origin <chr>, dest <chr>, air_time <dbl>,
#   hour <dbl>, minute <dbl>, time_hour <dttm>
\end{verbatim}

\begin{Shaded}
\begin{Highlighting}[]
\NormalTok{flights }\SpecialCharTok{|\textgreater{}}
  \FunctionTok{arrange}\NormalTok{(distance) }\SpecialCharTok{|\textgreater{}}
  \FunctionTok{relocate}\NormalTok{(distance) }\SpecialCharTok{|\textgreater{}}
  \FunctionTok{head}\NormalTok{(}\DecValTok{3}\NormalTok{)}
\end{Highlighting}
\end{Shaded}

\begin{verbatim}
# A tibble: 3 x 19
  distance  year month   day dep_time sched_dep_time dep_delay arr_time
     <dbl> <int> <int> <int>    <int>          <int>     <dbl>    <int>
1       17  2013     7    27       NA            106        NA       NA
2       80  2013     1     3     2127           2129        -2     2222
3       80  2013     1     4     1240           1200        40     1333
# i 11 more variables: sched_arr_time <int>, arr_delay <dbl>, carrier <chr>,
#   flight <int>, tailnum <chr>, origin <chr>, dest <chr>, air_time <dbl>,
#   hour <dbl>, minute <dbl>, time_hour <dttm>
\end{verbatim}

The shortest flight looks suspicious. Maybe there is something wrong
with this entry?

\subsection{Exercise 6}\label{exercise-6-2}

It does not matter in which order we \texttt{filter} or
\texttt{arrange}, but it is always preferable to \texttt{filter} as
early as possible. \texttt{filter} makes the data frame shorter and
means that the subsequent operations (function calls) will need to
process less data, and therefore be faster. However, if you are not
working with a lot of data, this will not matter much. Consider it a
good practice to adhere to though.

\section{4.3 Columns}\label{columns}

\subsection{Exercise 1}\label{exercise-1-11}

We would believe that \texttt{dep\_delay} = \texttt{dep\_time} -
\texttt{sched\_dep\_time}. We can test this by creating a new variable
called \texttt{my\_dep\_delay} and then compare it to
\texttt{dep\_delay}.

\begin{Shaded}
\begin{Highlighting}[]
\NormalTok{flights }\SpecialCharTok{|\textgreater{}}
  \FunctionTok{mutate}\NormalTok{(}
    \AttributeTok{my\_dep\_delay =}\NormalTok{ dep\_time }\SpecialCharTok{{-}}\NormalTok{ sched\_dep\_time,}
    \AttributeTok{dep\_delay =}\NormalTok{ dep\_delay,}
    \AttributeTok{.keep =} \StringTok{"used"}\NormalTok{) }\SpecialCharTok{|\textgreater{}}
  \FunctionTok{head}\NormalTok{(}\DecValTok{5}\NormalTok{)}
\end{Highlighting}
\end{Shaded}

\begin{verbatim}
# A tibble: 5 x 4
  dep_time sched_dep_time dep_delay my_dep_delay
     <int>          <int>     <dbl>        <int>
1      517            515         2            2
2      533            529         4            4
3      542            540         2            2
4      544            545        -1           -1
5      554            600        -6          -46
\end{verbatim}

We see that there is something wrong in our calculations. If we use
\texttt{?flights} we can see that \texttt{dep\_time} is stored in the
format HHMM and \texttt{dep\_delay} in minutes. There are several ways
we can solve this, see chapter 11.4 for another, similar solution. If
you are interested however, here is a solution:

\begin{Shaded}
\begin{Highlighting}[]
\NormalTok{flights }\SpecialCharTok{|\textgreater{}}
  \FunctionTok{mutate}\NormalTok{(}
    \AttributeTok{dep\_hours =} \FunctionTok{floor}\NormalTok{(dep\_time }\SpecialCharTok{/} \DecValTok{100}\NormalTok{),}
    \AttributeTok{dep\_minutes =}\NormalTok{ dep\_time }\SpecialCharTok{{-}}\NormalTok{ dep\_hours }\SpecialCharTok{*} \DecValTok{100}\NormalTok{,}
    \AttributeTok{dep\_time\_minutes =} \DecValTok{60} \SpecialCharTok{*}\NormalTok{ dep\_hours }\SpecialCharTok{+}\NormalTok{ dep\_minutes,}
    \AttributeTok{sched\_dep\_hours =} \FunctionTok{floor}\NormalTok{(sched\_dep\_time }\SpecialCharTok{/} \DecValTok{100}\NormalTok{),}
    \AttributeTok{sched\_dep\_minutes =}\NormalTok{ sched\_dep\_time }\SpecialCharTok{{-}}\NormalTok{ sched\_dep\_hours }\SpecialCharTok{*} \DecValTok{100}\NormalTok{,}
    \AttributeTok{sched\_dep\_time\_minutes =} \DecValTok{60} \SpecialCharTok{*}\NormalTok{ sched\_dep\_hours }\SpecialCharTok{+}\NormalTok{ sched\_dep\_minutes,}
    \AttributeTok{my\_dep\_delay =}\NormalTok{ dep\_time\_minutes }\SpecialCharTok{{-}}\NormalTok{ sched\_dep\_time\_minutes,}
    \AttributeTok{dep\_delay =}\NormalTok{ dep\_delay,}
    \AttributeTok{.keep =} \StringTok{"used"}\NormalTok{) }\SpecialCharTok{|\textgreater{}}
  \FunctionTok{relocate}\NormalTok{(dep\_delay, my\_dep\_delay) }\SpecialCharTok{|\textgreater{}}
  \FunctionTok{head}\NormalTok{(}\DecValTok{5}\NormalTok{)}
\end{Highlighting}
\end{Shaded}

\begin{verbatim}
# A tibble: 5 x 10
  dep_delay my_dep_delay dep_time sched_dep_time dep_hours dep_minutes
      <dbl>        <dbl>    <int>          <int>     <dbl>       <dbl>
1         2            2      517            515         5          17
2         4            4      533            529         5          33
3         2            2      542            540         5          42
4        -1           -1      544            545         5          44
5        -6           -6      554            600         5          54
# i 4 more variables: dep_time_minutes <dbl>, sched_dep_hours <dbl>,
#   sched_dep_minutes <dbl>, sched_dep_time_minutes <dbl>
\end{verbatim}

where the \texttt{floor} function rounds down a number.

\subsection{Exercise 2}\label{exercise-2-8}

Here are a two suggestions where we use \texttt{names} which extracts
the column names of a data frame.

\begin{Shaded}
\begin{Highlighting}[]
\CommentTok{\# Using select}
\NormalTok{flights }\SpecialCharTok{|\textgreater{}}
  \FunctionTok{select}\NormalTok{(dep\_time, dep\_delay, arr\_time, arr\_delay) }\SpecialCharTok{|\textgreater{}}
  \FunctionTok{names}\NormalTok{()}
\end{Highlighting}
\end{Shaded}

\begin{verbatim}
[1] "dep_time"  "dep_delay" "arr_time"  "arr_delay"
\end{verbatim}

\begin{Shaded}
\begin{Highlighting}[]
\CommentTok{\# Using starts\_with}
\NormalTok{flights }\SpecialCharTok{|\textgreater{}}
  \FunctionTok{select}\NormalTok{(}\FunctionTok{starts\_with}\NormalTok{(}\StringTok{"dep"}\NormalTok{) }\SpecialCharTok{|} \FunctionTok{starts\_with}\NormalTok{(}\StringTok{"arr"}\NormalTok{)) }\SpecialCharTok{|\textgreater{}}
  \FunctionTok{names}\NormalTok{()}
\end{Highlighting}
\end{Shaded}

\begin{verbatim}
[1] "dep_time"  "dep_delay" "arr_time"  "arr_delay"
\end{verbatim}

\begin{Shaded}
\begin{Highlighting}[]
\CommentTok{\# A unintuitive and not recommended solution, using mutate and .keep = used}
\NormalTok{flights }\SpecialCharTok{|\textgreater{}}
  \FunctionTok{mutate}\NormalTok{(}
    \AttributeTok{dep\_time =}\NormalTok{ dep\_time, }
    \AttributeTok{dep\_delay =}\NormalTok{ dep\_delay, }
    \AttributeTok{arr\_time =}\NormalTok{ arr\_time, }
    \AttributeTok{arr\_delay =}\NormalTok{ arr\_delay,}
    \AttributeTok{.keep =} \StringTok{"used"}\NormalTok{) }\SpecialCharTok{|\textgreater{}}
  \FunctionTok{names}\NormalTok{()}
\end{Highlighting}
\end{Shaded}

\begin{verbatim}
[1] "dep_time"  "dep_delay" "arr_time"  "arr_delay"
\end{verbatim}

\subsection{Exercise 3}\label{exercise-3-7}

Nothing seems to happen\ldots{}

\begin{Shaded}
\begin{Highlighting}[]
\NormalTok{flights }\SpecialCharTok{|\textgreater{}}
  \FunctionTok{select}\NormalTok{(dep\_time, dep\_time) }\SpecialCharTok{|\textgreater{}}
  \FunctionTok{head}\NormalTok{(}\DecValTok{3}\NormalTok{)}
\end{Highlighting}
\end{Shaded}

\begin{verbatim}
# A tibble: 3 x 1
  dep_time
     <int>
1      517
2      533
3      542
\end{verbatim}

\subsection{Exercise 4}\label{exercise-4-7}

\texttt{any\_of} can be used to select any matching variables in the
vector.

\begin{Shaded}
\begin{Highlighting}[]
\NormalTok{variables }\OtherTok{\textless{}{-}} \FunctionTok{c}\NormalTok{(}\StringTok{"year"}\NormalTok{, }\StringTok{"month"}\NormalTok{, }\StringTok{"day"}\NormalTok{, }\StringTok{"dep\_delay"}\NormalTok{, }\StringTok{"arr\_delay"}\NormalTok{)}
\NormalTok{flights }\SpecialCharTok{|\textgreater{}}
  \FunctionTok{select}\NormalTok{(}\FunctionTok{any\_of}\NormalTok{(variables)) }\SpecialCharTok{|\textgreater{}}
  \FunctionTok{head}\NormalTok{(}\DecValTok{3}\NormalTok{)}
\end{Highlighting}
\end{Shaded}

\begin{verbatim}
# A tibble: 3 x 5
   year month   day dep_delay arr_delay
  <int> <int> <int>     <dbl>     <dbl>
1  2013     1     1         2        11
2  2013     1     1         4        20
3  2013     1     1         2        33
\end{verbatim}

produces the same result as

\begin{Shaded}
\begin{Highlighting}[]
\NormalTok{variables }\OtherTok{\textless{}{-}} \FunctionTok{c}\NormalTok{(}\StringTok{"bill\_depth\_mm"}\NormalTok{, }\StringTok{"year"}\NormalTok{, }\StringTok{"month"}\NormalTok{, }\StringTok{"day"}\NormalTok{, }\StringTok{"dep\_delay"}\NormalTok{, }\StringTok{"arr\_delay"}\NormalTok{)}
\NormalTok{flights }\SpecialCharTok{|\textgreater{}}
  \FunctionTok{select}\NormalTok{(}\FunctionTok{any\_of}\NormalTok{(variables)) }\SpecialCharTok{|\textgreater{}}
  \FunctionTok{head}\NormalTok{(}\DecValTok{3}\NormalTok{)}
\end{Highlighting}
\end{Shaded}

\begin{verbatim}
# A tibble: 3 x 5
   year month   day dep_delay arr_delay
  <int> <int> <int>     <dbl>     <dbl>
1  2013     1     1         2        11
2  2013     1     1         4        20
3  2013     1     1         2        33
\end{verbatim}

while \texttt{all\_of} would produce an error

\begin{Shaded}
\begin{Highlighting}[]
\NormalTok{variables }\OtherTok{\textless{}{-}} \FunctionTok{c}\NormalTok{(}\StringTok{"bill\_depth\_mm"}\NormalTok{, }\StringTok{"year"}\NormalTok{, }\StringTok{"month"}\NormalTok{, }\StringTok{"day"}\NormalTok{, }\StringTok{"dep\_delay"}\NormalTok{, }\StringTok{"arr\_delay"}\NormalTok{)}
\NormalTok{flights }\SpecialCharTok{|\textgreater{}}
  \FunctionTok{select}\NormalTok{(}\FunctionTok{all\_of}\NormalTok{(variables)) }\SpecialCharTok{|\textgreater{}}
  \FunctionTok{head}\NormalTok{(}\DecValTok{3}\NormalTok{)}
\end{Highlighting}
\end{Shaded}

\begin{verbatim}
Error in `select()`:
i In argument: `all_of(variables)`.
Caused by error in `all_of()`:
! Can't subset elements that don't exist.
x Element `bill_depth_mm` doesn't exist.
\end{verbatim}

These functions may be useful if you select variables across multiple
data frames.

\subsection{Exercise 5}\label{exercise-5-3}

\begin{Shaded}
\begin{Highlighting}[]
\NormalTok{flights }\SpecialCharTok{|\textgreater{}} 
  \FunctionTok{select}\NormalTok{(}\FunctionTok{contains}\NormalTok{(}\StringTok{"TIME"}\NormalTok{)) }\SpecialCharTok{|\textgreater{}}
  \FunctionTok{head}\NormalTok{(}\DecValTok{3}\NormalTok{)}
\end{Highlighting}
\end{Shaded}

\begin{verbatim}
# A tibble: 3 x 6
  dep_time sched_dep_time arr_time sched_arr_time air_time time_hour          
     <int>          <int>    <int>          <int>    <dbl> <dttm>             
1      517            515      830            819      227 2013-01-01 05:00:00
2      533            529      850            830      227 2013-01-01 05:00:00
3      542            540      923            850      160 2013-01-01 05:00:00
\end{verbatim}

The default behavior is to ignore case, \texttt{ignore.case\ =\ TRUE}.
Set it to \texttt{FALSE} to match case

\begin{Shaded}
\begin{Highlighting}[]
\NormalTok{flights }\SpecialCharTok{|\textgreater{}} 
  \FunctionTok{select}\NormalTok{(}\FunctionTok{contains}\NormalTok{(}\StringTok{"TIME"}\NormalTok{, }\AttributeTok{ignore.case =} \ConstantTok{FALSE}\NormalTok{)) }\SpecialCharTok{|\textgreater{}}
  \FunctionTok{head}\NormalTok{(}\DecValTok{3}\NormalTok{)}
\end{Highlighting}
\end{Shaded}

\begin{verbatim}
# A tibble: 3 x 0
\end{verbatim}

\subsection{Exercise 6}\label{exercise-6-3}

\begin{Shaded}
\begin{Highlighting}[]
\NormalTok{flights }\SpecialCharTok{|\textgreater{}} 
  \FunctionTok{rename}\NormalTok{(}\AttributeTok{air\_time\_min =}\NormalTok{ air\_time) }\SpecialCharTok{|\textgreater{}}
  \FunctionTok{relocate}\NormalTok{(air\_time\_min) }\SpecialCharTok{|\textgreater{}}
  \FunctionTok{head}\NormalTok{(}\DecValTok{3}\NormalTok{)}
\end{Highlighting}
\end{Shaded}

\begin{verbatim}
# A tibble: 3 x 19
  air_time_min  year month   day dep_time sched_dep_time dep_delay arr_time
         <dbl> <int> <int> <int>    <int>          <int>     <dbl>    <int>
1          227  2013     1     1      517            515         2      830
2          227  2013     1     1      533            529         4      850
3          160  2013     1     1      542            540         2      923
# i 11 more variables: sched_arr_time <int>, arr_delay <dbl>, carrier <chr>,
#   flight <int>, tailnum <chr>, origin <chr>, dest <chr>, distance <dbl>,
#   hour <dbl>, minute <dbl>, time_hour <dttm>
\end{verbatim}

\subsection{Exercise 7}\label{exercise-7-2}

By \texttt{select}-ing \texttt{tailnum} we drop all other columns. Thus
there is no column called \texttt{arr\_delay} anymore.

\begin{Shaded}
\begin{Highlighting}[]
\NormalTok{flights }\SpecialCharTok{|\textgreater{}} 
  \FunctionTok{select}\NormalTok{(tailnum) }\SpecialCharTok{|\textgreater{}}
  \FunctionTok{head}\NormalTok{(}\DecValTok{3}\NormalTok{)}
\end{Highlighting}
\end{Shaded}

\begin{verbatim}
# A tibble: 3 x 1
  tailnum
  <chr>  
1 N14228 
2 N24211 
3 N619AA 
\end{verbatim}

\section{4.5 Groups}\label{groups}

\subsection{Exercise 1}\label{exercise-1-12}

First we will look at the carriers with the worst delays. Either by
using \texttt{group\_by}

\begin{Shaded}
\begin{Highlighting}[]
\NormalTok{flights }\SpecialCharTok{|\textgreater{}}
  \FunctionTok{group\_by}\NormalTok{(carrier) }\SpecialCharTok{|\textgreater{}}
  \FunctionTok{summarize}\NormalTok{(}\AttributeTok{avg\_dep\_delay =} \FunctionTok{mean}\NormalTok{(dep\_delay, }\AttributeTok{na.rm =} \ConstantTok{TRUE}\NormalTok{))}
\end{Highlighting}
\end{Shaded}

or by using \texttt{.by}

\begin{Shaded}
\begin{Highlighting}[]
\NormalTok{flights }\SpecialCharTok{|\textgreater{}}
  \FunctionTok{summarize}\NormalTok{(}
    \AttributeTok{dep\_delay\_avg =} \FunctionTok{mean}\NormalTok{(dep\_delay, }\AttributeTok{na.rm =} \ConstantTok{TRUE}\NormalTok{),}
    \AttributeTok{.by =}\NormalTok{ carrier) }\SpecialCharTok{|\textgreater{}}
  \FunctionTok{head}\NormalTok{(}\DecValTok{3}\NormalTok{)}
\end{Highlighting}
\end{Shaded}

\begin{verbatim}
# A tibble: 3 x 2
  carrier dep_delay_avg
  <chr>           <dbl>
1 UA              12.1 
2 AA               8.59
3 B6              13.0 
\end{verbatim}

Note: we recommend using \texttt{.by}, at least when grouping over more
variables to keep grouping tied to specific summarize. Performing
multiple \texttt{summarize} with one \texttt{group\_by} makes it harder
to follow the code, and there is a risk that the end result is a grouped
tibble, which is hard to keep track of the next time it is used.

Now we'll try the challenge. First we look at the hint provided, where
we use \texttt{group\_by} in conjunction with the \texttt{.group}
argument set to \texttt{drop}, and then try to visualize how many
flights each carrier has to each destination.

\begin{Shaded}
\begin{Highlighting}[]
\NormalTok{flights }\SpecialCharTok{|\textgreater{}}
  \FunctionTok{group\_by}\NormalTok{(carrier, dest) }\SpecialCharTok{|\textgreater{}}
  \FunctionTok{summarize}\NormalTok{(}
    \AttributeTok{n =} \FunctionTok{n}\NormalTok{(),}
    \AttributeTok{.groups =} \StringTok{"drop"}\NormalTok{) }\SpecialCharTok{|\textgreater{}}
  \FunctionTok{ggplot}\NormalTok{(}\FunctionTok{aes}\NormalTok{(}\AttributeTok{x =}\NormalTok{ dest, }\AttributeTok{y =}\NormalTok{ carrier, }\AttributeTok{col =}\NormalTok{ n)) }\SpecialCharTok{+} 
  \FunctionTok{geom\_tile}\NormalTok{()}
\end{Highlighting}
\end{Shaded}

\includegraphics{solutions_r4ds_data_transformation_files/figure-pdf/unnamed-chunk-23-1.pdf}

Since several carriers only fly to a few destinations, it may be hard to
disentangle if the fault is with the carrier or the airport. Do you have
a better solution? I'm sure there is one! Also, wouldn't it be more
reasonable to blame the \texttt{origin}? Hmm\ldots{}

\subsection{Exercise 2}\label{exercise-2-9}

\begin{Shaded}
\begin{Highlighting}[]
\NormalTok{flights }\SpecialCharTok{|\textgreater{}}
  \FunctionTok{slice\_max}\NormalTok{(}
\NormalTok{    dep\_delay,}
    \AttributeTok{by =}\NormalTok{ origin) }\SpecialCharTok{|\textgreater{}}
  \FunctionTok{relocate}\NormalTok{(origin, dep\_delay) }\SpecialCharTok{|\textgreater{}}
  \FunctionTok{head}\NormalTok{(}\DecValTok{3}\NormalTok{)}
\end{Highlighting}
\end{Shaded}

\begin{verbatim}
# A tibble: 3 x 19
  origin dep_delay  year month   day dep_time sched_dep_time arr_time
  <chr>      <dbl> <int> <int> <int>    <int>          <int>    <int>
1 EWR         1126  2013     1    10     1121           1635     1239
2 LGA          911  2013     3    17     2321            810      135
3 JFK         1301  2013     1     9      641            900     1242
# i 11 more variables: sched_arr_time <int>, arr_delay <dbl>, carrier <chr>,
#   flight <int>, tailnum <chr>, dest <chr>, air_time <dbl>, distance <dbl>,
#   hour <dbl>, minute <dbl>, time_hour <dttm>
\end{verbatim}

Alternatively

\begin{Shaded}
\begin{Highlighting}[]
\NormalTok{flights }\SpecialCharTok{|\textgreater{}}
\FunctionTok{group\_by}\NormalTok{(origin) }\SpecialCharTok{|\textgreater{}}
  \FunctionTok{slice\_max}\NormalTok{(dep\_delay) }\SpecialCharTok{|\textgreater{}}
  \FunctionTok{relocate}\NormalTok{(origin, dep\_delay) }\SpecialCharTok{|\textgreater{}}
  \FunctionTok{head}\NormalTok{(}\DecValTok{3}\NormalTok{)}
\end{Highlighting}
\end{Shaded}

\subsection{Exercise 3}\label{exercise-3-8}

First we average \texttt{dep\_delay} grouped by \texttt{hour}. Remember
to use \texttt{na.rm\ =TRUE} to remove any missing values

\begin{Shaded}
\begin{Highlighting}[]
\NormalTok{avg\_delays }\OtherTok{\textless{}{-}}\NormalTok{ flights }\SpecialCharTok{|\textgreater{}}
  \FunctionTok{summarize}\NormalTok{(}
    \AttributeTok{avg\_dep\_delay =} \FunctionTok{mean}\NormalTok{(dep\_delay, }\AttributeTok{na.rm =} \ConstantTok{TRUE}\NormalTok{), }
    \AttributeTok{.by =}\NormalTok{ hour) }\SpecialCharTok{|\textgreater{}}
  \FunctionTok{arrange}\NormalTok{(}\FunctionTok{desc}\NormalTok{(avg\_dep\_delay))}
\NormalTok{avg\_delays }\SpecialCharTok{|\textgreater{}} \FunctionTok{head}\NormalTok{(}\DecValTok{5}\NormalTok{)}
\end{Highlighting}
\end{Shaded}

\begin{verbatim}
# A tibble: 5 x 2
   hour avg_dep_delay
  <dbl>         <dbl>
1    19          24.8
2    20          24.3
3    21          24.2
4    18          21.1
5    17          21.1
\end{verbatim}

Looking at the full tibble, we can see that no flights depart at hours
2, 3, and 4, and that there seems to be something wrong with hour 1
since we have a \emph{Not-a-Number} or \texttt{NaN}. We will ignore this
now, but mention that \texttt{NaN} is used when some computation doesn't
have a meaningful result, such as division by zero, e.g.~1 / 0.
\texttt{NaN} are sometimes used for missing values, but \texttt{NA} is
preferable in that case.

Since both \texttt{hour} and \texttt{dep\_delay} are numerical
variables, we may use a scatter plot, \texttt{geom\_point}, but in this
case it may be more appealing to visualize the delays using a bar chart,
\texttt{geom\_bar}, or histogram, \texttt{geom\_histogram} with
\texttt{stat\ =\ "identity"} and \texttt{avg\_dep\_delay} selected as
the y aesthetic. We see that this works with \texttt{geom\_freqpoly} too
(and \texttt{geom\_density}).

\begin{Shaded}
\begin{Highlighting}[]
\NormalTok{avg\_delays }\SpecialCharTok{|\textgreater{}}
  \FunctionTok{arrange}\NormalTok{(hour) }\SpecialCharTok{|\textgreater{}}
  \FunctionTok{ggplot}\NormalTok{(}\FunctionTok{aes}\NormalTok{(}\AttributeTok{x =}\NormalTok{ hour, }\AttributeTok{y =}\NormalTok{ avg\_dep\_delay)) }\SpecialCharTok{+} 
  \FunctionTok{geom\_histogram}\NormalTok{(}\AttributeTok{stat =} \StringTok{"identity"}\NormalTok{) }\SpecialCharTok{+} 
  \FunctionTok{geom\_freqpoly}\NormalTok{(}\AttributeTok{stat =} \StringTok{"identity"}\NormalTok{)}
\end{Highlighting}
\end{Shaded}

\begin{verbatim}
Warning in geom_histogram(stat = "identity"): Ignoring unknown parameters:
`binwidth`, `bins`, and `pad`
\end{verbatim}

\begin{verbatim}
Warning: Removed 1 row containing missing values or values outside the scale range
(`geom_bar()`).
\end{verbatim}

\begin{verbatim}
Warning: Removed 1 row containing missing values or values outside the scale range
(`geom_path()`).
\end{verbatim}

\includegraphics{solutions_r4ds_data_transformation_files/figure-pdf/unnamed-chunk-26-1.pdf}

We see that the most delays are around 19-21.

\subsection{Exercise 4}\label{exercise-4-8}

When we \texttt{slice} using e negative value we \emph{drop} rows
instead of selecting them. \texttt{slice\_min} drops from the bottom and
\texttt{slice\_max} drops from the top, in an ordered tibble.

\begin{Shaded}
\begin{Highlighting}[]
\NormalTok{avg\_delays }\SpecialCharTok{|\textgreater{}}
  \FunctionTok{slice\_min}\NormalTok{(}
\NormalTok{    avg\_dep\_delay, }
    \AttributeTok{n =} \SpecialCharTok{{-}}\DecValTok{17}\NormalTok{)}
\end{Highlighting}
\end{Shaded}

\begin{verbatim}
# A tibble: 3 x 2
   hour avg_dep_delay
  <dbl>         <dbl>
1     5         0.688
2     6         1.64 
3     7         1.91 
\end{verbatim}

\subsection{Exercise 5}\label{exercise-5-4}

\texttt{count} counts the number of occurences of a value of a variable.
If we set the \texttt{sort} parameter to \texttt{TRUE} the resulting
tibble will be sorted according to the counts. Compare with
\texttt{distinct}. Below we count the number of flights departing at
each \texttt{hour}.

\begin{Shaded}
\begin{Highlighting}[]
\NormalTok{flights }\SpecialCharTok{|\textgreater{}}
  \FunctionTok{count}\NormalTok{(}
\NormalTok{    hour, }
    \AttributeTok{sort =} \ConstantTok{TRUE}\NormalTok{) }\SpecialCharTok{|\textgreater{}}
  \FunctionTok{head}\NormalTok{(}\DecValTok{3}\NormalTok{)}
\end{Highlighting}
\end{Shaded}

\begin{verbatim}
# A tibble: 3 x 2
   hour     n
  <dbl> <int>
1     8 27242
2     6 25951
3    17 24426
\end{verbatim}

We see that most flights depart in the morning. An interesting question
is why the most delays are in the evening when most flights depart in
the morning. Maybe we can look at when there are most arrivals?

\begin{Shaded}
\begin{Highlighting}[]
\NormalTok{flights }\SpecialCharTok{|\textgreater{}}
  \FunctionTok{mutate}\NormalTok{(}\AttributeTok{arr\_hour =} \FunctionTok{floor}\NormalTok{(arr\_time }\SpecialCharTok{/} \DecValTok{100}\NormalTok{)) }\SpecialCharTok{|\textgreater{}}
  \FunctionTok{count}\NormalTok{(}
\NormalTok{    arr\_hour, }
    \AttributeTok{sort =} \ConstantTok{TRUE}\NormalTok{) }\SpecialCharTok{|\textgreater{}}
  \FunctionTok{head}\NormalTok{(}\DecValTok{5}\NormalTok{)}
\end{Highlighting}
\end{Shaded}

\begin{verbatim}
# A tibble: 5 x 2
  arr_hour     n
     <dbl> <int>
1       10 25034
2       18 22367
3       20 22123
4       19 21017
5       21 20866
\end{verbatim}

Here we see what we expect, the hours with the most delays coincide with
top 5 hours of most arrivals. It still doesn't explain why we don't have
more delays at 10. To investigate this we would have to explore our data
further, and possibly conclude that the answer may lie beyond the
recorded data.

\subsection{Exercise 6}\label{exercise-6-4}

\begin{Shaded}
\begin{Highlighting}[]
\NormalTok{df }\OtherTok{\textless{}{-}} \FunctionTok{tibble}\NormalTok{(}
  \AttributeTok{x =} \DecValTok{1}\SpecialCharTok{:}\DecValTok{5}\NormalTok{,}
  \AttributeTok{y =} \FunctionTok{c}\NormalTok{(}\StringTok{"a"}\NormalTok{, }\StringTok{"b"}\NormalTok{, }\StringTok{"a"}\NormalTok{, }\StringTok{"a"}\NormalTok{, }\StringTok{"b"}\NormalTok{),}
  \AttributeTok{z =} \FunctionTok{c}\NormalTok{(}\StringTok{"K"}\NormalTok{, }\StringTok{"K"}\NormalTok{, }\StringTok{"L"}\NormalTok{, }\StringTok{"L"}\NormalTok{, }\StringTok{"K"}\NormalTok{)}
\NormalTok{)}
\end{Highlighting}
\end{Shaded}

\subsubsection{a.}\label{a.}

\begin{Shaded}
\begin{Highlighting}[]
\NormalTok{df }\SpecialCharTok{|\textgreater{}}
  \FunctionTok{group\_by}\NormalTok{(y)}
\end{Highlighting}
\end{Shaded}

\begin{verbatim}
# A tibble: 5 x 3
# Groups:   y [2]
      x y     z    
  <int> <chr> <chr>
1     1 a     K    
2     2 b     K    
3     3 a     L    
4     4 a     L    
5     5 b     K    
\end{verbatim}

Observe that the output is now grouped, which may be confusing if other
operations are applied to \texttt{df} in a new code segment. Compare
with 6 e. below.

\subsubsection{b.}\label{b.}

\begin{Shaded}
\begin{Highlighting}[]
\NormalTok{df }\SpecialCharTok{|\textgreater{}}
  \FunctionTok{arrange}\NormalTok{(y)}
\end{Highlighting}
\end{Shaded}

\begin{verbatim}
# A tibble: 5 x 3
      x y     z    
  <int> <chr> <chr>
1     1 a     K    
2     3 a     L    
3     4 a     L    
4     2 b     K    
5     5 b     K    
\end{verbatim}

The tibble is sorted according to \texttt{y} in alphabetical order.

\subsubsection{c.}\label{c.}

\begin{Shaded}
\begin{Highlighting}[]
\NormalTok{df }\SpecialCharTok{|\textgreater{}}
  \FunctionTok{group\_by}\NormalTok{(y) }\SpecialCharTok{|\textgreater{}}
  \FunctionTok{summarize}\NormalTok{(}\AttributeTok{mean\_x =} \FunctionTok{mean}\NormalTok{(x))}
\end{Highlighting}
\end{Shaded}

\begin{verbatim}
# A tibble: 2 x 2
  y     mean_x
  <chr>  <dbl>
1 a       2.67
2 b       3.5 
\end{verbatim}

The mean of \texttt{x} is calculated for all values of \texttt{y}.

\subsubsection{d.}\label{d.}

\begin{Shaded}
\begin{Highlighting}[]
\NormalTok{df }\SpecialCharTok{|\textgreater{}}
  \FunctionTok{group\_by}\NormalTok{(y, z) }\SpecialCharTok{|\textgreater{}}
  \FunctionTok{summarize}\NormalTok{(}\AttributeTok{mean\_x =} \FunctionTok{mean}\NormalTok{(x))}
\end{Highlighting}
\end{Shaded}

\begin{verbatim}
`summarise()` has grouped output by 'y'. You can override using the `.groups`
argument.
\end{verbatim}

\begin{verbatim}
# A tibble: 3 x 3
# Groups:   y [2]
  y     z     mean_x
  <chr> <chr>  <dbl>
1 a     K        1  
2 a     L        3.5
3 b     K        3.5
\end{verbatim}

The mean of \texttt{x} is calculated for all combinations of \texttt{y}
and \texttt{z}. The output is \emph{still} grouped by \texttt{y}.

\subsubsection{e.}\label{e.}

\begin{Shaded}
\begin{Highlighting}[]
\NormalTok{df }\SpecialCharTok{|\textgreater{}}
  \FunctionTok{group\_by}\NormalTok{(y, z) }\SpecialCharTok{|\textgreater{}}
  \FunctionTok{summarize}\NormalTok{(}
    \AttributeTok{mean\_x =} \FunctionTok{mean}\NormalTok{(x),}
    \AttributeTok{.groups =} \StringTok{"drop"}\NormalTok{)}
\end{Highlighting}
\end{Shaded}

\begin{verbatim}
# A tibble: 3 x 3
  y     z     mean_x
  <chr> <chr>  <dbl>
1 a     K        1  
2 a     L        3.5
3 b     K        3.5
\end{verbatim}

The mean of \texttt{x} is calculated for all combinations of \texttt{y}
and \texttt{z}. The output is \emph{not} grouped by \texttt{y}.

\subsubsection{f.}\label{f.}

\begin{Shaded}
\begin{Highlighting}[]
\NormalTok{df }\SpecialCharTok{|\textgreater{}}
  \FunctionTok{group\_by}\NormalTok{(y, z) }\SpecialCharTok{|\textgreater{}}
  \FunctionTok{summarize}\NormalTok{(}\AttributeTok{mean\_x =} \FunctionTok{mean}\NormalTok{(x))}
\end{Highlighting}
\end{Shaded}

\begin{verbatim}
`summarise()` has grouped output by 'y'. You can override using the `.groups`
argument.
\end{verbatim}

\begin{verbatim}
# A tibble: 3 x 3
# Groups:   y [2]
  y     z     mean_x
  <chr> <chr>  <dbl>
1 a     K        1  
2 a     L        3.5
3 b     K        3.5
\end{verbatim}

\begin{Shaded}
\begin{Highlighting}[]
\NormalTok{df }\SpecialCharTok{|\textgreater{}}
  \FunctionTok{group\_by}\NormalTok{(y, z) }\SpecialCharTok{|\textgreater{}}
  \FunctionTok{mutate}\NormalTok{(}\AttributeTok{mean\_x =} \FunctionTok{mean}\NormalTok{(x))}
\end{Highlighting}
\end{Shaded}

\begin{verbatim}
# A tibble: 5 x 4
# Groups:   y, z [3]
      x y     z     mean_x
  <int> <chr> <chr>  <dbl>
1     1 a     K        1  
2     2 b     K        3.5
3     3 a     L        3.5
4     4 a     L        3.5
5     5 b     K        3.5
\end{verbatim}

In the first piece of code we \texttt{summarize} over \texttt{z} and the
resulting tibble only has the unique values of \texttt{z}. In the second
piece of code we create a new variable with the mean value of \texttt{x}
for each combination of \texttt{y} and \texttt{z}.

\chapter{R4DS: Data tidying}\label{r4ds-data-tidying}

To run these solutions you need to load the necessary libraries.

\begin{Shaded}
\begin{Highlighting}[]
\FunctionTok{library}\NormalTok{(tidyverse)}
\end{Highlighting}
\end{Shaded}

\begin{verbatim}
-- Attaching core tidyverse packages ------------------------ tidyverse 2.0.0 --
v dplyr     1.1.4     v readr     2.1.5
v forcats   1.0.0     v stringr   1.5.1
v ggplot2   3.5.1     v tibble    3.2.1
v lubridate 1.9.3     v tidyr     1.3.1
v purrr     1.0.2     
-- Conflicts ------------------------------------------ tidyverse_conflicts() --
x dplyr::filter() masks stats::filter()
x dplyr::lag()    masks stats::lag()
i Use the conflicted package (<http://conflicted.r-lib.org/>) to force all conflicts to become errors
\end{verbatim}

\section{6.2 Tidy data}\label{tidy-data}

\subsection{Exercise 1}\label{exercise-1-13}

\begin{itemize}
\tightlist
\item
  \textbf{table1}

  \begin{itemize}
  \tightlist
  \item
    Observation: Tuberculosis and population for one country and one
    year.
  \item
    Columns:

    \begin{itemize}
    \tightlist
    \item
      country, year, tuberculosis cases, population
    \end{itemize}
  \end{itemize}
\item
  \textbf{table2}

  \begin{itemize}
  \tightlist
  \item
    Observation: Tuberculosis or population for one country and one
    year.
  \item
    Columns:

    \begin{itemize}
    \tightlist
    \item
      country, year, type of observation, count of the type of
      observation
    \end{itemize}
  \end{itemize}
\item
  \textbf{table3}

  \begin{itemize}
  \tightlist
  \item
    Observation: Tuberculosis per capita as a string for a country and
    one year.
  \item
    Columns:

    \begin{itemize}
    \tightlist
    \item
      country, year, tuberculosis per capita as a string
    \end{itemize}
  \end{itemize}
\end{itemize}

Observe that the rate in table three is stored as a string, which is not
the value of tuberculosis per capita.

\subsection{Exercise 2}\label{exercise-2-10}

Not recommended, but here is a solution for table2, if you are
interested. You do not need to be able to reproduce this.

\begin{Shaded}
\begin{Highlighting}[]
\NormalTok{table2 }\OtherTok{\textless{}{-}} \FunctionTok{tibble}\NormalTok{(}
  \AttributeTok{country =} \FunctionTok{c}\NormalTok{(}\StringTok{"Afghanistan"}\NormalTok{, }\StringTok{"Afghanistan"}\NormalTok{, }\StringTok{"Afghanistan"}\NormalTok{, }\StringTok{"Afghanistan"}\NormalTok{, }\StringTok{"Brazil"}\NormalTok{, }\StringTok{"Brazil"}\NormalTok{),}
  \AttributeTok{year =} \FunctionTok{c}\NormalTok{(}\DecValTok{1999}\NormalTok{, }\DecValTok{1999}\NormalTok{, }\DecValTok{2000}\NormalTok{, }\DecValTok{2000}\NormalTok{, }\DecValTok{1999}\NormalTok{, }\DecValTok{1999}\NormalTok{),}
  \AttributeTok{type =} \FunctionTok{c}\NormalTok{(}\StringTok{"cases"}\NormalTok{, }\StringTok{"population"}\NormalTok{, }\StringTok{"cases"}\NormalTok{, }\StringTok{"population"}\NormalTok{, }\StringTok{"cases"}\NormalTok{, }\StringTok{"population"}\NormalTok{),}
  \AttributeTok{count =} \FunctionTok{c}\NormalTok{(}\DecValTok{745}\NormalTok{, }\DecValTok{19987071}\NormalTok{, }\DecValTok{2666}\NormalTok{, }\DecValTok{20595360}\NormalTok{, }\DecValTok{37737}\NormalTok{, }\DecValTok{172006362}\NormalTok{)}
\NormalTok{)}

\NormalTok{cases }\OtherTok{\textless{}{-}}\NormalTok{ table2 }\SpecialCharTok{|\textgreater{}} 
  \FunctionTok{filter}\NormalTok{(type }\SpecialCharTok{==} \StringTok{"cases"}\NormalTok{) }\SpecialCharTok{|\textgreater{}}
  \FunctionTok{pull}\NormalTok{(count)}
\NormalTok{population }\OtherTok{\textless{}{-}}\NormalTok{ table2 }\SpecialCharTok{|\textgreater{}} 
  \FunctionTok{filter}\NormalTok{(type }\SpecialCharTok{==} \StringTok{"population"}\NormalTok{) }\SpecialCharTok{|\textgreater{}}
  \FunctionTok{pull}\NormalTok{(count)}
\NormalTok{rate }\OtherTok{\textless{}{-}}\NormalTok{ (cases }\SpecialCharTok{/}\NormalTok{ population) }\SpecialCharTok{*} \DecValTok{10000}
\NormalTok{table2}\SpecialCharTok{$}\NormalTok{rate }\OtherTok{\textless{}{-}} \FunctionTok{rep}\NormalTok{(rate, }\AttributeTok{each =} \DecValTok{2}\NormalTok{)}
\NormalTok{table2}
\end{Highlighting}
\end{Shaded}

\begin{verbatim}
# A tibble: 6 x 5
  country      year type           count  rate
  <chr>       <dbl> <chr>          <dbl> <dbl>
1 Afghanistan  1999 cases            745 0.373
2 Afghanistan  1999 population  19987071 0.373
3 Afghanistan  2000 cases           2666 1.29 
4 Afghanistan  2000 population  20595360 1.29 
5 Brazil       1999 cases          37737 2.19 
6 Brazil       1999 population 172006362 2.19 
\end{verbatim}

\chapter{R4DS: Exploratory data
analysis}\label{r4ds-exploratory-data-analysis}

\section{Notes for assignment}\label{notes-for-assignment}

\begin{itemize}
\tightlist
\item
  Asking questions to explore your data

  \begin{itemize}
  \tightlist
  \item
    What type of variation occurs within my variables?
  \item
    What type of covariation occurs between my variables?
  \end{itemize}
\item
  Visualize distributions to study variation within a variable.
\item
  Ask about

  \begin{itemize}
  \tightlist
  \item
    Typical values: Most common values, rare values, unusual patterns.
  \item
    Clusters: How are similar observations similar, how can clusters be
    explained, can clusters be misleading?
  \item
    Outliers
  \end{itemize}
\end{itemize}

Suggestions for exercises:

\begin{itemize}
\tightlist
\item
  Visualize distributions, find and extract interesting
  variables/features that represents an insight in data. 11.5.1 Exercise
  2
\end{itemize}

To run these solutions you need to load the necessary libraries.

\begin{Shaded}
\begin{Highlighting}[]
\FunctionTok{library}\NormalTok{(tidyverse)}
\end{Highlighting}
\end{Shaded}

\begin{verbatim}
-- Attaching core tidyverse packages ------------------------ tidyverse 2.0.0 --
v dplyr     1.1.4     v readr     2.1.5
v forcats   1.0.0     v stringr   1.5.1
v ggplot2   3.5.1     v tibble    3.2.1
v lubridate 1.9.3     v tidyr     1.3.1
v purrr     1.0.2     
-- Conflicts ------------------------------------------ tidyverse_conflicts() --
x dplyr::filter() masks stats::filter()
x dplyr::lag()    masks stats::lag()
i Use the conflicted package (<http://conflicted.r-lib.org/>) to force all conflicts to become errors
\end{verbatim}

\section{11.3 Variation}\label{variation}

We take a \texttt{glimpse} at the data before we begin

\begin{Shaded}
\begin{Highlighting}[]
\FunctionTok{glimpse}\NormalTok{(diamonds)}
\end{Highlighting}
\end{Shaded}

\begin{verbatim}
Rows: 53,940
Columns: 10
$ carat   <dbl> 0.23, 0.21, 0.23, 0.29, 0.31, 0.24, 0.24, 0.26, 0.22, 0.23, 0.~
$ cut     <ord> Ideal, Premium, Good, Premium, Good, Very Good, Very Good, Ver~
$ color   <ord> E, E, E, I, J, J, I, H, E, H, J, J, F, J, E, E, I, J, J, J, I,~
$ clarity <ord> SI2, SI1, VS1, VS2, SI2, VVS2, VVS1, SI1, VS2, VS1, SI1, VS1, ~
$ depth   <dbl> 61.5, 59.8, 56.9, 62.4, 63.3, 62.8, 62.3, 61.9, 65.1, 59.4, 64~
$ table   <dbl> 55, 61, 65, 58, 58, 57, 57, 55, 61, 61, 55, 56, 61, 54, 62, 58~
$ price   <int> 326, 326, 327, 334, 335, 336, 336, 337, 337, 338, 339, 340, 34~
$ x       <dbl> 3.95, 3.89, 4.05, 4.20, 4.34, 3.94, 3.95, 4.07, 3.87, 4.00, 4.~
$ y       <dbl> 3.98, 3.84, 4.07, 4.23, 4.35, 3.96, 3.98, 4.11, 3.78, 4.05, 4.~
$ z       <dbl> 2.43, 2.31, 2.31, 2.63, 2.75, 2.48, 2.47, 2.53, 2.49, 2.39, 2.~
\end{verbatim}

That's a lot of diamonds! Anyway\ldots{}

\subsection{Exercise 1}\label{exercise-1-14}

We start by visualizing the distributions using boxplots.

\begin{Shaded}
\begin{Highlighting}[]
\FunctionTok{library}\NormalTok{(patchwork)}

\NormalTok{p\_x }\OtherTok{\textless{}{-}}\NormalTok{ diamonds }\SpecialCharTok{|\textgreater{}}
  \FunctionTok{ggplot}\NormalTok{(}\FunctionTok{aes}\NormalTok{(}\AttributeTok{y =}\NormalTok{ x)) }\SpecialCharTok{+}
  \FunctionTok{geom\_boxplot}\NormalTok{()}

\NormalTok{p\_y }\OtherTok{\textless{}{-}}\NormalTok{ diamonds }\SpecialCharTok{|\textgreater{}}
  \FunctionTok{ggplot}\NormalTok{(}\FunctionTok{aes}\NormalTok{(}\AttributeTok{y =}\NormalTok{ y)) }\SpecialCharTok{+}
  \FunctionTok{geom\_boxplot}\NormalTok{()}

\NormalTok{p\_z }\OtherTok{\textless{}{-}}\NormalTok{ diamonds }\SpecialCharTok{|\textgreater{}}
  \FunctionTok{ggplot}\NormalTok{(}\FunctionTok{aes}\NormalTok{(}\AttributeTok{y =}\NormalTok{ z)) }\SpecialCharTok{+}
  \FunctionTok{geom\_boxplot}\NormalTok{()}

\NormalTok{p\_x }\SpecialCharTok{+}\NormalTok{ p\_y }\SpecialCharTok{+}\NormalTok{ p\_z}
\end{Highlighting}
\end{Shaded}

\includegraphics{solutions_r4ds_eda_files/figure-pdf/unnamed-chunk-3-1.pdf}

We see that there are several outliers, as noted in the book. We filter
the data to remove the outliers to visualize the distributions. Make
sure you save the data in a new tibble so you don't remove any variables
from our data! Just because the dimensions (x, y, z) may be misspecified
we may still be interested in other variables for these observations.
See 11.4 in textbook for a better way to deal with faulty values.

\begin{Shaded}
\begin{Highlighting}[]
\NormalTok{diamonds\_filtered }\OtherTok{\textless{}{-}}\NormalTok{ diamonds }\SpecialCharTok{|\textgreater{}}
  \FunctionTok{filter}\NormalTok{(x }\SpecialCharTok{\textgreater{}} \DecValTok{0} \SpecialCharTok{\&}\NormalTok{ y }\SpecialCharTok{\textgreater{}} \DecValTok{0} \SpecialCharTok{\&}\NormalTok{ z }\SpecialCharTok{\textgreater{}} \DecValTok{0}\NormalTok{) }\SpecialCharTok{|\textgreater{}}
  \FunctionTok{filter}\NormalTok{(y }\SpecialCharTok{\textless{}} \DecValTok{20} \SpecialCharTok{\&}\NormalTok{ z }\SpecialCharTok{\textless{}} \DecValTok{20}\NormalTok{)}

\NormalTok{p\_x }\OtherTok{\textless{}{-}}\NormalTok{ diamonds\_filtered }\SpecialCharTok{|\textgreater{}}
  \FunctionTok{ggplot}\NormalTok{(}\FunctionTok{aes}\NormalTok{(}\AttributeTok{y =}\NormalTok{ x)) }\SpecialCharTok{+}
  \FunctionTok{geom\_boxplot}\NormalTok{()}

\NormalTok{p\_y }\OtherTok{\textless{}{-}}\NormalTok{ diamonds\_filtered }\SpecialCharTok{|\textgreater{}}
  \FunctionTok{ggplot}\NormalTok{(}\FunctionTok{aes}\NormalTok{(}\AttributeTok{y =}\NormalTok{ y)) }\SpecialCharTok{+}
  \FunctionTok{geom\_boxplot}\NormalTok{()}

\NormalTok{p\_z }\OtherTok{\textless{}{-}}\NormalTok{ diamonds\_filtered }\SpecialCharTok{|\textgreater{}}
  \FunctionTok{ggplot}\NormalTok{(}\FunctionTok{aes}\NormalTok{(}\AttributeTok{y =}\NormalTok{ z)) }\SpecialCharTok{+}
  \FunctionTok{geom\_boxplot}\NormalTok{()}

\NormalTok{p\_x }\SpecialCharTok{+}\NormalTok{ p\_y }\SpecialCharTok{+}\NormalTok{ p\_z}
\end{Highlighting}
\end{Shaded}

\includegraphics{solutions_r4ds_eda_files/figure-pdf/unnamed-chunk-4-1.pdf}

We see that the distribution of the z-values are slightly higher than x
and y. Thinking about how a diamond is shaped, the height is usually
larger than the width and length, and z is commonly used to label
height. This seems like a plausible explanation for the different
distributions.

\subsection{Exercise 2}\label{exercise-2-11}

\begin{Shaded}
\begin{Highlighting}[]
\NormalTok{diamonds }\SpecialCharTok{|\textgreater{}}
  \FunctionTok{ggplot}\NormalTok{(}\FunctionTok{aes}\NormalTok{(}\AttributeTok{x =}\NormalTok{ price)) }\SpecialCharTok{+} 
  \FunctionTok{geom\_histogram}\NormalTok{(}\AttributeTok{binwidth =} \DecValTok{100}\NormalTok{)}
\end{Highlighting}
\end{Shaded}

\includegraphics{solutions_r4ds_eda_files/figure-pdf/unnamed-chunk-5-1.pdf}

There seems almost like there is a small gap in prices. Let's take a
closer look at that.

\begin{Shaded}
\begin{Highlighting}[]
\NormalTok{diamonds }\SpecialCharTok{|\textgreater{}}
  \FunctionTok{ggplot}\NormalTok{(}\FunctionTok{aes}\NormalTok{(}\AttributeTok{x =}\NormalTok{ price)) }\SpecialCharTok{+} 
  \FunctionTok{geom\_histogram}\NormalTok{(}\AttributeTok{binwidth =} \DecValTok{10}\NormalTok{) }\SpecialCharTok{+} 
  \FunctionTok{coord\_cartesian}\NormalTok{(}\AttributeTok{xlim =} \FunctionTok{c}\NormalTok{(}\DecValTok{1400}\NormalTok{, }\DecValTok{1600}\NormalTok{))}
\end{Highlighting}
\end{Shaded}

\includegraphics{solutions_r4ds_eda_files/figure-pdf/unnamed-chunk-6-1.pdf}

There are no diamonds priced around 1500!

\subsection{Exercise 3}\label{exercise-3-9}

we summarize all diamonds that are \texttt{carat\ ==\ 0.99} and
\texttt{carat\ ==\ 1}.

\begin{Shaded}
\begin{Highlighting}[]
\CommentTok{\# you can do this in many ways, as long as you achieve the same result!}
\NormalTok{diamonds }\SpecialCharTok{|\textgreater{}}
  \FunctionTok{filter}\NormalTok{(carat }\SpecialCharTok{\%in\%} \FunctionTok{c}\NormalTok{(}\FloatTok{0.99}\NormalTok{, }\DecValTok{1}\NormalTok{)) }\SpecialCharTok{|\textgreater{}}
  \FunctionTok{summarize}\NormalTok{(}
    \AttributeTok{count =} \FunctionTok{n}\NormalTok{(),}
    \AttributeTok{.by =}\NormalTok{ carat)}
\end{Highlighting}
\end{Shaded}

\begin{verbatim}
# A tibble: 2 x 2
  carat count
  <dbl> <int>
1  1     1558
2  0.99    23
\end{verbatim}

We see that there are very few diamonds of 0.99 carat compared to 1.
Would you expect the price of diamonds to go up as the carat increases?
We can investigate the mean price of the two groups.

\begin{Shaded}
\begin{Highlighting}[]
\CommentTok{\# you can do this in many ways, as long as you achieve the same result!}
\NormalTok{diamonds }\SpecialCharTok{|\textgreater{}}
  \FunctionTok{filter}\NormalTok{(carat }\SpecialCharTok{\%in\%} \FunctionTok{c}\NormalTok{(}\FloatTok{0.99}\NormalTok{, }\DecValTok{1}\NormalTok{)) }\SpecialCharTok{|\textgreater{}}
  \FunctionTok{summarize}\NormalTok{(}
    \AttributeTok{mean\_price =} \FunctionTok{mean}\NormalTok{(price),}
    \AttributeTok{.by =}\NormalTok{ carat)}
\end{Highlighting}
\end{Shaded}

\begin{verbatim}
# A tibble: 2 x 2
  carat mean_price
  <dbl>      <dbl>
1  1         5242.
2  0.99      4406.
\end{verbatim}

It seems that cutting a diamond so it is less than carat 1 is costly, so
it's better to try to keep it at 1! FIY
\href{https://en.wikipedia.org/wiki/Carat_(mass)}{carat} is a measure of
weight. The more you know about your data, the better questions (and
answers) you get!

\subsection{Exercise 4}\label{exercise-4-9}

We see that \texttt{coord\_cartesian()} works as zooming, while
\texttt{xlim} alters our data! Remember that \texttt{xlim}
\emph{removes} the data outside the specified values.

\begin{Shaded}
\begin{Highlighting}[]
\CommentTok{\# binwidth not specified}
\NormalTok{p1 }\OtherTok{\textless{}{-}}\NormalTok{ diamonds }\SpecialCharTok{|\textgreater{}}
  \FunctionTok{ggplot}\NormalTok{(}\FunctionTok{aes}\NormalTok{(}\AttributeTok{x =}\NormalTok{ price)) }\SpecialCharTok{+} 
  \FunctionTok{geom\_histogram}\NormalTok{() }\SpecialCharTok{+} 
  \FunctionTok{coord\_cartesian}\NormalTok{(}\AttributeTok{xlim =} \FunctionTok{c}\NormalTok{(}\DecValTok{1400}\NormalTok{, }\DecValTok{1600}\NormalTok{))}

\NormalTok{p2 }\OtherTok{\textless{}{-}}\NormalTok{ diamonds }\SpecialCharTok{|\textgreater{}}
  \FunctionTok{ggplot}\NormalTok{(}\FunctionTok{aes}\NormalTok{(}\AttributeTok{x =}\NormalTok{ price)) }\SpecialCharTok{+} 
  \FunctionTok{geom\_histogram}\NormalTok{() }\SpecialCharTok{+} 
  \FunctionTok{xlim}\NormalTok{(}\DecValTok{1400}\NormalTok{, }\DecValTok{1600}\NormalTok{)}
\NormalTok{p1 }\SpecialCharTok{/}\NormalTok{ p2}
\end{Highlighting}
\end{Shaded}

\begin{verbatim}
`stat_bin()` using `bins = 30`. Pick better value with `binwidth`.
`stat_bin()` using `bins = 30`. Pick better value with `binwidth`.
\end{verbatim}

\begin{verbatim}
Warning: Removed 52871 rows containing non-finite outside the scale range
(`stat_bin()`).
\end{verbatim}

\begin{verbatim}
Warning: Removed 2 rows containing missing values or values outside the scale range
(`geom_bar()`).
\end{verbatim}

\includegraphics{solutions_r4ds_eda_files/figure-pdf/unnamed-chunk-9-1.pdf}

\begin{Shaded}
\begin{Highlighting}[]
\CommentTok{\# Setting binwidth to the same size as the interval}
\NormalTok{p1 }\OtherTok{\textless{}{-}}\NormalTok{ diamonds }\SpecialCharTok{|\textgreater{}}
  \FunctionTok{ggplot}\NormalTok{(}\FunctionTok{aes}\NormalTok{(}\AttributeTok{x =}\NormalTok{ price)) }\SpecialCharTok{+} 
  \FunctionTok{geom\_histogram}\NormalTok{(}\AttributeTok{binwidth =} \DecValTok{200}\NormalTok{) }\SpecialCharTok{+} 
  \FunctionTok{coord\_cartesian}\NormalTok{(}\AttributeTok{xlim =} \FunctionTok{c}\NormalTok{(}\DecValTok{1400}\NormalTok{, }\DecValTok{1600}\NormalTok{))}

\NormalTok{p2 }\OtherTok{\textless{}{-}}\NormalTok{ diamonds }\SpecialCharTok{|\textgreater{}}
  \FunctionTok{ggplot}\NormalTok{(}\FunctionTok{aes}\NormalTok{(}\AttributeTok{x =}\NormalTok{ price)) }\SpecialCharTok{+} 
  \FunctionTok{geom\_histogram}\NormalTok{(}\AttributeTok{binwidth =} \DecValTok{200}\NormalTok{) }\SpecialCharTok{+} 
  \FunctionTok{xlim}\NormalTok{(}\DecValTok{1400}\NormalTok{, }\DecValTok{1600}\NormalTok{)}
\NormalTok{p1 }\SpecialCharTok{/}\NormalTok{ p2}
\end{Highlighting}
\end{Shaded}

\begin{verbatim}
Warning: Removed 52871 rows containing non-finite outside the scale range
(`stat_bin()`).
\end{verbatim}

\begin{verbatim}
Warning: Removed 2 rows containing missing values or values outside the scale range
(`geom_bar()`).
\end{verbatim}

\includegraphics{solutions_r4ds_eda_files/figure-pdf/unnamed-chunk-10-1.pdf}

\section{11.4 Unusual values}\label{unusual-values}

\subsection{Exercise 1}\label{exercise-1-15}

Let's take a look by making all prices less than 1000 \texttt{NA} and
plotting a histogram.

\begin{Shaded}
\begin{Highlighting}[]
\NormalTok{diamonds }\SpecialCharTok{|\textgreater{}}
  \FunctionTok{mutate}\NormalTok{(}\AttributeTok{price =} \FunctionTok{ifelse}\NormalTok{(price }\SpecialCharTok{\textless{}} \DecValTok{1000}\NormalTok{, }\ConstantTok{NA}\NormalTok{, price)) }\SpecialCharTok{|\textgreater{}}
  \FunctionTok{ggplot}\NormalTok{(}\FunctionTok{aes}\NormalTok{(}\AttributeTok{x =}\NormalTok{ price)) }\SpecialCharTok{+} 
  \FunctionTok{geom\_histogram}\NormalTok{()}
\end{Highlighting}
\end{Shaded}

\begin{verbatim}
`stat_bin()` using `bins = 30`. Pick better value with `binwidth`.
\end{verbatim}

\begin{verbatim}
Warning: Removed 14499 rows containing non-finite outside the scale range
(`stat_bin()`).
\end{verbatim}

\includegraphics{solutions_r4ds_eda_files/figure-pdf/unnamed-chunk-11-1.pdf}

The \texttt{NA} values are ignored when we plot a histogram. Now let's
see what happens if we set all Fair cuts to \texttt{NA} and then plot a
bar chart.

\begin{Shaded}
\begin{Highlighting}[]
\FunctionTok{glimpse}\NormalTok{(diamonds)}
\end{Highlighting}
\end{Shaded}

\begin{verbatim}
Rows: 53,940
Columns: 10
$ carat   <dbl> 0.23, 0.21, 0.23, 0.29, 0.31, 0.24, 0.24, 0.26, 0.22, 0.23, 0.~
$ cut     <ord> Ideal, Premium, Good, Premium, Good, Very Good, Very Good, Ver~
$ color   <ord> E, E, E, I, J, J, I, H, E, H, J, J, F, J, E, E, I, J, J, J, I,~
$ clarity <ord> SI2, SI1, VS1, VS2, SI2, VVS2, VVS1, SI1, VS2, VS1, SI1, VS1, ~
$ depth   <dbl> 61.5, 59.8, 56.9, 62.4, 63.3, 62.8, 62.3, 61.9, 65.1, 59.4, 64~
$ table   <dbl> 55, 61, 65, 58, 58, 57, 57, 55, 61, 61, 55, 56, 61, 54, 62, 58~
$ price   <int> 326, 326, 327, 334, 335, 336, 336, 337, 337, 338, 339, 340, 34~
$ x       <dbl> 3.95, 3.89, 4.05, 4.20, 4.34, 3.94, 3.95, 4.07, 3.87, 4.00, 4.~
$ y       <dbl> 3.98, 3.84, 4.07, 4.23, 4.35, 3.96, 3.98, 4.11, 3.78, 4.05, 4.~
$ z       <dbl> 2.43, 2.31, 2.31, 2.63, 2.75, 2.48, 2.47, 2.53, 2.49, 2.39, 2.~
\end{verbatim}

\begin{Shaded}
\begin{Highlighting}[]
\NormalTok{diamonds }\SpecialCharTok{|\textgreater{}}
  \FunctionTok{mutate}\NormalTok{(}\AttributeTok{cut =} \FunctionTok{case\_when}\NormalTok{( }\CommentTok{\# Here we use case\_when instead of ifelse. ifelse turns the variable into an integer variable}
\NormalTok{    cut }\SpecialCharTok{==} \StringTok{"Fair"} \SpecialCharTok{\textasciitilde{}} \ConstantTok{NA}\NormalTok{, }
    \ConstantTok{TRUE} \SpecialCharTok{\textasciitilde{}}\NormalTok{ cut)) }\SpecialCharTok{|\textgreater{}}
  \FunctionTok{ggplot}\NormalTok{(}\FunctionTok{aes}\NormalTok{(}\AttributeTok{x =}\NormalTok{ cut)) }\SpecialCharTok{+} 
  \FunctionTok{geom\_bar}\NormalTok{()}
\end{Highlighting}
\end{Shaded}

\includegraphics{solutions_r4ds_eda_files/figure-pdf/unnamed-chunk-12-1.pdf}

The bar chart shows the number of missing values as it's own category.
This is a reasonable way of showing missing values, compared to the
histogram that shows the counts of bins. Where would the \texttt{NA} bin
go?

\subsection{Exercise 2}\label{exercise-2-12}

It means that the \texttt{NA} variables are removed. These functions
throws an error and doesn't return a value if there are missing values
and \texttt{na.rm} is not set to \texttt{TRUE}!

\subsection{Exercise 3}\label{exercise-3-10}

We copy the code from the book and use \texttt{scales\ =\ "free\_y"} to
create two y-axes.

\begin{Shaded}
\begin{Highlighting}[]
\NormalTok{nycflights13}\SpecialCharTok{::}\NormalTok{flights }\SpecialCharTok{|\textgreater{}} 
  \FunctionTok{mutate}\NormalTok{(}
    \AttributeTok{cancelled =} \FunctionTok{is.na}\NormalTok{(dep\_time),}
    \AttributeTok{sched\_hour =}\NormalTok{ sched\_dep\_time }\SpecialCharTok{\%/\%} \DecValTok{100}\NormalTok{,}
    \AttributeTok{sched\_min =}\NormalTok{ sched\_dep\_time }\SpecialCharTok{\%\%} \DecValTok{100}\NormalTok{,}
    \AttributeTok{sched\_dep\_time =}\NormalTok{ sched\_hour }\SpecialCharTok{+}\NormalTok{ (sched\_min }\SpecialCharTok{/} \DecValTok{60}\NormalTok{)}
\NormalTok{  ) }\SpecialCharTok{|\textgreater{}} 
  \FunctionTok{ggplot}\NormalTok{(}\FunctionTok{aes}\NormalTok{(}\AttributeTok{x =}\NormalTok{ sched\_dep\_time)) }\SpecialCharTok{+} 
  \FunctionTok{geom\_freqpoly}\NormalTok{(}\FunctionTok{aes}\NormalTok{(}\AttributeTok{color =}\NormalTok{ cancelled), }\AttributeTok{binwidth =} \DecValTok{1}\SpecialCharTok{/}\DecValTok{4}\NormalTok{) }\SpecialCharTok{+} 
  \FunctionTok{facet\_wrap}\NormalTok{(}\SpecialCharTok{\textasciitilde{}}\NormalTok{ cancelled, }\AttributeTok{scales =} \StringTok{"free\_y"}\NormalTok{)}
\end{Highlighting}
\end{Shaded}

\includegraphics{solutions_r4ds_eda_files/figure-pdf/unnamed-chunk-13-1.pdf}

Note the different scales of the y-axes! If the variables are of similar
scale it is useful to plot them against the same axis, using the default
`scales = ``fixed'', but in this case it may be a better choice to show
the y-axes independently. It is important to think about what you want
the plot to convey, but also to not be misleading. See the plot below
where we instead create two plots, but don't show the second y-axis.

\begin{Shaded}
\begin{Highlighting}[]
\NormalTok{p1 }\OtherTok{\textless{}{-}}\NormalTok{ nycflights13}\SpecialCharTok{::}\NormalTok{flights }\SpecialCharTok{|\textgreater{}} 
  \FunctionTok{mutate}\NormalTok{(}
    \AttributeTok{cancelled =} \FunctionTok{is.na}\NormalTok{(dep\_time),}
    \AttributeTok{sched\_hour =}\NormalTok{ sched\_dep\_time }\SpecialCharTok{\%/\%} \DecValTok{100}\NormalTok{,}
    \AttributeTok{sched\_min =}\NormalTok{ sched\_dep\_time }\SpecialCharTok{\%\%} \DecValTok{100}\NormalTok{,}
    \AttributeTok{sched\_dep\_time =}\NormalTok{ sched\_hour }\SpecialCharTok{+}\NormalTok{ (sched\_min }\SpecialCharTok{/} \DecValTok{60}\NormalTok{)}
\NormalTok{  ) }\SpecialCharTok{|\textgreater{}} 
  \FunctionTok{filter}\NormalTok{(cancelled }\SpecialCharTok{==} \ConstantTok{FALSE}\NormalTok{) }\SpecialCharTok{|\textgreater{}}
  \FunctionTok{ggplot}\NormalTok{(}\FunctionTok{aes}\NormalTok{(}\AttributeTok{x =}\NormalTok{ sched\_dep\_time)) }\SpecialCharTok{+} 
  \FunctionTok{geom\_freqpoly}\NormalTok{(}\AttributeTok{color =} \StringTok{"red"}\NormalTok{, }\AttributeTok{binwidth =} \DecValTok{1}\SpecialCharTok{/}\DecValTok{4}\NormalTok{) }\SpecialCharTok{+} 
  \FunctionTok{ggtitle}\NormalTok{(}\StringTok{"Departed"}\NormalTok{)}

\NormalTok{p2 }\OtherTok{\textless{}{-}}\NormalTok{ nycflights13}\SpecialCharTok{::}\NormalTok{flights }\SpecialCharTok{|\textgreater{}} 
  \FunctionTok{mutate}\NormalTok{(}
    \AttributeTok{cancelled =} \FunctionTok{is.na}\NormalTok{(dep\_time),}
    \AttributeTok{sched\_hour =}\NormalTok{ sched\_dep\_time }\SpecialCharTok{\%/\%} \DecValTok{100}\NormalTok{,}
    \AttributeTok{sched\_min =}\NormalTok{ sched\_dep\_time }\SpecialCharTok{\%\%} \DecValTok{100}\NormalTok{,}
    \AttributeTok{sched\_dep\_time =}\NormalTok{ sched\_hour }\SpecialCharTok{+}\NormalTok{ (sched\_min }\SpecialCharTok{/} \DecValTok{60}\NormalTok{)}
\NormalTok{  ) }\SpecialCharTok{|\textgreater{}} 
  \FunctionTok{filter}\NormalTok{(cancelled }\SpecialCharTok{==} \ConstantTok{TRUE}\NormalTok{) }\SpecialCharTok{|\textgreater{}}
  \FunctionTok{ggplot}\NormalTok{(}\FunctionTok{aes}\NormalTok{(}\AttributeTok{x =}\NormalTok{ sched\_dep\_time)) }\SpecialCharTok{+} 
  \FunctionTok{geom\_freqpoly}\NormalTok{(}\AttributeTok{color =} \StringTok{"blue"}\NormalTok{, }\AttributeTok{binwidth =} \DecValTok{1}\SpecialCharTok{/}\DecValTok{4}\NormalTok{) }\SpecialCharTok{+} 
  \FunctionTok{theme}\NormalTok{(}
    \AttributeTok{axis.title.y =} \FunctionTok{element\_blank}\NormalTok{(),}
    \AttributeTok{axis.text.y =} \FunctionTok{element\_blank}\NormalTok{(),}
    \AttributeTok{axis.ticks.y =} \FunctionTok{element\_blank}\NormalTok{()) }\SpecialCharTok{+}
  \FunctionTok{ggtitle}\NormalTok{(}\StringTok{"Cancelled"}\NormalTok{)}
  
\NormalTok{p1 }\SpecialCharTok{+}\NormalTok{ p2}
\end{Highlighting}
\end{Shaded}

\includegraphics{solutions_r4ds_eda_files/figure-pdf/unnamed-chunk-14-1.pdf}

Notice how the plots look very similar, but now it seems there are about
the same amounts of departed and cancelled flight!

\begin{quote}
With great power comes great responsibility\ldots{} Be transparent with
how you create your data visualizations.
\end{quote}

\section{11.5.1 Covariation - A categorical and numerical
variable}\label{covariation---a-categorical-and-numerical-variable}

\subsection{Exercise 1}\label{exercise-1-16}

We can use boxplots and density plots to capture key values of the
distributions, the quartiles and medians in the boxplot, as well as
visualizing the distribution as a whole using the density plot.

\begin{Shaded}
\begin{Highlighting}[]
\NormalTok{p1 }\OtherTok{\textless{}{-}}\NormalTok{ nycflights13}\SpecialCharTok{::}\NormalTok{flights }\SpecialCharTok{|\textgreater{}} 
  \FunctionTok{mutate}\NormalTok{(}
    \AttributeTok{cancelled =} \FunctionTok{ifelse}\NormalTok{(}\FunctionTok{is.na}\NormalTok{(dep\_time), }\StringTok{"Cancelled"}\NormalTok{, }\StringTok{"Departed"}\NormalTok{), }\CommentTok{\# Improving labelling for visualization}
    \AttributeTok{sched\_hour =}\NormalTok{ sched\_dep\_time }\SpecialCharTok{\%/\%} \DecValTok{100}\NormalTok{,}
    \AttributeTok{sched\_min =}\NormalTok{ sched\_dep\_time }\SpecialCharTok{\%\%} \DecValTok{100}\NormalTok{,}
    \AttributeTok{sched\_dep\_time =}\NormalTok{ sched\_hour }\SpecialCharTok{+}\NormalTok{ (sched\_min }\SpecialCharTok{/} \DecValTok{60}\NormalTok{)}
\NormalTok{  ) }\SpecialCharTok{|\textgreater{}} 
  \FunctionTok{ggplot}\NormalTok{(}\FunctionTok{aes}\NormalTok{(}\AttributeTok{y =}\NormalTok{ sched\_dep\_time, }\AttributeTok{fill =}\NormalTok{ cancelled)) }\SpecialCharTok{+} 
  \FunctionTok{geom\_boxplot}\NormalTok{(}\AttributeTok{alpha =} \FloatTok{0.3}\NormalTok{) }\SpecialCharTok{+} 
  \FunctionTok{labs}\NormalTok{(}
    \AttributeTok{y =} \StringTok{"Scheduled departure time"}\NormalTok{,}
    \AttributeTok{fill =} \StringTok{""}\NormalTok{) }\SpecialCharTok{+}
  \FunctionTok{theme}\NormalTok{(}
    \AttributeTok{legend.position =} \StringTok{"none"}\NormalTok{,}
    \AttributeTok{axis.title.x =} \FunctionTok{element\_blank}\NormalTok{(),}
    \AttributeTok{axis.text.x =} \FunctionTok{element\_blank}\NormalTok{(),}
    \AttributeTok{axis.ticks.x =} \FunctionTok{element\_blank}\NormalTok{())}

\NormalTok{p2 }\OtherTok{\textless{}{-}}\NormalTok{ nycflights13}\SpecialCharTok{::}\NormalTok{flights }\SpecialCharTok{|\textgreater{}} 
  \FunctionTok{mutate}\NormalTok{(}
    \AttributeTok{cancelled =} \FunctionTok{ifelse}\NormalTok{(}\FunctionTok{is.na}\NormalTok{(dep\_time), }\StringTok{"Cancelled"}\NormalTok{, }\StringTok{"Departed"}\NormalTok{), }\CommentTok{\# Improving labelling for visualization}
    \AttributeTok{sched\_hour =}\NormalTok{ sched\_dep\_time }\SpecialCharTok{\%/\%} \DecValTok{100}\NormalTok{,}
    \AttributeTok{sched\_min =}\NormalTok{ sched\_dep\_time }\SpecialCharTok{\%\%} \DecValTok{100}\NormalTok{,}
    \AttributeTok{sched\_dep\_time =}\NormalTok{ sched\_hour }\SpecialCharTok{+}\NormalTok{ (sched\_min }\SpecialCharTok{/} \DecValTok{60}\NormalTok{)}
\NormalTok{  ) }\SpecialCharTok{|\textgreater{}} 
  \FunctionTok{ggplot}\NormalTok{(}
    \FunctionTok{aes}\NormalTok{(}\AttributeTok{y =}\NormalTok{ sched\_dep\_time, }\AttributeTok{fill =}\NormalTok{ cancelled)) }\SpecialCharTok{+} 
  \FunctionTok{geom\_density}\NormalTok{(}\AttributeTok{alpha =} \FloatTok{0.3}\NormalTok{) }\SpecialCharTok{+} 
  \FunctionTok{labs}\NormalTok{(}
    \AttributeTok{x =} \StringTok{"Scheduled departure time"}\NormalTok{,}
    \AttributeTok{fill =} \StringTok{""}\NormalTok{) }\SpecialCharTok{+}
  \FunctionTok{theme}\NormalTok{(}
    \AttributeTok{axis.title.y =} \FunctionTok{element\_blank}\NormalTok{(),}
    \AttributeTok{axis.text.y =} \FunctionTok{element\_blank}\NormalTok{(),}
    \AttributeTok{axis.ticks.y =} \FunctionTok{element\_blank}\NormalTok{(),}
    \AttributeTok{axis.title.x =} \FunctionTok{element\_blank}\NormalTok{(),}
    \AttributeTok{axis.text.x =} \FunctionTok{element\_blank}\NormalTok{(),}
    \AttributeTok{axis.ticks.x =} \FunctionTok{element\_blank}\NormalTok{())}

\NormalTok{p1 }\SpecialCharTok{+}\NormalTok{ p2}
\end{Highlighting}
\end{Shaded}

\includegraphics{solutions_r4ds_eda_files/figure-pdf/unnamed-chunk-15-1.pdf}

\subsection{Exercise 2}\label{exercise-2-13}

We can investigate the four Cs: \texttt{carat}, \texttt{cut},
\texttt{color}, and \texttt{clarity}. \texttt{carat} is a continuous
variable, while the other are (ordered) factors, see \texttt{?diamonds}
for a description. Here we will study the relationship between
\texttt{carat} and \texttt{price}, starting by visualizing their
relationship using a scatter plot with a line.

\begin{Shaded}
\begin{Highlighting}[]
\NormalTok{diamonds }\SpecialCharTok{|\textgreater{}}
  \FunctionTok{ggplot}\NormalTok{(}\FunctionTok{aes}\NormalTok{(}\AttributeTok{x =}\NormalTok{ carat, }\AttributeTok{y =}\NormalTok{ price)) }\SpecialCharTok{+}
  \FunctionTok{geom\_point}\NormalTok{(}\AttributeTok{alpha =} \FloatTok{0.1}\NormalTok{) }\SpecialCharTok{+} 
  \FunctionTok{geom\_smooth}\NormalTok{(}\AttributeTok{method =} \StringTok{"lm"}\NormalTok{, }\AttributeTok{se =}\NormalTok{ F) }\SpecialCharTok{+}
  \FunctionTok{coord\_cartesian}\NormalTok{(}\AttributeTok{ylim =} \FunctionTok{c}\NormalTok{(}\FunctionTok{min}\NormalTok{(diamonds}\SpecialCharTok{$}\NormalTok{price), }\FunctionTok{max}\NormalTok{(diamonds}\SpecialCharTok{$}\NormalTok{price)))}
\end{Highlighting}
\end{Shaded}

\begin{verbatim}
`geom_smooth()` using formula = 'y ~ x'
\end{verbatim}

\includegraphics{solutions_r4ds_eda_files/figure-pdf/unnamed-chunk-16-1.pdf}

We can clearly see that \texttt{carat} is important for the price of a
diamond, but what is the relationship between carat and cuts?

\begin{Shaded}
\begin{Highlighting}[]
\NormalTok{p1 }\OtherTok{\textless{}{-}}\NormalTok{ diamonds }\SpecialCharTok{|\textgreater{}}
  \FunctionTok{ggplot}\NormalTok{(}\FunctionTok{aes}\NormalTok{(}\AttributeTok{x =}\NormalTok{ carat, }\AttributeTok{fill =}\NormalTok{ cut)) }\SpecialCharTok{+} 
  \FunctionTok{geom\_density}\NormalTok{(}\AttributeTok{alpha =} \FloatTok{0.2}\NormalTok{) }\SpecialCharTok{+} 
  \FunctionTok{coord\_cartesian}\NormalTok{(}\AttributeTok{xlim =} \FunctionTok{c}\NormalTok{(}\DecValTok{0}\NormalTok{, }\DecValTok{2}\NormalTok{)) }\SpecialCharTok{+}
  \FunctionTok{labs}\NormalTok{(}\AttributeTok{x =} \StringTok{"Carat"}\NormalTok{) }\SpecialCharTok{+}
  \FunctionTok{theme}\NormalTok{(}
    \AttributeTok{legend.position =} \StringTok{"bottom"}\NormalTok{,}
    \AttributeTok{axis.title.y =} \FunctionTok{element\_blank}\NormalTok{(),}
    \AttributeTok{axis.ticks.y =} \FunctionTok{element\_blank}\NormalTok{(),}
    \AttributeTok{axis.text.y =} \FunctionTok{element\_blank}\NormalTok{())}
\NormalTok{p2 }\OtherTok{\textless{}{-}}\NormalTok{ diamonds }\SpecialCharTok{|\textgreater{}}
  \FunctionTok{ggplot}\NormalTok{(}\FunctionTok{aes}\NormalTok{(}\AttributeTok{x =} \FunctionTok{fct\_reorder}\NormalTok{(cut, carat, median), }\AttributeTok{y =}\NormalTok{ carat, }\AttributeTok{fill =}\NormalTok{ cut)) }\SpecialCharTok{+} 
  \FunctionTok{geom\_boxplot}\NormalTok{(}\AttributeTok{alpha =} \FloatTok{0.3}\NormalTok{) }\SpecialCharTok{+}
  \FunctionTok{labs}\NormalTok{(}
    \AttributeTok{x =} \StringTok{"Cut"}\NormalTok{,}
    \AttributeTok{y =} \StringTok{"Carat"}\NormalTok{) }\SpecialCharTok{+}
  \FunctionTok{theme}\NormalTok{(}\AttributeTok{legend.position =} \StringTok{"none"}\NormalTok{)}

\NormalTok{p1 }\SpecialCharTok{/}\NormalTok{ p2}
\end{Highlighting}
\end{Shaded}

\includegraphics{solutions_r4ds_eda_files/figure-pdf/unnamed-chunk-17-1.pdf}

We see that there are more diamonds with a fair cut and higher carat,
this could possibly explain the unexpected order of median price.
Observe that we have zoomed in on the density plot to better see the
median values!

If you find another variable that better explains the counter-intuitive
result, feel free to share!

\subsection{Exercise 3}\label{exercise-3-11}

In the previous exercise we used a boxplot along with a density plot,
however the exes were not align, Carat being on the x-axis in the
density plot and the y-axis in the boxplot. We recreate the plot from
before, using \texttt{coord\_flip} on the boxplot to make it correspond
to the density plot (we also make minor adjustments, dropping the zoom
and moving the legend).

\begin{Shaded}
\begin{Highlighting}[]
\NormalTok{p1 }\OtherTok{\textless{}{-}}\NormalTok{ diamonds }\SpecialCharTok{|\textgreater{}}
  \FunctionTok{ggplot}\NormalTok{(}\FunctionTok{aes}\NormalTok{(}\AttributeTok{x =}\NormalTok{ carat, }\AttributeTok{fill =}\NormalTok{ cut)) }\SpecialCharTok{+} 
  \FunctionTok{geom\_density}\NormalTok{(}\AttributeTok{alpha =} \FloatTok{0.2}\NormalTok{) }\SpecialCharTok{+} 
  \FunctionTok{labs}\NormalTok{(}\AttributeTok{x =} \StringTok{"Carat"}\NormalTok{) }\SpecialCharTok{+}
  \FunctionTok{theme}\NormalTok{(}
    \AttributeTok{legend.position =} \StringTok{"top"}\NormalTok{,}
    \AttributeTok{axis.title.y =} \FunctionTok{element\_blank}\NormalTok{(),}
    \AttributeTok{axis.ticks.y =} \FunctionTok{element\_blank}\NormalTok{(),}
    \AttributeTok{axis.text.y =} \FunctionTok{element\_blank}\NormalTok{())}
\NormalTok{p2 }\OtherTok{\textless{}{-}}\NormalTok{ diamonds }\SpecialCharTok{|\textgreater{}}
  \FunctionTok{ggplot}\NormalTok{(}\FunctionTok{aes}\NormalTok{(}\AttributeTok{x =} \FunctionTok{fct\_reorder}\NormalTok{(cut, carat, median), }\AttributeTok{y =}\NormalTok{ carat, }\AttributeTok{fill =}\NormalTok{ cut)) }\SpecialCharTok{+} 
  \FunctionTok{geom\_boxplot}\NormalTok{(}\AttributeTok{alpha =} \FloatTok{0.3}\NormalTok{) }\SpecialCharTok{+}
  \FunctionTok{labs}\NormalTok{(}
    \AttributeTok{x =} \StringTok{"Cut"}\NormalTok{,}
    \AttributeTok{y =} \StringTok{"Carat"}\NormalTok{) }\SpecialCharTok{+}
  \FunctionTok{theme}\NormalTok{(}\AttributeTok{legend.position =} \StringTok{"none"}\NormalTok{) }\SpecialCharTok{+}
  \FunctionTok{coord\_flip}\NormalTok{()}

\NormalTok{p1 }\SpecialCharTok{/}\NormalTok{ p2}
\end{Highlighting}
\end{Shaded}

\includegraphics{solutions_r4ds_eda_files/figure-pdf/unnamed-chunk-18-1.pdf}

It is usually a good idea to represent the same variable on the same
axis if you display multiple plots. One benefit of using
\texttt{coord\_flip} is that you don't have to respecify the labels,
\texttt{coord\_flip} takes care of that for you. However, if you have
defined axis related parameters in \texttt{theme}, you will have to
change them manually.

\subsection{Exercise 4}\label{exercise-4-10}

We install the \texttt{lvplot} package and compare it with a boxplot for
one cut first.

\begin{Shaded}
\begin{Highlighting}[]
\CommentTok{\#install.packages("lvplot")}
\FunctionTok{library}\NormalTok{(lvplot)}
\NormalTok{p1 }\OtherTok{\textless{}{-}}\NormalTok{ diamonds }\SpecialCharTok{|\textgreater{}}
  \FunctionTok{filter}\NormalTok{(cut }\SpecialCharTok{==} \StringTok{"Fair"}\NormalTok{) }\SpecialCharTok{|\textgreater{}} 
  \FunctionTok{ggplot}\NormalTok{(}\FunctionTok{aes}\NormalTok{(}\AttributeTok{x =}\NormalTok{ cut, }\AttributeTok{y =}\NormalTok{ price)) }\SpecialCharTok{+}
  \FunctionTok{geom\_lv}\NormalTok{() }\SpecialCharTok{+}
  \FunctionTok{labs}\NormalTok{(}
    \AttributeTok{x =} \StringTok{"Cut"}\NormalTok{,}
    \AttributeTok{y =} \StringTok{"Price"}\NormalTok{)}

\NormalTok{p2 }\OtherTok{\textless{}{-}}\NormalTok{ diamonds }\SpecialCharTok{|\textgreater{}}
  \FunctionTok{filter}\NormalTok{(cut }\SpecialCharTok{==} \StringTok{"Fair"}\NormalTok{) }\SpecialCharTok{|\textgreater{}} 
  \FunctionTok{ggplot}\NormalTok{(}\FunctionTok{aes}\NormalTok{(}\AttributeTok{y =}\NormalTok{ price)) }\SpecialCharTok{+} 
  \FunctionTok{geom\_boxplot}\NormalTok{(}\AttributeTok{alpha =} \FloatTok{0.3}\NormalTok{) }\SpecialCharTok{+}
  \FunctionTok{labs}\NormalTok{(}
    \AttributeTok{x =} \StringTok{"Cut"}\NormalTok{,}
    \AttributeTok{y =} \StringTok{"Price"}\NormalTok{) }\SpecialCharTok{+}
  \FunctionTok{theme}\NormalTok{(}\AttributeTok{legend.position =} \StringTok{"none"}\NormalTok{)}

\NormalTok{p1 }\SpecialCharTok{+}\NormalTok{ p2}
\end{Highlighting}
\end{Shaded}

\includegraphics{solutions_r4ds_eda_files/figure-pdf/unnamed-chunk-19-1.pdf}

The lower limit of the box in the boxplot is the first quartile (25th
percentile) and the upper limit is the third quartile (75th percentile).
The letter-value plot, \texttt{lvplot}, works in a similar way, where
the widest box is the same as the boxplot, and the following boxes are
the octiles (the lower limit is the value which 1/8 of all values are
lower than and the upper limit is the value value which 1/8 of all
values are higher than). It is important to note that the width of the
boxes \emph{do not} represent how many observations fall into that box,
but is in fact completely arbitrary, just as for a boxplot. For more
information about the letter-value plot see
\href{https://www.tandfonline.com/doi/full/10.1080/10618600.2017.1305277}{this
paper}.

Here is the letter-value plot of price vs.~cut for all cuts.

\begin{Shaded}
\begin{Highlighting}[]
\NormalTok{diamonds }\SpecialCharTok{|\textgreater{}}
  \FunctionTok{ggplot}\NormalTok{(}\FunctionTok{aes}\NormalTok{(}\AttributeTok{x =}\NormalTok{ cut, }\AttributeTok{y =}\NormalTok{ price)) }\SpecialCharTok{+}
  \FunctionTok{geom\_lv}\NormalTok{() }\SpecialCharTok{+}
  \FunctionTok{labs}\NormalTok{(}
    \AttributeTok{x =} \StringTok{"Cut"}\NormalTok{,}
    \AttributeTok{y =} \StringTok{"Price"}\NormalTok{)}
\end{Highlighting}
\end{Shaded}

\includegraphics{solutions_r4ds_eda_files/figure-pdf/unnamed-chunk-20-1.pdf}

There is no right or wrong when picking a specific plot (as long as it
shows what you intend to show). Which plot do you like more?

\subsection{Exercise 5}\label{exercise-5-5}

Observe the different use of axes, where some plots use counts or
density, others require a y-variable. Violin plot

\begin{Shaded}
\begin{Highlighting}[]
\NormalTok{diamonds }\SpecialCharTok{|\textgreater{}}
  \FunctionTok{ggplot}\NormalTok{(}\FunctionTok{aes}\NormalTok{(}\AttributeTok{x =}\NormalTok{ cut, }\AttributeTok{y =}\NormalTok{ price)) }\SpecialCharTok{+}
  \FunctionTok{geom\_violin}\NormalTok{() }\SpecialCharTok{+}
  \FunctionTok{labs}\NormalTok{(}
    \AttributeTok{x =} \StringTok{"Cut"}\NormalTok{,}
    \AttributeTok{y =} \StringTok{"Price"}\NormalTok{)}
\end{Highlighting}
\end{Shaded}

\includegraphics{solutions_r4ds_eda_files/figure-pdf/unnamed-chunk-21-1.pdf}

Histogram with \texttt{facet\_grid} and \texttt{scales\ =\ "free\_y"}.

\begin{Shaded}
\begin{Highlighting}[]
\NormalTok{diamonds }\SpecialCharTok{|\textgreater{}}
  \FunctionTok{ggplot}\NormalTok{(}\FunctionTok{aes}\NormalTok{(}\AttributeTok{x =}\NormalTok{ price)) }\SpecialCharTok{+}
  \FunctionTok{geom\_histogram}\NormalTok{() }\SpecialCharTok{+}
  \FunctionTok{labs}\NormalTok{(}
    \AttributeTok{x =} \StringTok{"Price"}\NormalTok{,}
    \AttributeTok{y =} \StringTok{"Count"}\NormalTok{) }\SpecialCharTok{+}
  \FunctionTok{facet\_grid}\NormalTok{(}\AttributeTok{rows =} \FunctionTok{vars}\NormalTok{(cut), }\AttributeTok{scales =} \StringTok{"free\_y"}\NormalTok{)}
\end{Highlighting}
\end{Shaded}

\begin{verbatim}
`stat_bin()` using `bins = 30`. Pick better value with `binwidth`.
\end{verbatim}

\includegraphics{solutions_r4ds_eda_files/figure-pdf/unnamed-chunk-22-1.pdf}

\texttt{freq\_poly}

\begin{Shaded}
\begin{Highlighting}[]
\NormalTok{diamonds }\SpecialCharTok{|\textgreater{}}
  \FunctionTok{ggplot}\NormalTok{(}\FunctionTok{aes}\NormalTok{(}\AttributeTok{x =}\NormalTok{ price, }\AttributeTok{col =}\NormalTok{ cut)) }\SpecialCharTok{+}
  \FunctionTok{geom\_freqpoly}\NormalTok{() }\SpecialCharTok{+}
  \FunctionTok{labs}\NormalTok{(}
    \AttributeTok{x =} \StringTok{"Price"}\NormalTok{,}
    \AttributeTok{y =} \StringTok{"Count"}\NormalTok{,}
    \AttributeTok{col =} \StringTok{"Cut"}\NormalTok{)}
\end{Highlighting}
\end{Shaded}

\begin{verbatim}
`stat_bin()` using `bins = 30`. Pick better value with `binwidth`.
\end{verbatim}

\includegraphics{solutions_r4ds_eda_files/figure-pdf/unnamed-chunk-23-1.pdf}

Density plot, using \texttt{fill} instead of \texttt{col} and
\texttt{alpha\ =\ 0.3} to make the fill transparent. It is possible to
use \texttt{col} as well, but personally I like \texttt{fill} more when
plotting densities.

\begin{Shaded}
\begin{Highlighting}[]
\NormalTok{diamonds }\SpecialCharTok{|\textgreater{}}
  \FunctionTok{ggplot}\NormalTok{(}\FunctionTok{aes}\NormalTok{(}\AttributeTok{x =}\NormalTok{ price, }\AttributeTok{fill =}\NormalTok{ cut)) }\SpecialCharTok{+}
  \FunctionTok{geom\_density}\NormalTok{(}\AttributeTok{alpha =} \FloatTok{0.3}\NormalTok{) }\SpecialCharTok{+}
  \FunctionTok{labs}\NormalTok{(}
    \AttributeTok{x =} \StringTok{"Price"}\NormalTok{,}
    \AttributeTok{y =} \StringTok{"Density"}\NormalTok{,}
    \AttributeTok{col =} \StringTok{"Cut"}\NormalTok{)}
\end{Highlighting}
\end{Shaded}

\includegraphics{solutions_r4ds_eda_files/figure-pdf/unnamed-chunk-24-1.pdf}

\subsection{Exercise 6}\label{exercise-6-5}

From the \href{https://github.com/eclarke/ggbeeswarm}{GitHub page} of
\texttt{ggbeeswarm}:

\begin{quote}
``Beeswarm plots (aka column scatter plots or violin scatter plots) are
a way of plotting points that would ordinarily overlap so that they fall
next to each other instead. In addition to reducing overplotting, it
helps visualize the density of the data at each point (similar to a
violin plot), while still showing each data point individually.''
\end{quote}

Here we will take a look at the violin bee swarm plot

\begin{Shaded}
\begin{Highlighting}[]
\FunctionTok{library}\NormalTok{(ggbeeswarm)}
\NormalTok{diamonds }\SpecialCharTok{|\textgreater{}}
  \FunctionTok{ggplot}\NormalTok{(}\FunctionTok{aes}\NormalTok{(}\AttributeTok{x =}\NormalTok{ cut, }\AttributeTok{y =}\NormalTok{ price)) }\SpecialCharTok{+}
  \FunctionTok{geom\_quasirandom}\NormalTok{(}\AttributeTok{alpha =} \FloatTok{0.1}\NormalTok{) }\SpecialCharTok{+}
  \FunctionTok{labs}\NormalTok{(}
    \AttributeTok{x =} \StringTok{"Cut"}\NormalTok{,}
    \AttributeTok{y =} \StringTok{"Price"}\NormalTok{)}
\end{Highlighting}
\end{Shaded}

\includegraphics{solutions_r4ds_eda_files/figure-pdf/unnamed-chunk-25-1.pdf}

You can compare this to using a scatter plot with jittering, which one
do you think is more informative? Is the bee swarm plot better than the
violin plot?

\begin{Shaded}
\begin{Highlighting}[]
\FunctionTok{library}\NormalTok{(ggbeeswarm)}
\NormalTok{diamonds }\SpecialCharTok{|\textgreater{}}
  \FunctionTok{ggplot}\NormalTok{(}\FunctionTok{aes}\NormalTok{(}\AttributeTok{x =}\NormalTok{ cut, }\AttributeTok{y =}\NormalTok{ price)) }\SpecialCharTok{+}
  \FunctionTok{geom\_point}\NormalTok{(}\AttributeTok{alpha =} \FloatTok{0.1}\NormalTok{, }\AttributeTok{position =} \StringTok{"jitter"}\NormalTok{) }\SpecialCharTok{+}
  \FunctionTok{labs}\NormalTok{(}
    \AttributeTok{x =} \StringTok{"Cut"}\NormalTok{,}
    \AttributeTok{y =} \StringTok{"Price"}\NormalTok{)}
\end{Highlighting}
\end{Shaded}

\includegraphics{solutions_r4ds_eda_files/figure-pdf/unnamed-chunk-26-1.pdf}

\section{11.5.2 Covariation - Two categorical
variables}\label{covariation---two-categorical-variables}

\subsection{Exercise 1}\label{exercise-1-17}

We can express the counts in percentage of the max count

\begin{Shaded}
\begin{Highlighting}[]
\NormalTok{diamonds }\SpecialCharTok{|\textgreater{}} 
  \FunctionTok{count}\NormalTok{(color, cut) }\SpecialCharTok{|\textgreater{}} 
  \FunctionTok{mutate}\NormalTok{(}\AttributeTok{n =} \DecValTok{100} \SpecialCharTok{*}\NormalTok{ (n }\SpecialCharTok{/} \FunctionTok{max}\NormalTok{(n))) }\SpecialCharTok{|\textgreater{}} 
  \FunctionTok{ggplot}\NormalTok{(}\FunctionTok{aes}\NormalTok{(}\AttributeTok{x =}\NormalTok{ color, }\AttributeTok{y =}\NormalTok{ cut)) }\SpecialCharTok{+}
  \FunctionTok{geom\_tile}\NormalTok{(}\FunctionTok{aes}\NormalTok{(}\AttributeTok{fill =}\NormalTok{ n)) }\SpecialCharTok{+}
  \FunctionTok{labs}\NormalTok{(}
    \AttributeTok{x =} \StringTok{"Color"}\NormalTok{,}
    \AttributeTok{y =} \StringTok{"Cut"}\NormalTok{,}
    \AttributeTok{fill =} \StringTok{"Percent"}\NormalTok{)}
\end{Highlighting}
\end{Shaded}

\includegraphics{solutions_r4ds_eda_files/figure-pdf/unnamed-chunk-27-1.pdf}

\subsection{Exercise 2}\label{exercise-2-14}

Using a segmented bar chart gives us insight about how the distribution
of cuts are within each color.

\begin{Shaded}
\begin{Highlighting}[]
\NormalTok{diamonds }\SpecialCharTok{|\textgreater{}} 
  \FunctionTok{ggplot}\NormalTok{(}\FunctionTok{aes}\NormalTok{(}\AttributeTok{x =}\NormalTok{ color, }\AttributeTok{fill =}\NormalTok{ cut)) }\SpecialCharTok{+}
  \FunctionTok{geom\_bar}\NormalTok{() }\SpecialCharTok{+}
  \FunctionTok{labs}\NormalTok{(}
    \AttributeTok{x =} \StringTok{"Color"}\NormalTok{,}
    \AttributeTok{y =} \StringTok{"Count"}\NormalTok{,}
    \AttributeTok{fill =} \StringTok{"Cut"}\NormalTok{)}
\end{Highlighting}
\end{Shaded}

\includegraphics{solutions_r4ds_eda_files/figure-pdf/unnamed-chunk-28-1.pdf}

\subsection{Exercise 3}\label{exercise-3-12}

First we take a \texttt{glimpse} at the data to remember what we're
dealing with since we've been obsessing about diamonds now for a while.

\begin{Shaded}
\begin{Highlighting}[]
\FunctionTok{glimpse}\NormalTok{(nycflights13}\SpecialCharTok{::}\NormalTok{flights)}
\end{Highlighting}
\end{Shaded}

\begin{verbatim}
Rows: 336,776
Columns: 19
$ year           <int> 2013, 2013, 2013, 2013, 2013, 2013, 2013, 2013, 2013, 2~
$ month          <int> 1, 1, 1, 1, 1, 1, 1, 1, 1, 1, 1, 1, 1, 1, 1, 1, 1, 1, 1~
$ day            <int> 1, 1, 1, 1, 1, 1, 1, 1, 1, 1, 1, 1, 1, 1, 1, 1, 1, 1, 1~
$ dep_time       <int> 517, 533, 542, 544, 554, 554, 555, 557, 557, 558, 558, ~
$ sched_dep_time <int> 515, 529, 540, 545, 600, 558, 600, 600, 600, 600, 600, ~
$ dep_delay      <dbl> 2, 4, 2, -1, -6, -4, -5, -3, -3, -2, -2, -2, -2, -2, -1~
$ arr_time       <int> 830, 850, 923, 1004, 812, 740, 913, 709, 838, 753, 849,~
$ sched_arr_time <int> 819, 830, 850, 1022, 837, 728, 854, 723, 846, 745, 851,~
$ arr_delay      <dbl> 11, 20, 33, -18, -25, 12, 19, -14, -8, 8, -2, -3, 7, -1~
$ carrier        <chr> "UA", "UA", "AA", "B6", "DL", "UA", "B6", "EV", "B6", "~
$ flight         <int> 1545, 1714, 1141, 725, 461, 1696, 507, 5708, 79, 301, 4~
$ tailnum        <chr> "N14228", "N24211", "N619AA", "N804JB", "N668DN", "N394~
$ origin         <chr> "EWR", "LGA", "JFK", "JFK", "LGA", "EWR", "EWR", "LGA",~
$ dest           <chr> "IAH", "IAH", "MIA", "BQN", "ATL", "ORD", "FLL", "IAD",~
$ air_time       <dbl> 227, 227, 160, 183, 116, 150, 158, 53, 140, 138, 149, 1~
$ distance       <dbl> 1400, 1416, 1089, 1576, 762, 719, 1065, 229, 944, 733, ~
$ hour           <dbl> 5, 5, 5, 5, 6, 5, 6, 6, 6, 6, 6, 6, 6, 6, 6, 5, 6, 6, 6~
$ minute         <dbl> 15, 29, 40, 45, 0, 58, 0, 0, 0, 0, 0, 0, 0, 0, 0, 59, 0~
$ time_hour      <dttm> 2013-01-01 05:00:00, 2013-01-01 05:00:00, 2013-01-01 0~
\end{verbatim}

We summarize the \texttt{mean} departure delay by destination and month
and use \texttt{geom\_tile} to visualize it.

\begin{Shaded}
\begin{Highlighting}[]
\NormalTok{nycflights13}\SpecialCharTok{::}\NormalTok{flights }\SpecialCharTok{|\textgreater{}} 
  \FunctionTok{summarize}\NormalTok{(}
    \AttributeTok{avg\_dep\_delay =} \FunctionTok{mean}\NormalTok{(dep\_delay, }\AttributeTok{na.rm =} \ConstantTok{TRUE}\NormalTok{),}
    \AttributeTok{.by =} \FunctionTok{c}\NormalTok{(dest, month)) }\SpecialCharTok{|\textgreater{}} 
  \FunctionTok{mutate}\NormalTok{(}\AttributeTok{month =} \FunctionTok{as.factor}\NormalTok{(month)) }\SpecialCharTok{|\textgreater{}} 
  \FunctionTok{ggplot}\NormalTok{(}\FunctionTok{aes}\NormalTok{(}\AttributeTok{x =}\NormalTok{ dest, }\AttributeTok{y =}\NormalTok{ month)) }\SpecialCharTok{+}
  \FunctionTok{geom\_tile}\NormalTok{(}\FunctionTok{aes}\NormalTok{(}\AttributeTok{fill =}\NormalTok{ avg\_dep\_delay)) }\SpecialCharTok{+}
  \FunctionTok{labs}\NormalTok{(}
    \AttributeTok{x =} \StringTok{"Destination"}\NormalTok{,}
    \AttributeTok{y =} \StringTok{"Month"}\NormalTok{,}
    \AttributeTok{fill =} \StringTok{"Avg. dep. delay"}\NormalTok{)}
\end{Highlighting}
\end{Shaded}

\includegraphics{solutions_r4ds_eda_files/figure-pdf/unnamed-chunk-30-1.pdf}

We see that we have \emph{a lot} of destinations, making it hard to see
any interesting information. To better visualize we should select
destinations of interest. We can imagine that we are working for some
decision maker who wants to implement a policy to reduce the average
delays. To better identify what to do it may be useful to look at
destinations with little delay and destinations with long delays. We can
use the \texttt{geom\_quasirandom} plot from before with
\texttt{geom\_label\_repel} from \texttt{ggrepel} to identify extreme
values, while still visualizing the distributions.
\texttt{geom\_label\_repel} will only show the labels where there is not
too much overlap. This will help us identify outliers with both short
and long average delays.

\begin{Shaded}
\begin{Highlighting}[]
\FunctionTok{library}\NormalTok{(ggrepel)}
\NormalTok{nycflights13}\SpecialCharTok{::}\NormalTok{flights }\SpecialCharTok{|\textgreater{}} 
  \FunctionTok{summarize}\NormalTok{(}
    \AttributeTok{avg\_dep\_delay =} \FunctionTok{mean}\NormalTok{(dep\_delay, }\AttributeTok{na.rm =} \ConstantTok{TRUE}\NormalTok{),}
    \AttributeTok{.by =} \FunctionTok{c}\NormalTok{(dest, month)) }\SpecialCharTok{|\textgreater{}} 
  \FunctionTok{mutate}\NormalTok{(}\AttributeTok{month =} \FunctionTok{as.factor}\NormalTok{(month)) }\SpecialCharTok{|\textgreater{}} 
  \FunctionTok{ggplot}\NormalTok{(}\FunctionTok{aes}\NormalTok{(}\AttributeTok{x =}\NormalTok{ month, }\AttributeTok{y =}\NormalTok{ avg\_dep\_delay, }\AttributeTok{label =}\NormalTok{ dest)) }\SpecialCharTok{+}
  \FunctionTok{geom\_quasirandom}\NormalTok{(}\AttributeTok{alpha =} \FloatTok{0.5}\NormalTok{) }\SpecialCharTok{+}
  \FunctionTok{geom\_label\_repel}\NormalTok{() }\SpecialCharTok{+}
  \FunctionTok{labs}\NormalTok{(}
    \AttributeTok{x =} \StringTok{"Month"}\NormalTok{,}
    \AttributeTok{y =} \StringTok{"Avg dep. delay"}\NormalTok{)}
\end{Highlighting}
\end{Shaded}

\includegraphics{solutions_r4ds_eda_files/figure-pdf/unnamed-chunk-31-1.pdf}

Now we have an idea of what destinations have the shortest and longest
average delays for each month, and we can investigate further why.

There are many other ways to visualize this, a good choice may be to
filter the dataset before, looking at destinations with short/long
average delays only or isolating some months. It all depends on the
questions we want to answer, and what information is useful for that
answer.

\section{11.5.3 Covariation - Two numerical
variables}\label{covariation---two-numerical-variables}

We first create the \texttt{smaller} dataset from the book

\begin{Shaded}
\begin{Highlighting}[]
\NormalTok{smaller }\OtherTok{\textless{}{-}}\NormalTok{ diamonds }\SpecialCharTok{|\textgreater{}} \FunctionTok{filter}\NormalTok{(carat }\SpecialCharTok{\textless{}} \DecValTok{3}\NormalTok{)}
\end{Highlighting}
\end{Shaded}

\subsection{Exercise 1}\label{exercise-1-18}

First, we try to color by \texttt{cut\_width(carat,\ 0.25)}.

\begin{Shaded}
\begin{Highlighting}[]
\NormalTok{smaller }\SpecialCharTok{|\textgreater{}} 
  \FunctionTok{ggplot}\NormalTok{(}\FunctionTok{aes}\NormalTok{(}\AttributeTok{x =}\NormalTok{ price)) }\SpecialCharTok{+} 
  \FunctionTok{geom\_freqpoly}\NormalTok{(}\FunctionTok{aes}\NormalTok{(}\AttributeTok{color =} \FunctionTok{cut\_width}\NormalTok{(carat, }\FloatTok{0.25}\NormalTok{))) }\SpecialCharTok{+}
  \FunctionTok{labs}\NormalTok{(}
    \AttributeTok{x =} \StringTok{"Price"}\NormalTok{,}
    \AttributeTok{y =} \StringTok{"Count"}\NormalTok{)}
\end{Highlighting}
\end{Shaded}

\begin{verbatim}
`stat_bin()` using `bins = 30`. Pick better value with `binwidth`.
\end{verbatim}

\includegraphics{solutions_r4ds_eda_files/figure-pdf/unnamed-chunk-33-1.pdf}

This plot is not very informative as several groups have too low counts
to even show up in our plot. For this, it can be useful to use
\texttt{cut\_number}, that bins the data into groups that have roughly
the sa,e number of observations.

\begin{Shaded}
\begin{Highlighting}[]
\NormalTok{smaller }\SpecialCharTok{|\textgreater{}} 
  \FunctionTok{ggplot}\NormalTok{(}\FunctionTok{aes}\NormalTok{(}\AttributeTok{x =}\NormalTok{ price)) }\SpecialCharTok{+} 
  \FunctionTok{geom\_freqpoly}\NormalTok{(}\FunctionTok{aes}\NormalTok{(}\AttributeTok{color =} \FunctionTok{cut\_number}\NormalTok{(carat, }\DecValTok{5}\NormalTok{))) }\SpecialCharTok{+}
  \FunctionTok{labs}\NormalTok{(}
    \AttributeTok{x =} \StringTok{"Price"}\NormalTok{,}
    \AttributeTok{y =} \StringTok{"Count"}\NormalTok{)}
\end{Highlighting}
\end{Shaded}

\begin{verbatim}
`stat_bin()` using `bins = 30`. Pick better value with `binwidth`.
\end{verbatim}

\includegraphics{solutions_r4ds_eda_files/figure-pdf/unnamed-chunk-34-1.pdf}

It also allows us to specify the number of bins in a convenient way, not
having to think about the range of the variable.

\subsection{Exercise 2}\label{exercise-2-15}

Here we group by using \texttt{cut\_number} and plot the distribution of
carat using \texttt{geom\_density}.

\begin{Shaded}
\begin{Highlighting}[]
\NormalTok{smaller }\SpecialCharTok{|\textgreater{}} 
  \FunctionTok{ggplot}\NormalTok{(}\FunctionTok{aes}\NormalTok{(}\AttributeTok{x =}\NormalTok{ carat)) }\SpecialCharTok{+} 
  \FunctionTok{geom\_density}\NormalTok{(}
    \FunctionTok{aes}\NormalTok{(}\AttributeTok{fill =} \FunctionTok{cut\_number}\NormalTok{(price, }\DecValTok{3}\NormalTok{)),}
    \AttributeTok{alpha =} \FloatTok{0.3}\NormalTok{) }\SpecialCharTok{+}
  \FunctionTok{labs}\NormalTok{(}
    \AttributeTok{x =} \StringTok{"Carat"}\NormalTok{,}
    \AttributeTok{y =} \StringTok{"Density"}\NormalTok{)}
\end{Highlighting}
\end{Shaded}

\includegraphics{solutions_r4ds_eda_files/figure-pdf/unnamed-chunk-35-1.pdf}

\subsection{Exercise 3}\label{exercise-3-13}

Large diamonds, high values of carat, have a lot more variation than
small diamonds, which are very concentrated.

\begin{Shaded}
\begin{Highlighting}[]
\NormalTok{diamonds }\SpecialCharTok{|\textgreater{}} 
  \FunctionTok{ggplot}\NormalTok{(}\FunctionTok{aes}\NormalTok{(}\AttributeTok{x =}\NormalTok{ carat, }\AttributeTok{y =}\NormalTok{ price)) }\SpecialCharTok{+} 
  \FunctionTok{geom\_boxplot}\NormalTok{(}\FunctionTok{aes}\NormalTok{(}\AttributeTok{group =} \FunctionTok{cut\_width}\NormalTok{(carat, }\FloatTok{0.1}\NormalTok{))) }\SpecialCharTok{+}
  \FunctionTok{labs}\NormalTok{(}
    \AttributeTok{x =} \StringTok{"Carat"}\NormalTok{,}
    \AttributeTok{y =} \StringTok{"Price"}\NormalTok{)}
\end{Highlighting}
\end{Shaded}

\includegraphics{solutions_r4ds_eda_files/figure-pdf/unnamed-chunk-36-1.pdf}

\subsection{Exercise 4}\label{exercise-4-11}

\begin{Shaded}
\begin{Highlighting}[]
\NormalTok{smaller }\SpecialCharTok{|\textgreater{}} 
  \FunctionTok{ggplot}\NormalTok{(}\FunctionTok{aes}\NormalTok{(}\AttributeTok{x =}\NormalTok{ carat, }\AttributeTok{y =}\NormalTok{ price, }\AttributeTok{col =}\NormalTok{ cut)) }\SpecialCharTok{+}
  \FunctionTok{geom\_point}\NormalTok{(}\AttributeTok{alpha =} \FloatTok{0.1}\NormalTok{)}
\end{Highlighting}
\end{Shaded}

\includegraphics{solutions_r4ds_eda_files/figure-pdf/unnamed-chunk-37-1.pdf}

\subsection{Exercise 5}\label{exercise-5-6}

The scatter plot visualizes two variables of each observation
explicitly, letting the reader see interesting patterns of both
variables at the same time. In this plot we see that some diamonds are
cut asymmetrically, but most have the same x and y dimensions.
Interestingly the diamonds most off the diagonal have fine cuts
(Premium), so they must have some desirable symmetry. Remembering what
the \href{https://en.wikipedia.org/wiki/Diamond_cut}{cut} represents
explains this:

\begin{quote}
Cut does not refer to shape (pear, oval), but the symmetry,
proportioning and polish of a diamond.
\end{quote}

\begin{Shaded}
\begin{Highlighting}[]
\NormalTok{diamonds }\SpecialCharTok{|\textgreater{}} 
  \FunctionTok{filter}\NormalTok{(x }\SpecialCharTok{\textgreater{}=} \DecValTok{4}\NormalTok{) }\SpecialCharTok{|\textgreater{}} 
  \FunctionTok{ggplot}\NormalTok{(}\FunctionTok{aes}\NormalTok{(}\AttributeTok{x =}\NormalTok{ x, }\AttributeTok{y =}\NormalTok{ y, }\AttributeTok{col =}\NormalTok{ cut)) }\SpecialCharTok{+}
  \FunctionTok{geom\_point}\NormalTok{() }\SpecialCharTok{+}
  \FunctionTok{coord\_cartesian}\NormalTok{(}\AttributeTok{xlim =} \FunctionTok{c}\NormalTok{(}\DecValTok{4}\NormalTok{, }\DecValTok{11}\NormalTok{), }\AttributeTok{ylim =} \FunctionTok{c}\NormalTok{(}\DecValTok{4}\NormalTok{, }\DecValTok{11}\NormalTok{))}
\end{Highlighting}
\end{Shaded}

\includegraphics{solutions_r4ds_eda_files/figure-pdf/unnamed-chunk-38-1.pdf}

\subsection{Exercise 6}\label{exercise-6-6}

Dividing the bins in groups with roughly equal number of points can
inform us about the distribution of the x-variable, carat in this case.
The downside in this case is that some bins are very narrow, making them
hard to understand (see carat \textless{} 1). Further there is no
systematic or unique binning of the variables, making the plot less
transparent. It may be a better idea to visualize the distribution of
carat separately.

\begin{Shaded}
\begin{Highlighting}[]
\NormalTok{smaller }\SpecialCharTok{|\textgreater{}} 
  \FunctionTok{ggplot}\NormalTok{(}\FunctionTok{aes}\NormalTok{(}\AttributeTok{x =}\NormalTok{ carat, }\AttributeTok{y =}\NormalTok{ price)) }\SpecialCharTok{+} 
  \FunctionTok{geom\_boxplot}\NormalTok{(}\FunctionTok{aes}\NormalTok{(}\AttributeTok{group =} \FunctionTok{cut\_number}\NormalTok{(carat, }\DecValTok{20}\NormalTok{)))}
\end{Highlighting}
\end{Shaded}

\includegraphics{solutions_r4ds_eda_files/figure-pdf/unnamed-chunk-39-1.pdf}

Combining the boxplot with a density plot gives the reader an idea of
the distribution of carat and price.

\begin{Shaded}
\begin{Highlighting}[]
\NormalTok{p1 }\OtherTok{\textless{}{-}}\NormalTok{ smaller }\SpecialCharTok{|\textgreater{}} 
  \FunctionTok{ggplot}\NormalTok{(}\FunctionTok{aes}\NormalTok{(}\AttributeTok{x =}\NormalTok{ carat, }\AttributeTok{y =}\NormalTok{ price)) }\SpecialCharTok{+} 
  \FunctionTok{geom\_boxplot}\NormalTok{(}\FunctionTok{aes}\NormalTok{(}\AttributeTok{group =} \FunctionTok{cut\_width}\NormalTok{(carat, }\FloatTok{0.2}\NormalTok{))) }\SpecialCharTok{+}
  \FunctionTok{theme}\NormalTok{(}
    \AttributeTok{axis.title.x =} \FunctionTok{element\_blank}\NormalTok{(),}
    \AttributeTok{axis.text.x =} \FunctionTok{element\_blank}\NormalTok{(),}
    \AttributeTok{axis.ticks.x =} \FunctionTok{element\_blank}\NormalTok{()) }\SpecialCharTok{+}
  \FunctionTok{labs}\NormalTok{(}\AttributeTok{y =} \StringTok{"Price"}\NormalTok{)}
    

\NormalTok{p2 }\OtherTok{\textless{}{-}}\NormalTok{ smaller }\SpecialCharTok{|\textgreater{}} 
  \FunctionTok{ggplot}\NormalTok{(}\FunctionTok{aes}\NormalTok{(}\AttributeTok{x =}\NormalTok{ carat)) }\SpecialCharTok{+}
  \FunctionTok{geom\_density}\NormalTok{() }\SpecialCharTok{+}
  \FunctionTok{theme}\NormalTok{(}
    \AttributeTok{axis.title.y =} \FunctionTok{element\_blank}\NormalTok{(),}
    \AttributeTok{axis.text.y =} \FunctionTok{element\_blank}\NormalTok{(),}
    \AttributeTok{axis.ticks.y =} \FunctionTok{element\_blank}\NormalTok{()) }\SpecialCharTok{+}
  \FunctionTok{labs}\NormalTok{(}\AttributeTok{x =} \StringTok{"Carat"}\NormalTok{)}

\NormalTok{p1 }\SpecialCharTok{/}\NormalTok{ p2 }\SpecialCharTok{+} \FunctionTok{plot\_layout}\NormalTok{(}\AttributeTok{heights =} \FunctionTok{c}\NormalTok{(}\DecValTok{10}\NormalTok{, }\DecValTok{1}\NormalTok{))}
\end{Highlighting}
\end{Shaded}

\includegraphics{solutions_r4ds_eda_files/figure-pdf/unnamed-chunk-40-1.pdf}

We can also use a one-dimensional scatter plot to make a really
compressed visualization.

\begin{Shaded}
\begin{Highlighting}[]
\NormalTok{p1 }\OtherTok{\textless{}{-}}\NormalTok{ smaller }\SpecialCharTok{|\textgreater{}} 
  \FunctionTok{ggplot}\NormalTok{(}\FunctionTok{aes}\NormalTok{(}\AttributeTok{x =}\NormalTok{ carat, }\AttributeTok{y =}\NormalTok{ price)) }\SpecialCharTok{+} 
  \FunctionTok{geom\_boxplot}\NormalTok{(}\FunctionTok{aes}\NormalTok{(}\AttributeTok{group =} \FunctionTok{cut\_width}\NormalTok{(carat, }\FloatTok{0.2}\NormalTok{))) }\SpecialCharTok{+}
  \FunctionTok{theme}\NormalTok{(}
    \AttributeTok{axis.title.x =} \FunctionTok{element\_blank}\NormalTok{(),}
    \AttributeTok{axis.text.x =} \FunctionTok{element\_blank}\NormalTok{(),}
    \AttributeTok{axis.ticks.x =} \FunctionTok{element\_blank}\NormalTok{()) }\SpecialCharTok{+}
  \FunctionTok{labs}\NormalTok{(}\AttributeTok{y =} \StringTok{"Price"}\NormalTok{)}
    

\NormalTok{p2 }\OtherTok{\textless{}{-}}\NormalTok{ smaller }\SpecialCharTok{|\textgreater{}} 
  \FunctionTok{ggplot}\NormalTok{(}\FunctionTok{aes}\NormalTok{(}\AttributeTok{x =}\NormalTok{ carat, }\AttributeTok{y =} \DecValTok{0}\NormalTok{)) }\SpecialCharTok{+}
  \FunctionTok{geom\_point}\NormalTok{(}\AttributeTok{alpha =} \FloatTok{0.01}\NormalTok{) }\SpecialCharTok{+}
  \FunctionTok{theme}\NormalTok{(}
    \AttributeTok{axis.title.y =} \FunctionTok{element\_blank}\NormalTok{(),}
    \AttributeTok{axis.text.y =} \FunctionTok{element\_blank}\NormalTok{(),}
    \AttributeTok{axis.ticks.y =} \FunctionTok{element\_blank}\NormalTok{()) }\SpecialCharTok{+}
  \FunctionTok{labs}\NormalTok{(}\AttributeTok{x =} \StringTok{"Carat"}\NormalTok{)}

\NormalTok{p1 }\SpecialCharTok{/}\NormalTok{ p2 }\SpecialCharTok{+} \FunctionTok{plot\_layout}\NormalTok{(}\AttributeTok{heights =} \FunctionTok{c}\NormalTok{(}\DecValTok{10}\NormalTok{, }\DecValTok{1}\NormalTok{))}
\end{Highlighting}
\end{Shaded}

\includegraphics{solutions_r4ds_eda_files/figure-pdf/unnamed-chunk-41-1.pdf}

Which one do you prefer? Do you have any other ideas of how to visualize
two distributions at the same time?

\chapter{Hypothesis testing}\label{hypothesis-testing}

\section{Sample and population}\label{sample-and-population-1}

\subsection{Exercise 1}\label{exercise-1-19}

The help function tells us that the penguins are observed in the Palmer
Archipelago, Antarctica during 2007-2009. So the population is all
Adelie, Chinstrap, and Gentoo penguins in the Palmer Archipelago,
Antarctica during the years 2007-2009. Is it reasonable to believe that
this population is \emph{representative} of the current penguin
population? If we believe that it is, i.e.~that the measurements made
between 2007 and 2009 are still valid today, we can use our dataset to
say something about the population today.

\begin{Shaded}
\begin{Highlighting}[]
\FunctionTok{library}\NormalTok{(tidyverse)}
\FunctionTok{library}\NormalTok{(palmerpenguins)}
\NormalTok{?penguins}
\end{Highlighting}
\end{Shaded}

\section{Hypothesis testing
(proportions)}\label{hypothesis-testing-proportions-1}

First we need to recreate the data

\begin{Shaded}
\begin{Highlighting}[]
\FunctionTok{library}\NormalTok{(tidyverse)}

\NormalTok{recreate\_valjarbarometern }\OtherTok{\textless{}{-}} \ControlFlowTok{function}\NormalTok{(valjarbarometern\_aggregate, n)}
\NormalTok{\{}
  \CommentTok{\# Calculate number of observations of each party}
\NormalTok{  obs\_per\_party }\OtherTok{\textless{}{-}} \FunctionTok{round}\NormalTok{(n }\SpecialCharTok{*}\NormalTok{ valjarbarometern\_aggregate}\SpecialCharTok{$}\NormalTok{percentage }\SpecialCharTok{/} \DecValTok{100}\NormalTok{)}
  
  \CommentTok{\# Create a list of party observations}
\NormalTok{  party\_observations }\OtherTok{\textless{}{-}} \FunctionTok{mapply}\NormalTok{(rep, valjarbarometern\_aggregate}\SpecialCharTok{$}\NormalTok{party, obs\_per\_party)}
  
  \CommentTok{\# Flatten the list into a single vector}
\NormalTok{  sample\_data }\OtherTok{\textless{}{-}} \FunctionTok{unlist}\NormalTok{(party\_observations)}
  
  \CommentTok{\# Create a data frame with the samples}
\NormalTok{  valjarbarometern }\OtherTok{\textless{}{-}} \FunctionTok{tibble}\NormalTok{(}\AttributeTok{party =}\NormalTok{ sample\_data)}
  
  \FunctionTok{return}\NormalTok{(valjarbarometern)}
\NormalTok{\}}
\CommentTok{\# Number of people in survey}
\NormalTok{n }\OtherTok{=} \DecValTok{3072}

\NormalTok{parties }\OtherTok{\textless{}{-}} \FunctionTok{c}\NormalTok{(}\StringTok{"V"}\NormalTok{, }\StringTok{"S"}\NormalTok{, }\StringTok{"MP"}\NormalTok{, }\StringTok{"C"}\NormalTok{, }\StringTok{"L"}\NormalTok{, }\StringTok{"KD"}\NormalTok{, }\StringTok{"M"}\NormalTok{, }\StringTok{"SD"}\NormalTok{, }\StringTok{"Other"}\NormalTok{)}
\CommentTok{\# Aggregated sample}
\NormalTok{valjarbarometern\_aggregate }\OtherTok{\textless{}{-}} \FunctionTok{tibble}\NormalTok{(}
  \AttributeTok{party =} \FunctionTok{factor}\NormalTok{(parties, }\AttributeTok{levels =}\NormalTok{ parties), }\CommentTok{\# Using levels to preserve order}
  \AttributeTok{percentage =} \FunctionTok{c}\NormalTok{(}\FloatTok{7.6}\NormalTok{, }\FloatTok{37.6}\NormalTok{, }\FloatTok{4.3}\NormalTok{, }\FloatTok{3.9}\NormalTok{, }\DecValTok{3}\NormalTok{, }\FloatTok{3.1}\NormalTok{, }\DecValTok{19}\NormalTok{, }\FloatTok{19.7}\NormalTok{, }\FloatTok{1.7}\NormalTok{))}

\NormalTok{valjarbarometern }\OtherTok{\textless{}{-}} \FunctionTok{recreate\_valjarbarometern}\NormalTok{(}
\NormalTok{  valjarbarometern\_aggregate, n)}
\end{Highlighting}
\end{Shaded}

\subsection{Exercise 1}\label{exercise-1-20}

We simply replace MP with C:

\begin{quote}
\textbf{\(H_0\)}: The proportion of people who would vote for MP is
\emph{smaller} than 0.04.
\end{quote}

\begin{quote}
\textbf{\(H_a\)}: The proportion of people who would vote for MP is
\emph{larger} than 0.04.
\end{quote}

\subsection{Exercise 2}\label{exercise-2-16}

We do the test for the C party instead.

\begin{Shaded}
\begin{Highlighting}[]
\NormalTok{n\_participants }\OtherTok{\textless{}{-}} \DecValTok{3072} \CommentTok{\# Number of participants in survey}
\NormalTok{n\_c\_voters }\OtherTok{\textless{}{-}}\NormalTok{ valjarbarometern }\SpecialCharTok{|\textgreater{}} \CommentTok{\# Number of participants who answered C}
  \FunctionTok{filter}\NormalTok{(party }\SpecialCharTok{==} \StringTok{"C"}\NormalTok{) }\SpecialCharTok{|\textgreater{}}
  \FunctionTok{summarize}\NormalTok{(}\AttributeTok{n =} \FunctionTok{n}\NormalTok{()) }\SpecialCharTok{|\textgreater{}}
  \FunctionTok{pull}\NormalTok{(n)}
\NormalTok{p\_h0 }\OtherTok{\textless{}{-}} \FloatTok{0.04} \CommentTok{\# Our null hypothesis}
\CommentTok{\# Conducting hypothesis test of proportions}
\FunctionTok{prop.test}\NormalTok{(n\_c\_voters, n\_participants, p\_h0, }\AttributeTok{alternative =} \StringTok{"less"}\NormalTok{)}
\end{Highlighting}
\end{Shaded}

\begin{verbatim}

    1-sample proportions test with continuity correction

data:  n_c_voters out of n_participants, null probability p_h0
X-squared = 0.048018, df = 1, p-value = 0.4133
alternative hypothesis: true p is less than 0.04
95 percent confidence interval:
 0.0000000 0.0454036
sample estimates:
        p 
0.0390625 
\end{verbatim}

We can see that the result is \emph{not} statistically significant at
significance level 0.05.

\subsection{Exercise 3}\label{exercise-3-14}

We try one of the parties with the lowest percentage, KD.

\begin{Shaded}
\begin{Highlighting}[]
\CommentTok{\# Note: L would also be below 4\% with statistical significance}
\NormalTok{n\_participants }\OtherTok{\textless{}{-}} \DecValTok{3072} \CommentTok{\# Number of participants in survey}
\NormalTok{n\_kd\_voters }\OtherTok{\textless{}{-}}\NormalTok{ valjarbarometern }\SpecialCharTok{|\textgreater{}} \CommentTok{\# Number of participants who answered KD}
  \FunctionTok{filter}\NormalTok{(party }\SpecialCharTok{==} \StringTok{"KD"}\NormalTok{) }\SpecialCharTok{|\textgreater{}}
  \FunctionTok{summarize}\NormalTok{(}\AttributeTok{n =} \FunctionTok{n}\NormalTok{()) }\SpecialCharTok{|\textgreater{}}
  \FunctionTok{pull}\NormalTok{(n)}
\NormalTok{p\_h0 }\OtherTok{\textless{}{-}} \FloatTok{0.04} \CommentTok{\# Our null hypothesis}
\CommentTok{\# Conducting hypothesis test of proportions}
\FunctionTok{prop.test}\NormalTok{(n\_kd\_voters, n\_participants, p\_h0, }\AttributeTok{alternative =} \StringTok{"less"}\NormalTok{)}
\end{Highlighting}
\end{Shaded}

\begin{verbatim}

    1-sample proportions test with continuity correction

data:  n_kd_voters out of n_participants, null probability p_h0
X-squared = 6.355, df = 1, p-value = 0.005853
alternative hypothesis: true p is less than 0.04
95 percent confidence interval:
 0.0000000 0.0366646
sample estimates:
         p 
0.03092448 
\end{verbatim}

We see that the \(p\)-value is below 0.05, so the result is
statistically significant at significance level 0.05.

\section{Hypothesis test (mean)}\label{hypothesis-test-mean-1}

\subsection{Exercise 1}\label{exercise-1-21}

We write down the hypotheses with the population explicit:

\begin{quote}
\textbf{\(H_0\)}: The mean flipper length of all penguins in the Palmer
Archipelago during 2007-2009 was 200 mm.
\end{quote}

\begin{quote}
\textbf{\(H_a\)}: The mean flipper length of all penguins in the Palmer
Archipelago during 2007-2009 was \emph{greater than} 200 mm.
\end{quote}

We then perform our t-test.

\begin{Shaded}
\begin{Highlighting}[]
\FunctionTok{library}\NormalTok{(palmerpenguins)}
\FunctionTok{t.test}\NormalTok{(penguins}\SpecialCharTok{$}\NormalTok{flipper\_length\_mm, }\AttributeTok{mu =} \DecValTok{200}\NormalTok{, }\AttributeTok{alternative =} \StringTok{"greater"}\NormalTok{)}
\end{Highlighting}
\end{Shaded}

\begin{verbatim}

    One Sample t-test

data:  penguins$flipper_length_mm
t = 1.2036, df = 341, p-value = 0.1148
alternative hypothesis: true mean is greater than 200
95 percent confidence interval:
 199.6611      Inf
sample estimates:
mean of x 
 200.9152 
\end{verbatim}

Our test shows that we do not have evidence in the data that the mean
flipper length of all penguins in the Palmer Archipelago during
2007-2009 was not greater than 200 mm at significance level 0.05.

\subsection{Exercise 2}\label{exercise-2-17}

We first formulate our hypothesis, here we only do it for the test of
flipper length between Gentoo and Chinstrap penguins.

\begin{quote}
\textbf{\(H_0\)}: Gentoo and Chinstrap penguins have the same flipper
length.
\end{quote}

\begin{quote}
\textbf{\(H_a\)}: Gentoo and Chinstrap penguins \emph{do not} have the
same flipper length.
\end{quote}

We then perform all the tests.

\begin{Shaded}
\begin{Highlighting}[]
\NormalTok{gentoo\_flipper\_length }\OtherTok{\textless{}{-}}\NormalTok{ penguins }\SpecialCharTok{|\textgreater{}}
  \FunctionTok{filter}\NormalTok{(species }\SpecialCharTok{==} \StringTok{"Gentoo"}\NormalTok{) }\SpecialCharTok{|\textgreater{}}
  \FunctionTok{select}\NormalTok{(flipper\_length\_mm) }\SpecialCharTok{|\textgreater{}}
  \FunctionTok{drop\_na}\NormalTok{()}
\NormalTok{chinstrap\_flipper\_length }\OtherTok{\textless{}{-}}\NormalTok{ penguins }\SpecialCharTok{|\textgreater{}}
  \FunctionTok{filter}\NormalTok{(species }\SpecialCharTok{==} \StringTok{"Chinstrap"}\NormalTok{) }\SpecialCharTok{|\textgreater{}}
  \FunctionTok{select}\NormalTok{(flipper\_length\_mm) }\SpecialCharTok{|\textgreater{}}
  \FunctionTok{drop\_na}\NormalTok{()}
\NormalTok{adelie\_flipper\_length }\OtherTok{\textless{}{-}}\NormalTok{ penguins }\SpecialCharTok{|\textgreater{}}
  \FunctionTok{filter}\NormalTok{(species }\SpecialCharTok{==} \StringTok{"Adelie"}\NormalTok{) }\SpecialCharTok{|\textgreater{}}
  \FunctionTok{select}\NormalTok{(flipper\_length\_mm) }\SpecialCharTok{|\textgreater{}}
  \FunctionTok{drop\_na}\NormalTok{()}

\FunctionTok{t.test}\NormalTok{(gentoo\_flipper\_length, chinstrap\_flipper\_length)}
\end{Highlighting}
\end{Shaded}

\begin{verbatim}

    Welch Two Sample t-test

data:  gentoo_flipper_length and chinstrap_flipper_length
t = 20.463, df = 127.61, p-value < 2.2e-16
alternative hypothesis: true difference in means is not equal to 0
95 percent confidence interval:
 19.29770 23.42923
sample estimates:
mean of x mean of y 
 217.1870  195.8235 
\end{verbatim}

\begin{Shaded}
\begin{Highlighting}[]
\FunctionTok{t.test}\NormalTok{(gentoo\_flipper\_length, adelie\_flipper\_length)}
\end{Highlighting}
\end{Shaded}

\begin{verbatim}

    Welch Two Sample t-test

data:  gentoo_flipper_length and adelie_flipper_length
t = 34.445, df = 261.75, p-value < 2.2e-16
alternative hypothesis: true difference in means is not equal to 0
95 percent confidence interval:
 25.67652 28.79018
sample estimates:
mean of x mean of y 
 217.1870  189.9536 
\end{verbatim}

\begin{Shaded}
\begin{Highlighting}[]
\FunctionTok{t.test}\NormalTok{(chinstrap\_flipper\_length, adelie\_flipper\_length)}
\end{Highlighting}
\end{Shaded}

\begin{verbatim}

    Welch Two Sample t-test

data:  chinstrap_flipper_length and adelie_flipper_length
t = 5.7804, df = 119.68, p-value = 6.049e-08
alternative hypothesis: true difference in means is not equal to 0
95 percent confidence interval:
 3.859244 7.880530
sample estimates:
mean of x mean of y 
 195.8235  189.9536 
\end{verbatim}

We have evidence that the flipper lengths are different between all
species of penguins at significance level 0.05. Flipper length seem to
be a useful variable to distinguish between different penguin species.

\subsection{Exercise 3}\label{exercise-3-15}

We start by visualizing using a bar chart.

\begin{Shaded}
\begin{Highlighting}[]
\NormalTok{penguins }\SpecialCharTok{|\textgreater{}} 
  \FunctionTok{ggplot}\NormalTok{(}\FunctionTok{aes}\NormalTok{(}\AttributeTok{x =}\NormalTok{ island, }\AttributeTok{fill =}\NormalTok{ species)) }\SpecialCharTok{+} 
  \FunctionTok{geom\_bar}\NormalTok{()}
\end{Highlighting}
\end{Shaded}

\includegraphics{solutions_hypothesis_testing_files/figure-pdf/unnamed-chunk-6-1.pdf}

We see that there are a few more Chinstrap penguins than Adelie, but is
the difference significant? We first formulate the hypotheses:

\begin{quote}
\textbf{\(H_0\)}: The proportion of Chinstrap penguins at Dream island
is 0.5.
\end{quote}

\begin{quote}
\textbf{\(H_a\)}: The proportion of Chinstrap penguins at Dream island
is \emph{greater than} 0.5.
\end{quote}

Then we perform the proportions test using \texttt{prop.test} with the
number of chinstrap penguins and the total number of penguins at Dream
island, p = 0.5 (our null hypothesis), and alternative hypothesis
``greater''.

\begin{Shaded}
\begin{Highlighting}[]
\NormalTok{n\_total }\OtherTok{\textless{}{-}}\NormalTok{ penguins }\SpecialCharTok{|\textgreater{}} 
  \FunctionTok{filter}\NormalTok{(island }\SpecialCharTok{==} \StringTok{"Dream"}\NormalTok{) }\SpecialCharTok{|\textgreater{}}
  \FunctionTok{summarize}\NormalTok{(}\AttributeTok{n =} \FunctionTok{n}\NormalTok{()) }\SpecialCharTok{|\textgreater{}} 
  \FunctionTok{pull}\NormalTok{()}
  
\NormalTok{n\_chinstrap }\OtherTok{\textless{}{-}}\NormalTok{ penguins }\SpecialCharTok{|\textgreater{}} 
  \FunctionTok{filter}\NormalTok{(species }\SpecialCharTok{==} \StringTok{"Chinstrap"} \SpecialCharTok{\&}\NormalTok{ island }\SpecialCharTok{==} \StringTok{"Dream"}\NormalTok{) }\SpecialCharTok{|\textgreater{}}
  \FunctionTok{summarize}\NormalTok{(}\AttributeTok{n =} \FunctionTok{n}\NormalTok{()) }\SpecialCharTok{|\textgreater{}} 
  \FunctionTok{pull}\NormalTok{()}
  
\FunctionTok{prop.test}\NormalTok{(n\_chinstrap, n\_total, }\AttributeTok{p =} \FloatTok{0.5}\NormalTok{, }\AttributeTok{alternative =} \StringTok{"greater"}\NormalTok{)}
\end{Highlighting}
\end{Shaded}

\begin{verbatim}

    1-sample proportions test with continuity correction

data:  n_chinstrap out of n_total, null probability 0.5
X-squared = 0.97581, df = 1, p-value = 0.1616
alternative hypothesis: true p is greater than 0.5
95 percent confidence interval:
 0.4706263 1.0000000
sample estimates:
        p 
0.5483871 
\end{verbatim}

We see that the result is not significant. We do not have evidence in
the data suggesting that there are more Chinstrap penguins than Adelie
penguins at Dream island.

Can you think of any issue with this test and the sample? How can we be
sure that the researchers didn't try to observe the same amount of both
species? It is important to ask these questions about data collection
before you start your analysis.

\chapter{Linear regression}\label{linear-regression}

\section{Setup}\label{setup-1}

First we load the dataset and have a \texttt{glimpse} at it.

\begin{Shaded}
\begin{Highlighting}[]
\FunctionTok{library}\NormalTok{(tidyverse)}
\FunctionTok{library}\NormalTok{(palmerpenguins)}
\FunctionTok{glimpse}\NormalTok{(penguins)}
\end{Highlighting}
\end{Shaded}

\begin{verbatim}
Rows: 344
Columns: 8
$ species           <fct> Adelie, Adelie, Adelie, Adelie, Adelie, Adelie, Adel~
$ island            <fct> Torgersen, Torgersen, Torgersen, Torgersen, Torgerse~
$ bill_length_mm    <dbl> 39.1, 39.5, 40.3, NA, 36.7, 39.3, 38.9, 39.2, 34.1, ~
$ bill_depth_mm     <dbl> 18.7, 17.4, 18.0, NA, 19.3, 20.6, 17.8, 19.6, 18.1, ~
$ flipper_length_mm <int> 181, 186, 195, NA, 193, 190, 181, 195, 193, 190, 186~
$ body_mass_g       <int> 3750, 3800, 3250, NA, 3450, 3650, 3625, 4675, 3475, ~
$ sex               <fct> male, female, female, NA, female, male, female, male~
$ year              <int> 2007, 2007, 2007, 2007, 2007, 2007, 2007, 2007, 2007~
\end{verbatim}

\section{A linear relationship}\label{a-linear-relationship-1}

\subsection{Exercise 1}\label{exercise-1-22}

We see that the average body weight for a penguin with flipper length of
201 mm is 4206, an increase of 50 g from 200. Note we round the values
in the plot outputs.

\begin{Shaded}
\begin{Highlighting}[]
\NormalTok{x }\OtherTok{\textless{}{-}} \DecValTok{201}
\NormalTok{model }\OtherTok{\textless{}{-}}\NormalTok{ penguins }\SpecialCharTok{|\textgreater{}} 
  \FunctionTok{lm}\NormalTok{(}\AttributeTok{formula =}\NormalTok{ body\_mass\_g }\SpecialCharTok{\textasciitilde{}}\NormalTok{ flipper\_length\_mm)}
\NormalTok{prediction }\OtherTok{\textless{}{-}} \FunctionTok{data.frame}\NormalTok{(}
  \AttributeTok{flipper\_length\_mm =}\NormalTok{ x,}
  \AttributeTok{body\_mass\_g =}\NormalTok{ model }\SpecialCharTok{|\textgreater{}} \FunctionTok{predict}\NormalTok{(}\FunctionTok{data.frame}\NormalTok{(}\AttributeTok{flipper\_length\_mm =}\NormalTok{ x)))}

\NormalTok{penguins }\SpecialCharTok{|\textgreater{}} 
  \FunctionTok{ggplot}\NormalTok{(}\FunctionTok{aes}\NormalTok{(}\AttributeTok{x =}\NormalTok{ flipper\_length\_mm, }\AttributeTok{y =}\NormalTok{ body\_mass\_g)) }\SpecialCharTok{+} 
  \FunctionTok{geom\_point}\NormalTok{() }\SpecialCharTok{+} 
  \FunctionTok{geom\_smooth}\NormalTok{(}\AttributeTok{method =} \StringTok{"lm"}\NormalTok{, }\AttributeTok{se =} \ConstantTok{FALSE}\NormalTok{) }\SpecialCharTok{+} 
  \FunctionTok{geom\_segment}\NormalTok{(}\AttributeTok{x =}\NormalTok{ prediction}\SpecialCharTok{$}\NormalTok{flipper\_length\_mm, }\AttributeTok{xend =}\NormalTok{ prediction}\SpecialCharTok{$}\NormalTok{flipper\_length\_mm,  }\AttributeTok{y =} \DecValTok{0}\NormalTok{, }\AttributeTok{yend =}\NormalTok{ prediction}\SpecialCharTok{$}\NormalTok{body\_mass\_g, }\AttributeTok{color =} \StringTok{"black"}\NormalTok{, }\AttributeTok{linetype =} \DecValTok{2}\NormalTok{, }\AttributeTok{linewidth =} \FloatTok{0.2}\NormalTok{) }\SpecialCharTok{+}
  \FunctionTok{geom\_segment}\NormalTok{(}\AttributeTok{x =} \DecValTok{0}\NormalTok{, }\AttributeTok{xend =}\NormalTok{ prediction}\SpecialCharTok{$}\NormalTok{flipper\_length\_mm,  }\AttributeTok{y =}\NormalTok{ prediction}\SpecialCharTok{$}\NormalTok{body\_mass\_g, }\AttributeTok{yend =}\NormalTok{ prediction}\SpecialCharTok{$}\NormalTok{body\_mass\_g, }\AttributeTok{color =} \StringTok{"black"}\NormalTok{, }\AttributeTok{linetype =} \DecValTok{2}\NormalTok{, }\AttributeTok{linewidth =} \FloatTok{0.2}\NormalTok{) }\SpecialCharTok{+}
  \FunctionTok{scale\_x\_continuous}\NormalTok{(}\AttributeTok{breaks =} \FunctionTok{c}\NormalTok{(prediction}\SpecialCharTok{$}\NormalTok{flipper\_length\_mm)) }\SpecialCharTok{+}
  \FunctionTok{scale\_y\_continuous}\NormalTok{(}\AttributeTok{breaks =} \FunctionTok{round}\NormalTok{(}\FunctionTok{c}\NormalTok{(prediction}\SpecialCharTok{$}\NormalTok{body\_mass\_g))) }\SpecialCharTok{+}
  \FunctionTok{geom\_point}\NormalTok{(}\AttributeTok{data =}\NormalTok{ prediction, }\AttributeTok{col =} \StringTok{"red"}\NormalTok{, }\AttributeTok{size =} \DecValTok{3}\NormalTok{) }\SpecialCharTok{+} 
  \FunctionTok{labs}\NormalTok{(}
    \AttributeTok{x =} \StringTok{"Flipper length (mm)"}\NormalTok{, }
    \AttributeTok{y =} \StringTok{"Body mass (g)"}\NormalTok{)}
\end{Highlighting}
\end{Shaded}

\includegraphics{solutions_linear_regression_files/figure-pdf/unnamed-chunk-2-1.pdf}

\subsection{Exercise 2}\label{exercise-2-18}

The predicted average body mass is -812 g, which doesn't seem to make
any sense! When we are predicting based on values that are outside of
our sample we may get poor predictions, as in this case. \textbf{Always
be mindful when predicting outside your sample.}

\begin{Shaded}
\begin{Highlighting}[]
\NormalTok{x }\OtherTok{\textless{}{-}} \DecValTok{100}
\NormalTok{model }\OtherTok{\textless{}{-}}\NormalTok{ penguins }\SpecialCharTok{|\textgreater{}} 
  \FunctionTok{lm}\NormalTok{(}\AttributeTok{formula =}\NormalTok{ body\_mass\_g }\SpecialCharTok{\textasciitilde{}}\NormalTok{ flipper\_length\_mm)}
\NormalTok{prediction }\OtherTok{\textless{}{-}} \FunctionTok{data.frame}\NormalTok{(}
  \AttributeTok{flipper\_length\_mm =}\NormalTok{ x,}
  \AttributeTok{body\_mass\_g =}\NormalTok{ model }\SpecialCharTok{|\textgreater{}} \FunctionTok{predict}\NormalTok{(}\FunctionTok{data.frame}\NormalTok{(}\AttributeTok{flipper\_length\_mm =}\NormalTok{ x)))}

\NormalTok{penguins }\SpecialCharTok{|\textgreater{}} 
  \FunctionTok{ggplot}\NormalTok{(}\FunctionTok{aes}\NormalTok{(}\AttributeTok{x =}\NormalTok{ flipper\_length\_mm, }\AttributeTok{y =}\NormalTok{ body\_mass\_g)) }\SpecialCharTok{+} 
  \FunctionTok{geom\_point}\NormalTok{() }\SpecialCharTok{+} 
  \FunctionTok{geom\_smooth}\NormalTok{(}\AttributeTok{method =} \StringTok{"lm"}\NormalTok{, }\AttributeTok{se =} \ConstantTok{FALSE}\NormalTok{) }\SpecialCharTok{+} 
  \FunctionTok{geom\_segment}\NormalTok{(}\AttributeTok{x =}\NormalTok{ prediction}\SpecialCharTok{$}\NormalTok{flipper\_length\_mm, }\AttributeTok{xend =}\NormalTok{ prediction}\SpecialCharTok{$}\NormalTok{flipper\_length\_mm,  }\AttributeTok{y =} \DecValTok{0}\NormalTok{, }\AttributeTok{yend =}\NormalTok{ prediction}\SpecialCharTok{$}\NormalTok{body\_mass\_g, }\AttributeTok{color =} \StringTok{"black"}\NormalTok{, }\AttributeTok{linetype =} \DecValTok{2}\NormalTok{, }\AttributeTok{linewidth =} \FloatTok{0.2}\NormalTok{) }\SpecialCharTok{+}
  \FunctionTok{geom\_segment}\NormalTok{(}\AttributeTok{x =} \DecValTok{0}\NormalTok{, }\AttributeTok{xend =}\NormalTok{ prediction}\SpecialCharTok{$}\NormalTok{flipper\_length\_mm,  }\AttributeTok{y =}\NormalTok{ prediction}\SpecialCharTok{$}\NormalTok{body\_mass\_g, }\AttributeTok{yend =}\NormalTok{ prediction}\SpecialCharTok{$}\NormalTok{body\_mass\_g, }\AttributeTok{color =} \StringTok{"black"}\NormalTok{, }\AttributeTok{linetype =} \DecValTok{2}\NormalTok{, }\AttributeTok{linewidth =} \FloatTok{0.2}\NormalTok{) }\SpecialCharTok{+}
  \FunctionTok{scale\_x\_continuous}\NormalTok{(}\AttributeTok{breaks =} \FunctionTok{c}\NormalTok{(prediction}\SpecialCharTok{$}\NormalTok{flipper\_length\_mm)) }\SpecialCharTok{+}
  \FunctionTok{scale\_y\_continuous}\NormalTok{(}\AttributeTok{breaks =} \FunctionTok{round}\NormalTok{(}\FunctionTok{c}\NormalTok{(prediction}\SpecialCharTok{$}\NormalTok{body\_mass\_g))) }\SpecialCharTok{+}
  \FunctionTok{geom\_point}\NormalTok{(}\AttributeTok{data =}\NormalTok{ prediction, }\AttributeTok{col =} \StringTok{"red"}\NormalTok{, }\AttributeTok{size =} \DecValTok{3}\NormalTok{) }\SpecialCharTok{+} 
  \FunctionTok{labs}\NormalTok{(}
    \AttributeTok{x =} \StringTok{"Flipper length (mm)"}\NormalTok{, }
    \AttributeTok{y =} \StringTok{"Body mass (g)"}\NormalTok{)}
\end{Highlighting}
\end{Shaded}

\includegraphics{solutions_linear_regression_files/figure-pdf/unnamed-chunk-3-1.pdf}

\section{Fitting a linear model}\label{fitting-a-linear-model-1}

\subsection{Exercise 1}\label{exercise-1-23}

\begin{Shaded}
\begin{Highlighting}[]
\NormalTok{model }\OtherTok{\textless{}{-}}\NormalTok{ penguins }\SpecialCharTok{|\textgreater{}} 
  \FunctionTok{lm}\NormalTok{(}\AttributeTok{formula =}\NormalTok{ flipper\_length\_mm }\SpecialCharTok{\textasciitilde{}}\NormalTok{ body\_mass\_g)}
\FunctionTok{summary}\NormalTok{(model)}
\end{Highlighting}
\end{Shaded}

\begin{verbatim}

Call:
lm(formula = flipper_length_mm ~ body_mass_g, data = penguins)

Residuals:
     Min       1Q   Median       3Q      Max 
-23.7626  -4.9138   0.9891   5.1166  16.6392 

Coefficients:
             Estimate Std. Error t value Pr(>|t|)    
(Intercept) 1.367e+02  1.997e+00   68.47   <2e-16 ***
body_mass_g 1.528e-02  4.668e-04   32.72   <2e-16 ***
---
Signif. codes:  0 '***' 0.001 '**' 0.01 '*' 0.05 '.' 0.1 ' ' 1

Residual standard error: 6.913 on 340 degrees of freedom
  (2 observations deleted due to missingness)
Multiple R-squared:  0.759, Adjusted R-squared:  0.7583 
F-statistic:  1071 on 1 and 340 DF,  p-value: < 2.2e-16
\end{verbatim}

The linear model doesn't care which variable is the dependent and
independent variable, it will fit a model either way. The coefficients
will be very different, as the increase in body mass is much smaller per
mm flipper length than the other way around, but the relationship will
still be significant. The coefficient is 1.528e-02 which is the same as
\(1.528 \cdot 10^{-2} = 0.01528\). Thus the flipper length increases on
average by 0.015 cm per g the body mass increases. We can never draw any
causal conclusions without additional assumptions, and in this case it
doesn't seem appropriate, as it seems more like a ``chicken and egg''
dilemma than a causal effect. The model can still be useful for
predicting and to understand the relationship between the variables,
without suggesting a causal relationship.

\begin{Shaded}
\begin{Highlighting}[]
\NormalTok{penguins }\SpecialCharTok{|\textgreater{}} 
  \FunctionTok{ggplot}\NormalTok{(}\FunctionTok{aes}\NormalTok{(}\AttributeTok{x =}\NormalTok{ body\_mass\_g, }\AttributeTok{y =}\NormalTok{ flipper\_length\_mm)) }\SpecialCharTok{+} 
  \FunctionTok{geom\_point}\NormalTok{() }\SpecialCharTok{+} 
  \FunctionTok{geom\_smooth}\NormalTok{(}\AttributeTok{method =} \StringTok{"lm"}\NormalTok{, }\AttributeTok{se =} \ConstantTok{FALSE}\NormalTok{) }\SpecialCharTok{+} 
  \FunctionTok{labs}\NormalTok{(}
    \AttributeTok{x =} \StringTok{"Flipper length (mm)"}\NormalTok{, }
    \AttributeTok{y =} \StringTok{"Body mass (g)"}\NormalTok{)}
\end{Highlighting}
\end{Shaded}

\includegraphics{solutions_linear_regression_files/figure-pdf/unnamed-chunk-5-1.pdf}

\subsection{Exercise 2}\label{exercise-2-19}

From the visualization we see that the sales increase during the summer
months and decreases during winter, but the line is completely flat.
There is a clear relationship, but it is not a linear relationship! We
can see from how the data is generated that it depends on the
\texttt{sin} function, which is not linear.

\begin{Shaded}
\begin{Highlighting}[]
\CommentTok{\# Simulating data}
\NormalTok{day }\OtherTok{\textless{}{-}} \FunctionTok{seq}\NormalTok{(}\DecValTok{0}\NormalTok{, }\DecValTok{12}\NormalTok{, }\AttributeTok{length.out =} \DecValTok{356}\NormalTok{)}
\NormalTok{sales }\OtherTok{\textless{}{-}} \DecValTok{400} \SpecialCharTok{*} \FunctionTok{sin}\NormalTok{(pi }\SpecialCharTok{*}\NormalTok{ day }\SpecialCharTok{/} \DecValTok{12}\NormalTok{)}
\NormalTok{ice\_cream }\OtherTok{\textless{}{-}} \FunctionTok{tibble}\NormalTok{(}\AttributeTok{day =}\NormalTok{ day, }\AttributeTok{sales =}\NormalTok{ sales)}

\NormalTok{ice\_cream }\SpecialCharTok{|\textgreater{}} 
  \FunctionTok{ggplot}\NormalTok{(}\FunctionTok{aes}\NormalTok{(}\AttributeTok{x =}\NormalTok{ day, }\AttributeTok{y =}\NormalTok{ sales)) }\SpecialCharTok{+} 
  \FunctionTok{geom\_point}\NormalTok{() }\SpecialCharTok{+} 
  \FunctionTok{geom\_smooth}\NormalTok{(}\AttributeTok{method =} \StringTok{"lm"}\NormalTok{) }\SpecialCharTok{+} 
  \FunctionTok{scale\_x\_continuous}\NormalTok{(}
    \AttributeTok{breaks =} \FunctionTok{seq}\NormalTok{(}\DecValTok{0}\NormalTok{, }\DecValTok{12}\NormalTok{)) }\SpecialCharTok{+}
  \FunctionTok{labs}\NormalTok{(}
    \AttributeTok{x =} \StringTok{"Month"}\NormalTok{,}
    \AttributeTok{y =} \StringTok{"Ice cream sales"}\NormalTok{)}
\end{Highlighting}
\end{Shaded}

\begin{verbatim}
`geom_smooth()` using formula = 'y ~ x'
\end{verbatim}

\includegraphics{solutions_linear_regression_files/figure-pdf/unnamed-chunk-6-1.pdf}

We see that the \(p\)-value is 1, there is no significant relationship.
It is important to remember that we are looking for a \emph{linear}
relationship when performing linear regression. It will not detect any
other relationship.

\begin{Shaded}
\begin{Highlighting}[]
\NormalTok{ice\_cream }\SpecialCharTok{|\textgreater{}} 
  \FunctionTok{lm}\NormalTok{(}\AttributeTok{formula =}\NormalTok{ sales }\SpecialCharTok{\textasciitilde{}}\NormalTok{ day) }\SpecialCharTok{|\textgreater{}} 
  \FunctionTok{summary}\NormalTok{()}
\end{Highlighting}
\end{Shaded}

\begin{verbatim}

Call:
lm(formula = sales ~ day, data = ice_cream)

Residuals:
    Min      1Q  Median      3Q     Max 
-253.93 -102.08   28.28  115.45  146.06 

Coefficients:
             Estimate Std. Error t value Pr(>|t|)    
(Intercept) 2.539e+02  1.312e+01   19.36   <2e-16 ***
day         1.821e-14  1.892e+00    0.00        1    
---
Signif. codes:  0 '***' 0.001 '**' 0.01 '*' 0.05 '.' 0.1 ' ' 1

Residual standard error: 124 on 354 degrees of freedom
Multiple R-squared:  1.039e-30, Adjusted R-squared:  -0.002825 
F-statistic: 3.679e-28 on 1 and 354 DF,  p-value: 1
\end{verbatim}

\subsection{Exercise 3}\label{exercise-3-16}

We see that we get one coefficient for Chinstrap and one for Gentoo. The
value, under ``Estimate'' is the difference from the Adelie penguins in
average body weight, compare with the values computed below.

\begin{Shaded}
\begin{Highlighting}[]
\NormalTok{penguins }\SpecialCharTok{|\textgreater{}} 
  \FunctionTok{lm}\NormalTok{(}\AttributeTok{formula =}\NormalTok{ body\_mass\_g }\SpecialCharTok{\textasciitilde{}}\NormalTok{ species) }\SpecialCharTok{|\textgreater{}}
  \FunctionTok{summary}\NormalTok{()}
\end{Highlighting}
\end{Shaded}

\begin{verbatim}

Call:
lm(formula = body_mass_g ~ species, data = penguins)

Residuals:
     Min       1Q   Median       3Q      Max 
-1126.02  -333.09   -33.09   316.91  1223.98 

Coefficients:
                 Estimate Std. Error t value Pr(>|t|)    
(Intercept)       3700.66      37.62   98.37   <2e-16 ***
speciesChinstrap    32.43      67.51    0.48    0.631    
speciesGentoo     1375.35      56.15   24.50   <2e-16 ***
---
Signif. codes:  0 '***' 0.001 '**' 0.01 '*' 0.05 '.' 0.1 ' ' 1

Residual standard error: 462.3 on 339 degrees of freedom
  (2 observations deleted due to missingness)
Multiple R-squared:  0.6697,    Adjusted R-squared:  0.6677 
F-statistic: 343.6 on 2 and 339 DF,  p-value: < 2.2e-16
\end{verbatim}

We can check this

\begin{Shaded}
\begin{Highlighting}[]
\CommentTok{\# Calculating mean body weight for each species}
\NormalTok{adelie\_mean }\OtherTok{\textless{}{-}}\NormalTok{ penguins }\SpecialCharTok{|\textgreater{}} 
  \FunctionTok{filter}\NormalTok{(species }\SpecialCharTok{==} \StringTok{"Adelie"}\NormalTok{) }\SpecialCharTok{|\textgreater{}} 
  \FunctionTok{pull}\NormalTok{(body\_mass\_g) }\SpecialCharTok{|\textgreater{}}
  \FunctionTok{mean}\NormalTok{(}\AttributeTok{na.rm =} \ConstantTok{TRUE}\NormalTok{)}
\NormalTok{chinstrap\_mean }\OtherTok{\textless{}{-}}\NormalTok{ penguins }\SpecialCharTok{|\textgreater{}} 
  \FunctionTok{filter}\NormalTok{(species }\SpecialCharTok{==} \StringTok{"Chinstrap"}\NormalTok{) }\SpecialCharTok{|\textgreater{}} 
  \FunctionTok{pull}\NormalTok{(body\_mass\_g) }\SpecialCharTok{|\textgreater{}}
  \FunctionTok{mean}\NormalTok{(}\AttributeTok{na.rm =} \ConstantTok{TRUE}\NormalTok{)}
\NormalTok{gentoo\_mean }\OtherTok{\textless{}{-}}\NormalTok{ penguins }\SpecialCharTok{|\textgreater{}} 
  \FunctionTok{filter}\NormalTok{(species }\SpecialCharTok{==} \StringTok{"Gentoo"}\NormalTok{) }\SpecialCharTok{|\textgreater{}} 
  \FunctionTok{pull}\NormalTok{(body\_mass\_g) }\SpecialCharTok{|\textgreater{}}
  \FunctionTok{mean}\NormalTok{(}\AttributeTok{na.rm =} \ConstantTok{TRUE}\NormalTok{)}

 \CommentTok{\# Differences in mean}
\NormalTok{chinstrap\_mean }\SpecialCharTok{{-}}\NormalTok{ adelie\_mean}
\end{Highlighting}
\end{Shaded}

\begin{verbatim}
[1] 32.42598
\end{verbatim}

\begin{Shaded}
\begin{Highlighting}[]
\NormalTok{gentoo\_mean }\SpecialCharTok{{-}}\NormalTok{ adelie\_mean}
\end{Highlighting}
\end{Shaded}

\begin{verbatim}
[1] 1375.354
\end{verbatim}

In addition to comparing the means the linear regression outputs
\(p\)-values that tests if there is a significant difference between the
species. Only the difference between Adelie and Gentoos is significant.

\subsection{Exercise 4}\label{exercise-4-12}

We see that the \(p\)-values are the same and we can use linear
regression instead of the t-test. Linear regression is more versatile
and can accommodate many levels in the categorical predictor. If it is
possible to use a t-test it may be more straightforward and clear
though.

\begin{Shaded}
\begin{Highlighting}[]
\CommentTok{\# t{-}test}
\NormalTok{female\_gentoo\_body\_mass }\OtherTok{\textless{}{-}}\NormalTok{ penguins }\SpecialCharTok{|\textgreater{}}
  \FunctionTok{filter}\NormalTok{(species }\SpecialCharTok{==} \StringTok{"Gentoo"} \SpecialCharTok{\&}\NormalTok{ sex }\SpecialCharTok{==} \StringTok{"female"}\NormalTok{) }\SpecialCharTok{|\textgreater{}}
  \FunctionTok{select}\NormalTok{(body\_mass\_g) }\SpecialCharTok{|\textgreater{}}
  \FunctionTok{drop\_na}\NormalTok{()}
\NormalTok{male\_gentoo\_body\_mass }\OtherTok{\textless{}{-}}\NormalTok{ penguins }\SpecialCharTok{|\textgreater{}}
  \FunctionTok{filter}\NormalTok{(species }\SpecialCharTok{==} \StringTok{"Gentoo"} \SpecialCharTok{\&}\NormalTok{ sex }\SpecialCharTok{==} \StringTok{"male"}\NormalTok{) }\SpecialCharTok{|\textgreater{}}
  \FunctionTok{select}\NormalTok{(body\_mass\_g) }\SpecialCharTok{|\textgreater{}}
  \FunctionTok{drop\_na}\NormalTok{()}

\FunctionTok{t.test}\NormalTok{(male\_gentoo\_body\_mass, female\_gentoo\_body\_mass)}
\end{Highlighting}
\end{Shaded}

\begin{verbatim}

    Welch Two Sample t-test

data:  male_gentoo_body_mass and female_gentoo_body_mass
t = 14.761, df = 116.64, p-value < 2.2e-16
alternative hypothesis: true difference in means is not equal to 0
95 percent confidence interval:
 697.0763 913.1130
sample estimates:
mean of x mean of y 
 5484.836  4679.741 
\end{verbatim}

\begin{Shaded}
\begin{Highlighting}[]
\CommentTok{\# Linear regression}
\NormalTok{penguins }\SpecialCharTok{|\textgreater{}} 
  \FunctionTok{filter}\NormalTok{(species }\SpecialCharTok{==} \StringTok{"Gentoo"}\NormalTok{) }\SpecialCharTok{|\textgreater{}}
  \FunctionTok{lm}\NormalTok{(}\AttributeTok{formula =}\NormalTok{ body\_mass\_g }\SpecialCharTok{\textasciitilde{}}\NormalTok{ sex) }\SpecialCharTok{|\textgreater{}}
  \FunctionTok{summary}\NormalTok{()}
\end{Highlighting}
\end{Shaded}

\begin{verbatim}

Call:
lm(formula = body_mass_g ~ sex, data = filter(penguins, species == 
    "Gentoo"))

Residuals:
    Min      1Q  Median      3Q     Max 
-734.84 -207.29   20.26  215.16  815.16 

Coefficients:
            Estimate Std. Error t value Pr(>|t|)    
(Intercept)  4679.74      39.15  119.52   <2e-16 ***
sexmale       805.09      54.69   14.72   <2e-16 ***
---
Signif. codes:  0 '***' 0.001 '**' 0.01 '*' 0.05 '.' 0.1 ' ' 1

Residual standard error: 298.2 on 117 degrees of freedom
  (5 observations deleted due to missingness)
Multiple R-squared:  0.6494,    Adjusted R-squared:  0.6464 
F-statistic: 216.7 on 1 and 117 DF,  p-value: < 2.2e-16
\end{verbatim}

\section{Categorical predictors}\label{categorical-predictors-1}

\subsection{Exercise 1}\label{exercise-1-24}

We start with the linear model.

\begin{Shaded}
\begin{Highlighting}[]
\NormalTok{penguins }\SpecialCharTok{|\textgreater{}} 
  \FunctionTok{lm}\NormalTok{(}\AttributeTok{formula =}\NormalTok{ body\_mass\_g }\SpecialCharTok{\textasciitilde{}}\NormalTok{ island) }\SpecialCharTok{|\textgreater{}} 
  \FunctionTok{summary}\NormalTok{()}
\end{Highlighting}
\end{Shaded}

\begin{verbatim}

Call:
lm(formula = body_mass_g ~ island, data = penguins)

Residuals:
     Min       1Q   Median       3Q      Max 
-1866.02  -362.90    -6.37   433.98  1583.98 

Coefficients:
                Estimate Std. Error t value Pr(>|t|)    
(Intercept)      4716.02      48.47   97.30   <2e-16 ***
islandDream     -1003.11      74.25  -13.51   <2e-16 ***
islandTorgersen -1009.65     100.21  -10.08   <2e-16 ***
---
Signif. codes:  0 '***' 0.001 '**' 0.01 '*' 0.05 '.' 0.1 ' ' 1

Residual standard error: 626.3 on 339 degrees of freedom
  (2 observations deleted due to missingness)
Multiple R-squared:  0.3936,    Adjusted R-squared:   0.39 
F-statistic:   110 on 2 and 339 DF,  p-value: < 2.2e-16
\end{verbatim}

We see that the coefficients for both Dream and Torgersen are
significant, but let's have a look at a density plot to better
understand the result.

\begin{Shaded}
\begin{Highlighting}[]
\NormalTok{mean\_bws }\OtherTok{\textless{}{-}}\NormalTok{ penguins }\SpecialCharTok{|\textgreater{}} 
  \FunctionTok{summarize}\NormalTok{(}
    \AttributeTok{body\_mass\_g =} \FunctionTok{mean}\NormalTok{(body\_mass\_g, }\AttributeTok{na.rm =} \ConstantTok{TRUE}\NormalTok{),}
    \AttributeTok{.by =}\NormalTok{ island)}
\NormalTok{penguins }\SpecialCharTok{|\textgreater{}} 
  \FunctionTok{ggplot}\NormalTok{(}\FunctionTok{aes}\NormalTok{(}\AttributeTok{x =}\NormalTok{ body\_mass\_g, }\AttributeTok{fill =}\NormalTok{ island)) }\SpecialCharTok{+} 
  \FunctionTok{geom\_density}\NormalTok{(}\AttributeTok{alpha =} \FloatTok{0.3}\NormalTok{) }\SpecialCharTok{+}
  \FunctionTok{geom\_vline}\NormalTok{(}
    \FunctionTok{aes}\NormalTok{(}
      \AttributeTok{xintercept =}\NormalTok{ body\_mass\_g,}
      \AttributeTok{col =}\NormalTok{ island),}
    \AttributeTok{data =}\NormalTok{ mean\_bws) }\SpecialCharTok{+}
  \FunctionTok{labs}\NormalTok{(}
    \AttributeTok{x =} \StringTok{"Body mass (g)"}\NormalTok{,}
    \AttributeTok{y =} \StringTok{""}\NormalTok{,}
    \AttributeTok{col =} \StringTok{"Species"}\NormalTok{,}
    \AttributeTok{fill =} \StringTok{"Species"}\NormalTok{) }\SpecialCharTok{+}
  \FunctionTok{theme}\NormalTok{(}\AttributeTok{axis.text.y =} \FunctionTok{element\_blank}\NormalTok{(),}
        \AttributeTok{axis.ticks.y =} \FunctionTok{element\_blank}\NormalTok{())}
\end{Highlighting}
\end{Shaded}

\includegraphics{solutions_linear_regression_files/figure-pdf/unnamed-chunk-12-1.pdf}

We see that \texttt{lm} choose Biscoe as the reference island, making
the coefficients for both the other islands significant. The result
would have been different if Dream or Torgersen would've been the
reference.

\subsection{Exercise 2}\label{exercise-2-20}

We set Torgersen as reference and rerun the linear regression.

\begin{Shaded}
\begin{Highlighting}[]
\NormalTok{penguins\_re }\OtherTok{\textless{}{-}}\NormalTok{ penguins }\SpecialCharTok{|\textgreater{}} \FunctionTok{mutate}\NormalTok{(}\AttributeTok{island =} \FunctionTok{relevel}\NormalTok{(island, }\AttributeTok{ref =} \StringTok{"Torgersen"}\NormalTok{))}
\NormalTok{model\_summary }\OtherTok{\textless{}{-}}\NormalTok{ penguins\_re }\SpecialCharTok{|\textgreater{}} 
  \FunctionTok{lm}\NormalTok{(}\AttributeTok{formula =}\NormalTok{ body\_mass\_g }\SpecialCharTok{\textasciitilde{}}\NormalTok{ island) }\SpecialCharTok{|\textgreater{}} 
  \FunctionTok{summary}\NormalTok{()}
\NormalTok{model\_summary}\SpecialCharTok{$}\NormalTok{coefficients}
\end{Highlighting}
\end{Shaded}

We see that now only Biscoe is significant.

\subsection{Exercise 3}\label{exercise-3-17}

Gentoos only exists on Biscoe, which is the island with the significant
difference in body weight. It is fair to assume that species is more
responsible than island for the effect we see on body mass.

\begin{Shaded}
\begin{Highlighting}[]
\NormalTok{penguins }\SpecialCharTok{|\textgreater{}} 
  \FunctionTok{ggplot}\NormalTok{(}\FunctionTok{aes}\NormalTok{(}\AttributeTok{x =}\NormalTok{ island, }\AttributeTok{fill =}\NormalTok{ species)) }\SpecialCharTok{+}
  \FunctionTok{geom\_bar}\NormalTok{() }\SpecialCharTok{+}
  \FunctionTok{labs}\NormalTok{(}
    \AttributeTok{x =} \StringTok{"Island"}\NormalTok{,}
    \AttributeTok{y =} \StringTok{"Count"}\NormalTok{,}
    \AttributeTok{fill =} \StringTok{"Species"}\NormalTok{)}
\end{Highlighting}
\end{Shaded}

\includegraphics{solutions_linear_regression_files/figure-pdf/unnamed-chunk-14-1.pdf}

\subsection{Exercise 4}\label{exercise-4-13}

We see that island doesn't have a significant effect on the Adelie
penguins body mass and the coefficients are quite small in comparison to
the total body mass of the average penguin.

\begin{Shaded}
\begin{Highlighting}[]
\NormalTok{model\_summary }\OtherTok{\textless{}{-}}\NormalTok{ penguins }\SpecialCharTok{|\textgreater{}} 
  \FunctionTok{filter}\NormalTok{(species }\SpecialCharTok{==} \StringTok{"Adelie"}\NormalTok{) }\SpecialCharTok{|\textgreater{}}  
  \FunctionTok{lm}\NormalTok{(}\AttributeTok{formula =}\NormalTok{ body\_mass\_g }\SpecialCharTok{\textasciitilde{}}\NormalTok{ island) }\SpecialCharTok{|\textgreater{}} 
  \FunctionTok{summary}\NormalTok{()}
\NormalTok{model\_summary}\SpecialCharTok{$}\NormalTok{coefficients}
\end{Highlighting}
\end{Shaded}

\begin{verbatim}
                   Estimate Std. Error     t value     Pr(>|t|)
(Intercept)     3709.659091   69.58192 53.31355215 1.671273e-98
islandDream      -21.266234   92.98275 -0.22871161 8.194088e-01
islandTorgersen   -3.286542   94.96708 -0.03460717 9.724396e-01
\end{verbatim}

We can visualize the result to see that it seems to make sense.

\begin{Shaded}
\begin{Highlighting}[]
\NormalTok{mean\_bws }\OtherTok{\textless{}{-}}\NormalTok{ penguins }\SpecialCharTok{|\textgreater{}} 
  \FunctionTok{filter}\NormalTok{(species }\SpecialCharTok{==} \StringTok{"Adelie"}\NormalTok{) }\SpecialCharTok{|\textgreater{}}  
  \FunctionTok{summarize}\NormalTok{(}
    \AttributeTok{body\_mass\_g =} \FunctionTok{mean}\NormalTok{(body\_mass\_g, }\AttributeTok{na.rm =} \ConstantTok{TRUE}\NormalTok{),}
    \AttributeTok{.by =}\NormalTok{ island)}
\NormalTok{penguins }\SpecialCharTok{|\textgreater{}} 
  \FunctionTok{filter}\NormalTok{(species }\SpecialCharTok{==} \StringTok{"Adelie"}\NormalTok{) }\SpecialCharTok{|\textgreater{}}  
  \FunctionTok{ggplot}\NormalTok{(}\FunctionTok{aes}\NormalTok{(}\AttributeTok{x =}\NormalTok{ body\_mass\_g, }\AttributeTok{fill =}\NormalTok{ island, }\AttributeTok{col =}\NormalTok{ island)) }\SpecialCharTok{+} 
  \FunctionTok{geom\_density}\NormalTok{(}\AttributeTok{alpha =} \FloatTok{0.3}\NormalTok{) }\SpecialCharTok{+}
  \FunctionTok{geom\_vline}\NormalTok{(}
    \FunctionTok{aes}\NormalTok{(}\AttributeTok{xintercept =}\NormalTok{ body\_mass\_g,}
        \AttributeTok{col =}\NormalTok{ island),}
    \AttributeTok{data =}\NormalTok{ mean\_bws) }\SpecialCharTok{+}
  \FunctionTok{labs}\NormalTok{(}
    \AttributeTok{x =} \StringTok{"Body mass (g)"}\NormalTok{,}
    \AttributeTok{y =} \StringTok{""}\NormalTok{,}
    \AttributeTok{col =} \StringTok{"Island"}\NormalTok{,}
    \AttributeTok{fill =} \StringTok{"Island"}\NormalTok{) }\SpecialCharTok{+}
  \FunctionTok{theme}\NormalTok{(}\AttributeTok{axis.text.y =} \FunctionTok{element\_blank}\NormalTok{(),}
        \AttributeTok{axis.ticks.y =} \FunctionTok{element\_blank}\NormalTok{())}
\end{Highlighting}
\end{Shaded}

\begin{verbatim}
Warning: Removed 1 row containing non-finite outside the scale range
(`stat_density()`).
\end{verbatim}

\includegraphics{solutions_linear_regression_files/figure-pdf/unnamed-chunk-16-1.pdf}

\section{A surprising result (Simpson's
paradox)}\label{a-surprising-result-simpsons-paradox-1}

\subsection{Exercise 2}\label{exercise-2-21}

We see that the \(p\)-values are low for both species, even if it is
just above 0.05 for Gentoos. Even if the points seemingly line up well,
the difference is still statistically significant, at least for one
species. This shows that it is hard to determine statistical
significance by eye only.

\begin{Shaded}
\begin{Highlighting}[]
\NormalTok{model\_summary }\OtherTok{\textless{}{-}}\NormalTok{ penguins }\SpecialCharTok{|\textgreater{}} 
  \FunctionTok{lm}\NormalTok{(}\AttributeTok{formula =}\NormalTok{ body\_mass\_g }\SpecialCharTok{\textasciitilde{}}\NormalTok{ flipper\_length\_mm }\SpecialCharTok{+}\NormalTok{ species) }\SpecialCharTok{|\textgreater{}} 
  \FunctionTok{summary}\NormalTok{()}
\NormalTok{model\_summary}\SpecialCharTok{$}\NormalTok{coefficients}
\end{Highlighting}
\end{Shaded}

\begin{verbatim}
                    Estimate Std. Error   t value     Pr(>|t|)
(Intercept)       -4031.4769  584.15134 -6.901425 2.551453e-11
flipper_length_mm    40.7054    3.07102 13.254686 1.401600e-32
speciesChinstrap   -206.5101   57.73065 -3.577132 3.980986e-04
speciesGentoo       266.8096   95.26374  2.800747 5.391808e-03
\end{verbatim}

\begin{Shaded}
\begin{Highlighting}[]
\NormalTok{penguins }\SpecialCharTok{|\textgreater{}} 
  \FunctionTok{ggplot}\NormalTok{(}\FunctionTok{aes}\NormalTok{(}\AttributeTok{x =}\NormalTok{ flipper\_length\_mm, }\AttributeTok{y =}\NormalTok{ body\_mass\_g, }\AttributeTok{col =}\NormalTok{ species)) }\SpecialCharTok{+}
  \FunctionTok{geom\_point}\NormalTok{() }\SpecialCharTok{+} 
  \FunctionTok{geom\_smooth}\NormalTok{(}\AttributeTok{method =} \StringTok{"lm"}\NormalTok{, }\AttributeTok{se =} \ConstantTok{FALSE}\NormalTok{) }\SpecialCharTok{+}
  \FunctionTok{labs}\NormalTok{(}
    \AttributeTok{x =} \StringTok{"Flipper length (mm)"}\NormalTok{,}
    \AttributeTok{y =} \StringTok{"Body mass (g)"}\NormalTok{,}
    \AttributeTok{col =} \StringTok{"Species"}\NormalTok{)}
\end{Highlighting}
\end{Shaded}

\begin{verbatim}
`geom_smooth()` using formula = 'y ~ x'
\end{verbatim}

\begin{verbatim}
Warning: Removed 2 rows containing non-finite outside the scale range
(`stat_smooth()`).
\end{verbatim}

\begin{verbatim}
Warning: Removed 2 rows containing missing values or values outside the scale range
(`geom_point()`).
\end{verbatim}

\includegraphics{solutions_linear_regression_files/figure-pdf/unnamed-chunk-18-1.pdf}

\section{Prediction and prediction
interval}\label{prediction-and-prediction-interval-1}

\subsection{Exercise 1}\label{exercise-1-25}

\begin{Shaded}
\begin{Highlighting}[]
\NormalTok{model }\OtherTok{\textless{}{-}}\NormalTok{ penguins }\SpecialCharTok{|\textgreater{}}  \FunctionTok{lm}\NormalTok{(}\AttributeTok{formula =}\NormalTok{ body\_mass\_g }\SpecialCharTok{\textasciitilde{}}\NormalTok{ flipper\_length\_mm)}
\NormalTok{model }\SpecialCharTok{|\textgreater{}} \FunctionTok{predict}\NormalTok{(}\FunctionTok{data.frame}\NormalTok{(}\AttributeTok{flipper\_length\_mm =} \DecValTok{201}\NormalTok{), }\AttributeTok{interval =} \StringTok{"prediction"}\NormalTok{, }\AttributeTok{level =} \FloatTok{0.95}\NormalTok{)}
\end{Highlighting}
\end{Shaded}

\begin{verbatim}
       fit      lwr      upr
1 4205.967 3429.303 4982.632
\end{verbatim}

\begin{Shaded}
\begin{Highlighting}[]
\NormalTok{model }\SpecialCharTok{|\textgreater{}} \FunctionTok{predict}\NormalTok{(}\FunctionTok{data.frame}\NormalTok{(}\AttributeTok{flipper\_length\_mm =} \DecValTok{100}\NormalTok{), }\AttributeTok{interval =} \StringTok{"prediction"}\NormalTok{, }\AttributeTok{level =} \FloatTok{0.95}\NormalTok{)}
\end{Highlighting}
\end{Shaded}

\begin{verbatim}
        fit       lwr      upr
1 -812.2747 -1645.371 20.82124
\end{verbatim}

\subsection{Exercise 2}\label{exercise-2-22}

The interval becomes smaller. The confidence level specifies our
confidence in a new observation falling within the interval. If we make
the interval smaller, we have to reduce our confidence. On the other
hand, if we make our interval infidelity big, we can be 100\% confident
it will cover the next observation. Of course there is no use for such a
100\% interval in practice.

\begin{Shaded}
\begin{Highlighting}[]
\NormalTok{model }\SpecialCharTok{|\textgreater{}} \FunctionTok{predict}\NormalTok{(}\FunctionTok{data.frame}\NormalTok{(}\AttributeTok{flipper\_length\_mm =} \DecValTok{200}\NormalTok{), }\AttributeTok{interval =} \StringTok{"prediction"}\NormalTok{, }\AttributeTok{level =} \FloatTok{0.5}\NormalTok{)}
\end{Highlighting}
\end{Shaded}

\begin{verbatim}
       fit     lwr      upr
1 4156.282 3889.67 4422.894
\end{verbatim}

\chapter{Multiple linear regression}\label{multiple-linear-regression}

\section{Setup}\label{setup-2}

First we load the dataset and have a \texttt{glimpse} at it.

\begin{Shaded}
\begin{Highlighting}[]
\FunctionTok{library}\NormalTok{(tidyverse)}
\NormalTok{birth }\OtherTok{\textless{}{-}} \FunctionTok{read\_csv}\NormalTok{(}\StringTok{"datasets/birth.csv"}\NormalTok{) }\CommentTok{\# Remember to set the correct path}
\CommentTok{\# Recoding sex as factor}
\NormalTok{birth }\OtherTok{\textless{}{-}}\NormalTok{ birth }\SpecialCharTok{|\textgreater{}}  \FunctionTok{mutate}\NormalTok{(}\AttributeTok{sex =} \FunctionTok{as.factor}\NormalTok{(}\FunctionTok{ifelse}\NormalTok{(sex }\SpecialCharTok{==} \DecValTok{1}\NormalTok{, }\StringTok{"Male"}\NormalTok{, }\StringTok{"Female"}\NormalTok{)))}
\FunctionTok{glimpse}\NormalTok{(birth)}
\end{Highlighting}
\end{Shaded}

\begin{verbatim}
Rows: 230
Columns: 7
$ birth_weight_g         <dbl> 4073, 4060, 3071, 3546, 3222, 3420, 3106, 3274,~
$ pregnancy_duration_wks <dbl> 41, 40, 39, 39, 39, 41, 40, 40, 40, 41, 41, 38,~
$ sex                    <fct> Male, Female, Female, Male, Female, Male, Femal~
$ age                    <dbl> 42, 27, 37, 39, 31, 33, 38, 26, 30, 35, 29, 29,~
$ length_cm              <dbl> 176, 180, 171, 167, 158, 166, 169, 164, 167, 16~
$ weight_kg              <dbl> 65, 81, 70, 80, 54, 50, 50, 65, 80, 55, 57, 63,~
$ bmi                    <dbl> 20.98399, 25.00000, 23.93899, 28.68514, 21.6311~
\end{verbatim}

\section{Birth weight dataset}\label{birth-weight-dataset-1}

\subsection{Exercise 1}\label{exercise-1-26}

We don't know anything else than that the sample was collected in
Sweden, so our population is children born in Sweden. It would be useful
to know more about where the children was born to better specify the
population, and when they were born.

Information about the father, ethnicity, and health related variables
such as smoking and blood pressure would be interesting.

\subsection{Exercise 2}\label{exercise-2-23}

The null and alternative hypotheses are

\begin{quote}
\textbf{\(H_0\)}: There is \emph{no} linear relationship between the
mothers age and the birth weight of the child.
\end{quote}

\begin{quote}
\textbf{\(H_a\)}: There \emph{is} a linear relationship between the
mothers age and the birth weight of the child.
\end{quote}

We perform linear regression with \texttt{age} as our predictor.

\begin{Shaded}
\begin{Highlighting}[]
\NormalTok{model\_summary }\OtherTok{\textless{}{-}}\NormalTok{ birth }\SpecialCharTok{|\textgreater{}} 
  \FunctionTok{lm}\NormalTok{(}\AttributeTok{formula =}\NormalTok{ birth\_weight\_g }\SpecialCharTok{\textasciitilde{}}\NormalTok{ age) }\SpecialCharTok{|\textgreater{}} 
  \FunctionTok{summary}\NormalTok{()}
\NormalTok{model\_summary}\SpecialCharTok{$}\NormalTok{coefficients}
\end{Highlighting}
\end{Shaded}

\begin{verbatim}
               Estimate Std. Error   t value     Pr(>|t|)
(Intercept) 3453.230738 201.341654 17.151099 6.616019e-43
age            1.233859   6.336775  0.194714 8.457902e-01
\end{verbatim}

We see that the \(p\)-value is \(\sim0.85\), so we conclude that there
is no significant linear relationship between the birth weight and the
mothers age.

\subsection{Exercise 3}\label{exercise-3-18}

We fit the three linear models separately.

\begin{Shaded}
\begin{Highlighting}[]
\CommentTok{\# Mother length}
\NormalTok{model\_summary }\OtherTok{\textless{}{-}}\NormalTok{ birth }\SpecialCharTok{|\textgreater{}} 
  \FunctionTok{lm}\NormalTok{(}\AttributeTok{formula =}\NormalTok{ birth\_weight\_g }\SpecialCharTok{\textasciitilde{}}\NormalTok{ length\_cm) }\SpecialCharTok{|\textgreater{}} 
  \FunctionTok{summary}\NormalTok{()}
\NormalTok{model\_summary}\SpecialCharTok{$}\NormalTok{coefficients}
\end{Highlighting}
\end{Shaded}

\begin{verbatim}
             Estimate Std. Error   t value     Pr(>|t|)
(Intercept) 108.03363 767.697547 0.1407242 8.882121e-01
length_cm    20.25099   4.590836 4.4111770 1.585424e-05
\end{verbatim}

\begin{Shaded}
\begin{Highlighting}[]
\CommentTok{\# Mother weight}
\NormalTok{model\_summary }\OtherTok{\textless{}{-}}\NormalTok{ birth }\SpecialCharTok{|\textgreater{}} 
  \FunctionTok{lm}\NormalTok{(}\AttributeTok{formula =}\NormalTok{ birth\_weight\_g }\SpecialCharTok{\textasciitilde{}}\NormalTok{ weight\_kg) }\SpecialCharTok{|\textgreater{}} 
  \FunctionTok{summary}\NormalTok{()}
\NormalTok{model\_summary}\SpecialCharTok{$}\NormalTok{coefficients}
\end{Highlighting}
\end{Shaded}

\begin{verbatim}
               Estimate Std. Error   t value     Pr(>|t|)
(Intercept) 2974.057674  179.75748 16.544834 6.358329e-41
weight_kg      8.032416    2.74819  2.922803 3.818211e-03
\end{verbatim}

\begin{Shaded}
\begin{Highlighting}[]
\CommentTok{\# Mother BMI}
\NormalTok{model\_summary }\OtherTok{\textless{}{-}}\NormalTok{ birth }\SpecialCharTok{|\textgreater{}} 
  \FunctionTok{lm}\NormalTok{(}\AttributeTok{formula =}\NormalTok{ birth\_weight\_g }\SpecialCharTok{\textasciitilde{}}\NormalTok{ bmi) }\SpecialCharTok{|\textgreater{}} 
  \FunctionTok{summary}\NormalTok{()}
\NormalTok{model\_summary}\SpecialCharTok{$}\NormalTok{coefficients}
\end{Highlighting}
\end{Shaded}

\begin{verbatim}
               Estimate Std. Error   t value     Pr(>|t|)
(Intercept) 3294.148417 193.123860 17.057180 1.340283e-42
bmi            8.570779   8.260604  1.037549 3.005793e-01
\end{verbatim}

We see that the \(p\)-values for length and weight are below the
significance level 0.05 but not BMI.

\section{Comparing models}\label{comparing-models-1}

\subsection{Selecting between two simple linear regression
models}\label{selecting-between-two-simple-linear-regression-models-1}

The linear relationship was not significant so we would not expect
\texttt{age} to explain much of the variation of
\texttt{birth\_weight\_g}.

\begin{Shaded}
\begin{Highlighting}[]
\NormalTok{model\_summary }\OtherTok{\textless{}{-}}\NormalTok{ birth }\SpecialCharTok{|\textgreater{}} 
  \FunctionTok{lm}\NormalTok{(}\AttributeTok{formula =}\NormalTok{ birth\_weight\_g }\SpecialCharTok{\textasciitilde{}}\NormalTok{ age) }\SpecialCharTok{|\textgreater{}} 
  \FunctionTok{summary}\NormalTok{()}
\NormalTok{model\_summary}\SpecialCharTok{$}\NormalTok{r.squared}
\end{Highlighting}
\end{Shaded}

\begin{verbatim}
[1] 0.0001662599
\end{verbatim}

In our plot we can see that the line is almost completely flat.

\begin{Shaded}
\begin{Highlighting}[]
\NormalTok{birth }\SpecialCharTok{|\textgreater{}} 
  \FunctionTok{ggplot}\NormalTok{(}\FunctionTok{aes}\NormalTok{(}\AttributeTok{x =}\NormalTok{ age, }\AttributeTok{y =}\NormalTok{ birth\_weight\_g)) }\SpecialCharTok{+}
  \FunctionTok{geom\_point}\NormalTok{() }\SpecialCharTok{+} 
  \FunctionTok{geom\_smooth}\NormalTok{(}\AttributeTok{method =} \StringTok{"lm"}\NormalTok{, }\AttributeTok{se =} \ConstantTok{FALSE}\NormalTok{) }\SpecialCharTok{+}
  \FunctionTok{labs}\NormalTok{(}
    \AttributeTok{x =} \StringTok{"Age"}\NormalTok{,}
    \AttributeTok{y =} \StringTok{"Birth weight (g)"}\NormalTok{)}
\end{Highlighting}
\end{Shaded}

\begin{verbatim}
`geom_smooth()` using formula = 'y ~ x'
\end{verbatim}

\includegraphics{solutions_linear_regression_multiple_files/figure-pdf/unnamed-chunk-5-1.pdf}

We can also look at a few selected points and look at their explained
variation like in the chapter. We see that almost all variation is
unexplained. Note: You were not asked to do this visualization on your
own.

\begin{Shaded}
\begin{Highlighting}[]
\CommentTok{\# Limits}
\NormalTok{ylims }\OtherTok{\textless{}{-}} \FunctionTok{c}\NormalTok{(}\DecValTok{3000}\NormalTok{, }\DecValTok{4000}\NormalTok{)}
\CommentTok{\# Predictions}
\NormalTok{model }\OtherTok{\textless{}{-}}\NormalTok{ birth }\SpecialCharTok{|\textgreater{}} \FunctionTok{lm}\NormalTok{(}\AttributeTok{formula =}\NormalTok{ birth\_weight\_g }\SpecialCharTok{\textasciitilde{}}\NormalTok{ age)}
\CommentTok{\# Subset points to show}
\FunctionTok{set.seed}\NormalTok{(}\DecValTok{1337}\NormalTok{)}
\NormalTok{subset\_birth }\OtherTok{\textless{}{-}}\NormalTok{ birth }\SpecialCharTok{|\textgreater{}} 
  \FunctionTok{mutate}\NormalTok{(}
    \AttributeTok{pred =} \FunctionTok{predict}\NormalTok{(model),}
    \AttributeTok{resid =} \FunctionTok{resid}\NormalTok{(model)) }\SpecialCharTok{|\textgreater{}} 
  \FunctionTok{filter}\NormalTok{(}
\NormalTok{    (age }\SpecialCharTok{\textless{}} \FunctionTok{mean}\NormalTok{(birth}\SpecialCharTok{$}\NormalTok{age) }\SpecialCharTok{\&}\NormalTok{ birth\_weight\_g }\SpecialCharTok{\textless{}}\NormalTok{ pred }\SpecialCharTok{{-}} \DecValTok{100}\NormalTok{) }\SpecialCharTok{|}
\NormalTok{    (age }\SpecialCharTok{\textgreater{}} \FunctionTok{mean}\NormalTok{(birth}\SpecialCharTok{$}\NormalTok{age) }\SpecialCharTok{\&}\NormalTok{ birth\_weight\_g }\SpecialCharTok{\textgreater{}}\NormalTok{ pred }\SpecialCharTok{+} \DecValTok{100}\NormalTok{)) }\SpecialCharTok{|\textgreater{}}  
  \FunctionTok{filter}\NormalTok{(birth\_weight\_g }\SpecialCharTok{\textless{}} \DecValTok{4000} \SpecialCharTok{\&}\NormalTok{ birth\_weight\_g }\SpecialCharTok{\textgreater{}} \DecValTok{3000}\NormalTok{) }\SpecialCharTok{|\textgreater{}} 
  \FunctionTok{sample\_n}\NormalTok{(}\DecValTok{20}\NormalTok{)}
\CommentTok{\# Calculate mean of birth weight}
\NormalTok{mean\_bw }\OtherTok{\textless{}{-}}\NormalTok{ birth }\SpecialCharTok{|\textgreater{}} \FunctionTok{pull}\NormalTok{(birth\_weight\_g) }\SpecialCharTok{|\textgreater{}} \FunctionTok{mean}\NormalTok{()}

\NormalTok{birth }\SpecialCharTok{|\textgreater{}} 
  \FunctionTok{ggplot}\NormalTok{(}\FunctionTok{aes}\NormalTok{(}
    \AttributeTok{x =}\NormalTok{ age, }
    \AttributeTok{y =}\NormalTok{ birth\_weight\_g)) }\SpecialCharTok{+} 
  \FunctionTok{geom\_hline}\NormalTok{(}
    \AttributeTok{yintercept =}\NormalTok{ mean\_bw,}
    \AttributeTok{linetype =} \DecValTok{3}\NormalTok{,}
    \AttributeTok{linewidth =} \DecValTok{1}\NormalTok{) }\SpecialCharTok{+}
  \FunctionTok{geom\_point}\NormalTok{(}\AttributeTok{alpha =} \FloatTok{0.1}\NormalTok{) }\SpecialCharTok{+} 
  \FunctionTok{geom\_smooth}\NormalTok{(}\AttributeTok{method =} \StringTok{"lm"}\NormalTok{, }\AttributeTok{se =} \ConstantTok{FALSE}\NormalTok{) }\SpecialCharTok{+}
  \FunctionTok{geom\_segment}\NormalTok{( }\CommentTok{\# Total variation}
    \AttributeTok{data =}\NormalTok{ subset\_birth,}
    \FunctionTok{aes}\NormalTok{(}
      \AttributeTok{x =}\NormalTok{ age }\SpecialCharTok{+} \FloatTok{0.1}\NormalTok{,}
      \AttributeTok{xend =}\NormalTok{ age }\SpecialCharTok{+} \FloatTok{0.1}\NormalTok{,}
      \AttributeTok{y =}\NormalTok{  mean\_bw,}
      \AttributeTok{yend =}\NormalTok{ birth\_weight\_g,}
      \AttributeTok{col =} \StringTok{"Total"}\NormalTok{),}
    \AttributeTok{linetype =} \DecValTok{2}\NormalTok{,}
    \AttributeTok{linewidth =} \FloatTok{1.5}\NormalTok{,}
    \AttributeTok{alpha =} \FloatTok{0.5}\NormalTok{) }\SpecialCharTok{+}
    \FunctionTok{geom\_segment}\NormalTok{( }\CommentTok{\# Unexplained variation}
    \AttributeTok{data =}\NormalTok{ subset\_birth,}
    \FunctionTok{aes}\NormalTok{(}
      \AttributeTok{x =}\NormalTok{ age,}
      \AttributeTok{xend =}\NormalTok{ age,}
      \AttributeTok{y =}\NormalTok{  pred,}
      \AttributeTok{yend =}\NormalTok{ birth\_weight\_g,}
      \AttributeTok{col =} \StringTok{"Unexplained"}\NormalTok{),}
    \AttributeTok{linewidth =} \FloatTok{1.5}\NormalTok{,}
    \AttributeTok{alpha =} \FloatTok{0.7}
\NormalTok{    ) }\SpecialCharTok{+}
  \FunctionTok{geom\_segment}\NormalTok{( }\CommentTok{\# Explained variation}
    \AttributeTok{data =}\NormalTok{ subset\_birth,}
    \FunctionTok{aes}\NormalTok{(}
      \AttributeTok{x =}\NormalTok{ age,}
      \AttributeTok{xend =}\NormalTok{ age,}
      \AttributeTok{y =}\NormalTok{  mean\_bw,}
      \AttributeTok{yend =}\NormalTok{ pred,}
      \AttributeTok{col =} \StringTok{"Explained"}\NormalTok{),}
    \AttributeTok{linewidth =} \FloatTok{1.5}\NormalTok{,}
    \AttributeTok{alpha =} \FloatTok{0.7}
\NormalTok{    ) }\SpecialCharTok{+}
  \FunctionTok{geom\_point}\NormalTok{(}\AttributeTok{data =}\NormalTok{ subset\_birth, }\AttributeTok{size =} \DecValTok{4}\NormalTok{) }\SpecialCharTok{+} 
  \FunctionTok{labs}\NormalTok{(}
    \AttributeTok{col =} \StringTok{"Variation"}\NormalTok{) }\SpecialCharTok{+}
  \FunctionTok{scale\_color\_manual}\NormalTok{(}
    \AttributeTok{values =} \FunctionTok{c}\NormalTok{(}
      \StringTok{"Total"} \OtherTok{=} \StringTok{"black"}\NormalTok{,}
      \StringTok{"Explained"} \OtherTok{=} \StringTok{"darkgreen"}\NormalTok{, }
      \StringTok{"Unexplained"} \OtherTok{=} \StringTok{"darkorange"}\NormalTok{),}
    \AttributeTok{breaks =} \FunctionTok{c}\NormalTok{(}
      \StringTok{"Total"}\NormalTok{, }\StringTok{"Explained"}\NormalTok{, }\StringTok{"Unexplained"}\NormalTok{)) }\SpecialCharTok{+}
  \FunctionTok{coord\_cartesian}\NormalTok{(}\AttributeTok{ylim =}\NormalTok{ ylims)}
\end{Highlighting}
\end{Shaded}

\includegraphics{solutions_linear_regression_multiple_files/figure-pdf/unnamed-chunk-6-1.pdf}

\subsection{Exercise 2}\label{exercise-2-24}

We fit a multiple linear regression model using the provided formula.

\begin{Shaded}
\begin{Highlighting}[]
\NormalTok{model }\OtherTok{\textless{}{-}}\NormalTok{ birth }\SpecialCharTok{|\textgreater{}} \FunctionTok{lm}\NormalTok{(}\AttributeTok{formula =}\NormalTok{ birth\_weight\_g }\SpecialCharTok{\textasciitilde{}}\NormalTok{ weight\_kg }\SpecialCharTok{+}\NormalTok{ length\_cm)}
\NormalTok{model }\SpecialCharTok{|\textgreater{}} \FunctionTok{summary}\NormalTok{()}
\end{Highlighting}
\end{Shaded}

\begin{verbatim}

Call:
lm(formula = birth_weight_g ~ weight_kg + length_cm, data = birth)

Residuals:
     Min       1Q   Median       3Q      Max 
-1252.37  -286.14    11.38   271.66  1224.75 

Coefficients:
            Estimate Std. Error t value Pr(>|t|)    
(Intercept)  279.653    774.487   0.361 0.718374    
weight_kg      4.255      2.881   1.477 0.141141    
length_cm     17.582      4.923   3.572 0.000433 ***
---
Signif. codes:  0 '***' 0.001 '**' 0.01 '*' 0.05 '.' 0.1 ' ' 1

Residual standard error: 447 on 227 degrees of freedom
Multiple R-squared:  0.0874,    Adjusted R-squared:  0.07936 
F-statistic: 10.87 on 2 and 227 DF,  p-value: 3.103e-05
\end{verbatim}

Including both \texttt{weight\_kg} and \texttt{length\_cm} increase
R\(^2\) by a little, 0.0874 compared to 0.0786 when using only
\texttt{length\_cm}. It is possible that a lot of the variation
explained by weight is also explained by length, as it is reasonable to
believe that longer women also weigh more, on average.

\subsection{Exercise 3}\label{exercise-3-19}

We see that our regression model is visualize as a \emph{plane} in 3D
space, where we can vary both \texttt{weight\_kg} and
\texttt{length\_cm}. In general we will not visualize multiple linear
regression models, and in fact any model containing more than two
predictors is not possible to visualize. Understanding the visualization
requires some geometrical intuition, and if you struggle to understand
the visualization you can focus on interpreting the models instead, as
it is more important.

\textbf{Note:} We are working on rendering this for the published
lecture notes. You can see the visualization in the Multiple linear
regression module on Canvas.

\section{Fitting multiple linear regression
models}\label{fitting-multiple-linear-regression-models-1}

\subsection{Interpreting multiple linear
regression}\label{interpreting-multiple-linear-regression-1}

\subsubsection{Exercise 1}\label{exercise-1-27}

We first run the code to simulate the data.

\begin{Shaded}
\begin{Highlighting}[]
\CommentTok{\# Simulating data}
\FunctionTok{set.seed}\NormalTok{(}\DecValTok{42}\NormalTok{) }\CommentTok{\# Important to include this!}
\NormalTok{n\_days }\OtherTok{\textless{}{-}} \DecValTok{153} 
\NormalTok{n\_years }\OtherTok{\textless{}{-}} \DecValTok{5}
\NormalTok{month }\OtherTok{\textless{}{-}} \FunctionTok{rep}\NormalTok{(}\FunctionTok{seq}\NormalTok{(}\DecValTok{4}\NormalTok{, }\DecValTok{8}\NormalTok{, }\AttributeTok{length.out =}\NormalTok{ n\_days), n\_years)}
\NormalTok{ice\_cream\_sales }\OtherTok{\textless{}{-}} \DecValTok{2000} \SpecialCharTok{+} \DecValTok{3000} \SpecialCharTok{*} \FunctionTok{sin}\NormalTok{(pi }\SpecialCharTok{*}\NormalTok{ month }\SpecialCharTok{/} \DecValTok{14}\NormalTok{) }\SpecialCharTok{+} \FunctionTok{rnorm}\NormalTok{(}\FunctionTok{length}\NormalTok{(month), }\DecValTok{0}\NormalTok{, }\DecValTok{200}\NormalTok{)}
\NormalTok{snake\_bites }\OtherTok{\textless{}{-}} \DecValTok{20} \SpecialCharTok{*} \FunctionTok{sin}\NormalTok{(pi }\SpecialCharTok{*}\NormalTok{ month }\SpecialCharTok{/} \DecValTok{14}\NormalTok{) }\SpecialCharTok{+} \FunctionTok{rnorm}\NormalTok{(}\FunctionTok{length}\NormalTok{(month), }\DecValTok{0}\NormalTok{, }\DecValTok{8}\NormalTok{)}

\NormalTok{ice\_cream }\OtherTok{\textless{}{-}} \FunctionTok{tibble}\NormalTok{(}\AttributeTok{month =}\NormalTok{ month, }\AttributeTok{ice\_cream\_sales =}\NormalTok{ ice\_cream\_sales, }\AttributeTok{snake\_bites =}\NormalTok{ snake\_bites)}
\end{Highlighting}
\end{Shaded}

We then visualize the relationship and see that there may be a linear
relationship.

\begin{Shaded}
\begin{Highlighting}[]
\NormalTok{ice\_cream }\SpecialCharTok{|\textgreater{}} 
  \FunctionTok{ggplot}\NormalTok{(}\FunctionTok{aes}\NormalTok{(}\AttributeTok{x =}\NormalTok{ snake\_bites, }\AttributeTok{y =}\NormalTok{ ice\_cream\_sales)) }\SpecialCharTok{+} 
  \FunctionTok{geom\_point}\NormalTok{() }\SpecialCharTok{+} 
  \FunctionTok{geom\_smooth}\NormalTok{(}\AttributeTok{method =} \StringTok{"lm"}\NormalTok{, }\AttributeTok{se =} \ConstantTok{FALSE}\NormalTok{)}
\end{Highlighting}
\end{Shaded}

\begin{verbatim}
`geom_smooth()` using formula = 'y ~ x'
\end{verbatim}

\includegraphics{solutions_linear_regression_multiple_files/figure-pdf/unnamed-chunk-10-1.pdf}

We finally fit a regression model to see if the relationship is
significant.

\begin{Shaded}
\begin{Highlighting}[]
\NormalTok{ice\_cream }\SpecialCharTok{|\textgreater{}} 
  \FunctionTok{lm}\NormalTok{(}\AttributeTok{formula =}\NormalTok{ ice\_cream\_sales }\SpecialCharTok{\textasciitilde{}}\NormalTok{ snake\_bites) }\SpecialCharTok{|\textgreater{}} 
  \FunctionTok{summary}\NormalTok{()}
\end{Highlighting}
\end{Shaded}

\begin{verbatim}

Call:
lm(formula = ice_cream_sales ~ snake_bites, data = ice_cream)

Residuals:
    Min      1Q  Median      3Q     Max 
-844.14 -179.84   17.15  193.59  709.88 

Coefficients:
            Estimate Std. Error t value Pr(>|t|)    
(Intercept) 4726.072     24.170 195.536  < 2e-16 ***
snake_bites    4.907      1.183   4.148 3.74e-05 ***
---
Signif. codes:  0 '***' 0.001 '**' 0.01 '*' 0.05 '.' 0.1 ' ' 1

Residual standard error: 268.7 on 763 degrees of freedom
Multiple R-squared:  0.02205,   Adjusted R-squared:  0.02077 
F-statistic:  17.2 on 1 and 763 DF,  p-value: 3.735e-05
\end{verbatim}

We see that the \(p\)-value is less than 0.05 and the relationship is
significant with significance level 0.05.

\subsubsection{Exercise 2}\label{exercise-2-25}

It doesn't seem likely that snake bites increases ice cream sales,
i.e.~the effect doesn't seem intuitively causal. Instead there may be
another variable explaining why they correlate\ldots{}

\subsubsection{Exercise 3}\label{exercise-3-20}

We include \texttt{month} and fit a multiple linear regression model
with hypotheses

\begin{quote}
\textbf{\(H_0\)}: There is \emph{no} linear relationship between snake
bites and ice cream sales, given that month is part of the model.
\end{quote}

\begin{quote}
\textbf{\(H_a\)}: There \emph{is} a linear relationship between snake
bites and ice cream sales, given that month is part of the model.
\end{quote}

\begin{Shaded}
\begin{Highlighting}[]
\NormalTok{ice\_cream }\SpecialCharTok{|\textgreater{}} 
  \FunctionTok{lm}\NormalTok{(}\AttributeTok{formula =}\NormalTok{ ice\_cream\_sales }\SpecialCharTok{\textasciitilde{}}\NormalTok{ snake\_bites }\SpecialCharTok{+}\NormalTok{ month) }\SpecialCharTok{|\textgreater{}} 
  \FunctionTok{summary}\NormalTok{()}
\end{Highlighting}
\end{Shaded}

\begin{verbatim}

Call:
lm(formula = ice_cream_sales ~ snake_bites + month, data = ice_cream)

Residuals:
    Min      1Q  Median      3Q     Max 
-590.93 -138.02   -0.79  142.45  746.38 

Coefficients:
             Estimate Std. Error t value Pr(>|t|)    
(Intercept) 3936.8267    42.0421  93.640   <2e-16 ***
snake_bites    1.4352     0.9548   1.503    0.133    
month        142.3655     6.7454  21.106   <2e-16 ***
---
Signif. codes:  0 '***' 0.001 '**' 0.01 '*' 0.05 '.' 0.1 ' ' 1

Residual standard error: 213.6 on 762 degrees of freedom
Multiple R-squared:  0.3828,    Adjusted R-squared:  0.3812 
F-statistic: 236.3 on 2 and 762 DF,  p-value: < 2.2e-16
\end{verbatim}

We don't have evidence in our data that there is a linear relationship
between snake bites and ice cream sales at significance level 0.05,
given that the month is included in the model.

\textbf{Beyond the scope of the course}

If we look at the code that generated the data we can see that in fact
\texttt{snake\_bites} and \texttt{ice\_cream\_sales} are both generated
based on \texttt{month}, but are otherwise unrelated. You do not need to
understand the mathematical formula, just notice \texttt{month} appears
on the right side of \texttt{\textless{}-}, indicating that the values
of \texttt{ice\_cream\_sales} and \texttt{snake\_bites} depend on
\texttt{month}.

\begin{Shaded}
\begin{Highlighting}[]
\NormalTok{ice\_cream\_sales }\OtherTok{\textless{}{-}} \DecValTok{2000} \SpecialCharTok{+} \DecValTok{3000} \SpecialCharTok{*} \FunctionTok{sin}\NormalTok{(pi }\SpecialCharTok{*}\NormalTok{ month }\SpecialCharTok{/} \DecValTok{14}\NormalTok{) }\SpecialCharTok{+} \FunctionTok{rnorm}\NormalTok{(}\FunctionTok{length}\NormalTok{(month), }\DecValTok{0}\NormalTok{, }\DecValTok{200}\NormalTok{)}
\NormalTok{snake\_bites }\OtherTok{\textless{}{-}} \DecValTok{20} \SpecialCharTok{*} \FunctionTok{sin}\NormalTok{(pi }\SpecialCharTok{*}\NormalTok{ month }\SpecialCharTok{/} \DecValTok{14}\NormalTok{) }\SpecialCharTok{+} \FunctionTok{rnorm}\NormalTok{(}\FunctionTok{length}\NormalTok{(month), }\DecValTok{0}\NormalTok{, }\DecValTok{8}\NormalTok{)}
\end{Highlighting}
\end{Shaded}

\section{Model selection}\label{model-selection-2}

\subsection{Multicollinearity}\label{multicollinearity-1}

\subsubsection{Exercise 1}\label{exercise-1-28}

Visualizing shows the linear relationship.

\begin{Shaded}
\begin{Highlighting}[]
\NormalTok{birth }\SpecialCharTok{|\textgreater{}} 
  \FunctionTok{ggplot}\NormalTok{(}\FunctionTok{aes}\NormalTok{(}\AttributeTok{x=}\NormalTok{ weight\_kg, }\AttributeTok{y =}\NormalTok{ bmi)) }\SpecialCharTok{+} 
  \FunctionTok{geom\_point}\NormalTok{() }\SpecialCharTok{+}
  \FunctionTok{geom\_smooth}\NormalTok{(}\AttributeTok{method =} \StringTok{"lm"}\NormalTok{, }\AttributeTok{se =} \ConstantTok{FALSE}\NormalTok{)}
\end{Highlighting}
\end{Shaded}

\begin{verbatim}
`geom_smooth()` using formula = 'y ~ x'
\end{verbatim}

\includegraphics{solutions_linear_regression_multiple_files/figure-pdf/unnamed-chunk-14-1.pdf}

Unsurprisingly, the relationship is significant and R\(^2\) is as high
as 0.8.

\begin{Shaded}
\begin{Highlighting}[]
\NormalTok{birth }\SpecialCharTok{|\textgreater{}} 
  \FunctionTok{lm}\NormalTok{(}\AttributeTok{formula =}\NormalTok{ bmi }\SpecialCharTok{\textasciitilde{}}\NormalTok{ weight\_kg) }\SpecialCharTok{|\textgreater{}} 
  \FunctionTok{summary}\NormalTok{()}
\end{Highlighting}
\end{Shaded}

\begin{verbatim}

Call:
lm(formula = bmi ~ weight_kg, data = birth)

Residuals:
    Min      1Q  Median      3Q     Max 
-4.7078 -1.0939 -0.1947  0.9263  5.8122 

Coefficients:
            Estimate Std. Error t value Pr(>|t|)    
(Intercept) 3.572167   0.653213   5.469 1.19e-07 ***
weight_kg   0.302571   0.009987  30.298  < 2e-16 ***
---
Signif. codes:  0 '***' 0.001 '**' 0.01 '*' 0.05 '.' 0.1 ' ' 1

Residual standard error: 1.666 on 228 degrees of freedom
Multiple R-squared:  0.801, Adjusted R-squared:  0.8002 
F-statistic:   918 on 1 and 228 DF,  p-value: < 2.2e-16
\end{verbatim}

\subsubsection{Exercise 2}\label{exercise-2-26}

We can solve this in two ways, either by avoiding using \texttt{.} in
our formula or by removing \texttt{z} from \texttt{df}.

\begin{Shaded}
\begin{Highlighting}[]
\FunctionTok{set.seed}\NormalTok{(}\DecValTok{42}\NormalTok{) }\CommentTok{\# Setting seed for reproducibility}
\NormalTok{x }\OtherTok{\textless{}{-}} \FunctionTok{seq}\NormalTok{(}\DecValTok{1}\NormalTok{, }\DecValTok{100}\NormalTok{) }\SpecialCharTok{+} \FunctionTok{rnorm}\NormalTok{(}\DecValTok{100}\NormalTok{, }\DecValTok{0}\NormalTok{, }\DecValTok{2}\NormalTok{)}
\NormalTok{y }\OtherTok{\textless{}{-}} \DecValTok{3} \SpecialCharTok{*}\NormalTok{ x }\SpecialCharTok{+} \FunctionTok{rnorm}\NormalTok{(}\DecValTok{100}\NormalTok{, }\DecValTok{0}\NormalTok{, }\DecValTok{2}\NormalTok{)}
\NormalTok{z }\OtherTok{\textless{}{-}} \FunctionTok{rep}\NormalTok{(}\ConstantTok{NA}\NormalTok{, }\DecValTok{100}\NormalTok{)}

\NormalTok{df }\OtherTok{\textless{}{-}} \FunctionTok{data.frame}\NormalTok{(}\AttributeTok{x =}\NormalTok{ x, }\AttributeTok{y =}\NormalTok{ y, }\AttributeTok{z =}\NormalTok{ z)}

\NormalTok{model1 }\OtherTok{\textless{}{-}}\NormalTok{ df }\SpecialCharTok{|\textgreater{}} \FunctionTok{lm}\NormalTok{(}\AttributeTok{formula =}\NormalTok{ y }\SpecialCharTok{\textasciitilde{}}\NormalTok{ x)}
\NormalTok{model1\_summary }\OtherTok{\textless{}{-}}\NormalTok{ model1 }\SpecialCharTok{|\textgreater{}} \FunctionTok{summary}\NormalTok{()}
\NormalTok{model1\_summary}\SpecialCharTok{$}\NormalTok{r.squared}
\NormalTok{model2 }\OtherTok{\textless{}{-}}\NormalTok{ df }\SpecialCharTok{|\textgreater{}} 
  \FunctionTok{select}\NormalTok{(}\SpecialCharTok{!}\NormalTok{z) }\SpecialCharTok{|\textgreater{}} 
  \FunctionTok{lm}\NormalTok{(}\AttributeTok{formula =}\NormalTok{ y }\SpecialCharTok{\textasciitilde{}}\NormalTok{ .)}
\NormalTok{model2\_summary }\OtherTok{\textless{}{-}}\NormalTok{ model2 }\SpecialCharTok{|\textgreater{}} \FunctionTok{summary}\NormalTok{()}
\NormalTok{model2\_summary}\SpecialCharTok{$}\NormalTok{r.squared}
\end{Highlighting}
\end{Shaded}

\subsection{Model selection
strategies}\label{model-selection-strategies-1}

\subsubsection{Exercise 1}\label{exercise-1-29}

When performing forward selection we start by fitting simple linear
regression models for each of the potential predictors and add the
variable that improves the adjusted R\(^2\) the most to our model. It is
important to use adjusted R\(^2\) since adding any variable will
increase R\(^2\) but adjusted R\(^2\) will penalize adding more
variables to the model, balancing model fit and model complexity.

To keep the code as compact as possible, we will use the \texttt{;}
operator that has the same effect as a new line when extracting the
adjusted R\(^2\) value from each model.

\begin{Shaded}
\begin{Highlighting}[]
\NormalTok{ms }\OtherTok{\textless{}{-}}\NormalTok{ birth }\SpecialCharTok{|\textgreater{}} \FunctionTok{lm}\NormalTok{(}\AttributeTok{formula =}\NormalTok{ birth\_weight\_g }\SpecialCharTok{\textasciitilde{}}\NormalTok{ pregnancy\_duration\_wks) }\SpecialCharTok{|\textgreater{}} \FunctionTok{summary}\NormalTok{() ; ms}\SpecialCharTok{$}\NormalTok{adj.r.squared}
\end{Highlighting}
\end{Shaded}

\begin{verbatim}
[1] 0.2433635
\end{verbatim}

\begin{Shaded}
\begin{Highlighting}[]
\NormalTok{ms }\OtherTok{\textless{}{-}}\NormalTok{ birth }\SpecialCharTok{|\textgreater{}} \FunctionTok{lm}\NormalTok{(}\AttributeTok{formula =}\NormalTok{ birth\_weight\_g }\SpecialCharTok{\textasciitilde{}}\NormalTok{ sex) }\SpecialCharTok{|\textgreater{}} \FunctionTok{summary}\NormalTok{() ; ms}\SpecialCharTok{$}\NormalTok{adj.r.squared}
\end{Highlighting}
\end{Shaded}

\begin{verbatim}
[1] 0.05762742
\end{verbatim}

\begin{Shaded}
\begin{Highlighting}[]
\NormalTok{ms }\OtherTok{\textless{}{-}}\NormalTok{ birth }\SpecialCharTok{|\textgreater{}} \FunctionTok{lm}\NormalTok{(}\AttributeTok{formula =}\NormalTok{ birth\_weight\_g }\SpecialCharTok{\textasciitilde{}}\NormalTok{ age) }\SpecialCharTok{|\textgreater{}} \FunctionTok{summary}\NormalTok{() ; ms}\SpecialCharTok{$}\NormalTok{adj.r.squared}
\end{Highlighting}
\end{Shaded}

\begin{verbatim}
[1] -0.004218976
\end{verbatim}

\begin{Shaded}
\begin{Highlighting}[]
\NormalTok{ms }\OtherTok{\textless{}{-}}\NormalTok{ birth }\SpecialCharTok{|\textgreater{}} \FunctionTok{lm}\NormalTok{(}\AttributeTok{formula =}\NormalTok{ birth\_weight\_g }\SpecialCharTok{\textasciitilde{}}\NormalTok{ length\_cm) }\SpecialCharTok{|\textgreater{}} \FunctionTok{summary}\NormalTok{() ; ms}\SpecialCharTok{$}\NormalTok{adj.r.squared}
\end{Highlighting}
\end{Shaded}

\begin{verbatim}
[1] 0.07459224
\end{verbatim}

\begin{Shaded}
\begin{Highlighting}[]
\NormalTok{ms }\OtherTok{\textless{}{-}}\NormalTok{ birth }\SpecialCharTok{|\textgreater{}} \FunctionTok{lm}\NormalTok{(}\AttributeTok{formula =}\NormalTok{ birth\_weight\_g }\SpecialCharTok{\textasciitilde{}}\NormalTok{ weight\_kg) }\SpecialCharTok{|\textgreater{}} \FunctionTok{summary}\NormalTok{() ; ms}\SpecialCharTok{$}\NormalTok{adj.r.squared}
\end{Highlighting}
\end{Shaded}

\begin{verbatim}
[1] 0.03188758
\end{verbatim}

\begin{Shaded}
\begin{Highlighting}[]
\NormalTok{ms }\OtherTok{\textless{}{-}}\NormalTok{ birth }\SpecialCharTok{|\textgreater{}} \FunctionTok{lm}\NormalTok{(}\AttributeTok{formula =}\NormalTok{ birth\_weight\_g }\SpecialCharTok{\textasciitilde{}}\NormalTok{ bmi) }\SpecialCharTok{|\textgreater{}} \FunctionTok{summary}\NormalTok{() ; ms}\SpecialCharTok{$}\NormalTok{adj.r.squared}
\end{Highlighting}
\end{Shaded}

\begin{verbatim}
[1] 0.0003339815
\end{verbatim}

Unsurprisingly, \texttt{pregnancy\_duration\_wks} is the variable with
the highest adjusted R\(^2\). It seems reasonable since even when they
are born infants grow very fast. In the next step we fit multiple linear
regression models with two predictors, the first one being
\texttt{pregnancy\_duration\_wks}.

\begin{Shaded}
\begin{Highlighting}[]
\NormalTok{ms }\OtherTok{\textless{}{-}}\NormalTok{ birth }\SpecialCharTok{|\textgreater{}} \FunctionTok{lm}\NormalTok{(}\AttributeTok{formula =}\NormalTok{ birth\_weight\_g }\SpecialCharTok{\textasciitilde{}}\NormalTok{ pregnancy\_duration\_wks }\SpecialCharTok{+}\NormalTok{ sex) }\SpecialCharTok{|\textgreater{}} \FunctionTok{summary}\NormalTok{() ; ms}\SpecialCharTok{$}\NormalTok{adj.r.squared}
\end{Highlighting}
\end{Shaded}

\begin{verbatim}
[1] 0.2710353
\end{verbatim}

\begin{Shaded}
\begin{Highlighting}[]
\NormalTok{ms }\OtherTok{\textless{}{-}}\NormalTok{ birth }\SpecialCharTok{|\textgreater{}} \FunctionTok{lm}\NormalTok{(}\AttributeTok{formula =}\NormalTok{ birth\_weight\_g }\SpecialCharTok{\textasciitilde{}}\NormalTok{ pregnancy\_duration\_wks }\SpecialCharTok{+}\NormalTok{ age) }\SpecialCharTok{|\textgreater{}} \FunctionTok{summary}\NormalTok{() ; ms}\SpecialCharTok{$}\NormalTok{adj.r.squared}
\end{Highlighting}
\end{Shaded}

\begin{verbatim}
[1] 0.2462937
\end{verbatim}

\begin{Shaded}
\begin{Highlighting}[]
\NormalTok{ms }\OtherTok{\textless{}{-}}\NormalTok{ birth }\SpecialCharTok{|\textgreater{}} \FunctionTok{lm}\NormalTok{(}\AttributeTok{formula =}\NormalTok{ birth\_weight\_g }\SpecialCharTok{\textasciitilde{}}\NormalTok{ pregnancy\_duration\_wks }\SpecialCharTok{+}\NormalTok{ length\_cm) }\SpecialCharTok{|\textgreater{}} \FunctionTok{summary}\NormalTok{() ; ms}\SpecialCharTok{$}\NormalTok{adj.r.squared}
\end{Highlighting}
\end{Shaded}

\begin{verbatim}
[1] 0.314511
\end{verbatim}

\begin{Shaded}
\begin{Highlighting}[]
\NormalTok{ms }\OtherTok{\textless{}{-}}\NormalTok{ birth }\SpecialCharTok{|\textgreater{}} \FunctionTok{lm}\NormalTok{(}\AttributeTok{formula =}\NormalTok{ birth\_weight\_g }\SpecialCharTok{\textasciitilde{}}\NormalTok{ pregnancy\_duration\_wks }\SpecialCharTok{+}\NormalTok{ weight\_kg) }\SpecialCharTok{|\textgreater{}} \FunctionTok{summary}\NormalTok{() ; ms}\SpecialCharTok{$}\NormalTok{adj.r.squared}
\end{Highlighting}
\end{Shaded}

\begin{verbatim}
[1] 0.2789052
\end{verbatim}

\begin{Shaded}
\begin{Highlighting}[]
\NormalTok{ms }\OtherTok{\textless{}{-}}\NormalTok{ birth }\SpecialCharTok{|\textgreater{}} \FunctionTok{lm}\NormalTok{(}\AttributeTok{formula =}\NormalTok{ birth\_weight\_g }\SpecialCharTok{\textasciitilde{}}\NormalTok{ pregnancy\_duration\_wks }\SpecialCharTok{+}\NormalTok{ bmi) }\SpecialCharTok{|\textgreater{}} \FunctionTok{summary}\NormalTok{() ; ms}\SpecialCharTok{$}\NormalTok{adj.r.squared}
\end{Highlighting}
\end{Shaded}

\begin{verbatim}
[1] 0.2464091
\end{verbatim}

\texttt{length\_cm} is now the variable that increases adjusted R\(^2\)
the most, so we include it in our model and continue by adding a third
variable.

\begin{Shaded}
\begin{Highlighting}[]
\NormalTok{ms }\OtherTok{\textless{}{-}}\NormalTok{ birth }\SpecialCharTok{|\textgreater{}} \FunctionTok{lm}\NormalTok{(}\AttributeTok{formula =}\NormalTok{ birth\_weight\_g }\SpecialCharTok{\textasciitilde{}}\NormalTok{ pregnancy\_duration\_wks }\SpecialCharTok{+}\NormalTok{ length\_cm }\SpecialCharTok{+}\NormalTok{ sex) }\SpecialCharTok{|\textgreater{}} \FunctionTok{summary}\NormalTok{() ; ms}\SpecialCharTok{$}\NormalTok{adj.r.squared}
\end{Highlighting}
\end{Shaded}

\begin{verbatim}
[1] 0.3365938
\end{verbatim}

\begin{Shaded}
\begin{Highlighting}[]
\NormalTok{ms }\OtherTok{\textless{}{-}}\NormalTok{ birth }\SpecialCharTok{|\textgreater{}} \FunctionTok{lm}\NormalTok{(}\AttributeTok{formula =}\NormalTok{ birth\_weight\_g }\SpecialCharTok{\textasciitilde{}}\NormalTok{ pregnancy\_duration\_wks }\SpecialCharTok{+}\NormalTok{ length\_cm }\SpecialCharTok{+}\NormalTok{ age) }\SpecialCharTok{|\textgreater{}} \FunctionTok{summary}\NormalTok{() ; ms}\SpecialCharTok{$}\NormalTok{adj.r.squared}
\end{Highlighting}
\end{Shaded}

\begin{verbatim}
[1] 0.3123593
\end{verbatim}

\begin{Shaded}
\begin{Highlighting}[]
\NormalTok{ms }\OtherTok{\textless{}{-}}\NormalTok{ birth }\SpecialCharTok{|\textgreater{}} \FunctionTok{lm}\NormalTok{(}\AttributeTok{formula =}\NormalTok{ birth\_weight\_g }\SpecialCharTok{\textasciitilde{}}\NormalTok{ pregnancy\_duration\_wks }\SpecialCharTok{+}\NormalTok{ length\_cm }\SpecialCharTok{+}\NormalTok{ weight\_kg) }\SpecialCharTok{|\textgreater{}} \FunctionTok{summary}\NormalTok{() ; ms}\SpecialCharTok{$}\NormalTok{adj.r.squared}
\end{Highlighting}
\end{Shaded}

\begin{verbatim}
[1] 0.3223792
\end{verbatim}

\begin{Shaded}
\begin{Highlighting}[]
\NormalTok{ms }\OtherTok{\textless{}{-}}\NormalTok{ birth }\SpecialCharTok{|\textgreater{}} \FunctionTok{lm}\NormalTok{(}\AttributeTok{formula =}\NormalTok{ birth\_weight\_g }\SpecialCharTok{\textasciitilde{}}\NormalTok{ pregnancy\_duration\_wks }\SpecialCharTok{+}\NormalTok{ length\_cm }\SpecialCharTok{+}\NormalTok{ bmi) }\SpecialCharTok{|\textgreater{}} \FunctionTok{summary}\NormalTok{() ; ms}\SpecialCharTok{$}\NormalTok{adj.r.squared}
\end{Highlighting}
\end{Shaded}

\begin{verbatim}
[1] 0.3219845
\end{verbatim}

\texttt{sex} is now the variable that increases adjusted R\(^2\) the
most, so we include it in our model and continue with by a fourth
variable.

\begin{Shaded}
\begin{Highlighting}[]
\NormalTok{ms }\OtherTok{\textless{}{-}}\NormalTok{ birth }\SpecialCharTok{|\textgreater{}} \FunctionTok{lm}\NormalTok{(}\AttributeTok{formula =}\NormalTok{ birth\_weight\_g }\SpecialCharTok{\textasciitilde{}}\NormalTok{ pregnancy\_duration\_wks }\SpecialCharTok{+}\NormalTok{ length\_cm }\SpecialCharTok{+}\NormalTok{ sex }\SpecialCharTok{+}\NormalTok{ age) }\SpecialCharTok{|\textgreater{}} \FunctionTok{summary}\NormalTok{() ; ms}\SpecialCharTok{$}\NormalTok{adj.r.squared}
\end{Highlighting}
\end{Shaded}

\begin{verbatim}
[1] 0.3341412
\end{verbatim}

\begin{Shaded}
\begin{Highlighting}[]
\NormalTok{ms }\OtherTok{\textless{}{-}}\NormalTok{ birth }\SpecialCharTok{|\textgreater{}} \FunctionTok{lm}\NormalTok{(}\AttributeTok{formula =}\NormalTok{ birth\_weight\_g }\SpecialCharTok{\textasciitilde{}}\NormalTok{ pregnancy\_duration\_wks }\SpecialCharTok{+}\NormalTok{ length\_cm }\SpecialCharTok{+}\NormalTok{ sex }\SpecialCharTok{+}\NormalTok{ weight\_kg) }\SpecialCharTok{|\textgreater{}} \FunctionTok{summary}\NormalTok{() ; ms}\SpecialCharTok{$}\NormalTok{adj.r.squared}
\end{Highlighting}
\end{Shaded}

\begin{verbatim}
[1] 0.3449951
\end{verbatim}

\begin{Shaded}
\begin{Highlighting}[]
\NormalTok{ms }\OtherTok{\textless{}{-}}\NormalTok{ birth }\SpecialCharTok{|\textgreater{}} \FunctionTok{lm}\NormalTok{(}\AttributeTok{formula =}\NormalTok{ birth\_weight\_g }\SpecialCharTok{\textasciitilde{}}\NormalTok{ pregnancy\_duration\_wks }\SpecialCharTok{+}\NormalTok{ length\_cm }\SpecialCharTok{+}\NormalTok{ sex }\SpecialCharTok{+}\NormalTok{ bmi) }\SpecialCharTok{|\textgreater{}} \FunctionTok{summary}\NormalTok{() ; ms}\SpecialCharTok{$}\NormalTok{adj.r.squared}
\end{Highlighting}
\end{Shaded}

\begin{verbatim}
[1] 0.3443753
\end{verbatim}

\texttt{weight\_kg} is now the variable that increases adjusted R\(^2\)
the most and is also statistically significant at significance level
0.05, so we include it in our model and continue by adding a fifth
variable.

\begin{Shaded}
\begin{Highlighting}[]
\NormalTok{ms }\OtherTok{\textless{}{-}}\NormalTok{ birth }\SpecialCharTok{|\textgreater{}} \FunctionTok{lm}\NormalTok{(}\AttributeTok{formula =}\NormalTok{ birth\_weight\_g }\SpecialCharTok{\textasciitilde{}}\NormalTok{ pregnancy\_duration\_wks }\SpecialCharTok{+}\NormalTok{ length\_cm }\SpecialCharTok{+}\NormalTok{ sex }\SpecialCharTok{+}\NormalTok{ weight\_kg }\SpecialCharTok{+}\NormalTok{ age) }\SpecialCharTok{|\textgreater{}} \FunctionTok{summary}\NormalTok{() ; ms}\SpecialCharTok{$}\NormalTok{adj.r.squared}
\end{Highlighting}
\end{Shaded}

\begin{verbatim}
[1] 0.3426022
\end{verbatim}

\begin{Shaded}
\begin{Highlighting}[]
\NormalTok{ms }\OtherTok{\textless{}{-}}\NormalTok{ birth }\SpecialCharTok{|\textgreater{}} \FunctionTok{lm}\NormalTok{(}\AttributeTok{formula =}\NormalTok{ birth\_weight\_g }\SpecialCharTok{\textasciitilde{}}\NormalTok{ pregnancy\_duration\_wks }\SpecialCharTok{+}\NormalTok{ length\_cm }\SpecialCharTok{+}\NormalTok{ sex }\SpecialCharTok{+}\NormalTok{ weight\_kg }\SpecialCharTok{+}\NormalTok{ bmi) }\SpecialCharTok{|\textgreater{}} \FunctionTok{summary}\NormalTok{() ; ms}\SpecialCharTok{$}\NormalTok{adj.r.squared}
\end{Highlighting}
\end{Shaded}

\begin{verbatim}
[1] 0.3435731
\end{verbatim}

Adding either \texttt{age} or \texttt{bmi} decreased our adjusted
R\(^2\) and we have reached our stopping criteria. Using forward
selection we found the same model as when we looked at multicollinearity
and removed \texttt{bmi} based on domain knowledge.

\begin{Shaded}
\begin{Highlighting}[]
\CommentTok{\# Forward selection (fs) model}
\NormalTok{model\_fs }\OtherTok{\textless{}{-}}\NormalTok{ birth }\SpecialCharTok{|\textgreater{}} \FunctionTok{lm}\NormalTok{(}\AttributeTok{formula =}\NormalTok{ birth\_weight\_g }\SpecialCharTok{\textasciitilde{}}\NormalTok{ pregnancy\_duration\_wks }\SpecialCharTok{+}\NormalTok{ length\_cm }\SpecialCharTok{+}\NormalTok{ sex }\SpecialCharTok{+}\NormalTok{ weight\_kg)}
\NormalTok{model\_fs }\SpecialCharTok{|\textgreater{}} \FunctionTok{summary}\NormalTok{()}
\end{Highlighting}
\end{Shaded}

\begin{verbatim}

Call:
lm(formula = birth_weight_g ~ pregnancy_duration_wks + length_cm + 
    sex + weight_kg, data = birth)

Residuals:
     Min       1Q   Median       3Q      Max 
-1000.23  -249.24    -6.36   249.70  1102.63 

Coefficients:
                        Estimate Std. Error t value Pr(>|t|)    
(Intercept)            -6504.850   1026.365  -6.338 1.26e-09 ***
pregnancy_duration_wks   175.099     20.167   8.683 7.97e-16 ***
length_cm                 15.842      4.163   3.806 0.000182 ***
sexMale                  149.730     50.465   2.967 0.003332 ** 
weight_kg                  4.800      2.431   1.975 0.049545 *  
---
Signif. codes:  0 '***' 0.001 '**' 0.01 '*' 0.05 '.' 0.1 ' ' 1

Residual standard error: 377 on 225 degrees of freedom
Multiple R-squared:  0.3564,    Adjusted R-squared:  0.345 
F-statistic: 31.15 on 4 and 225 DF,  p-value: < 2.2e-16
\end{verbatim}

\subsubsection{Exercise 2}\label{exercise-2-27}

We can use our results from the previous exercise to see if mixed
selection would have given us a different model. After each forward step
we check the \(p\)-values. If any variable is not significant we can
conclude that mixed selection would have given a different result.

\begin{Shaded}
\begin{Highlighting}[]
\CommentTok{\# First forward step}
\NormalTok{ms }\OtherTok{\textless{}{-}}\NormalTok{ birth }\SpecialCharTok{|\textgreater{}} \FunctionTok{lm}\NormalTok{(}\AttributeTok{formula =}\NormalTok{ birth\_weight\_g }\SpecialCharTok{\textasciitilde{}}\NormalTok{ pregnancy\_duration\_wks) }\SpecialCharTok{|\textgreater{}} \FunctionTok{summary}\NormalTok{() ; ms}\SpecialCharTok{$}\NormalTok{coefficients}
\end{Highlighting}
\end{Shaded}

\begin{verbatim}
                        Estimate Std. Error   t value     Pr(>|t|)
(Intercept)            -3867.089  852.13015 -4.538144 9.183294e-06
pregnancy_duration_wks   185.083   21.42084  8.640327 9.938670e-16
\end{verbatim}

\begin{Shaded}
\begin{Highlighting}[]
\CommentTok{\# Second forward step}
\NormalTok{ms }\OtherTok{\textless{}{-}}\NormalTok{ birth }\SpecialCharTok{|\textgreater{}} \FunctionTok{lm}\NormalTok{(}\AttributeTok{formula =}\NormalTok{ birth\_weight\_g }\SpecialCharTok{\textasciitilde{}}\NormalTok{ pregnancy\_duration\_wks }\SpecialCharTok{+}\NormalTok{ length\_cm) }\SpecialCharTok{|\textgreater{}} \FunctionTok{summary}\NormalTok{() ; ms}\SpecialCharTok{$}\NormalTok{coefficients}
\end{Highlighting}
\end{Shaded}

\begin{verbatim}
                         Estimate  Std. Error   t value     Pr(>|t|)
(Intercept)            -7075.6886 1036.945658 -6.823587 7.987029e-11
pregnancy_duration_wks   183.3006   20.392024  8.988839 9.983322e-17
length_cm                 19.6258    3.951779  4.966319 1.341270e-06
\end{verbatim}

\begin{Shaded}
\begin{Highlighting}[]
\CommentTok{\# Third forward step}
\NormalTok{ms }\OtherTok{\textless{}{-}}\NormalTok{ birth }\SpecialCharTok{|\textgreater{}} \FunctionTok{lm}\NormalTok{(}\AttributeTok{formula =}\NormalTok{ birth\_weight\_g }\SpecialCharTok{\textasciitilde{}}\NormalTok{ pregnancy\_duration\_wks }\SpecialCharTok{+}\NormalTok{ length\_cm }\SpecialCharTok{+}\NormalTok{ sex) }\SpecialCharTok{|\textgreater{}} \FunctionTok{summary}\NormalTok{() ; ms}\SpecialCharTok{$}\NormalTok{coefficients}
\end{Highlighting}
\end{Shaded}

\begin{verbatim}
                          Estimate  Std. Error   t value     Pr(>|t|)
(Intercept)            -6670.10127 1029.486946 -6.479054 5.704978e-10
pregnancy_duration_wks   174.36396   20.292184  8.592666 1.411882e-15
length_cm                 18.86122    3.896383  4.840700 2.397739e-06
sexMale                  148.54634   50.783535  2.925089 3.794448e-03
\end{verbatim}

\begin{Shaded}
\begin{Highlighting}[]
\CommentTok{\# Fourth forward step}
\NormalTok{ms }\OtherTok{\textless{}{-}}\NormalTok{ birth }\SpecialCharTok{|\textgreater{}} \FunctionTok{lm}\NormalTok{(}\AttributeTok{formula =}\NormalTok{ birth\_weight\_g }\SpecialCharTok{\textasciitilde{}}\NormalTok{ pregnancy\_duration\_wks }\SpecialCharTok{+}\NormalTok{ length\_cm }\SpecialCharTok{+}\NormalTok{ sex }\SpecialCharTok{+}\NormalTok{ weight\_kg) }\SpecialCharTok{|\textgreater{}} \FunctionTok{summary}\NormalTok{() ; ms}\SpecialCharTok{$}\NormalTok{coefficients}
\end{Highlighting}
\end{Shaded}

\begin{verbatim}
                           Estimate  Std. Error   t value     Pr(>|t|)
(Intercept)            -6504.849584 1026.365436 -6.337752 1.257868e-09
pregnancy_duration_wks   175.098880   20.166721  8.682566 7.969632e-16
length_cm                 15.841621    4.162720  3.805594 1.823702e-04
sexMale                  149.730071   50.464514  2.967037 3.331594e-03
weight_kg                  4.799832    2.430882  1.974523 4.954531e-02
\end{verbatim}

No variable became not significant after each forward step. We conclude
that mixed selection would have given us the same model as forward
selection in this case.

\subsubsection{Exercise 3}\label{exercise-3-21}

We first remember the variables in the \texttt{penguins} dataset.

\begin{Shaded}
\begin{Highlighting}[]
\FunctionTok{library}\NormalTok{(palmerpenguins)}
\FunctionTok{glimpse}\NormalTok{(penguins)}
\end{Highlighting}
\end{Shaded}

\begin{verbatim}
Rows: 344
Columns: 8
$ species           <fct> Adelie, Adelie, Adelie, Adelie, Adelie, Adelie, Adel~
$ island            <fct> Torgersen, Torgersen, Torgersen, Torgersen, Torgerse~
$ bill_length_mm    <dbl> 39.1, 39.5, 40.3, NA, 36.7, 39.3, 38.9, 39.2, 34.1, ~
$ bill_depth_mm     <dbl> 18.7, 17.4, 18.0, NA, 19.3, 20.6, 17.8, 19.6, 18.1, ~
$ flipper_length_mm <int> 181, 186, 195, NA, 193, 190, 181, 195, 193, 190, 186~
$ body_mass_g       <int> 3750, 3800, 3250, NA, 3450, 3650, 3625, 4675, 3475, ~
$ sex               <fct> male, female, female, NA, female, male, female, male~
$ year              <int> 2007, 2007, 2007, 2007, 2007, 2007, 2007, 2007, 2007~
\end{verbatim}

We then start by fitting a regression model using all potential
predictors.

\begin{Shaded}
\begin{Highlighting}[]
\NormalTok{penguins }\SpecialCharTok{|\textgreater{}} \FunctionTok{lm}\NormalTok{(}\AttributeTok{formula =}\NormalTok{ body\_mass\_g }\SpecialCharTok{\textasciitilde{}}\NormalTok{ .) }\SpecialCharTok{|\textgreater{}} \FunctionTok{summary}\NormalTok{()}
\end{Highlighting}
\end{Shaded}

\begin{verbatim}

Call:
lm(formula = body_mass_g ~ ., data = penguins)

Residuals:
    Min      1Q  Median      3Q     Max 
-809.70 -180.87   -6.25  176.76  864.22 

Coefficients:
                   Estimate Std. Error t value Pr(>|t|)    
(Intercept)       84087.945  41912.019   2.006  0.04566 *  
speciesChinstrap   -282.539     88.790  -3.182  0.00160 ** 
speciesGentoo       890.958    144.563   6.163 2.12e-09 ***
islandDream         -21.180     58.390  -0.363  0.71704    
islandTorgersen     -58.777     60.852  -0.966  0.33482    
bill_length_mm       18.964      7.112   2.667  0.00805 ** 
bill_depth_mm        60.798     20.002   3.040  0.00256 ** 
flipper_length_mm    18.504      3.128   5.915 8.46e-09 ***
sexmale             378.977     48.074   7.883 4.95e-14 ***
year                -42.785     20.949  -2.042  0.04194 *  
---
Signif. codes:  0 '***' 0.001 '**' 0.01 '*' 0.05 '.' 0.1 ' ' 1

Residual standard error: 286.5 on 323 degrees of freedom
  (11 observations deleted due to missingness)
Multiple R-squared:  0.8768,    Adjusted R-squared:  0.8734 
F-statistic: 255.4 on 9 and 323 DF,  p-value: < 2.2e-16
\end{verbatim}

We see that both coefficients relating to \texttt{island} are not
significant, so we remove \texttt{island} from our model.

\begin{Shaded}
\begin{Highlighting}[]
\NormalTok{penguins }\SpecialCharTok{|\textgreater{}} \FunctionTok{lm}\NormalTok{(}\AttributeTok{formula =}\NormalTok{ body\_mass\_g }\SpecialCharTok{\textasciitilde{}}\NormalTok{ . }\SpecialCharTok{{-}}\NormalTok{ island) }\SpecialCharTok{|\textgreater{}} \FunctionTok{summary}\NormalTok{() }
\end{Highlighting}
\end{Shaded}

\begin{verbatim}

Call:
lm(formula = body_mass_g ~ . - island, data = penguins)

Residuals:
    Min      1Q  Median      3Q     Max 
-809.15 -179.43   -3.42  183.60  866.03 

Coefficients:
                   Estimate Std. Error t value Pr(>|t|)    
(Intercept)       80838.770  41677.817   1.940 0.053292 .  
speciesChinstrap   -274.496     81.558  -3.366 0.000855 ***
speciesGentoo       929.275    136.036   6.831 4.16e-11 ***
bill_length_mm       18.997      7.086   2.681 0.007718 ** 
bill_depth_mm        60.749     19.926   3.049 0.002487 ** 
flipper_length_mm    18.064      3.088   5.849 1.20e-08 ***
sexmale             382.168     47.797   7.996 2.28e-14 ***
year                -41.139     20.832  -1.975 0.049131 *  
---
Signif. codes:  0 '***' 0.001 '**' 0.01 '*' 0.05 '.' 0.1 ' ' 1

Residual standard error: 286.1 on 325 degrees of freedom
  (11 observations deleted due to missingness)
Multiple R-squared:  0.8764,    Adjusted R-squared:  0.8738 
F-statistic: 329.3 on 7 and 325 DF,  p-value: < 2.2e-16
\end{verbatim}

\section{Prediction}\label{prediction-3}

\subsection{Exercise 1}\label{exercise-1-30}

We first note that \texttt{model\_fs} and \texttt{model\_de} are the
same models.

\begin{Shaded}
\begin{Highlighting}[]
\NormalTok{model\_de }\OtherTok{\textless{}{-}}\NormalTok{ birth }\SpecialCharTok{|\textgreater{}} 
  \FunctionTok{lm}\NormalTok{(}\AttributeTok{formula =}\NormalTok{ birth\_weight\_g }\SpecialCharTok{\textasciitilde{}}\NormalTok{ . }\SpecialCharTok{{-}}\NormalTok{ bmi }\SpecialCharTok{{-}}\NormalTok{ age)}
\NormalTok{prediction }\OtherTok{\textless{}{-}} \FunctionTok{predict}\NormalTok{(}
\NormalTok{  model\_de, }
  \FunctionTok{data.frame}\NormalTok{(}
    \AttributeTok{sex =} \StringTok{"Female"}\NormalTok{, }
    \AttributeTok{pregnancy\_duration\_wks =} \DecValTok{40}\NormalTok{, }
    \AttributeTok{weight\_kg =} \DecValTok{68}\NormalTok{, }
    \AttributeTok{bmi =} \DecValTok{24}\NormalTok{, }
    \AttributeTok{length\_cm =} \DecValTok{170}\NormalTok{, }
    \AttributeTok{age =} \DecValTok{0}\NormalTok{))}
\NormalTok{prediction}
\end{Highlighting}
\end{Shaded}

\begin{verbatim}
      1 
3518.57 
\end{verbatim}

We see that the prediction differs slightly compared to
\texttt{model\_bs} that predicted \(\sim 3492\) grams compared to
\texttt{model\_de}s 3518 g. Would you say that the difference is
relevant?

\subsection{Exercise 2}\label{exercise-2-28}

\begin{Shaded}
\begin{Highlighting}[]
\CommentTok{\# Prediction interval by backward selection model}
\NormalTok{model\_bs }\OtherTok{\textless{}{-}}\NormalTok{ birth }\SpecialCharTok{|\textgreater{}} \FunctionTok{lm}\NormalTok{(}\AttributeTok{formula =}\NormalTok{ birth\_weight\_g }\SpecialCharTok{\textasciitilde{}}\NormalTok{ pregnancy\_duration\_wks }\SpecialCharTok{+}\NormalTok{ sex }\SpecialCharTok{+}\NormalTok{ weight\_kg }\SpecialCharTok{+}\NormalTok{ bmi)}
\FunctionTok{predict}\NormalTok{(}
\NormalTok{  model\_bs, }
  \FunctionTok{data.frame}\NormalTok{(}
    \AttributeTok{sex =} \StringTok{"Female"}\NormalTok{, }
    \AttributeTok{pregnancy\_duration\_wks =} \DecValTok{40}\NormalTok{, }
    \AttributeTok{weight\_kg =} \DecValTok{68}\NormalTok{, }
    \AttributeTok{bmi =} \DecValTok{24}\NormalTok{),}
  \AttributeTok{interval =} \StringTok{"prediction"}\NormalTok{)}
\end{Highlighting}
\end{Shaded}

\begin{verbatim}
       fit     lwr      upr
1 3491.656 2745.76 4237.552
\end{verbatim}

\begin{Shaded}
\begin{Highlighting}[]
\CommentTok{\# Prediction interval by backward selection model}
\NormalTok{model\_de }\OtherTok{\textless{}{-}}\NormalTok{ birth }\SpecialCharTok{|\textgreater{}} 
  \FunctionTok{lm}\NormalTok{(}\AttributeTok{formula =}\NormalTok{ birth\_weight\_g }\SpecialCharTok{\textasciitilde{}}\NormalTok{ . }\SpecialCharTok{{-}}\NormalTok{ bmi }\SpecialCharTok{{-}}\NormalTok{ age)}
\FunctionTok{predict}\NormalTok{(}
\NormalTok{  model\_de, }
  \FunctionTok{data.frame}\NormalTok{(}
    \AttributeTok{sex =} \StringTok{"Female"}\NormalTok{, }
    \AttributeTok{pregnancy\_duration\_wks =} \DecValTok{40}\NormalTok{, }
    \AttributeTok{weight\_kg =} \DecValTok{68}\NormalTok{, }
    \AttributeTok{bmi =} \DecValTok{24}\NormalTok{, }
    \AttributeTok{length\_cm =} \DecValTok{170}\NormalTok{, }
    \AttributeTok{age =} \DecValTok{0}\NormalTok{),}
  \AttributeTok{interval =} \StringTok{"prediction"}\NormalTok{)}
\end{Highlighting}
\end{Shaded}

\begin{verbatim}
      fit      lwr      upr
1 3518.57 2771.581 4265.558
\end{verbatim}

The predictions overlap quite a lot. In this case the models seem very
similar and it shouldn't matter much which one we use, at least for this
particular prediction.

\chapter{Logistic regression}\label{logistic-regression-1}

\section{Setup}\label{setup-3}

First we load the dataset and have a \texttt{glimpse} at it.

\begin{Shaded}
\begin{Highlighting}[]
\FunctionTok{library}\NormalTok{(tidyverse)}
\CommentTok{\#install.packages("mlbench") \# Install to get PimaIndiansDiabetes2 dataset}
\FunctionTok{library}\NormalTok{(mlbench)}
\FunctionTok{data}\NormalTok{(}\StringTok{"PimaIndiansDiabetes2"}\NormalTok{) }\CommentTok{\# Using 2 because it is a corrected version.}
\FunctionTok{glimpse}\NormalTok{(PimaIndiansDiabetes2)}
\end{Highlighting}
\end{Shaded}

\begin{verbatim}
Rows: 768
Columns: 9
$ pregnant <dbl> 6, 1, 8, 1, 0, 5, 3, 10, 2, 8, 4, 10, 10, 1, 5, 7, 0, 7, 1, 1~
$ glucose  <dbl> 148, 85, 183, 89, 137, 116, 78, 115, 197, 125, 110, 168, 139,~
$ pressure <dbl> 72, 66, 64, 66, 40, 74, 50, NA, 70, 96, 92, 74, 80, 60, 72, N~
$ triceps  <dbl> 35, 29, NA, 23, 35, NA, 32, NA, 45, NA, NA, NA, NA, 23, 19, N~
$ insulin  <dbl> NA, NA, NA, 94, 168, NA, 88, NA, 543, NA, NA, NA, NA, 846, 17~
$ mass     <dbl> 33.6, 26.6, 23.3, 28.1, 43.1, 25.6, 31.0, 35.3, 30.5, NA, 37.~
$ pedigree <dbl> 0.627, 0.351, 0.672, 0.167, 2.288, 0.201, 0.248, 0.134, 0.158~
$ age      <dbl> 50, 31, 32, 21, 33, 30, 26, 29, 53, 54, 30, 34, 57, 59, 51, 3~
$ diabetes <fct> pos, neg, pos, neg, pos, neg, pos, neg, pos, pos, neg, pos, n~
\end{verbatim}

\section{The Pima Indians dataset}\label{the-pima-indians-dataset-1}

\subsection{Exercise 1}\label{exercise-1-31}

Our population is all women of the Pima Indian population aged 21 or
older. Looking at the original paper, we can see that the study started
in 1965, which we may also want to specify.

\subsection{Exercise 2}\label{exercise-2-29}

We see that the blood pressure seems higher in the group that were
diagnosed with diabetes. In the next chapter we will see if it is a
significant relationship.

\begin{Shaded}
\begin{Highlighting}[]
\FunctionTok{library}\NormalTok{(patchwork)}
\NormalTok{p1 }\OtherTok{\textless{}{-}}\NormalTok{ PimaIndiansDiabetes2 }\SpecialCharTok{|\textgreater{}}
  \FunctionTok{ggplot}\NormalTok{(}\FunctionTok{aes}\NormalTok{(}\AttributeTok{y =}\NormalTok{ pressure)) }\SpecialCharTok{+} 
  \FunctionTok{facet\_grid}\NormalTok{(}\SpecialCharTok{\textasciitilde{}}\NormalTok{ diabetes) }\SpecialCharTok{+}
  \FunctionTok{geom\_boxplot}\NormalTok{() }\SpecialCharTok{+}
  \FunctionTok{theme}\NormalTok{(}
    \AttributeTok{axis.text.x =} \FunctionTok{element\_blank}\NormalTok{(),  }\CommentTok{\# Remove axis text}
    \AttributeTok{axis.ticks.x =} \FunctionTok{element\_blank}\NormalTok{()  }\CommentTok{\# Remove axis title}
\NormalTok{  ) }\SpecialCharTok{+}
  \FunctionTok{labs}\NormalTok{(}\AttributeTok{y =} \StringTok{"Blood pressure"}\NormalTok{)}
\NormalTok{p2 }\OtherTok{\textless{}{-}}\NormalTok{ PimaIndiansDiabetes2 }\SpecialCharTok{|\textgreater{}}
  \FunctionTok{ggplot}\NormalTok{(}\FunctionTok{aes}\NormalTok{(}\AttributeTok{x =}\NormalTok{ pressure, }\AttributeTok{fill =}\NormalTok{ diabetes)) }\SpecialCharTok{+} 
  \FunctionTok{geom\_density}\NormalTok{(}\AttributeTok{alpha =} \FloatTok{0.5}\NormalTok{) }\SpecialCharTok{+}
  \FunctionTok{theme}\NormalTok{(}
    \AttributeTok{legend.position =} \StringTok{"top"}\NormalTok{,}
    \AttributeTok{axis.text.y =} \FunctionTok{element\_blank}\NormalTok{(),  }\CommentTok{\# Remove axis text}
    \AttributeTok{axis.ticks.y =} \FunctionTok{element\_blank}\NormalTok{()  }\CommentTok{\# Remove axis title}
\NormalTok{  ) }\SpecialCharTok{+}
  \FunctionTok{labs}\NormalTok{(}
    \AttributeTok{x =} \StringTok{"Blood pressure"}\NormalTok{,}
    \AttributeTok{y =} \StringTok{""}\NormalTok{,}
    \AttributeTok{fill =} \StringTok{"Diabetes"}
\NormalTok{    )}

\NormalTok{p1 }\SpecialCharTok{+}\NormalTok{ p2}
\end{Highlighting}
\end{Shaded}

\begin{verbatim}
Warning: Removed 35 rows containing non-finite outside the scale range
(`stat_boxplot()`).
\end{verbatim}

\begin{verbatim}
Warning: Removed 35 rows containing non-finite outside the scale range
(`stat_density()`).
\end{verbatim}

\includegraphics{solutions_logistic_regression_files/figure-pdf/unnamed-chunk-2-1.pdf}

\section{Using linear regression}\label{using-linear-regression-1}

\subsection{Exercise 1}\label{exercise-1-32}

First we formulate our hypotheses:

\begin{quote}
\textbf{H\_0}: There is \emph{no} linear relationship between mass and
blood pressure in the women in the Pima Indian tribe.
\end{quote}

\begin{quote}
\textbf{H\_a}: There \emph{is} a linear relationship between mass and
blood pressure in the women in the Pima Indian tribe.
\end{quote}

Then we fit our model.

\begin{Shaded}
\begin{Highlighting}[]
\NormalTok{model }\OtherTok{\textless{}{-}} \FunctionTok{lm}\NormalTok{(pressure }\SpecialCharTok{\textasciitilde{}}\NormalTok{ mass, }\AttributeTok{data =}\NormalTok{ PimaIndiansDiabetes2)}
\FunctionTok{summary}\NormalTok{(model)}
\end{Highlighting}
\end{Shaded}

\begin{verbatim}

Call:
lm(formula = pressure ~ mass, data = PimaIndiansDiabetes2)

Residuals:
    Min      1Q  Median      3Q     Max 
-54.081  -7.628  -0.331   7.262  54.868 

Coefficients:
            Estimate Std. Error t value Pr(>|t|)    
(Intercept) 55.48694    2.11810  26.197  < 2e-16 ***
mass         0.51989    0.06382   8.147 1.63e-15 ***
---
Signif. codes:  0 '***' 0.001 '**' 0.01 '*' 0.05 '.' 0.1 ' ' 1

Residual standard error: 11.86 on 727 degrees of freedom
  (39 observations deleted due to missingness)
Multiple R-squared:  0.08365,   Adjusted R-squared:  0.08239 
F-statistic: 66.37 on 1 and 727 DF,  p-value: 1.63e-15
\end{verbatim}

From the summary output we can see that the \(p\)-value is very small.
According to our model there is a linear relationship between mass and
blood pressure. We can also visualize this to see that our results are
reasonable.

\begin{Shaded}
\begin{Highlighting}[]
\NormalTok{PimaIndiansDiabetes2 }\SpecialCharTok{|\textgreater{}} 
  \FunctionTok{ggplot}\NormalTok{(}\FunctionTok{aes}\NormalTok{(}\AttributeTok{x =}\NormalTok{ mass, }\AttributeTok{y =}\NormalTok{ pressure)) }\SpecialCharTok{+}
  \FunctionTok{geom\_point}\NormalTok{() }\SpecialCharTok{+}
  \FunctionTok{geom\_smooth}\NormalTok{(}\AttributeTok{method =} \StringTok{"lm"}\NormalTok{, }\AttributeTok{se =} \ConstantTok{FALSE}\NormalTok{)}
\end{Highlighting}
\end{Shaded}

\begin{verbatim}
`geom_smooth()` using formula = 'y ~ x'
\end{verbatim}

\begin{verbatim}
Warning: Removed 39 rows containing non-finite outside the scale range
(`stat_smooth()`).
\end{verbatim}

\begin{verbatim}
Warning: Removed 39 rows containing missing values or values outside the scale range
(`geom_point()`).
\end{verbatim}

\includegraphics{solutions_logistic_regression_files/figure-pdf/unnamed-chunk-4-1.pdf}

We see a clear slope of the line.

\subsection{Exercise 2}\label{exercise-2-30}

We formulate the hypotheses:

\begin{quote}
\textbf{H\_0}: There is \emph{no} linear relationship between glucose
and the onset of diabetes in the women in the Pima Indian tribe.
\end{quote}

\begin{quote}
\textbf{H\_a}: There \emph{is} a linear relationship between glucose
levels and the onset of diabetes in the women in the Pima Indian tribe.
\end{quote}

As discussed in the chapter, a \emph{linear} relationship is not right
for this problem, since the outcome is binary.

\section{Using logistic regression}\label{using-logistic-regression-1}

\subsection{Exercise 1 (a)}\label{exercise-1-a}

We start by formulating the hypotheses

\begin{quote}
\textbf{H\_0}: There is \emph{no} relationship between blood pressure
and being diagnosed with diabetes.
\end{quote}

\begin{quote}
\textbf{H\_a}: There \emph{is} a relationship between blood pressure
level and being diagnosed with diabetes.
\end{quote}

then fit our model

\begin{Shaded}
\begin{Highlighting}[]
\NormalTok{model }\OtherTok{\textless{}{-}} \FunctionTok{glm}\NormalTok{(diabetes }\SpecialCharTok{\textasciitilde{}}\NormalTok{ pressure, }\AttributeTok{data =}\NormalTok{ PimaIndiansDiabetes2, }\AttributeTok{family =} \StringTok{"binomial"}\NormalTok{)}
\FunctionTok{summary}\NormalTok{(model)}
\end{Highlighting}
\end{Shaded}

\begin{verbatim}

Call:
glm(formula = diabetes ~ pressure, family = "binomial", data = PimaIndiansDiabetes2)

Coefficients:
             Estimate Std. Error z value Pr(>|z|)    
(Intercept) -2.830259   0.491750  -5.755 8.64e-09 ***
pressure     0.029884   0.006587   4.537 5.72e-06 ***
---
Signif. codes:  0 '***' 0.001 '**' 0.01 '*' 0.05 '.' 0.1 ' ' 1

(Dispersion parameter for binomial family taken to be 1)

    Null deviance: 943.40  on 732  degrees of freedom
Residual deviance: 921.78  on 731  degrees of freedom
  (35 observations deleted due to missingness)
AIC: 925.78

Number of Fisher Scoring iterations: 4
\end{verbatim}

We see that the \(p\)-value is low and that there is a significant
relationship between blood pressure and diabetes at significance level
0.05.

\subsection{Exercise 1 (b)}\label{exercise-1-b}

We start by formulating the hypotheses

\begin{quote}
\textbf{H\_0}: There is \emph{no} relationship between body mass index
and being diagnosed with diabetes.
\end{quote}

\begin{quote}
\textbf{H\_a}: There \emph{is} a relationship between body mass index
level and being diagnosed with diabetes.
\end{quote}

then fit our model

\begin{Shaded}
\begin{Highlighting}[]
\NormalTok{model }\OtherTok{\textless{}{-}} \FunctionTok{glm}\NormalTok{(diabetes }\SpecialCharTok{\textasciitilde{}}\NormalTok{ mass, }\AttributeTok{data =}\NormalTok{ PimaIndiansDiabetes2, }\AttributeTok{family =} \StringTok{"binomial"}\NormalTok{)}
\FunctionTok{summary}\NormalTok{(model)}
\end{Highlighting}
\end{Shaded}

\begin{verbatim}

Call:
glm(formula = diabetes ~ mass, family = "binomial", data = PimaIndiansDiabetes2)

Coefficients:
            Estimate Std. Error z value Pr(>|z|)    
(Intercept) -3.99682    0.42885   -9.32  < 2e-16 ***
mass         0.10250    0.01261    8.13 4.31e-16 ***
---
Signif. codes:  0 '***' 0.001 '**' 0.01 '*' 0.05 '.' 0.1 ' ' 1

(Dispersion parameter for binomial family taken to be 1)

    Null deviance: 981.53  on 756  degrees of freedom
Residual deviance: 904.89  on 755  degrees of freedom
  (11 observations deleted due to missingness)
AIC: 908.89

Number of Fisher Scoring iterations: 4
\end{verbatim}

We see that the \(p\)-value is low and that there is a significant
relationship between BMI and diabetes at significance level 0.05.

\section{Prediction}\label{prediction-4}

\subsection{Exercise 1}\label{exercise-1-33}

Here are two examples of predicting the probability of being diagnosed
with diabetes based on \texttt{pressure} and \texttt{mass}

\begin{Shaded}
\begin{Highlighting}[]
\NormalTok{pressure\_model }\OtherTok{\textless{}{-}} \FunctionTok{glm}\NormalTok{(diabetes }\SpecialCharTok{\textasciitilde{}}\NormalTok{ pressure, }\AttributeTok{data =}\NormalTok{ PimaIndiansDiabetes2, }\AttributeTok{family =} \StringTok{"binomial"}\NormalTok{)}
\FunctionTok{predict}\NormalTok{(pressure\_model, }\FunctionTok{data.frame}\NormalTok{(}\AttributeTok{pressure =} \DecValTok{90}\NormalTok{), }\AttributeTok{type =} \StringTok{"response"}\NormalTok{)}
\end{Highlighting}
\end{Shaded}

\begin{verbatim}
        1 
0.4648821 
\end{verbatim}

\begin{Shaded}
\begin{Highlighting}[]
\NormalTok{mass\_model }\OtherTok{\textless{}{-}} \FunctionTok{glm}\NormalTok{(diabetes }\SpecialCharTok{\textasciitilde{}}\NormalTok{ mass, }\AttributeTok{data =}\NormalTok{ PimaIndiansDiabetes2, }\AttributeTok{family =} \StringTok{"binomial"}\NormalTok{)}
\FunctionTok{predict}\NormalTok{(mass\_model, }\FunctionTok{data.frame}\NormalTok{(}\AttributeTok{mass =} \DecValTok{25}\NormalTok{), }\AttributeTok{type =} \StringTok{"response"}\NormalTok{)}
\end{Highlighting}
\end{Shaded}

\begin{verbatim}
        1 
0.1924252 
\end{verbatim}

\subsection{Exercise 2}\label{exercise-2-31}

We try to find the x value (glucose level) that corresponds to 0.9
probability of being diagnosed with diabetes. According to the plot this
roughly coincides with 195 mg/dl. Our prediction confirms this.

\begin{Shaded}
\begin{Highlighting}[]
\NormalTok{model }\OtherTok{\textless{}{-}} \FunctionTok{glm}\NormalTok{(diabetes }\SpecialCharTok{\textasciitilde{}}\NormalTok{ glucose, }\AttributeTok{data =}\NormalTok{ PimaIndiansDiabetes2, }\AttributeTok{family =} \StringTok{"binomial"}\NormalTok{)}
\FunctionTok{predict}\NormalTok{(model, }\FunctionTok{data.frame}\NormalTok{(}\AttributeTok{glucose =} \DecValTok{195}\NormalTok{), }\AttributeTok{type =} \StringTok{"response"}\NormalTok{)}
\end{Highlighting}
\end{Shaded}

\begin{verbatim}
        1 
0.9010074 
\end{verbatim}

\subsection{Exercise 3}\label{exercise-3-22}

Our model predicts a \(\sim0.92\) probability of the onset of diabetes
given a glucose level of 200 mg/dl. Our model was not trained on any
data of women with diabetes, i.e.~the model was not trained on any
glucose levels 200 or above and is not suited for predictions outside of
what was observed in the dataset. Further, the model doesn't ``know''
about the classification of diabetes, i.e.~that glucose levels over 200
are classified as diabetes.

\begin{Shaded}
\begin{Highlighting}[]
\NormalTok{model }\OtherTok{\textless{}{-}} \FunctionTok{glm}\NormalTok{(diabetes }\SpecialCharTok{\textasciitilde{}}\NormalTok{ glucose, }\AttributeTok{data =}\NormalTok{ PimaIndiansDiabetes2, }\AttributeTok{family =} \StringTok{"binomial"}\NormalTok{)}
\FunctionTok{predict}\NormalTok{(model, }\FunctionTok{data.frame}\NormalTok{(}\AttributeTok{glucose =} \DecValTok{200}\NormalTok{), }\AttributeTok{type =} \StringTok{"response"}\NormalTok{)}
\end{Highlighting}
\end{Shaded}

\begin{verbatim}
        1 
0.9177104 
\end{verbatim}

\chapter{Multiple logistic
regression}\label{multiple-logistic-regression-1}

\section{Setup}\label{setup-4}

First we load the dataset and have a \texttt{glimpse} at it.

\begin{Shaded}
\begin{Highlighting}[]
\FunctionTok{library}\NormalTok{(tidyverse)}
\CommentTok{\#install.packages("mlbench") \# Install to get PimaIndiansDiabetes2 dataset}
\FunctionTok{library}\NormalTok{(mlbench)}
\FunctionTok{data}\NormalTok{(}\StringTok{"PimaIndiansDiabetes2"}\NormalTok{) }\CommentTok{\# Using 2 because it is a corrected version.}
\FunctionTok{glimpse}\NormalTok{(PimaIndiansDiabetes2)}
\end{Highlighting}
\end{Shaded}

\begin{verbatim}
Rows: 768
Columns: 9
$ pregnant <dbl> 6, 1, 8, 1, 0, 5, 3, 10, 2, 8, 4, 10, 10, 1, 5, 7, 0, 7, 1, 1~
$ glucose  <dbl> 148, 85, 183, 89, 137, 116, 78, 115, 197, 125, 110, 168, 139,~
$ pressure <dbl> 72, 66, 64, 66, 40, 74, 50, NA, 70, 96, 92, 74, 80, 60, 72, N~
$ triceps  <dbl> 35, 29, NA, 23, 35, NA, 32, NA, 45, NA, NA, NA, NA, 23, 19, N~
$ insulin  <dbl> NA, NA, NA, 94, 168, NA, 88, NA, 543, NA, NA, NA, NA, 846, 17~
$ mass     <dbl> 33.6, 26.6, 23.3, 28.1, 43.1, 25.6, 31.0, 35.3, 30.5, NA, 37.~
$ pedigree <dbl> 0.627, 0.351, 0.672, 0.167, 2.288, 0.201, 0.248, 0.134, 0.158~
$ age      <dbl> 50, 31, 32, 21, 33, 30, 26, 29, 53, 54, 30, 34, 57, 59, 51, 3~
$ diabetes <fct> pos, neg, pos, neg, pos, neg, pos, neg, pos, pos, neg, pos, n~
\end{verbatim}

\section{Model selection}\label{model-selection-3}

\subsection{Exercise 1}\label{exercise-1-34}

We fit the two models and use the \texttt{AIC} function

\begin{Shaded}
\begin{Highlighting}[]
\CommentTok{\# Pressure model}
\NormalTok{model }\OtherTok{\textless{}{-}} \FunctionTok{glm}\NormalTok{(diabetes }\SpecialCharTok{\textasciitilde{}}\NormalTok{ pressure, }\AttributeTok{data =}\NormalTok{ PimaIndiansDiabetes2, }\AttributeTok{family =} \StringTok{"binomial"}\NormalTok{)}
\FunctionTok{AIC}\NormalTok{(model)}
\end{Highlighting}
\end{Shaded}

\begin{verbatim}
[1] 925.7838
\end{verbatim}

\begin{Shaded}
\begin{Highlighting}[]
\CommentTok{\# Triceps model}
\NormalTok{model }\OtherTok{\textless{}{-}} \FunctionTok{glm}\NormalTok{(diabetes }\SpecialCharTok{\textasciitilde{}}\NormalTok{ triceps, }\AttributeTok{data =}\NormalTok{ PimaIndiansDiabetes2, }\AttributeTok{family =} \StringTok{"binomial"}\NormalTok{)}
\FunctionTok{AIC}\NormalTok{(model)}
\end{Highlighting}
\end{Shaded}

\begin{verbatim}
[1] 654.857
\end{verbatim}

We see that \texttt{triceps} has much lower AIC and comparing to the
result in the lecture notes including \texttt{pressure} barely improves
the model!

\subsection{Exercise 2}\label{exercise-2-32}

We fit the model with all three predictors and look at the AIC value and
the VIF value

\begin{Shaded}
\begin{Highlighting}[]
\NormalTok{model }\OtherTok{\textless{}{-}} \FunctionTok{glm}\NormalTok{(diabetes }\SpecialCharTok{\textasciitilde{}}\NormalTok{ pressure }\SpecialCharTok{+}\NormalTok{ mass }\SpecialCharTok{+}\NormalTok{ triceps, }\AttributeTok{data =}\NormalTok{ PimaIndiansDiabetes2, }\AttributeTok{family =} \StringTok{"binomial"}\NormalTok{)}
\FunctionTok{summary}\NormalTok{(model)}
\end{Highlighting}
\end{Shaded}

\begin{verbatim}

Call:
glm(formula = diabetes ~ pressure + mass + triceps, family = "binomial", 
    data = PimaIndiansDiabetes2)

Coefficients:
             Estimate Std. Error z value Pr(>|z|)    
(Intercept) -5.105936   0.719401  -7.097 1.27e-12 ***
pressure     0.018727   0.008321   2.251 0.024416 *  
mass         0.071762   0.019412   3.697 0.000218 ***
triceps      0.021741   0.012031   1.807 0.070760 .  
---
Signif. codes:  0 '***' 0.001 '**' 0.01 '*' 0.05 '.' 0.1 ' ' 1

(Dispersion parameter for binomial family taken to be 1)

    Null deviance: 683.62  on 536  degrees of freedom
Residual deviance: 623.13  on 533  degrees of freedom
  (231 observations deleted due to missingness)
AIC: 631.13

Number of Fisher Scoring iterations: 3
\end{verbatim}

We see that the AIC is lower than using only \texttt{pressure} and
\texttt{triceps}, but the \(p\)-value for \texttt{triceps} is now over
0.05, indicating that the relationship is not statistically significant
at significance level 0.05. However, the \(p\)-value is low so we can
keep it in our model, but it may be reasonable to use only one of BMI
and tricep skin fold thickness as both are measurements relating to body
fat.

We also calculate the VIF value to see if we have multicollinearity

\begin{Shaded}
\begin{Highlighting}[]
\CommentTok{\# warning: false}
\FunctionTok{library}\NormalTok{(car)}
\end{Highlighting}
\end{Shaded}

\begin{verbatim}
Loading required package: carData
\end{verbatim}

\begin{verbatim}

Attaching package: 'car'
\end{verbatim}

\begin{verbatim}
The following object is masked from 'package:dplyr':

    recode
\end{verbatim}

\begin{verbatim}
The following object is masked from 'package:purrr':

    some
\end{verbatim}

\begin{Shaded}
\begin{Highlighting}[]
\FunctionTok{vif}\NormalTok{(model)}
\end{Highlighting}
\end{Shaded}

\begin{verbatim}
pressure     mass  triceps 
1.054880 1.594327 1.549618 
\end{verbatim}

We see that all values are below 5, so we do not need to remove any
variable because of multicollinearity.

\subsection{Model selection
strategies}\label{model-selection-strategies-2}

\subsubsection{Exercise 1}\label{exercise-1-35}

See Assignment 2.

\section{Accuracy}\label{accuracy-1}

\subsection{Exercise 1}\label{exercise-1-36}

We see that the accuracy is the same, with or without \texttt{pressure}.
In this case the simpler model, not including \texttt{pressure} is the
better option.

\begin{Shaded}
\begin{Highlighting}[]
\NormalTok{model }\OtherTok{\textless{}{-}}\NormalTok{ PimaIndiansDiabetes2 }\SpecialCharTok{|\textgreater{}} 
  \FunctionTok{glm}\NormalTok{(}\AttributeTok{formula =}\NormalTok{ diabetes }\SpecialCharTok{\textasciitilde{}}\NormalTok{ ., }\AttributeTok{family =} \StringTok{"binomial"}\NormalTok{)}

\CommentTok{\# Calculating the accuracy}
\NormalTok{PimaIndiansDiabetes2 }\SpecialCharTok{|\textgreater{}} 
  \CommentTok{\# Calculated the predicted probabilities for all observations}
  \FunctionTok{mutate}\NormalTok{(}\AttributeTok{predicted\_prob =}\NormalTok{ model }\SpecialCharTok{|\textgreater{}} \FunctionTok{predict}\NormalTok{(PimaIndiansDiabetes2, }\AttributeTok{type =} \StringTok{"response"}\NormalTok{)) }\SpecialCharTok{|\textgreater{}} 
  \CommentTok{\# Predicting the outcome}
  \FunctionTok{mutate}\NormalTok{(}\AttributeTok{predicted\_outcome =} \FunctionTok{ifelse}\NormalTok{(predicted\_prob }\SpecialCharTok{\textgreater{}} \FloatTok{0.5}\NormalTok{, }\StringTok{"pos"}\NormalTok{, }\StringTok{"neg"}\NormalTok{)) }\SpecialCharTok{|\textgreater{}} 
  \CommentTok{\# Dropping predictions with missing values}
  \FunctionTok{filter}\NormalTok{(}\SpecialCharTok{!}\FunctionTok{is.na}\NormalTok{(predicted\_outcome)) }\SpecialCharTok{|\textgreater{}}
  \CommentTok{\# Summarizing (Accuracy = \# Correct predictions / \# Total predictions}
  \FunctionTok{summarize}\NormalTok{(}\AttributeTok{accuracy =} \FunctionTok{mean}\NormalTok{(predicted\_outcome }\SpecialCharTok{==}\NormalTok{ diabetes))}
\end{Highlighting}
\end{Shaded}

\begin{verbatim}
   accuracy
1 0.7831633
\end{verbatim}

\subsection{Exercise 2}\label{exercise-2-33}

We leave this one to you.

\subsection{Exercise 3}\label{exercise-3-23}

Using this seed the accuracy is the same.

\begin{Shaded}
\begin{Highlighting}[]
\CommentTok{\# code{-}fold: true}
\CommentTok{\# comment: Code}
\CommentTok{\# Set seed for reproducibility}
\FunctionTok{set.seed}\NormalTok{(}\DecValTok{42}\NormalTok{) }

\CommentTok{\# Create an index to keep track of subjects}
\NormalTok{pima }\OtherTok{\textless{}{-}}\NormalTok{ PimaIndiansDiabetes2 }\SpecialCharTok{|\textgreater{}} 
  \FunctionTok{drop\_na}\NormalTok{() }\SpecialCharTok{|\textgreater{}} \CommentTok{\# Using only complete observations}
  \FunctionTok{mutate}\NormalTok{(}\AttributeTok{index =} \DecValTok{1}\SpecialCharTok{:}\FunctionTok{n}\NormalTok{())}

\CommentTok{\# Create training and testing set}
\NormalTok{training\_set }\OtherTok{\textless{}{-}}\NormalTok{ pima }\SpecialCharTok{|\textgreater{}} 
  \FunctionTok{sample\_frac}\NormalTok{(}\FloatTok{0.8}\NormalTok{) }\CommentTok{\# Use 80\% of the data in training set}
\NormalTok{test\_set }\OtherTok{\textless{}{-}}\NormalTok{ pima }\SpecialCharTok{|\textgreater{}} 
  \FunctionTok{anti\_join}\NormalTok{(training\_set, }\AttributeTok{by =} \StringTok{"index"}\NormalTok{)}

\CommentTok{\# Train model on training set}
\NormalTok{model }\OtherTok{\textless{}{-}}\NormalTok{ training\_set }\SpecialCharTok{|\textgreater{}} 
  \FunctionTok{glm}\NormalTok{(}\AttributeTok{formula =}\NormalTok{ diabetes }\SpecialCharTok{\textasciitilde{}}\NormalTok{ ., }\AttributeTok{family =} \StringTok{"binomial"}\NormalTok{)}

\CommentTok{\# Evaluate model by accuracy on test set}
\NormalTok{predictions }\OtherTok{\textless{}{-}}\NormalTok{ model }\SpecialCharTok{|\textgreater{}} 
  \FunctionTok{predict}\NormalTok{(test\_set, }\AttributeTok{type =} \StringTok{"response"}\NormalTok{)}
\NormalTok{test\_set }\SpecialCharTok{|\textgreater{}} 
  \FunctionTok{mutate}\NormalTok{(}\AttributeTok{prediction =} \FunctionTok{ifelse}\NormalTok{(predictions }\SpecialCharTok{\textgreater{}} \FloatTok{0.5}\NormalTok{, }\StringTok{"pos"}\NormalTok{, }\StringTok{"neg"}\NormalTok{)) }\SpecialCharTok{|\textgreater{}} 
  \FunctionTok{drop\_na}\NormalTok{() }\SpecialCharTok{|\textgreater{}} \CommentTok{\# Dropping predictions with missing values}
  \FunctionTok{summarize}\NormalTok{(}\AttributeTok{accuracy =} \FunctionTok{sum}\NormalTok{(prediction }\SpecialCharTok{==}\NormalTok{ diabetes) }\SpecialCharTok{/} \FunctionTok{n}\NormalTok{())}
\end{Highlighting}
\end{Shaded}

\begin{verbatim}
   accuracy
1 0.8076923
\end{verbatim}

\begin{Shaded}
\begin{Highlighting}[]
\NormalTok{model }\OtherTok{\textless{}{-}}\NormalTok{ training\_set }\SpecialCharTok{|\textgreater{}} 
  \FunctionTok{glm}\NormalTok{(}\AttributeTok{formula =}\NormalTok{ diabetes }\SpecialCharTok{\textasciitilde{}}\NormalTok{ . }\SpecialCharTok{{-}}\NormalTok{ pressure, }\AttributeTok{family =} \StringTok{"binomial"}\NormalTok{)}

\CommentTok{\# Evaluate model by accuracy on test set}
\NormalTok{predictions }\OtherTok{\textless{}{-}}\NormalTok{ model }\SpecialCharTok{|\textgreater{}} 
  \FunctionTok{predict}\NormalTok{(test\_set, }\AttributeTok{type =} \StringTok{"response"}\NormalTok{)}
\NormalTok{test\_set }\SpecialCharTok{|\textgreater{}} 
  \FunctionTok{mutate}\NormalTok{(}\AttributeTok{prediction =} \FunctionTok{ifelse}\NormalTok{(predictions }\SpecialCharTok{\textgreater{}} \FloatTok{0.5}\NormalTok{, }\StringTok{"pos"}\NormalTok{, }\StringTok{"neg"}\NormalTok{)) }\SpecialCharTok{|\textgreater{}} 
  \FunctionTok{drop\_na}\NormalTok{() }\SpecialCharTok{|\textgreater{}} \CommentTok{\# Dropping predictions with missing values}
  \FunctionTok{summarize}\NormalTok{(}\AttributeTok{accuracy =} \FunctionTok{mean}\NormalTok{(prediction }\SpecialCharTok{==}\NormalTok{ diabetes))}
\end{Highlighting}
\end{Shaded}

\begin{verbatim}
   accuracy
1 0.8076923
\end{verbatim}

With a different seed the full model has a higher accuracy.

\begin{Shaded}
\begin{Highlighting}[]
\CommentTok{\# code{-}fold: true}
\CommentTok{\# comment: Code}
\CommentTok{\# Set seed for reproducibility}
\FunctionTok{set.seed}\NormalTok{(}\DecValTok{44}\NormalTok{) }

\CommentTok{\# Create an index to keep track of subjects}
\NormalTok{pima }\OtherTok{\textless{}{-}}\NormalTok{ PimaIndiansDiabetes2 }\SpecialCharTok{|\textgreater{}} 
  \FunctionTok{drop\_na}\NormalTok{() }\SpecialCharTok{|\textgreater{}} \CommentTok{\# Using only complete observations}
  \FunctionTok{mutate}\NormalTok{(}\AttributeTok{index =} \DecValTok{1}\SpecialCharTok{:}\FunctionTok{n}\NormalTok{())}

\CommentTok{\# Create training and testing set}
\NormalTok{training\_set }\OtherTok{\textless{}{-}}\NormalTok{ pima }\SpecialCharTok{|\textgreater{}} 
  \FunctionTok{sample\_frac}\NormalTok{(}\FloatTok{0.8}\NormalTok{) }\CommentTok{\# Use 80\% of the data in training set}
\NormalTok{test\_set }\OtherTok{\textless{}{-}}\NormalTok{ pima }\SpecialCharTok{|\textgreater{}} 
  \FunctionTok{anti\_join}\NormalTok{(training\_set, }\AttributeTok{by =} \StringTok{"index"}\NormalTok{)}

\CommentTok{\# Train model on training set}
\NormalTok{model }\OtherTok{\textless{}{-}}\NormalTok{ training\_set }\SpecialCharTok{|\textgreater{}} 
  \FunctionTok{glm}\NormalTok{(}\AttributeTok{formula =}\NormalTok{ diabetes }\SpecialCharTok{\textasciitilde{}}\NormalTok{ ., }\AttributeTok{family =} \StringTok{"binomial"}\NormalTok{)}

\CommentTok{\# Evaluate model by accuracy on test set}
\NormalTok{predictions }\OtherTok{\textless{}{-}}\NormalTok{ model }\SpecialCharTok{|\textgreater{}} 
  \FunctionTok{predict}\NormalTok{(test\_set, }\AttributeTok{type =} \StringTok{"response"}\NormalTok{)}
\NormalTok{test\_set }\SpecialCharTok{|\textgreater{}} 
  \FunctionTok{mutate}\NormalTok{(}\AttributeTok{prediction =} \FunctionTok{ifelse}\NormalTok{(predictions }\SpecialCharTok{\textgreater{}} \FloatTok{0.5}\NormalTok{, }\StringTok{"pos"}\NormalTok{, }\StringTok{"neg"}\NormalTok{)) }\SpecialCharTok{|\textgreater{}} 
  \FunctionTok{summarize}\NormalTok{(}\AttributeTok{accuracy =} \FunctionTok{sum}\NormalTok{(prediction }\SpecialCharTok{==}\NormalTok{ diabetes) }\SpecialCharTok{/} \FunctionTok{n}\NormalTok{())}
\end{Highlighting}
\end{Shaded}

\begin{verbatim}
   accuracy
1 0.7820513
\end{verbatim}

\begin{Shaded}
\begin{Highlighting}[]
\NormalTok{model }\OtherTok{\textless{}{-}}\NormalTok{ training\_set }\SpecialCharTok{|\textgreater{}} 
  \FunctionTok{glm}\NormalTok{(}\AttributeTok{formula =}\NormalTok{ diabetes }\SpecialCharTok{\textasciitilde{}}\NormalTok{ . }\SpecialCharTok{{-}}\NormalTok{ pressure, }\AttributeTok{family =} \StringTok{"binomial"}\NormalTok{)}

\CommentTok{\# Evaluate model by accuracy on test set}
\NormalTok{predictions }\OtherTok{\textless{}{-}}\NormalTok{ model }\SpecialCharTok{|\textgreater{}} 
  \FunctionTok{predict}\NormalTok{(test\_set, }\AttributeTok{type =} \StringTok{"response"}\NormalTok{)}
\NormalTok{test\_set }\SpecialCharTok{|\textgreater{}} 
  \FunctionTok{mutate}\NormalTok{(}\AttributeTok{prediction =} \FunctionTok{ifelse}\NormalTok{(predictions }\SpecialCharTok{\textgreater{}} \FloatTok{0.5}\NormalTok{, }\StringTok{"pos"}\NormalTok{, }\StringTok{"neg"}\NormalTok{)) }\SpecialCharTok{|\textgreater{}} 
  \FunctionTok{summarize}\NormalTok{(}\AttributeTok{accuracy =} \FunctionTok{sum}\NormalTok{(prediction }\SpecialCharTok{==}\NormalTok{ diabetes) }\SpecialCharTok{/} \FunctionTok{n}\NormalTok{())}
\end{Highlighting}
\end{Shaded}

\begin{verbatim}
   accuracy
1 0.7692308
\end{verbatim}

To properly study accuracy it is good to use more than one train-test
combination or using k-fold validation
\href{'https://en.wikipedia.org/wiki/Cross-validation_(statistics)'}{cross-validation}.
However, this is beyond the scope of this course.

We see that including \texttt{pressure} doesn't seem to improve our
accuracy much, suggesting it is better to go for the smaller, less
complex model.

\subsection{Exercise 4}\label{exercise-4-14}

The model would achieve an accuracy of 90\% since
\(\text{Accuracy} = \#\text{Correct predictions} / \#\text{Total predictions} = 90 / 100 = 0.9 = 90 \%\).



\end{document}
